% 本文件由 LaTeX 排版 Agent 根据有效 translation.json 自动生成。
% author=Sozomen (c. 375-c. 447); work_id=2602; title=教会史

\chapter{教会史}

% structure: p0001 source=h1 target=omitted reason=page-title-already-rendered-by-parent

\Emph{致辞与提议}\newline{}\Emph{卷一}\newline{}\Emph{卷二}\newline{}\Emph{卷三}\newline{}\Emph{卷四}\newline{}\Emph{卷五}\newline{}\Emph{卷六}\newline{}\Emph{卷七}\newline{}\Emph{卷八}\newline{}\Emph{卷九}\par

\SourceRecord{Translated by Chester D. Hartranft.；Nicene and Post-Nicene Fathers, Second Series；Vol. 2.；Edited by Philip Schaff and Henry Wace.；Buffalo, NY: Christian Literature Publishing Co.,；1890.；<http://www.newadvent.org/fathers/2602.htm>.}

% structure: p0001 source=h1 target=section reason=page-title

\section{呈文与建议}

致\Person{狄奥多西}皇帝的呈文及\WorkTitle{教会史}建议书，\Person{撒拉米尼乌斯·赫尔米亚斯·索佐曼}著。\par

民间常言，先前的皇帝们各有所好：有的喜爱装饰，关心皇家紫袍、王冠之类；有的好学，撰写能令读者着迷的神话作品或论文；有的习武，力求投枪正中目标、猎杀野兽、掷矛或跃马。凡献身于某一取悦君主之技艺者，皆入宫自荐。有人带来一颗难以打磨的宝石；有人致力于调制比紫袍更鲜亮的颜色；有人献上诗歌或论文；有人展示一种精巧而奇特的铠甲。\par

人们认为，对全体人民的统治者来说，至少拥有一种家庭美德是最伟大的、君王般的事；但他们对虔诚却未作如此高的评价，而虔诚终究是帝国真正的装饰。然而，最强大的皇帝\Person{狄奥多西}啊，您一句话，靠天主的助佑，培养了一切美德。身披紫袍和冠冕，作为旁观者眼中您尊严的象征，您内心常佩戴着统治权柄的真正装饰：虔诚与爱人之心。因此，诗人和作家，以及您的大部分官员和其余臣民，无时无刻不关注您和您的事迹。当您作为竞赛的裁判和演讲的评判者主持时，任何矫饰的声音和形式都无法剥夺您的准确性，而是真诚地颁发奖品，观察措辞是否适合作品的构思；同样，对于词语的形式、划分、秩序、统一、用语、结构、论证、思想和叙事，您也如此。您用赞许的判断和掌声、金像、竖立雕像、礼物以及各种荣誉来酬报演讲者。您对演讲者表现出的个人恩宠超过了古代克里特人对广为赞颂的\Person{荷马}，或\Person{阿勒瓦达}对\Person{西摩尼得斯}，或西西里暴君\Person{狄奥尼修}对\Person{苏格拉底}的同伴\Person{柏拉图}，或马其顿的\Person{腓力}对历史学家\Person{泰奥彭波斯}，或\Person{塞维鲁}皇帝对以诗歌叙述鱼类种类、特性和捕捞的\Person{奥皮安}。因为克里特人用一千\OriginalTerm{努米}{nummi}奖赏\Person{荷马}后，他们将奖赏的数额刻在公共柱子上，仿佛要夸耀他们过分的慷慨。\Person{阿勒瓦达}、\Person{狄奥尼修}和\Person{腓力}并不比克里特人更谨慎，克里特人夸耀自己温和而哲学化的治理，但他们很快模仿了克里特人的柱子，以免在捐赠上逊色。但当\Person{塞维鲁}赐予\Person{奥皮安}为他的每行中等诗句一枚金礼时，他的慷慨使所有人惊讶，以至于\Person{奥皮安}的诗作至今被称为\Quote{金言}。这就是过往学问与演讲爱好者的捐赠。但是，皇帝啊，您在对文学的慷慨上超越了任何古人，在我看来，您这样做并非没有道理。因为您不仅努力以您的美德征服一切，而且根据您对古代事迹（希腊人和罗马人如此成功地指导这些事迹）的透彻了解，您也成功处理了自己的事务。传闻说，白天您进行军事和身体锻炼，通过司法裁决、记录必要事项以及观察公私应做之事来安排国事；晚上则忙于读书。有一个说法：有一盏灯为您研读这些著作服务，这盏灯通过某种机制使油自动流入灯芯，这样王宫中没有一个仆人被迫分担您的辛劳，且因与睡眠斗争而违背本性。因此，您既对近臣也对所有人仁慈温和，因为您效法作为您楷模的天上君王；祂喜爱降雨，使日光照耀义人和不义之人，并毫不吝啬地赐予其他恩惠。自然，我也听说，凭借您的渊博学识，您对石头的性质、草根的功效和药物的力量，绝不亚于\Person{达味}最智慧的儿子\Person{撒罗满}；而在美德上您超过他，因为\Person{撒罗满}成为情欲的奴隶，没有保持那曾为他带来繁荣和智慧的虔诚到底。但是，最强大的皇帝啊，因为您用您的约束理性对抗轻浮，您不仅是人的独裁者，也是灵魂和身体情欲的独裁者，这是人们自然会认为的。还有这一点也值得注意：我了解到您克服了对一切饮食的欲望；无论更甜的无花果（诗意地说），还是其他任何当季水果，都不能俘虏您，除了您在感谢万有的创造者后少量触摸和品尝的。您习惯通过日常锻炼征服口渴、闷热和寒冷，以至于您似乎将自制作为第二天性。最近，众所周知，您急于访问\Place{本都}的\Place{赫拉克利亚}城，并重建被时间摧毁的它，您在夏季取道\Place{比提尼亚}。当中午太阳非常炽热时，一名亲卫看见您因大量汗水和尘土而发热，仿佛为了施恩于您，他预先递给您一个碗，碗反射出太阳的光芒；他倒入一些甜饮料，并加入冷水。但是，最强大的皇帝，您接过了它，称赞此人的好意，并明确表示您将很快以王室的慷慨酬报他这精心制作的举动。但是，当所有士兵张着嘴惊奇地看着碗，并认为那喝饮料的人有福时，您，尊贵的皇帝，将饮料还给他，命令他随意使用。因此，在我看来，马其顿的\Person{腓力}之子\Person{亚历山大}也被您的美德超越了；他的崇拜者传言，当他和马其顿人经过一个无水之地时，一名焦急的士兵找到了水，舀了它，献给亚历山大；他不肯喝，却倒掉了饮料。因此，简言之，按照\Person{荷马}的说法，称您比在您之前的国王更为君王是恰当的；因为我们听说过有些人没有获得任何值得钦佩的东西，另一些人则用一两件事迹装饰他们的统治。但是，最强大的皇帝啊，您聚集了所有美德，在虔诚、爱人之心、勇敢、明智、正义、慷慨以及与君王尊严相称的豪迈上超越了每个人。每个时代将夸耀您的统治是唯一清白、纯洁、没有谋杀的，超出所有曾存在的政府。您教导您的臣民愉快地追求严肃的事务，使他们为您和公共事务表现出热情、善意和尊重。因此，基于所有这些原因，作为\WorkTitle{教会史}的著者，我认为有必要向您致辞。因为我即将叙述许多虔诚之人的美德和\OriginalTerm{天主教会}{Catholic Church}的事件，还能向谁更合适地这样做呢？既然她与众多敌人的斗争引导我来到您和您的祖先的门前？来吧，您，无所不知，拥有一切美德，特别是虔诚——圣言说：\BibleQuote{虔诚是智慧的开端}——从我这里接受这部著作，并用您的辛劳调整其事实，凭您的准确知识加以纯化，无论是增添还是删减。因为任何您看来合意的方式，对读者都将完全有益和辉煌，在您认可之后，将无人再改动它。我的历史始于凯撒\Person{克里斯普斯}和\Person{君士坦丁}的第三度执政官任期，延伸到您的第十七度执政官任期。我决定将整部作品分为九部分：第一和第二卷将涵盖\Person{君士坦丁}治下的教会事务；第三和第四卷，他儿子们治下的；第五和第六卷，伟大的\Person{君士坦丁}的儿子们的堂弟\Person{尤利安}、\Person{约维安}，以及\Person{瓦伦提尼安}和\Person{瓦伦斯}治下的；第七和第八卷，最强大的皇帝啊，将揭示\Person{格拉提安}和\Person{瓦伦提尼安}兄弟治下的事务，直到您的神圣祖父\Person{狄奥多西}被宣布，以及您著名的父亲\Person{阿卡狄乌斯}与您最虔诚和敬神的叔父\Person{霍诺留}共同接受父权统治并参与管理罗马世界；第九卷我献给热爱基督、最纯洁的陛下，愿\Person{天主}永远保佑您在完整的善意中，大大战胜仇敌，使万有在您脚下，并将神圣帝国传承给您的子子孙孙，蒙\Person{基督}的赞许，藉着祂、并与祂一起，愿光荣归于\Person{天主}、\Person{圣父}及\Person{圣神}，直到永远。阿们。\par

\SourceRecord{Translated by Chester D. Hartranft.；Nicene and Post-Nicene Fathers, Second Series；Vol. 2.；Edited by Philip Schaff and Henry Wace.；Buffalo, NY: Christian Literature Publishing Co.,；1890.；<http://www.newadvent.org/fathers/26020.htm>.}

% structure: p0001 source=h1 target=section reason=page-title

\section{\WorkTitle{教会史（第二卷）}}

% structure: p0002 source=h2 target=subsection reason=relative-heading-rank; base=h2

\subsection{第一章 本书序言：探讨犹太民族的历史；提及开始此类著作的人；如何以及从哪些来源搜集历史；如何致力于真理；以及历史还将包含哪些细节。}

我常思考：为何其他人很容易相信天主圣言，而犹太人却如此不信，尽管自起初，关于天主事物的教导就由先知传授给他们，并且先知们也使他们预知基督来临之前的事件。此外，他们民族的创始人和割礼的建立者亚巴郎，曾蒙恩亲眼目睹并款待了天主之子。\EditorialNote{创世纪第十八章}他的儿子依撒格，则被尊为十字架祭献的预像，因为他被他父亲捆绑着牵到祭台上；正如精通圣经的人所肯定的，基督的苦难也以类似的方式发生。雅各伯预言万民的期望将是基督，正如现在这样；他也预言了基督来临的时间，说：\Quote{犹大支派的领袖，即部落首领，将要衰败。}\BibleRef{创 49:10}\par

这明确指的是黑落德的统治，他父系是厄东人，母系是阿拉伯人；犹太民族由罗马元老院和奥古斯都·凯撒交给了他。其余的先知中，有些预言了基督的诞生、祂不可言喻的受孕、祂出生后母亲仍保持童贞、祂的民族和故乡；有些预言了祂神圣而奇妙的事迹；另一些则预言了祂的苦难、祂从死者中复活、祂升天，以及伴随每件事的事件。如果有人对这些事实无知，阅读圣经便不难知晓。玛塔提雅的儿子若瑟夫，他是一位司祭，在犹太人和罗马人中最为杰出，可作为关于基督真理的重要见证；因为他犹豫是否称祂为人，因祂行了奇妙的事迹，并是真理教义的导师，却公开称祂为基督；祂被判处十字架之死，并在第三日重新显现。若瑟夫对圣先知们事先预言关于基督的无数其他奇妙之事也并非无知。他还见证基督吸引了许多希腊人和犹太人归向祂，这些人持续爱祂，并且以祂命名的民族并未灭绝。在我看来，在叙述这些事情时，他几乎是在宣告，通过比较作为，基督就是天主。仿佛被奇迹所震撼，他某种程度上走了一条中间路线，既不攻击那些信仰耶稣的人，反而赞同他们。\par

当我思考此事时，在我看来相当值得注意：希伯来人竟没有预先料到并在其他人之前立即皈依基督教；因为尽管西碧尔和一些神谕曾预先宣布关于基督的将来事件，但我们不能因此将不信归咎于所有希腊人。因为只有少数人，在学识上优越，才能理解这类预言，而这些预言大部分以诗歌形式并用隐晦的词语向民众宣布。因此，在我看来，这是出于天主的预知，为了将来事件的一致性，即将发生的事实不仅要通过祂自己的先知，也要部分通过外邦人来显明。就像一个音乐家，在陌生旋律的压力下，可能用拨子轻拂琴弦的多余音调，或在已有的音符上添加其他音符。\par

既然现在已表明，希伯来人虽拥有众多且更明确的关于基督来临的预言，却比希腊人更不愿接受对祂的信仰，关于此事的论述就到此为止。然而，决不可因此认为教会主要通过其他民族的皈依而建立是不合理的；因为首先，显而易见，在神圣而伟大的事务中，天主喜欢以奇妙的方式促成变化；其次，要记住，正是凭借非凡的德行，那些在起初负责宗教事务的人维持了他们的影响。如果他们确实没有具备锐利的语言表达能力或优美的措辞，也没有通过言语或数学证明说服听众的能力，他们却仍然完成了他们承担的工作。他们放弃财产，忽视亲属，被钉在十字架上，仿佛拥有不属于自己的身躯，忍受了许多惨烈的折磨；既不被任何城市的人民和统治者的奉承所诱惑，也不被他们的恐吓所吓倒，他们通过自己的行为清楚地证明，在斗争中他们得到了丰厚奖赏的盼望的支持。因此，他们实际上无需诉诸言语论证；因为他们毫不费力，他们的行为本身迫使每个家庭和每个城市的居民相信他们的见证，甚至在知道其内容之前。\par

既然人类的状况发生了如此神圣而奇妙的转变，以致古代的崇拜和民族的法律遭到蔑视；既然许多著名的希腊作家曾竭尽其雄辩之才描述卡吕冬野猪、马拉松公牛以及其他类似奇事——这些事确实发生在某些国家或城市，或有神秘的起源——为什么我不超越自我，撰写一部教会史呢？因为我深信，既然主题不是人的成就，我写这样一部历史似乎几乎难以置信；但是，在天主，没有不可能的事。\par

我起初强烈倾向于从事件的开端追溯历程；但考虑到那些最有智慧的人——\Person{克雷孟}、\Person{赫格西普}，宗徒的继承者，史学家\Person{阿非利加努}，以及以\Person{旁非鲁}著称的\Person{欧瑟伯}，一位深谙\WorkTitle{圣经}及希腊诗人与史学家著作的人——已经编纂了直至他们那个时代的类似历史记录，我便只以两卷简述从\Person{基督}升天到\Person{利西尼}被废黜期间教会所发生的一切。如今，因\Person{天主}的助佑，我也将尽力叙述此后的事件。\par

我将记录我所参与的事件，以及从那些亲身知晓或目睹本时代或前一代事务的人那里听闻的事。但对于更早的事件，我查阅了与宗教相关的既定法律、同时期公会议的进程、出现的创新，以及君王与司铎的书信。这些文献有些保存在宫廷与教会中，有些则散佚于学者手中。我曾多次考虑全文转录，但进一步思量后，鉴于文献数量庞大，我决定仅给出其内容的简要概述；然而每当引入争议话题时，我将随时自由转录任何有助于阐明真理的作品。若有对过往事件无知者，因见其他著作中的矛盾记载而断定我的历史为伪，请其知晓：自从\Person{亚略}的教义及其他更新的假说被提出后，教会领袖们彼此意见分歧，为自己的追随者留下书面的独特见解；而且，请谨记，这些领袖召开公会议并随心所欲地颁布法令，往往未经听证就谴责那些信条不同于他们的人，并竭尽所能劝说当时的在位君王与贵族支持他们。各派系的党徒为维护自己教义的正统性，分别收集利于自己异端的书信，而省略所有相反倾向的文件。这便是我们在试图就此事得出结论时所面临的障碍！尽管如此，为保持历史准确性，必须对获取真理的手段予以最严格关注，我仍感到必须尽力检视这一类所有著作。\par

愿无人指责我怀有冒昧或恶意，因我详述了圣职人员之间关于首席地位及自身异端优越性的争论。首先，如我所言，历史学家应将真理视为重于一切；此外，天主教会的教义被证明尤其最为纯正，因为它屡次经受对立思想者的阴谋考验；然而，\Person{天主}的命定使天主教会保持了自身的优越，重掌其权能，并引导所有教会与人民接受其真理。\par

我曾思考是否应仅限于叙述罗马政府统治下教会相关的事件；但看来更明智的做法是尽量纳入\Place{波斯}与野蛮人之间与宗教相关的事迹。在本著作中引入所谓隐修生活的始祖与创始人及其直接继承者——这些人因观察或传闻而声名显赫——的记述，也并非与教会史无关。因为我不愿被视为对他们忘恩负义，不愿将他们的德行投入遗忘，也不愿被认为对他们的历史无知；我愿留下他们生活方式的记录，使他人效法其榜样，获得有福而圆满的结局。随着著作进行，这些主题将尽可能地有所记述。\par

呼求\Person{天主}的帮助与慈恩，我现在开始叙述事件；本历史将从这一点起始。\par

% structure: p0013 source=h2 target=subsection reason=relative-heading-rank; base=h2

\subsection{第二章　\Person{君士坦丁}在位时各大城的主教；以及因恐惧\Person{利西尼}，东方直至\Place{利比亚}的基督教如何谨慎地宣认，而在西方，因\Person{君士坦丁}的恩惠，基督教得以自由宣认。}

在\Person{君士坦丁}凯撒与\Person{克里斯普斯}凯撒担任执政官期间，\Person{西尔维斯特}治理罗马教会；\Person{亚历山大}治理亚历山大教会；\Person{玛加略}治理耶路撒冷教会。自\Person{罗马努}以来，无人被任命管理\Place{奥伦特河畔安提约基雅}教会；显然，迫害阻碍了授任礼仪的举行。但不久后在\Place{尼西亚}召开的公会议的主教们深知\Person{欧斯塔修}的品行与教义之纯洁，判定他配得上宗徒之座；尽管他当时是邻近的\Place{贝洛亚}的主教，他们仍将他调任至\Place{安提约基雅}。\par

东方直至\Place{埃及}边境\Place{利比亚}的基督徒，不敢公开以教会形式聚会；因为\Person{利西尼}已收回对他们的恩惠。但西方的基督徒——\Place{希腊}人、\Place{马其顿}人、\Place{伊利里亚}人——在\Person{君士坦丁}的保护下安全地聚会敬拜，君士坦丁当时是\Place{罗马帝国}的元首。\par

% structure: p0016 source=h2 target=subsection reason=relative-heading-rank; base=h2

\subsection{第三章　藉十字架异象与\Person{基督}显现，\Person{君士坦丁}被引导信奉基督教。——他从我们的弟兄接受宗教训诲。}

据我们所知，\Person{君士坦丁}因若干不同事件的共同作用，特别因一个来自天上的征兆显现，而受到激励尊崇基督教。\par

当他最初下定决心与马克森提乌斯开战时，他对如何实施军事行动以及从何处寻求援助感到困惑。在他不知所措之际，他在异象中看见十字架的光芒照耀在天上。他对这景象感到惊奇，但一些站在旁边的圣天使喊道：\Quote{噢，君士坦丁！凭此标记，征服吧！}据说基督亲自向他显现，并向他展示了十字架的标记，命令他制作一个同样的标记，并在战斗中将其作为他的帮助，因为它将确保胜利。\par

欧瑟伯（别号庞菲鲁斯）断言，他曾听皇帝发誓说，当太阳在正午时分即将偏移时，他与随行的士兵看见在天上有一个由光组成的十字架凯旋标记，并被以下话语环绕：\Quote{凭此标记，征服。}\par

这个异象在他正困惑该将军队引向何处时出现在他路上。当他思考这可能意味着什么时，夜晚来临；他入睡时，基督带着他曾在天上见过的标记显现，并命令他制作一个该标记的复制品，并在敌对交锋中将其作为他的帮助。无需进一步说明；因为皇帝清楚地领悟到事奉天主的必要性。\par

天亮时，他召集了基督的司铎，询问他们的教义。他们翻开圣经，阐释了关于基督的真理，并从先知书中向他指出，那些曾预言过的迹象如何已经应验。他们说，向他显现的标记是战胜地狱的象征；因为基督来到人间，被钉在十字架上，死亡，并在第三日复活。因此，他们说，有希望在当前时代结束时，将有普遍的死人复活，并进入不朽，那时行善者将得相应赏报，作恶者将受惩罚。然而，他们继续说，救恩和从罪中得洁净的方法已被提供；即，对于未入门者，按教会规章施行入门礼；对于已入门者，则要避免再犯罪。但即使是在圣人中，也很少有人能遵守这后一条件，因此设立了另一种洁净的方法，即悔改；因为天主出于其对人类的爱，将宽恕赐予那些犯罪后悔改，并以善行证实其悔改的人。\par

% structure: p0022 source=h2 target=subsection reason=relative-heading-rank; base=h2

\subsection{第四章 君士坦丁命令在战斗中把十字架标记抬在他前面；关于十字架标记抬举者的非凡记述.}

皇帝对司铎们为他阐释的关于基督的预言感到惊奇，便召来一些熟练的工匠，命令他们改造罗马人称为\Emph{Labarum}的军旗，将其改造为十字架的形状，并用黄金和宝石装饰它。这个战争标记比其他所有标记都更受珍视；因为它总是被抬在皇帝前面，并被士兵崇敬。我认为君士坦丁将罗马权力中最尊贵的象征改为基督的标记，主要是为了让士兵们通过习惯经常看到它并崇敬它，从而放弃他们古老的迷信形式，并承认皇帝所敬拜的真天主为他们的领袖和战斗中的帮助；因为这个标记总是被抬在他的军队前面，并且奉皇帝的命令，由一队杰出的持矛者在战斗最激烈时抬着穿过方阵，每个人轮流将标记扛在肩上，在队列中巡行。据说有一次，在敌军意外调动时，手持标记的人因恐惧将标记交给另一个人，并秘密逃离战场。当他逃出敌人武器射程之外时，突然中箭倒下，而那个站在神圣标记旁边的人却安然无恙，尽管许多武器瞄准了他；因为敌军的箭矢被神奇地引导到标记上，而抬着标记的人在危险中得以保全。\par

还传说，没有哪个在战斗中抬着这个标记的士兵曾遭遇到战争中常发生的黑暗灾祸而倒下，或受伤，或被俘。\par

% structure: p0025 source=h2 target=subsection reason=relative-heading-rank; base=h2

\subsection{第五章 驳斥君士坦丁因杀害其子克里斯普斯而成为基督徒的断言。}

我知道异教徒中流传着一种说法：\Person{君士坦丁}在杀死自己的一些近亲，特别是同意处死其子\Person{克里斯普斯}之后，对自己的恶行感到懊悔，便向当时执掌\Person{普罗提诺}学派的哲学家\Person{索帕特}请教除去罪污的方法。据传说，哲学家回答说，\Quote{这样的道德污秽无法得到净化}。皇帝因此沮丧，但偶然遇到几位主教，他们告诉他，\Quote{只要悔改并接受圣洗，就能从罪中获得洁净}。他听了这些话很高兴，钦佩他们的教义，于是成为基督徒，并引导臣民信奉同一信仰。在我看来，这故事是那些想要诋毁基督宗教的人的捏造。据说\Person{君士坦丁}需要净化的原因——\Person{克里斯普斯}，直到他父亲统治的第二十年才去世；他享有帝国的第二高位，拥有凯撒的称号，并且许多由他批准的支持基督宗教的法令至今犹存。通过查考这些法令的日期和立法者名单，可以证明这一点。\Person{索帕特}似乎不太可能与\Person{君士坦丁}有任何往来，因为当时\Person{君士坦丁}的政府位于\Place{大洋}和\Place{莱茵河}附近地区；他与\Place{意大利}统治者\Person{马克森提乌斯}的争执在\Place{罗马}领土上造成了如此大的分裂，以至于当时在\Place{高卢}、\Place{不列颠}或邻近地区定居并非易事，而普遍承认\Person{君士坦丁}在与\Person{马克森提乌斯}作战之前、返回\Place{罗马}和\Place{意大利}之前，就已经接受了基督宗教；这从他颁布的支持宗教的法令的日期可以得到证明。但即便假设\Person{索帕特}偶然遇到了皇帝，或与他有书信往来，也不能想象这位哲学家不知道\Person{阿尔克墨涅}之子\Person{赫拉克勒斯}在杀死自己的孩子以及他的客人兼朋友\Person{伊菲托斯}之后，通过参加\Person{谷神}的神秘仪式在\Place{雅典}获得了净化。我刚刚举的例子就足以说明，\Place{希腊}人认为这种性质的罪污可以得到净化；因此，声称\Person{索帕特}教导相反论调的人是在诽谤。\par

我不能认为这位哲学家可能对这些事实一无所知；因为他在当时被视为\Place{希腊}最有学问的人。\par

% structure: p0028 source=h2 target=subsection reason=relative-heading-rank; base=h2

\subsection{第六章 \Person{君士坦丁}的父亲容许\Person{基督}的名号得以扩展；\Person{君士坦丁大帝}预备它渗透各处。}

在\Person{君士坦丁}的统治下，教会兴旺发达，信徒人数日益增多，因为它们受到一位仁慈善良的皇帝的善行所尊荣，而天主也在其他方面保护它们免受先前所遭遇的迫害和骚扰。当世界其他地区的教会遭受迫害时，唯有\Person{君士坦丁}的父亲\Person{君士坦提乌斯}允许基督徒无所畏惧地敬拜天主。我知道他做了一件了不起的事，值得记载。他想考验某些与他宫廷关系密切的优秀善良基督徒的忠诚。他把他们全部召来，告诉他们，如果他们既拜偶像又侍奉天主，就可以留在他的宫廷并保留职位；但如果拒绝服从他的意愿，就会被赶出宫廷，并且几乎难逃他的报复。当意见分歧使他们分成两派，即那些同意放弃信仰的人和那些宁愿以现世利益换取天主荣耀的人，皇帝决定保留那些坚持信仰的人作为他的朋友和顾问；但他转而远离其他人，视他们为懦夫和骗子，并将他们逐出宫廷，断定那些轻易背叛天主的人永远不会忠于他们的君王。因此，很可能在\Person{君士坦提乌斯}在世时，\Place{意大利}以外的地区居民，即\Place{高卢}、\Place{不列颠}或\Place{比利牛斯山脉}地区直至\Place{西海}，公开宣认基督宗教并不违反法律。当\Person{君士坦丁}继承同一政府时，教会的事务变得更加辉煌；因为\Person{赫库利乌斯}之子\Person{马克森提乌斯}被杀后，他的份额也归\Person{君士坦丁}所有；\Place{台伯河}和\Place{埃里达努斯河}（当地人称为\Place{波河}）沿岸的居民，以及\Place{阿奎利斯河}（据说\WorkTitle{阿尔戈号}曾被拖曳至此）沿岸的居民，和\Place{第勒尼安海}沿岸的居民，都被允许不受干扰地行使他们的宗教。\par

当\Person{阿尔戈英雄}逃离\Person{埃厄忒斯}时，他们从另一条路线返回故乡，跨越\Place{斯基泰海}，穿过那里的几条河流，最终抵达\Place{意大利}海岸，在那里过冬并建立了一座城市，他们称之为\Place{埃莫纳}。次年夏天，在当地人的帮助下，他们借助机械将\WorkTitle{阿尔戈号}拖曳了四百\OriginalTerm{斯塔迪亚}{stadia}，从而到达\Place{阿奎利斯河}，该河流入\Place{埃里达努斯河}；\Place{埃里达努斯河}本身流入\Place{意大利海}。\par

\Place{西巴利斯}战役后，\Person{达尔达尼亚人}和\Person{马其顿人}、\Place{伊斯特河}沿岸居民、\Person{希腊人}以及整个\Person{伊利里亚}民族都归顺了\Person{君士坦丁}。\par

% structure: p0032 source=h2 target=subsection reason=relative-heading-rank; base=h2

\subsection{第七章 关于\Person{君士坦丁}与其妹夫\Person{李锡尼}因基督徒引起的争端，以及\Person{李锡尼}如何被武力征服并处死。}

遭受这次挫折之后，先前尊重基督徒的\Person{里奇纽斯}改变了看法，苛待了许多在他统治下的司铎；他也迫害了许多其他人，尤其是士兵。因为与\Person{君士坦丁}的分歧，他对基督徒深怀愤怒，想借着他们为宗教所受的苦难来伤害他；此外，他怀疑教会祈祷并热心希望\Person{君士坦丁}独自享有至高统治权。除此之外，当与\Person{君士坦丁}再次交战前夕，\Person{里奇纽斯}如同惯常所做的那样，通过祭献和神谕对预期的战争进行了预测，并因被胜利的许诺所欺骗，他回到了异教徒的宗教。\par

异教徒们自己也叙述说，大约在这个时候，他在\Place{米利都}向\Person{阿波罗}·\Place{狄迪默斯}的神谕求问，并从魔鬼那里得到了关于战争结果的回答，以\Person{荷马}的以下诗句表述：\par

\Quote{老者啊，年轻人使你备受困扰，与你争战！\newline{}你的力量已变得衰弱，但你的老迈却仍将坚毅。}\par

从许多事实中，我常常看出基督徒的教导得到天主眷顾的支持，其推进也由此保障；而尤以当时所发生的事为最：就在\Person{里奇纽斯}即将迫害他治下所有教会之时，\Place{比提尼亚}的战争爆发了，最终导致了他与\Person{君士坦丁}之间的战争，而\Person{君士坦丁}因天主的援助大为增强，以致在海陆两方都战胜了他的敌人。在他的舰队和军队覆灭后，\Person{里奇纽斯}逃入\Place{尼科米底亚}，之后作为平民在\Place{塞萨洛尼卡}居住了一段时间，并最终在那里被杀害。这就是这样一位统治者的结局：他在统治之初曾在战争与和平中卓著功勋，并曾因娶得\Person{君士坦丁}之妹而受尊荣。\par

% structure: p0037 source=h2 target=subsection reason=relative-heading-rank; base=h2

\subsection{第八章 \Person{君士坦丁}在基督徒自由和建造教堂方面所赐恩惠的清单；以及其他为公共福祉所行的事。}

一旦罗马帝国的唯一统治权归于\Person{君士坦丁}，他便发布了一道公开敕令，命令他在东方的所有臣民尊崇基督宗教，谨慎敬拜天主，并承认唯有天主才是本质神圣且拥有永永远远能力的：因为他乐于毫不吝惜地将一切美善赐给那些热诚拥抱真理的人；他以最好的期盼对待他们的努力，而无论是和平还是战争、公共还是私人生活中的不幸，则临到违命者身上。\Person{君士坦丁}接着补充道，但并非出于虚夸：既然天主视他为合用的仆人，堪当统治，他便从\Place{不列颠}海域被引领至东方诸省，为的是基督宗教得以扩展，而那些因敬拜天主在宣认或殉道中坚持的人得以被提升至公开的尊荣。在作了这些声明之后，他又提出了无数其他细节，以为可使他的臣民被吸引归向宗教。他下令撤销所有教会迫害者对基督信仰所通过的法令和判决；并命令所有因宣认基督而被流放——无论是到海岛还是其他地方，违背他们自己意愿的——以及所有被判在矿场、公共工程、后宫、亚麻布厂劳动，或被编入公职人员名册的人，都应恢复自由。他从那些被加诸耻辱的人身上除去耻辱的污名，并允许那些被剥夺军中高职的人，要么重新担任原有职位，要么以荣誉退役的方式，按自己选择享受宽松的安闲生活；当他召回了所有人享受原有的自由和惯常的荣誉之后，他也归还了他们的财产。对于那些被杀且财产被充公的人，他规定遗产应转交给最近的亲属，若无继承人，则转交给该地产所在地区的教会；当遗产已转入他人之手，并已成为私有或国有财产时，他命令将其归还。他同样承诺，当财产已由国库购买或从国库接收作为赠与时，将采取最恰当和最好的安排。如上所述，这些措施经皇帝颁布并由法律批准后，立即付诸实施。基督徒因此几乎被安置在罗马政府的所有主要职位上；假神的敬拜被普遍禁止；占卜之术、雕像的奉献以及异教节日的庆祝均被取缔。在众多城市中保有的许多最古老习俗逐渐废止：在埃及人中，用来标示\Place{尼罗河}水上涨的尺度不再被抬入异教庙宇，而是抬入教堂。角斗士的表演随后在罗马人中也被禁止；在\Place{黎巴嫩}和\Place{赫利奥波利斯}的腓尼基人中，婚前使处女卖淫的习俗——她们在经历一次非法交合后便惯于合法结婚——也被废除。关于祈祷之所，皇帝修缮了一些规模足够的；另一些则因增加长度和宽度而变得辉煌，他还在这类建筑先前未曾存在的地方建造了新的建筑。他从皇室国库提供了必要的供给，并致信各城的主教和各行省的总督，要求他们贡献任何所愿之物，并命令他们服从并热诚服从于司铎。\par

宗教的繁荣与帝国的日益繁荣同步发展。与\Person{李锡尼}的战争之后，皇帝在与外国的战斗中取得了胜利；他征服了\Place{萨尔马特人}和称为\Place{哥特人}的民族，并与他们订立了有利的条约。这些民族居住在\Place{伊斯特河}畔；由于他们非常尚武，且因人数众多和身体魁梧而常备不懈，他们使其他蛮族部落感到敬畏，并只与罗马人抗衡。据说，在这次战争中，\Person{君士坦丁}通过征兆和梦境清楚地看到，\OriginalTerm{天主眷顾}{Divine Providence}的特殊保护已临于他。因此，当他战胜了所有在战场上起来反对他的人，他通过热切关注宗教事务表达了对\Person{基督}的感激，并劝勉总督们认信唯一的真信仰和救恩之路。他制定法律，规定从附属国征收的部分资金应由各城市转交给主教和神职人员，无论他们居住何处，并命令这项赠予的法律应永为成规。为了训练士兵像他一样敬拜天主，他在他们的武器上刻上十字架的标记，并在皇宫内建造了一座祈祷所。当他参战时，他让人在他面前抬着一个建成教堂形状的帐篷，以便他或他的军队若被领入旷野时，能有一座神圣建筑在其中赞美和朝拜天主，并参与奥迹。司铎和执事跟随帐篷，按照教会法律执行这些事宜的命令。从那时起，罗马军团（当时以编号称呼）各自提供自己的帐篷，并配有随行司铎和执事。他还命令遵守称为主日的日子（犹太人称为一周的第一天，异教徒献给太阳神的日子），以及第七日之前一日，并下令在这些日子不得处理司法或其他事务，而应以祈祷和恳求侍奉天主。他尊崇主日，因为\Person{基督}在这一天从死者中复活；也尊崇上述的日子，因为他在这一天被钉十字架。他特别崇敬十字架，既因为它在与敌人的战斗中赋予他力量，也因为这一象征曾以神圣的方式向他显现。他以法律废除了罗马法庭惯用的十字架刑罚。他命令，每当铸造钱币和雕像时，都应始终刻印这一神圣标记，至今存在的以这种形式制成的他的雕像仍证明着这一命令。事实上，他在一切事情上，尤其是在制定法律时，都努力侍奉天主。而且，他似乎还禁止了许多此前未被禁止的可耻和放荡的结合；关心此事的人从这些少数实例中便可一目了然他就这些点所制定的法律是什么；在我看来，现在详尽讨论它们是不合理的。然而，我认为有必要提及为宗教的荣誉和巩固而制定的法律，因为它们构成了教会历史的重要组成部分。因此，我将继续叙述。\par

% structure: p0040 source=h2 target=subsection reason=relative-heading-rank; base=h2

\subsection{第九章 \Person{君士坦丁}制定法律，优待独身者和神职人员。}

有一条古罗马法律，规定二十五岁以上未婚者不享有与已婚者相同的特权；该法律的其他条款规定，非至亲者不能从遗嘱中获得任何财产；此外，无子女者将被剥夺继承财产的一半。这条古罗马法律的目的是增加罗马和属民的人口，这些人口在此法律颁布前不久的内战中大为减少。皇帝察觉到这一规定损害了那些为天主而保持独身和无子女者的利益，并认为试图依靠人的关心和热忱来增加人类繁衍是荒谬的（因为自然总是按照上主的命令而增减），因此制定法律，规定独身和无子女者享有与已婚者相同的权益。他甚至赐予那些选择持守节操和童贞生活的人特殊特权，并允许他们（违反罗马帝国普遍习俗）在达到青春期前立遗嘱；因为他相信那些献身于侍奉天主和追求哲学的人，在任何情况下都会正确判断。出于类似原因，古代罗马人允许维斯塔贞女一满六岁就可立遗嘱。这是对宗教崇高尊崇的最大证明。\Person{君士坦丁}免除各地神职人员的赋税，并允许诉讼当事人如果更愿意选择主教而非国家统治者，可向主教提出上诉，请求其裁决。他规定他们的判决有效，并且如同皇帝亲自宣判一样，优于其他法官的判决；总督和下属军官应负责执行这些判决；主教会议的决议不得更改。\par

在我的历史叙述进行到此时，若不提及那些为教会中获得自由的人而通过的法律，将是不恰当的。由于法律的严格和主人的不愿，获得这种更好的自由——即罗马公民权——有许多困难。因此，\Person{君士坦丁}制定了三条法律，规定凡在教会内，其自由身份经司铎证实的人，都应获得罗马公民权。\par

这些敬虔规条的记录仍然存在，因为当时习惯将一切有关释放奴隶的法律刻在石板上。这便是君士坦丁的立法；他事事致力促进宗教的尊荣；而宗教之所以受到重视，不仅因其本身，也因当时参与其中者的德行。\par

% structure: p0044 source=h2 target=subsection reason=relative-heading-rank; base=h2

\subsection{第十章　论幸存的伟大宣信者}

由于迫害刚刚停止，许多杰出的基督徒以及许多幸存的宣信者装饰了各教会：其中有科尔多瓦主教奥西乌斯、奇里乞亚的厄皮法尼亚主教安菲翁、继玛加略之后治理耶路撒冷教会的马克西穆斯，以及一位埃及人帕弗努提乌斯。据说天主借后者行了许多奇迹，制伏魔鬼，并赐给他恩宠医治各种疾病。这位帕弗努提乌斯和方才提到的马克西穆斯，曾属于马克西米努斯判处到矿场劳役的宣信者之列，他剜去了他们的右眼，并弄瘸了他们的左腿。\par

% structure: p0046 source=h2 target=subsection reason=relative-heading-rank; base=h2

\subsection{第十一章　圣斯彼里迪翁小传：他的谦逊与坚定}

塞浦路斯特里米通主教斯彼里迪翁于此时兴盛。要彰显他的德行，我认为至今仍流传的声誉已足够。他借天主助佑所行的奇事，似乎为同地区居民所共知。我不会隐瞒我所知的那些事迹。\newline{}\par

他是一个农夫，已婚，并有子女；然而他并不因此而缺乏灵性上的成就。据说某夜，有恶人进入他的羊圈，正在偷他的羊时，忽然被捆绑，却无人捆绑他们。次日，当他来到羊圈，发现他们被缚，便解开那无形的束缚；但他责备他们宁愿偷窃本可合法赚取或取得之物，又如此在夜间大费周章；然而他怜悯他们，渴望教导他们，以引导他们度更善的生活，便对他们说：\Quote{去吧，把这公羊也带去；因为你们守夜已疲乏，不应叫你们的劳作如此受责，以致空手离开我的羊圈。}这一举动值得万分钦佩，但下面我要述说的同样令人钦佩。有一个人将其存款托付给他的女儿，她是一位童贞女，名叫依勒内。为更安全起见，她将存款埋藏起来；但不久后她去世了，未向任何人提及此事。存款人来讨要存款。斯彼里迪翁不知如何回答，便搜遍全屋寻找；但找不到，那人便哭泣，撕扯头发，几乎要断气。斯彼里迪翁动了怜悯之心，来到坟墓前，呼唤那女孩的名字；她回答后，他便询问存款的事。得到所需信息后，他返回，在所指明的地方找到了宝藏，并交给了物主。既然我已开始讲述此事，再补充这件事或许也不为过。\par

斯彼里迪翁有一个习惯：他将自己的一部分出产分给穷人，另一部分则借给那些想要的人作为馈赠；但无论是给予还是收回，他从不亲手分发或接收：他只指点仓库，告诉来求他的人随意取用所需，或归还所借之物。有一个人曾这样借了东西，后假装前来归还；当他照常被指示将所借之物放回仓库时，他看到了不义的机会；以为此事会隐瞒过去，他便不清偿债务，而是欺诈地装作已履行义务，假装已归还便离开了。然而此事不能长久隐藏。过了一段时间，那人又来借取，被派往仓库，允许他自行量取所需。他发现仓库空无一物，便去告知斯彼里迪翁；斯彼里迪翁对他说：\Quote{人哪，我奇怪怎么唯独你发现仓库空空如也，没有你所需的物品：你想想是否没有归还第一次的借款，以致你再次需要？否则，你所求的不会短缺。去吧，信靠，你必将找到。}那人感到被责备，承认了自己的过错。这位天主之人在教会事务管理上的坚定与精确值得钦佩。据说此后有一次，塞浦路斯的主教们集会商议某件紧急事宜。斯彼里迪翁出席，莱德里主教特里菲利乌斯也在场，他是一位善于辞令的人，因曾研习法律，在贝里图斯独居过。\par

当集会召开时，特里菲利乌斯应邀向众人讲道，在讲道中途引用了经文：\Quote{拿起你的床走吧}\BibleRef{玛 9:6}，但将\Quote{床}（\OriginalTerm{床}{\GreekText{κράββατος}}）一词改成了\Quote{榻}（\OriginalTerm{榻}{\GreekText{σκίμπους}}）。斯彼里迪翁非常恼怒，喊道：\Quote{难道你比说出\Quote{床}字的那位更大，以致以用祂的话为羞耻？}他说完这话，便从司铎的宝座转过身来，面向众人；以此行动教导他们约束那以口才自豪的人；他配作这样的责备，因为他因自身的行为而受尊敬且极其显赫；同时他比那位长老年长，且在司铎职位上更高。\par

斯彼里东接待陌生人的情形可从以下事件看出。在四旬期内，有一位旅客在那些他习惯与家人连续禁食的日子里前来拜访他，而在规定尝食物的那一天，他直到中午都不进食。斯彼里东看出那旅客十分疲惫，便对他的女儿说：\Quote{来，给他洗脚并摆上肉。}那贞女回答说家里既无面包也无大麦食物，因为在禁食期间准备这些东西是多余的；斯彼里东先祈祷并求宽恕，然后吩咐她煮一些碰巧在家里的咸猪肉。当肉煮好后，他坐下到桌旁与那客人一同坐席，吃了肉，并叫他照样做。但那客人以自己是基督徒为由推辞，斯彼里东对他说：\Quote{正因如此，你更不可推辞吃肉；因为天主的话显示：\BibleQuote{对洁净的人，一切都是洁净的。}}\BibleRef{铎 1:15}这就是我所能讲述的关于斯彼里东的细节。\par

% structure: p0052 source=h2 target=subsection reason=relative-heading-rank; base=h2

\subsection{第十二章 论隐修士的组织：其起源与创始人。}

在这时期，那些接受了隐修生活的人，在彰显教会的光荣与以他们的善行证明其教义真理方面，并非最不重要。的确，人从天主那里所领受的最有益的事，就是他们的\OriginalTerm{哲学}{philosophy}。他们忽略数学的许多分支和辩证法的技巧，因为认为这些研究是多余的，是对时间的无用耗费，因为它们对正确的生活毫无贡献。他们专务于培养自然和有用的学问，为能减轻，若非完全根除，邪恶。他们始终坚持，不把任何处于德行与恶习之间的行为或原则视为善，因为他们只以善为乐。他们认为，凡是虽不做恶却不行善的人都是恶人。因为他们不是通过辩论，而是通过实践来证明德行，视世人的荣誉为无物。他们勇敢地克胜灵魂的激情，既不屈服于本性的需要，也不屈从于身体的软弱。既已拥有天主之心的力量，他们总是仰望万有的创造者，日夜朝拜祂，以祈祷和恳求平息祂。藉着灵魂的纯洁和善行的一生，他们洁净地进入宗教礼仪，蔑视净化、净水器皿及此类仪式；因为他们认为唯有罪才是污点。他们超越了我们所易受的外在际遇，仿佛掌控着一切；因此，他们不为支配人生的灾难或所迫而偏离所选的道路。受辱时不觉痛苦，遭恶意时不自卫，面对疾病或缺乏必需品时不气馁，反而以忍耐和温良在此类考验中喜乐。他们在整个生命中习惯于知足，并尽可能接近天主，达到人性所能及的程度。他们认为现世生命只是一段旅程，因此不关心积累财富，仅为紧迫的需要而预备现世所需。他们赞叹自然的美丽与纯朴，但他们的希望寄托于天堂和未来的福乐。他们完全沉浸于对天主的崇拜中，厌恶污秽的言语；既然已驱逐恶行，他们甚至不允许提及此类事情。他们尽可能地限制本性的需求，迫使身体满足于适度的供给。他们以节制克服放纵，以正义克服不义，以真理克服谎言，并在一切事上达到适中。他们与邻人和睦共处，彼此共融。他们为朋友和陌生人提供所需，按需要施舍给贫困者，安慰受苦者。由于他们在一切事上勤奋，热切寻求至高之善，他们的训导虽谦逊审慎、缺乏虚浮华丽的辞藻，却如同良药，具有医治听众道德疾病的效力；他们说话也怀着敬畏与尊敬，避免一切争辩、嘲讽和愤怒。的确，抑制一切非理性的情感，克胜肉体和本性的激情，是合理的。如有些人所说，\Person{厄里亚}先知和\Person{若翰洗者}是这崇高哲学的创始人。毕达哥拉斯派的\Person{斐洛}叙述，在他那个时代，最虔诚的希伯来人从世界各地聚集，定居在玛雷奥蒂斯湖附近一座山丘上的地区，为能以哲学家的方式生活。他描述了他们的住所、生活方式和习俗，与我们现在在埃及隐修士中所见到的相似。他说，自从他们开始从事哲学研究起，他们就将财产交给亲属，放弃事务和社交，住在城墙之外、田野和花园中。他还告诉我们，他们有称为\OriginalTerm{隐修院}{monasteries}的圣殿，在其中他们独自居住，专心举行神圣奥秘，并以圣咏和赞歌勤勉地朝拜天主。他们日落前从不进食，有些人每隔三天甚至更长时间才进食一次。最后，他说，在某些日子，他们躺在地上，戒除酒和肉食；他们的食物是面包、盐和牛膝草，饮料是水；其中还有保持童贞到老年的妇女，她们为了对哲学的热爱并出于自愿，实践独身。在这段叙述中，\Person{斐洛}似乎描述了某些已接受了基督信仰却仍保留本民族习俗的犹太人；因为这种生活方式在其他任何地方都找不到踪迹；因此我得出结论，这种哲学从那时起就在埃及兴盛。然而，另一些人断言，这种生活方式源于不时发生的、为了宗教而受的迫害，许多人被迫逃往山岭、旷野和森林，并习惯了这种生活。\par

% structure: p0054 source=h2 target=subsection reason=relative-heading-rank; base=h2

\subsection{第十三章 论\Person{圣安当}与\Person{圣保禄（简朴者）}}

无论这哲学的开创者是埃及人还是其他人，普遍公认的是伟大的隐修士\Person{安当}通过道德和适宜的操练，将这种生活方式发展到了精确和完美的顶峰。他的名声传遍埃及的旷野，以至于皇帝\Person{君士坦丁}因其美德的声誉寻求他的友谊，以书信尊崇他，并敦促他写信告知可能需要什么。他种族上是埃及人，出身于\Place{科马}的一个显赫家庭，该地位于埃及边境附近的\Place{赫拉克利亚}附近。他年幼时失去了父母，将父亲的遗产分给了同村的人，卖掉了其余财产，将所得分给穷人；因为他知道哲学不仅仅在于放弃财产，而在于妥善分配财产。他结识了当时的热心人士，并效仿所有人的美德。他相信行善的习惯会变得令人愉快，尽管起初很艰难，他反思了更严格的苦修方法，并日复一日地通过自制来增强它，仿佛他总是在重新开始他的事业。他通过劳动抑制身体的享乐，并借助天主上智控制灵魂的激情。他的食物是面包和盐，饮料是水，他从未在日落前开斋。他经常两天或更长时间不吃东西。他几乎整夜不睡，持续祈祷直到天亮。如果他有时睡觉，也只是在一张短席上小睡片刻；但通常光秃秃的地面就是他的床铺。他拒绝用油膏抹身体，不使用浴池和类似可能通过湿气使身体松弛的奢侈品；据说他从未在任何时候看到自己赤裸。他既没有拥有也不欣赏学问，但他重视良好的理解力，认为它先于文字，并且是文字的发现者。他非常温和、仁慈、谨慎、刚毅；谈话愉快，辩论友好，即使别人将争论的话题作为争吵的借口。他通过自己的习惯和一种智慧平息了日益增长的争论，使他们恢复温和；他也缓和了与他交谈者的热情，并规范了他们的举止。虽然因他非凡的德行，他充满了天主的预知，但他不把预知未来视为美德，也不劝告别人轻易寻求这一恩赐，因为他认为没有人会因对未来无知或知晓而受到惩罚或奖赏；因为真正的幸福在于事奉天主和遵守祂的诫命。\Quote{但是，}他说，\Quote{如果有人想知道未来，让他不断净化自己的灵魂，因为他将有能力行在光明中，理解将要发生的事，因为天主将向他启示未来。}他从不让自己闲散，而是劝勉所有似乎愿意过好生活的人勤于劳作，自我反省，并在创造昼夜的天主面前承认自己的罪；当他们犯错误时，他敦促他们将过犯记录下来，以便他们为自己的罪感到羞愧，并害怕有人会发现许多记录的事情；因为如果文件被追溯到他们，他们会害怕自己在别人面前暴露为堕落的人。他比任何人都更勇敢、最热心地为受伤害者辩护，并为此常去城市；因为许多人来见他，迫使他为他们向统治者和掌权者求情。所有的人都以见到他为荣，热切地听他的讲论，并同意他的论点；但他更喜欢在旷野中不为人知、隐藏自己。当被迫访问城市时，他一旦完成所承担的工作，就立刻返回旷野；因为，他说，正如鱼在水中生长一样，旷野是为隐修士预备的世界；正如鱼被抛在干地上就死亡一样，隐修士进入城市就会失去庄重。他对所有见到他的人都顺从和和蔼，并注意不拥有或似乎拥有傲慢的本性。我对\Person{安当}的行为作了这个简短的描述，以便通过类比，从对他旷野行为的描述中形成对他哲学的概念。\par

他有许多著名的门徒，其中一些在埃及和利比亚兴盛，另一些在巴勒斯坦、叙利亚和阿拉伯；每个门徒都不逊于他们的老师，与同住的人一起过生活，规范自己的行为，教导许多人，并使他们与同类的美德和哲学结合。但任何人很难通过仔细走访城市和村庄来找到\Person{安当}的同伴或他们的继承者，因为他们寻求隐藏比许多野心家通过炫耀和表演寻求声望和名声更为热切。\par

我们必须按年代顺序叙述\Person{安当}最著名门徒的历史，特别是那位被称为\Quote{质朴的\Person{保禄}}的\Person{保禄}。据说他住在乡间，娶了一位美丽的女子，后来撞见她通奸，他平静地笑了笑，并发誓不再与她同住；他把她留给了奸夫，立刻去旷野投奔\Person{安当}。还据说他极其温和忍耐；由于他年事已高，不习惯隐修院的严厉，\Person{安当}通过种种考验测试他的力量，因为他是新来的，但发现他并无卑劣之处；并且，在证明了自己有完美的哲学之后，他被派去独居，因为他不再需要老师了。天主亲自证实了\Person{安当}的见证，并显示此人在行为上最为卓越，甚至比他的老师更善于折磨和驱逐魔鬼。\par

% structure: p0058 source=h2 target=subsection reason=relative-heading-rank; base=h2

\subsection{第十四章。关于圣\Person{安蒙}和\Place{奥林匹斯}的\Person{欧提基乌斯}的记述。}

大约在此时，\Place{埃及}人\Person{阿孟}开始了隐修生活。据说他的家人强迫他结婚，但他的妻子从未在肉身上认识他；因为在婚礼当天，当他们独处时，作为新郎的他引领她作为新娘到他的床前，他对她说：\Quote{妇人啊！我们的婚姻确实已经缔结，但并未圆满}；然后他从\WorkTitle{圣经}中向她指出，保持童贞是她的首要益处，并恳求他们分开生活。她被他关于童贞的论点说服，但想到要与他分离，感到十分痛苦；因此，虽然分床而睡，他与她共同生活了十八年，在此期间他并未忽略隐修功课。在这段时期结束时，那位妻子的好胜心因丈夫的德行而被强烈激发，她深信让这样一个人因她而生活在世俗领域是不公平的；她认为为了隐修生活，他们有必要彼此分开生活；她向丈夫恳求了此事。因此，他离开了，感谢天主给了他妻子的建议，并对她说：\Quote{你保留这房子，我将为自己另造一间。}他退隐到\Place{马雷奥提斯湖}以南、\Place{西提斯}与称为\Place{尼特里亚}的山之间的一处旷野；在此地，他致力于隐修生活二十二年，每年看望妻子两次。这位神圣的人是那里隐修院的创始人，并聚集了许多著名的门徒，正如继承记录所显示的那样。许多非凡的事件发生在他身上，这些事件已被\Place{埃及}的隐修士们准确记载，他们非常仔细地保存了更古老苦修者的德行，通过口传传统传承。我将叙述那些为我们所知的事情。\par

\Person{阿孟}和他的门徒\Person{狄奥多若}有一次需要出行，路上发现必须渡过一条称为\Place{吕科斯}的运河。\Person{阿孟}命令\Person{狄奥多若}背对河水过河，以免他们看到彼此的裸体，而由于他自己也羞于看到自己赤身，他突然被一股神圣的力量抓住并带过河，降落在对岸。当\Person{狄奥多若}过河后，他发觉长者的衣服和脚都没有湿，便询问原因；没有得到回答，他强烈质问这件事；最后\Person{阿孟}在约定有生之年不提此事后，承认了事实。\par

以下是同一性质的另一个奇迹。一些邪恶的父亲带着一个被疯狗咬伤、濒临死亡的儿子来到他面前，哀哭恳求他医治。他对他们说：\Quote{你们的儿子不需要我的医治，但如果你们愿意将你们偷来的牛归还给主人，他会立刻痊愈。}结果正如所预言的那样：牛被归还，孩子的病也消除了。据说，当\Person{阿孟}去世时，\Person{安当}看见他的灵魂升入天堂，因为天使伴随他唱着\BibleBook{}{圣咏}；当同伴问他为何明显惊愕时，他没有向他们隐瞒；因为人们看到他专注地凝视天空，因所见奇妙景象而感到惊异。不久之后，有人从\Place{西提斯}来，宣布\Person{阿孟}去世的时刻，\Person{安当}预言的真相就显明了。因此，正如所有善人所见证的，这些圣者每个人都以特殊的方式蒙福：一人因脱离此生而蒙福，另一人因被堪当目睹天主向他显示如此奇迹的景象而蒙福；因为\Person{安当}和\Person{阿孟}彼此相距多日的路程，且上述事件得到与二人皆熟识之人的证实。\par

我相信同样是在这朝统治期间，\Person{欧提基安}开始了隐修生活。他将居所定在\Place{比提尼亚}，靠近\Place{奥林普斯}山。他属于\OriginalTerm{诺瓦替安派}{Novatians}，并分享了天主的恩宠；他医治疾病、施行奇迹，他德行的名声使得\Person{君士坦丁}与他保持亲密和友谊。恰在这时期，一位皇家卫兵因涉嫌谋反君主而逃跑，经搜索后在\Place{奥林普斯}附近被捕。\Person{欧提基安}被那人的亲戚恳求，为他向皇帝求情，同时指示松开囚犯的锁链，以免他因锁链过重而死亡。据传，\Person{欧提基安}因此派人去见看守囚犯的官员，要求他们松开锁链；他们拒绝后，他亲自前往监狱，尽管门锁着，却自动打开，囚犯的锁链也脱落了。\Person{欧提基安}随后前往当时居住在\Place{拜占庭}的皇帝那里，轻易获得了赦免，因为\Person{君士坦丁}通常不拒绝他的请求，非常尊敬此人。\par

我现在简要叙述了隐修哲学生活中最杰出的代表人物。如果有人想获得关于这些人的更详细信息，可以在为其中许多人撰写的传记中找到。\par

% structure: p0064 source=h2 target=subsection reason=relative-heading-rank; base=h2

\subsection{第十五章　亚略异端，其起源、发展，以及由此在主教们中引起的争论。}

尽管，如我们所展示的，其时宗教正处于繁盛状态，然而教会却被严重的纷争所扰乱；因为借着虔诚和寻求对天主的更深认识为借口，某些此前未被审视的问题被激发出来。\Person{阿里乌}是这些争论的始作俑者。他是埃及\Place{亚历山大里亚}教会的一位司铎，起初对教义有热切的思想，并支持\Person{梅利提乌斯}的革新。然而，他最终放弃了后一意见，并由\Place{亚历山大里亚}主教\Person{伯多禄}祝圣为执事，后来\Person{伯多禄}将他逐出教会，因为当\Person{伯多禄}绝罚\Person{梅利提乌斯}的狂热者并拒绝他们的圣洗时，\Person{阿里乌}为这些行为攻击他，且不能安静地受约束。在\Person{伯多禄}殉道之后，\Person{阿里乌}向\Person{阿基拉}请求宽恕，并恢复执事之职，后来晋升为司铎。此后，\Person{亚历山大}也对他高度尊重，因为他是一位极其熟练的逻辑学家；据说他并不缺乏此类知识。他陷入荒谬的言论，以至于他敢于在教会中宣讲此前无人提及的事，即：天主子是从无先有之物受造的，存在过一个他不存在的时段；因拥有自由意志，他有能力行恶与行善，并且他是受造的和被造的：除此之外，随着他深入论证和探究细节，还添加了许多其他类似的断言。那些听到这些学说被提出的人，责备\Person{亚历山大}没有反对违背教义的革新。但这主教认为更明智的做法是让各方自由讨论有疑点的问题，以便通过劝说而非强制来停止纷争；因此他与一些神职人员坐下作为裁判，引导双方进行讨论。但这次讨论中，如通常言辞之争中所发生的，各方都声称胜利。\Person{阿里乌}为他的断言辩护，但其他人争辩说圣子与圣父\OriginalTerm{同体}{consubstantial}且\OriginalTerm{同永}{co-eternal}。会议再次召开，同样的问题被争论，但他们之间未能达成一致。在辩论中，\Person{亚历山大}似乎先倾向一方，然后又倾向另一方；最终，他表明自己支持那些肯定圣子与圣父同体同永的人，并命令\Person{阿里乌}接受这一教义，摒弃他先前的观点。然而，\Person{阿里乌}不肯被说服而服从，许多主教和神职人员认为他的教义陈述是正确的。因此，\Person{亚历山大}将他以及与他见解一致的神职人员逐出教会。\Place{亚历山大里亚}教区中接受他意见的有司铎\Person{艾塔拉斯}、\Person{阿基拉}、\Person{卡尔波内斯}、\Person{萨尔马特斯}和\Person{阿里乌}，以及执事\Person{尤佐伊乌斯}、\Person{玛卡里乌斯}、\Person{犹利乌}、\Person{梅纳斯}和\Person{赫拉迪乌斯}。同样，许多民众也站在他们一边：有些人因为他们认为这些教义来自天主；另一些人，如类似情况中常发生的，因为他们相信这些人受到虐待和不公正的绝罚。\Place{亚历山大里亚}的局势如此，\Person{阿里乌}的党派认为明智之举是寻求其他城市主教的青睐，于是派遣使团到他们那里；他们将自己的教义书面陈述送给他们，请求他们如果认为这些思想来自天主，就通知\Person{亚历山大}不应打扰他们；但如果他们不认可这些教义，就应教导他们必须持有什么意见。这一预防措施对他们不无益处；因为他们的主张因此普遍传播，他们提出的问题成为所有主教辩论的主题。有些人写信给\Person{亚历山大}，恳求他不要接纳\Person{阿里乌}的党派共融，除非他们放弃自己的意见，而另一些人则写信催促采取相反的行动。当\Person{亚历山大}观察到许多以良好品行外表为尊、以雄辩说服力为重的教士——特别是\Place{尼科美底亚}教会领袖\Person{欧瑟伯}，一位颇有学识并在宫廷中享有盛誉的人——支持\Person{阿里乌}派时，他写信给各教会的主教，要求他们不要与这些人共融。这一措施更加激发了各方的热忱，并且可以预料，争论越来越激烈。\Person{欧瑟伯}和他的党派曾多次请求\Person{亚历山大}，但未能说服他；因此，他们感到受辱而愤慨，更加坚决地支持\Person{阿里乌}的教义。在\Place{彼提尼亚}召开了一次主教会议后，他们写信给所有主教，要求他们与\Person{阿里乌}派共融，如同与持正确宣认的人一样，并同样要求\Person{亚历山大}与他们共融。由于无法从\Person{亚历山大}那里得到服从，\Person{阿里乌}派派遣使者到\Place{提洛}主教\Person{保利纳斯}、\Place{凯撒勒雅}（在\Place{巴勒斯坦}）教会领袖\Person{欧瑟伯·旁非禄}和\Place{斯基托波利斯}主教\Person{帕特罗斐禄}那里，请求允许他自己和他的追随者——因为他们先前已达到司铎级别——将跟随他们的人组成一个教会。因为在\Place{亚历山大里亚}，当时和如今仍然有这样的习俗：所有教会应归一位主教管辖，但每位司铎应有自己的教堂，在其中召集民众。这三位主教，连同在\Place{巴勒斯坦}聚集的其他主教，同意了\Person{阿里乌}的请求，并允许他像以前一样召集民众；但要求服从\Person{亚历山大}，并命令\Person{阿里乌}不断努力，以恢复与他的和平与共融。\par

% structure: p0066 source=h2 target=subsection reason=relative-heading-rank; base=h2

\subsection{第十六章．\Person{君士坦丁}听闻主教们的纷争及关于逾越节的意见分歧后，大为忧虑，并派遣西班牙人、\Place{科尔多瓦}主教\Person{何西乌}前往\Place{亚历山大里亚}，以消除主教们的分歧，并解决关于逾越节的争端。}

在埃及举行了许多会议之后，争论仍持续加剧，分歧的报告传到了皇宫，\Person{君士坦丁}因此大为烦恼；因为正是在这个时期，当宗教开始更广泛地传播时，许多人因教义分歧而被劝阻，不愿信奉基督信仰。皇帝公开指责\Person{阿略}和\Person{亚历山大}是这场骚动的始作俑者，并写信责备他们，因为他们本可以隐瞒争议，却将其公开化；他们激烈地争论了一个本不应提出，或者至少应安静提出自己观点的问题。他告诉他们，不应当因对某些教义观点持有不同意见而与他人分离。\par

关于上智的安排，人们必须持守同一信仰；但对这一问题所做的细微探究，尤其是当它们不能使人达成同一意见时，理应按照一切道理隐藏起来。他劝诫他们摒弃关于这些问题的所有闲谈，并要同心合意；因为他感到十分痛心，为此他放弃了他访问东方城市的计划。他就是以这种语气写信给\Person{亚历山大}和\Person{阿略}，同时责备和劝诫他们二人。\par

君士坦丁也因有关逾越节庆祝方式的分歧而深感忧虑；因为东方的一些城市在这点上有所不同，尽管它们没有彼此拒绝共融；他们更按照犹太人的习俗过节，而由于这种分歧，自然就减损了节日祭献的光彩。皇帝热切地努力从教会中消除这两个导致分裂的原因；他认为在祸患扩大前能够除去它，便派遣了一位因其信德、德行备受尊崇，且在其先前时代因对这教义的宣认而极受认可的人，去调解埃及因教义而分裂的人，以及东方在逾越节上意见不同的人。此人就是\Place{科尔多瓦}主教\Person{何西乌}。\par

% structure: p0070 source=h2 target=subsection reason=relative-heading-rank; base=h2

\subsection{第十七章 因\Person{阿略}而召集的\Place{尼西亚}会议}

当发现事情并未如皇帝所愿，反而争论太大无法和解，以至于被派去缔造和平的人未完成任务而归时，\Person{君士坦丁}在\Place{彼提尼亚}的\Place{尼西亚}召集了一次会议，并写信给各地教会中最杰出的人士，指示他们于指定日期到场。在那些占据宗座的人中，以下几位出席了这次会议：\Person{耶路撒冷的玛加略}、当时已主持\Place{奥龙特斯河畔安提约基雅}教会的\Person{欧斯塔修}，以及\Place{马里奥特湖}附近\Place{亚历山大城}的\Person{亚历山大}。\Person{罗马主教朱利乌斯}因年事极高无法出席，但由他教会的司铎\Person{维托}和\Person{味增爵}代表。来自不同民族的许多其他优秀人士聚集在一起，其中有些人以学识、口才和精通圣经及其他学科而闻名；有些人以德行高尚著称；还有些人兼有这些优点。大约有三百二十位主教到场，还有众多司铎和执事陪同。同样，还有擅长辩证法的人在场，准备协助讨论。而且，如同这类场合常见的情况，许多司铎也因处理自己的私事而赴会；因为他们认为这是一个纠正冤屈的好机会，每个人在哪些方面对他人有不满，就向皇帝呈递一份文件，在其中指出对自己所犯的过犯。这种做法日复一日，皇帝便指定了一天，所有控诉都要在这一天呈到他面前。到了预定时间，他拿起呈递给他的请愿书，说道：\Quote{所有这些指控到了适当的时候，都会在伟大的审判日被提出，并由众人的大审判者审判；至于我，我不过是一个人，若我受理此类事务，实属不当，因为控告者和被告都是司铎；而司铎的行为应当永远不会受他人审判。因此，你们要效法天主的圣爱与慈悲，彼此和好；撤回彼此间的控告；让我们听从劝告，专心致力于我们集会所关注的信仰问题。}这番讲话之后，为使每个人的文件无效，皇帝命令将请愿书烧毁，然后指定了一天来解决疑问。但在预定时间到来之前，主教们聚集在一起，召\Person{阿略}到场，开始审查有争议的话题，每人提出自己的意见。然而，可以预料的是，调查引出了许多不同的问题：有些主教反对引入与从起初传授给他们的信仰相悖的新奇事物。特别是那些坚持教义简朴的人主张，天主的信仰应当不加好奇探究地接受；而另一些人则主张，不应不加审查地遵循古旧的看法。当时聚集的许多主教以及陪同他们的神职人员，因擅长辩证技巧并娴熟于这类修辞方法而引人注目，获得了皇帝和宫廷的注意。其中，当时是\Place{亚历山大城}执事、陪同其主教\Person{亚历山大}的\Person{亚他那修}，似乎在这些问题的商议中起到了最大的作用。\par

% structure: p0072 source=h2 target=subsection reason=relative-heading-rank; base=h2

\subsection{第十八章 两位哲学家因与他们辩论的两位老人的纯朴而归信}

当这些辩论进行时，某些外邦哲学家渴望参与其中；有些人是想了解所教导的教义，另一些人则因最近外邦宗教被压制而对基督徒感到愤怒，他们想把关于教义的探究变成言辞之争，以便在他们中间引起分歧，并使他们显得持有矛盾的意见。据说，其中一位哲学家，以其公认的出众口才而自傲，开始嘲弄司铎，从而激怒了一位纯朴的老人，他被称为精修者而备受尊敬，虽然不谙逻辑的精细和冗词，却决意与他争辩。那些认识这位精修者的人中，较不严肃的嘲笑他承担如此重任；但较有思想的人则担心，在同如此雄辩之人争辩时，他只会使自己成为笑柄；然而他的影响力如此之大，声望如此之高，以致他们无法阻止他参加辩论；于是他以下面的话发言：\Quote{奉耶稣基督之名，哲学家啊，请听我说。有一位天主，是天地及一切有形无形之物的创造者。祂藉圣言的能力创造万物，并藉祂的圣神的圣德确立万物。圣言，我们称之为天主子，见人类陷于谬误，如禽兽般生活，便怜悯他们，屈尊由女人所生，与人交往，并为他们而死。祂还要再来，按此生的行为审判我们每个人。我们以极纯朴的心相信这些事是真实的。因此，不要徒劳费力去反驳那些只能藉信德理解的事实，或去探究这些事是否真正发生的方式。回答我，你信吗？} 这位哲学家对所发生的事感到惊讶，回答说：\Quote{我信}；并感谢这位老人以辩论胜过了他，便开始向其他人教导同样的教义。他劝诫那些仍持他先前主张的人接受他所采纳的观点，并起誓保证，他因一种无法解释的冲动而被迫接受了基督信仰。\par

据说，管理\Place{君士坦丁堡}教会的\Person{亚历山大}也行过类似的奇迹。当\Person{君士坦丁}返回\Place{拜占庭}时，某些哲学家来到他面前，抱怨宗教上的革新，特别是他在国家中引入了一种新的敬拜形式，违背了他的祖先以及所有从前掌权者——无论希腊人或罗马人——所遵循的方式。他们也希望与主教\Person{亚历山大}就教义进行辩论。\Person{亚历山大}虽然不擅长此类辩驳，或许因他的生活而被说服，深知他是一位卓越而良善的人，便遵皇帝之命接受了这场较量。当哲学家们集合时，由于每个人都想参与讨论，他请求他们推选一位他们认为配得的人作代言人，其他人保持沉默。当其中一位哲学家开始展开辩论时，\Person{亚历山大}对他说：\Quote{我奉耶稣基督之名，命令你住口。}那人立刻哑口无言。那么，应该思考一下，是一个被一句话轻易止住的人——而且是一位哲学家——更伟大，还是被话的力量劈开石墙更伟大？我听说有人将后一奇迹归于称为加色丁人的\Person{犹利安}。据我所知，这些事正如上文所述发生。\par

% structure: p0075 source=h2 target=subsection reason=relative-heading-rank; base=h2

\subsection{第十九章  当大公会议召集时，皇帝发表公开演说。}

主教们进行了长时间的磋商；在召来\Person{亚略}到他们面前之后，他们对他的命题进行了准确的检验；他们谨慎提防，不投票支持任何一方。当最终指定了解决疑难之处的日期到来时，他们聚集在皇宫内，因为皇帝已表示他打算参与审议。当他与司铎们同在一处时，他穿过会议室的上座，坐在为他预备的宝座上，然后命令大公会议就坐；因为座位已沿着宫殿房间的墙壁两边排列，那房间是最大且优于其他厅堂的。\par

他们就坐后，\Person{欧瑟伯·庞非录}起身发表了一篇颂扬皇帝的演说，并因皇帝而感谢天主。他讲完后，恢复了肃静，皇帝以如下言辞发言：\Quote{我为一切事感谢天主，但尤其，朋友们，因得以看见你们的会议而感谢。这事件超乎我的祈祷，因如此多的基督司铎被领到同一地点；现在，我渴望你们同心合意，拥有一致的判断，因为我认为天主的教会中的分裂比任何其他邪恶都更危险。因此，当有人报告说你们有分歧时，我听到这有害之事，心中深感痛苦；这对你们最无益处，因为你们是神圣敬礼的引导者和和平的仲裁者。正是为此，我召集你们举行神圣的大公会议，我既是你们的皇帝，也是你们的同僚医师，我向你们寻求一件恩惠，这恩惠为我们共同的主所悦纳，我接受它如同你们赐予它一样光荣。我所寻求的恩惠是，你们调查纷争的原因，并寻求一致与和平的解决，使我与你们一同战胜那嫉妒的魔鬼，它见我们脚下的外敌与暴君被征服，又妒忌我们的美好光景，便激起了这场内乱。}皇帝用拉丁语发表了这篇演说，由他身旁的人提供翻译。\par

% structure: p0078 source=h2 target=subsection reason=relative-heading-rank; base=h2

\subsection{第二十章  在听取双方意见后，皇帝谴责了亚略的追随者并将他们放逐。}

随后的司铎辩论转向了教义。皇帝耐心地听取了双方的发言；他赞扬了言辞得当的人，斥责了那些表现出争论倾向的人，并根据他对所听内容的理解（因为他并非完全不谙希腊语）和蔼地对每个人讲话。终于，所有司铎都彼此达成一致，承认圣子与圣父\OriginalTerm{同体}{consubstantial}。在会议开始时，只有十七人称赞\Person{亚略}的观点，但最终这些人中的大多数都同意了普遍的见解。皇帝也赞同这一判断，因为他认为会议的一致是天主的认可；他下令，任何人若对此有违抗，应立即流放，作为企图推翻神圣定义的罪人。我曾以为有必要把那关于这事的文件本身复制出来，作为真理的范例，以便后代能以固定而清晰的形式拥有那在当时促成和好的信仰象征；但因某些理解此类事务的虔诚朋友建议，这些真理应由\OriginalTerm{已入教者}{initiated}及其引导者来谈论和听取，我便同意了他们的意见；因为很可能有\OriginalTerm{未入教者}{uninitiated}会读到这本书。虽然我隐藏了那些我应当保持沉默的禁止材料，但我并未完全让读者对会议所持的意见一无所知。\par

% structure: p0080 source=h2 target=subsection reason=relative-heading-rank; base=h2

\subsection{第21章 会议对\Person{亚略}的决定；其追随者的谴责；他的著作须被烧毁；某些大司祭与会议意见分歧；逾越节的确定。}

应当知道，他们断言圣子与圣父\OriginalTerm{同体}{consubstantial}；那些声称圣子曾有一段时间不存在，在受生之前他不存在，他是由虚无之物所造，他与圣父具有不同的\OriginalTerm{本体}{hypostasis}或\OriginalTerm{实体}{substance}，以及他易于变化或改变的人，应被逐出教会，被视为与天主教会隔绝者。这个决定得到了\Place{尼科米底亚}主教\Person{欧瑟伯}、\Place{尼西亚}主教\Person{泰奥格尼斯}、\Place{加采东}主教\Person{玛里斯}、\Place{西徐波利斯}主教\Person{帕特罗斐鲁}、\Place{利比亚} \Place{托勒迈}主教\Person{塞贡多}的认可。然而，\Person{欧瑟伯·潘斐鲁}暂时未表示同意，但经进一步考察后同意了。会议逐出了\Person{亚略}及其追随者，并禁止他进入\Place{亚历山大里亚}。表达他观点的措辞也被谴责，其中包括他就此问题所写的一部名为\WorkTitle{塔利亚}的著作。我没有读过这本书，但我知道它内容松散，其放荡程度类似于\Person{索塔德斯}的作品。应当知道，虽然\Place{尼科米底亚}主教\Person{欧瑟伯}和\Place{尼西亚}主教\Person{泰奥格尼斯}同意了会议所制定的这份信仰文件，但他们既不同意也未签署对\Person{亚略}的罢免。皇帝以流放惩罚\Person{亚略}，并向各地的主教和人民发送敕令，谴责他及其追随者为不敬神之人，并命令销毁他们的书籍，以免留下对他本人或其宣扬的教义的任何记忆。凡被发现藏匿他的著作且在被指控时不立即烧毁者，应处以死刑，承受极刑。皇帝致信每个城市反对\Person{亚略}及其接受其教义的人，并命令\Person{欧瑟伯}和\Person{泰奥格尼斯}离开他们担任主教的城邑；他特别致信\Place{尼科米底亚}教会，敦促其持守会议所制定的信仰，选举正统主教，服从他们，并使过去的事被遗忘；他威胁要惩罚那些胆敢为被流放的主教说好话或采纳其意见的人。在这些信件及其他信件中，他对\Person{欧瑟伯}表现出愤怒，因为他先前曾附和暴君的意见并参与其阴谋。根据皇帝敕令，\Person{欧瑟伯}和\Person{泰奥格尼斯}被逐出他们所管辖的教会，\Person{安菲翁}接管了\Place{尼科米底亚}教会，\Person{克雷斯特}接管了\Place{尼西亚}教会。这场教义争论结束后，会议决定各地应在同一时间庆祝逾越节。\par

% structure: p0082 source=h2 target=subsection reason=relative-heading-rank; base=h2

\subsection{第22章 \OriginalTerm{诺瓦替安派}{Novatians}主教\Person{阿开修}应皇帝之邀出席第一次大公会议。}

据说，皇帝出于渴望看到基督徒之间重建和谐的强烈愿望，传唤\OriginalTerm{诺瓦替安派}{Novatians}主教\Person{阿开修}到会，向他展示了已经由主教们签署确认的信仰和节期的定义，询问他是否能同意。\Person{阿开修}回答说，他们的阐述没有规定新教义，他赞同会议的意见，并且他从一开始就对信仰和节期持有这些观点。皇帝说：\Quote{那么，如果你与他们意见一致，为何要与他人保持共融距离？}他回答说，分歧最初是在\Person{德西乌斯}统治期间在\Person{诺瓦替安}和\Person{科尔乃略}之间爆发的，并且他认为那些在领洗后犯下圣经中称为\OriginalTerm{至于死的罪}{sins unto death}\BibleRef{若一 5:16}的人不配领受共融；因为他认为这些罪的赦免只取决于天主的权威，而不是司铎的。皇帝回答说：\Quote{\Person{阿开修}，拿一把梯子，独自爬到天堂去吧。}通过这番话，我认为皇帝并非意在赞扬\Person{阿开修}，而是责备他，因为他作为一个人，竟以为自己无罪。\par

% structure: p0084 source=h2 target=subsection reason=relative-heading-rank; base=h2

\subsection{第23章 会议制定的法令；一位名叫\Person{帕夫努提乌}的\OriginalTerm{宣信者}{Confessor}阻止会议制定一项关于即将被授予司铎尊荣者须守\OriginalTerm{独身}{Celibacy}的法令。}

热切于改革那些从事教会事务之人的性生活，该会议颁布了被称为\OriginalTerm{法规}{canons}的法律。当他们在讨论此事时，有些人认为当通过一项法律，规定主教、长老、执事及副执事不得与他们晋铎前所娶的妻子交往；但\OriginalTerm{宣信者}{confessor}帕弗努提乌斯站起来反对这一提案；他说婚姻是尊贵而纯洁的，与自己的妻子同居便是贞洁，并建议会议不要制定这样的法律，因为这将难以承受，并可能成为他们及其妻子失节的场合；他提醒他们说，按照教会的古老传统，那些在领受神品共融时未婚的人，必须保持独身，而已婚者则不可休妻。这就是帕弗努提乌斯的建议，尽管他本人未婚，会议依此采纳了他的意见，没有制定任何法律，而是将此事留给个人判断决定，而非强制。然而，会议制定了其他法律来管理教会；这些法律很容易找到，因为许多人都有保存。\par

% structure: p0086 source=h2 target=subsection reason=relative-heading-rank; base=h2

\subsection{第二十四章 关于梅利提乌斯；神圣议会在其纠纷中所作的卓越指示。}

在审查了梅利提乌斯在埃及的行为之后，会议判决他居住在吕科斯，仅保留主教之名；并禁止他在任何城市或乡村施行圣秩。此前由他所祝圣的人，依此法可保留共融与职务，但在尊严上应次于教会及堂区的圣职人员。当某职位因死亡而空缺时，他们若被认为相称，可根据大众的选举继任，但在此情况下，应由亚历山大里亚教会的主教祝圣，因为他们被禁止在选举中行使任何权力或影响。这一规定在会议看来是公正的，因为梅利提乌斯及其追随者在施行圣秩时表现出极大的鲁莽和轻率；因此它也剥夺了所有与伯多禄不同的圣秩的考虑。伯多禄在治理亚历山大里亚教会时，因当时迫害猖獗而逃亡，后来殉道。\par

% structure: p0088 source=h2 target=subsection reason=relative-heading-rank; base=h2

\subsection{第二十五章 皇帝为会议预备了公宴，邀请其成员至君士坦丁堡，以礼物尊荣他们，劝勉众人一心，并将神圣会议的法令送达亚历山大里亚及各地。}

就在会议通过这些法令之时，庆祝了君士坦丁在位二十周年；因为罗马习俗是在每个统治的第十年举行庆典。皇帝因此认为时机适宜，邀请会议参加庆典，并向他们赠送了合适的礼物；当他们准备返乡时，他召集所有人，劝勉他们在信仰上同心合意，彼此和睦，以免日后纷争在他们中间滋生。在许多类似的劝勉之后，他最后命令他们勤于祈祷，常为他自己、他的子女及帝国祈求天主，在这样对前来尼西亚的人讲话后，他便告别了他们。他致信各城的教会，以便使那些未出席的人明了会议所纠正的事项；尤其对亚历山大里亚教会，他写了更多的内容，敦促他们摒弃一切分歧，在会议颁布的信仰上保持一致；因为这只能是天主的判断，既然它是由圣神藉着如此众多且如此卓越的至高司祭的共识而确立，并经准确查考和验证所有疑难之点后所批准的。\par

\SourceRecord{Translated by Chester D. Hartranft.；Nicene and Post-Nicene Fathers, Second Series；Vol. 2.；Edited by Philip Schaff and Henry Wace.；Buffalo, NY: Christian Literature Publishing Co.,；1890.；<http://www.newadvent.org/fathers/26021.htm>.}

% structure: p0001 source=h1 target=section reason=page-title

\section{\WorkTitle{教会史（第二卷）}}

% structure: p0002 source=h2 target=subsection reason=relative-heading-rank; base=h2

\subsection{第一章 发现赐予生命的十字架与圣钉}

当尼西亚的事务如上所述处理完毕后，司铎们返回家中。皇帝因天主教会意见一致的恢复而极其喜乐，并渴望代表自己、他的子女及帝国，向天主表达主教们的一致同意所激发的感恩，于是下令在耶路撒冷靠近称为加尔瓦略的地方为天主建造一所祈祷之所。同时，他的母亲海伦前往该城，为献上祈祷并探访圣地。她对基督信仰的热诚使她急切地寻找曾构成可敬十字架的木头。但发现这件遗物或主的坟墓并非易事；因为外教人先前曾迫害教会，并在基督信仰初传时采用一切诡计企图消灭它，他们将那地方用大量泥土掩盖，并抬高了原本低洼的地面，正如现今所见，为更有效地隐藏它们，将整个复活之地和加尔瓦略山用墙围起，并且装饰了整个地点，铺上石头。他们还建造了一座阿佛洛狄忒神庙，并设立了一尊小像，以致那些前往那里敬拜基督的人会显得向阿佛洛狄忒屈膝，从而使在那地方敬拜的真正原因随着时间被遗忘；而且，由于基督徒不敢无所畏惧地经常前往那地方或向他人指明，这神庙和雕像将逐渐被视为专属于外教人。然而，那地方最终被发现，关于它的欺诈也被热心地揭露出来；有人说事实首先由一位住在东方的希伯来人披露，他从祖传的一些文件中获得信息；但更符合真理的看法是，天主通过征兆和梦境启示了事实；因为我认为，当天主认为应显明这事时，不需要人的信息。当奉皇帝之命深挖那地方时，发现了我们的主从死者中复活的洞穴；在距离不远处，找到了三个十字架和另一块单独的木头，上面用希伯来文、希腊文和拉丁文的白字刻着以下文字：\Quote{纳匝肋人耶稣，犹太人的君王。}这些文字，正如福音的神圣经卷所述，是由犹太总督比拉多下令安放在基督头上的。然而，仍有困难区分神圣的十字架与其他十字架；因为铭文已从它上面被扭下并丢弃，十字架本身也与其他十字架一起被随意丢弃，当时被钉者的遗体被取下。据记载，士兵们发现耶稣死在十字架上，便将他取下，交给人埋葬；而为了加速钉在两边的那两个强盗的死亡，他们打断了他们的腿，然后取下十字架，把它们扔到一边。他们并不关心按原样放置十字架；因为天色已晚，且那些人已经死了，他们认为不值得留下处理十字架。因此，需要比人所能提供的更神圣的启示来区分神圣的十字架与其他十字架；这个启示以如下方式赐下：耶路撒冷有一位贵族妇女，患有极其严重且无法治愈的疾病；耶路撒冷主教玛加略，由皇帝的母亲及其随从陪同，前往她的病榻旁。祈祷后，玛加略用手势向旁观者示意，那神圣的十字架将是在接触病人后去除疾病的那一个。他依次将每个十字架靠近她；但当两个十字架放在她身上时，对她而言，这似乎只是愚蠢和嘲弄，因为她已临近死亡之门。然而，当第三个十字架同样被带到她面前时，她突然睁开眼睛，恢复体力，立刻从床上跳起来，痊愈了。据说，有一个人也以同样的方式从死者中复活。可敬的木头就这样被确认后，大部分被存放在一个银匣中，至今仍保存在耶路撒冷；但皇后将一部分送给了她的儿子君士坦丁，连同钉住基督身体的钉子。据说，皇帝用这些钉子为他的马制作了一个头盔和马嚼子，依照匝加利亚的预言，他提及这个时期时说：\Quote{那在马辔头上的，当为全能的上主所圣。}\BibleRef{匝 14:20}这些事，从前确为圣先知所知并预言过，最后，当天主认为它们应被彰显时，便以奇事证实了。这并不显得如此惊人，因为甚至在外教人中，也承认西彼拉曾预言将如此——\par

\Quote{最蒙祝福的树木，我们的主曾悬挂其上。}\par

我们最热忱的敌人无法否认这件事的真实性，由此可见，十字架的木料及其所受的\OriginalTerm{朝拜}{\GreekText{σέβας}}早有预显。\par

以上事件，正如我们从非常严谨的人那里所得知的那样，准确讲述，这些人通过父子相传的继承获得信息；其他人也已将这些事件记录下来，以利后人。\par

% structure: p0007 source=h2 target=subsection reason=relative-heading-rank; base=h2

\subsection{第二章 关于皇帝的母亲海伦；她访问耶路撒冷，在那城里建造圣殿，并行了其他虔诚善工；她的逝世}

大约在同一时期，皇帝决意建造一座圣殿以尊崇天主，遂命令总督们务必以最宏伟、最昂贵的方式完成这项工程。他的母亲赫肋纳也建造了两座圣殿：一座在白冷靠近基督诞生的山洞处，另一座在橄榄山的山脊上，基督正是从那里被接升天堂的。赫肋纳的虔诚与宗教热忱还体现在其他许多事迹中，以下所述尤为突出：据说她在耶路撒冷居住期间，曾召集奉献贞女参加宴会，在晚餐时服侍她们，亲自递上食物，为她们倒水洗手，并履行其他类似服侍宾客的惯例。当她访问东方各城时，向每座城镇的教会赐予适当的礼物，使那些被剥夺财产的人富裕起来，毫不吝啬地供应穷人的必需品，并释放了那些长期被监禁、或被判处流放或矿役的人。在我看来，如此众多的神圣行为理应得到报偿；事实上，即使在今生，她也已被提升到辉煌荣耀的顶峰：她被宣告为奥古斯塔，她的肖像被印在金质钱币上，她的儿子授予她掌管帝国国库的权力，可按照她的判断支配。她的死亡也是荣耀的：当她八十岁离世时，留下了她的儿子和后代（像她一样属于凯撒的世系）成为罗马世界的主人。如果这样的名声有任何益处——遗忘并未因她去世而隐藏她——未来世代将有她永恒记忆的保证：因为两座城市以她的名字命名，一座在彼提尼雅，另一座在巴勒斯坦。这就是赫肋纳的历史。\par

% structure: p0009 source=h2 target=subsection reason=relative-heading-rank; base=h2

\subsection{第三章 君士坦丁大帝建造的圣殿；以他名字命名的城市；其建立；城内的建筑；索斯特尼雍的弥额尔总领天使圣殿及在那里发生的奇迹。}

皇帝始终致力于推动宗教发展，在各地建造了最美丽的天主圣殿，尤其是在大都市中，例如彼提尼雅的尼科美底亚、奥龙特斯河畔的安提约基雅，以及拜占庭。他大力改善拜占庭，使其在权力和政府参与方面与罗马平等；因为当他按照自己的心意整顿帝国事务，并通过战争和条约处理外交之后，他决意建立一座应以他自己的名字命名的城市，并在名声上与罗马相等。为此，他前往特洛伊山麓的平原，靠近赫勒斯滂，在埃阿斯墓之上——据说阿开亚人围困特洛伊时曾在那里设置舰队和营帐——在此规划了一座宏伟美丽的城市，并在高地上建造了城门，至今仍可从海上被航行过去的人看见。但他进展至此，天主夜间显现给他，命令他另寻地点。在天主之手的引领下，他到达了色雷斯的拜占庭，在彼提尼雅的加采东对面，在此他被要求建造他的城市，使其无愧于君士坦丁之名。他于是顺从天主的话语，扩建了原称拜占庭的城市，用高墙环绕。他还在向南的区域建造了宏伟的住宅。他意识到原有人口不足以支撑如此巨大的城市，便以权贵及其家室充实之，将他们从旧罗马和其他国家召来此处。他征收税款以支付建造和装饰城市的费用，并供应居民食物，以及提供城市所需的一切其他物资。他用赛马场、喷泉、柱廊和其他建筑豪华地装饰它。他称它为\Quote{新罗马}和\Quote{君士坦丁堡}，并使之成为所有北、南、东和地中海沿岸居民的帝国首都，从伊斯特河上的城市、埃皮达姆努斯和爱奥尼亚湾，直到库勒纳和被称为波里姆的利比亚地区。\par

他建造了另一个议事厅，他们称之为元老院；他下令给予与罗马其他城市相同的荣誉和节庆日，并且他勤奋地使以他名字命名的城市在各方面与意大利的罗马相等，没有失败；他的愿望也没有受挫；因为在天主的助佑下，不得不承认它在人口和财富上都是伟大的。我不知道有什么理由可以解释这种非凡的壮大，除非是建造者和居民的虔诚，以及他们对穷人的怜悯和慷慨。他们对基督信仰表现出的热忱如此之大，以至于许多犹太居民和大部分希腊人皈依了。由于这座城市在宗教繁荣时期成为帝国的首都，它没有被祭台、希腊神庙和祭献所玷污；尽管\Person{尤利安}曾短期准许引入偶像崇拜，但不久之后它又消失了。\Person{君士坦丁}还通过用许多宏伟的祈祷所装饰它，进一步尊荣这座以他自己名字命名的、新组成的属于基督的城市。天主也与皇帝的精神合作，并通过神圣的显示说服人们，城里的这些祈祷所是神圣而拯救人的。根据外国人和公民的普遍看法，最著名的圣堂是建在以前称为\Place{赫斯提}的地方的那座。这个地方，现在称为\Place{弥额尔堂}，位于从\Place{本都}航行到\Place{君士坦丁堡}的人的右侧，离后者水路大约三十五斯塔迪亚，但如果绕海湾走，则距离为七十斯塔迪亚以上。这个地方得到了现在流行的名称，因为人们相信神圣的总领天使\Person{弥额尔}曾在那里显现。我也肯定这是真的，因为我本人得到了极大的恩惠，而且许多其他人确实得到帮助的经历证明了这一点。因为有些人陷入了可怕的逆境或不可避免的危险，有些人患有疾病和未知的苦难，他们在那里向天主祈祷，他们的不幸得到了改变。如果我详细叙述情况和人物，我会很冗长。但我不能省略提到\Person{阿奎利诺}的例子，他甚至现在仍和我们住在一起，他是我们所属于的同一个法庭的辩护人。我将叙述我从他那里听到的关于这件事的情况以及我所看到的。他因黄胆汁而患上严重的热病，医生给他喝了一些外来的药物。他吐了出来，并且由于呕吐的发力，胆汁扩散开来，使他的面容染上了黄色。因此他不得不呕吐出所有饮食。他长时间处于这种状态；由于营养物在他体内无法安定，医生的技艺也对这种痛苦束手无策。发现他已经半死不活，他命令仆人把他抬到祈祷所；因为他认真地说，在那里他要么死，要么摆脱疾病。当他躺在那里时，夜间有一种神圣的力量向他显现，并命令他将脚浸入一种由蜂蜜、酒和胡椒制成的合剂中。那人照做了，就摆脱了疾病，尽管这个处方违背了医生的职业规则，因为如此热的性质的合剂被认为对胆汁失调不利。我也听说，宫廷医生之一的\Person{普罗比亚诺}，曾患脚病严重，也在这个地方摆脱了疾病，并被认为配得上被一个奇妙而神圣的异象所造访。他以前信奉外教迷信，但后来成为基督徒；然而，尽管他以某种方式承认了我们其余教义的可能性，但他无法理解如何通过神圣的十字架实现所有人的救恩。当他对这个问题心存疑虑时，这座圣堂祭台上的十字架标记在神圣异象中向他指出，他听到一个声音公开宣告：正如基督被钉在十字架上，人类或个人的一切需要，无论是什么，都不能通过神圣天使或虔诚善良的人的服侍来满足；因为除了受人尊敬的十字架之外，没有其他纠正的能力。我只记录了我知道发生在这座圣堂中的几件事，因为没有时间全部叙述。\par

% structure: p0012 source=h2 target=subsection reason=relative-heading-rank; base=h2

\subsection{第四章。\Person{君士坦丁}大帝在\Place{玛默勒}橡树处所做之事；他也建造了一座圣殿。}

我认为有必要详述\Person{君士坦丁}关于称为\Place{玛默勒橡树}之地的作为。此地如今称为\Place{特雷宾托斯}，距南方的\Place{赫贝龙}约十五斯塔狄亚，距离\Place{耶路撒冷}则有二百五十斯塔狄亚。据记载，天主圣子曾在此地与两位奉命前往\Place{索多玛}的天使一同显现给\Person{亚巴郎}，并预言了他儿子的诞生。此地居民以及\Place{巴勒斯坦}周边地区的人、\Person{腓尼基人}和\Person{阿拉伯人}，每年夏季聚集在此举行盛大的节庆；还有许多商贩因市集而前往。实际上，所有民族都殷勤参与这个节庆：犹太人因以圣祖\Person{亚巴郎}的后裔自夸；异教徒因天使曾在此向人显现；基督徒因那为了人类的救恩而由童贞女所生者，后来在此地显现给一位虔诚的人。此外，这地方也以宗教仪式合宜地受到尊敬。有人在此向万有的天主祈祷；有人呼求天使，奠酒、焚香，或献上牛、公山羊、绵羊或公鸡。各人献上自己劳动的美果，经过整年精心保存，按照许诺作为节庆的食物，为自己和家属预备。或因尊敬圣地，或因畏惧天主的义怒，他们都远离自己的妻子，尽管在节庆期间她们更加刻意修饰装扮。即使她们偶然出现并参与公共游行，也绝无放荡行为。在其他方面，他们也未曾轻率行事，虽然帐幕彼此相连，众人混杂共处。该地是开阔的田野，可耕种，没有房屋，只有\Person{亚巴郎}古橡树周围和他所掘的井旁的建筑。节庆期间无人从那口井中取水；因为按异教徒的习俗，有人在井旁放置点燃的灯盏；有人奠酒或投入饼糕；还有人投入钱币、没药或乳香。因此，我想，水因与投入之物混合而变得无用。有一次，异教徒按上述方式和既定习俗欢庆时，\Person{君士坦丁}的岳母在场祈祷，并将所行之事告知皇帝。得到消息后，他毫不客气地斥责了\Place{巴勒斯坦}的主教们，因为他们玩忽职守，容许圣地被不洁的奠祭和祭献所玷污；并在他就此写给\Place{耶路撒冷}主教\Person{玛卡略}、\Person{欧瑟比·庞非勒}和\Place{巴勒斯坦}诸位主教的信函中，表达了虔诚的谴责。他命令这些主教与\Place{腓尼基}的主教们就此问题举行会议，并发出指令：将从前在那里建立的祭台彻底拆毁，将雕像用火销毁，并建造一座与如此古老神圣之地相称的圣堂。皇帝最后命令，在那地点不得再举行任何奠祭或祭献，应完全按照教会的法律专用于敬拜天主；若有人试图恢复先前的仪式，主教们应检举违法者，使其受最严厉的惩罚。总督和\Person{基督}的司铎们严格执行了皇帝信件中的命令。\par

% structure: p0014 source=h2 target=subsection reason=relative-heading-rank; base=h2

\subsection{第五章 \Person{君士坦丁}摧毁了供奉偶像的地方，并劝说人民偏好基督教。}

在他所有臣民的整个王国中，许多民族和城市仍然对他们的虚妄偶像怀有敬畏之心，这导致他们忽视基督徒的教义，而关心他们祖先的古旧习俗、生活方式和节庆。因此，皇帝认为有必要教导总督们压制他们迷信的崇拜仪式。他想，如果能让他们藐视自己的神庙及其中的雕像，这件事就容易完成。实施这一计划时，他并不需要军事援助；因为宫廷中的基督徒男子带着皇帝的诏书走遍各城。受胁迫的人民因担心若抗拒这些法令，自己和妻子儿女会遭受祸害，而保持沉默。看守和祭司因缺乏民众支持，交出他们最宝贵的财富，以及被称为\OriginalTerm{从天而降的}{\GreekText{διοπετῆ}}的偶像；通过这些仆役，礼物从神龛和庙宇的隐秘之处被取出。先前无法进入、只有祭司知道的地点，变得对所有愿意进入的人开放。那些用贵重材料制成的雕像以及其他有价值的东西，被火净化，成为公共财产。工艺精湛的铜像被运到以皇帝命名的城市，作为装饰品放置在那里，如今仍可在公共场所如街道、赛马场和宫殿中见到。其中就有位于\Person{皮提亚}女巫神谕所中的\Person{阿波罗}雕像，还有来自\Place{赫利孔山}的\Person{缪斯}雕像、来自\Place{德尔斐}的鼎，以及备受赞颂的\Person{潘神}雕像，这些都是\Person{保萨尼阿斯}（拉栖第梦人）和希腊城邦在\Place{米底}战争之后奉献的。\par

至于这些神庙，有些被夺去了门，有些被掀去了屋顶，其他的则被废弃，任其倾颓或拆毁。\Place{基里基亚}的\Place{埃吉斯}城的\Person{阿斯克勒庇俄斯}神庙和靠近\Place{黎巴嫩山}与\Place{阿多尼斯河}的\Place{阿法卡}的\Person{维纳斯}神庙，当时都被掘毁，彻底消灭。这两座神庙最受古人尊崇和敬拜；因为\Place{埃吉斯}人常说，他们中间身体虚弱的人得以摆脱疾病，是因为那鬼神在夜间显现并治愈了他们。而在\Place{阿法卡}，人们相信，在某个特定的日子念出某篇祷词后，一团像星一样的火焰从\Place{黎巴嫩山}顶降下，沉入附近的河中；他们宣称这是\Person{乌拉尼亚}，因为他们用这个名字称呼\Person{阿佛洛狄忒}。皇帝的作为达到了他预期的最大成功；因为看到他们以前所敬畏和恐惧的对象被大胆推倒，塞满稻草和干草，人民遂开始鄙视他们先前所崇敬之物，并责备祖先的错误见解。另一些人，嫉妒基督徒在皇帝那里所获得的尊荣，认为有必要效法君主的作为；还有些人专心考察基督教，藉着记号、异梦或与主教和隐修士的会谈，而深信成为基督徒是更好的。从这时起，各族和市民自发地抛弃了他们从前的见解。那时，\Place{加沙}的一个港口\Place{马尤马}，那里迷信和古老的礼仪直到当时仍受推崇，全体居民一致转而信奉基督教。皇帝为酬报他们的虔诚，认为他们堪受最大的光荣，将那地方提升为城市——它此前无此身份——并命名为\Place{君士坦提亚}，以此因其虔诚而尊荣那地，将他最亲爱的子女的名字赐给它。同样，\Place{腓尼基}的\Place{君士坦丁}城也已知是从皇帝得名。但纪录每一此类事件并不适宜，因为此时还有许多别的城市皈依了宗教，并且未经皇帝任何命令，自行拆毁了邻近的神庙和雕像，建造了祈祷之所。\par

% structure: p0017 source=h2 target=subsection reason=relative-heading-rank; base=h2

\subsection{第6章  在君士坦丁治下，基督的圣名何以传遍整个世界。}

教会就这样传遍了整个罗马世界，宗教甚至传入了野蛮人中间。\Place{莱茵河}两岸的部落都归依了基督教，同样还有住在遥远海岸的\Person{凯尔特人}和\Person{高卢人}；此外\Person{哥特人}以及与他们相邻的部落，他们从前住在\Place{多瑙河}两岸的高岸上，早已分享了基督信仰，并转变为更温和、更理性的遵守。几乎所有的野蛮人，从\Person{加里恩努斯}及继他之后的皇帝统治时期罗马人与外族之间的战争开始，都宣称尊崇基督的教导。因为当一群无法言说、混杂各族的民众从\Place{色雷斯}渡海进入\Place{亚洲}并蹂躏它，而其他地区来的野蛮人也对邻近的罗马人做了同样的事时，许多被俘的基督的\OriginalTerm{司铎}{priests}就住在这些部落中；在他们居留期间，只用基督的名和呼求\Person{圣子}，就治愈了病人，洁净了附魔的人；此外，他们过着无可指责的生活，并以他们的德行引起羡慕。野蛮人对这些人的行为和奇妙作为感到惊异，认为效法他们所见的更好的人，并像他们一样敬拜天主，对他们来说是明智的，也是中悦天主的。当他们被告知当行的教导后，民众受教、受洗，后来被聚集到教会中。\par

% structure: p0019 source=h2 target=subsection reason=relative-heading-rank; base=h2

\subsection{第7章  \Person{伊比利亚人}如何接受基督信仰。}

据说，在该朝期间，伊比利亚人——一个庞大而好战的蛮族——承认了基督。他们居住在亚美尼亚以北。一位被俘的基督徒妇女促使他们摒弃了祖传的宗教。她非常忠信而虔诚，在异邦人中也不松懈她惯常的宗教义务。禁食、日夜祈祷、赞美天主，是她的喜乐。蛮族人询问她忍耐的缘由：她简单地回答，必须这样敬拜天主之子；但祂的名号以及敬拜的方式，在他们看来是陌生的。恰巧当地一个男孩患病，他的母亲按照伊比利亚人的习俗，带他挨家挨户走访，希望有人能治愈这病，使患者脱离痛苦。当找不到能医治他的人时，男孩被带到女俘虏面前，她说：\Quote{至于药物，我既无经验也无知识，也不懂得涂抹油膏或药膏的方法；但是，妇人啊，我相信我所敬拜的基督，那真实而伟大的天主，会拯救你的孩子。}于是她立刻为他祈祷，解除了他的疾病，尽管此前人们认为他快要死了。不久之后，该民族总督的妻子因不治之症濒临死亡，然而她也以同样的方式得救。就这样，这位女俘虏通过介绍基督作为健康的赐予者、生命的主宰、君权的主宰和万有的主宰，教导了关于基督的知识。总督的妻子因亲身体验而相信了女俘虏的话，接受了基督信仰，并极为尊敬这妇人。国王对治愈的速度、信德的神奇与治愈力量感到惊奇，从妻子那里得知缘由后，便命令用礼物赏赐女俘虏。\Quote{礼物，}王后说，\Quote{无论价值如何，她看得很轻；她只重视对她天主的侍奉。因此，若我们想取悦她，或想做稳妥而正确的事，让我们也敬拜那大能而拯救的天主吧，祂按自己的意愿赐予君王延续，贬抑高傲者，使显赫者卑微，拯救危难中的人。}王后继续以如此卓越的方式论辩，但伊比利亚的君主仍犹豫不决，未被说服，因为他思忖这些事的新奇，也尊重祖传的宗教。不久后，他带着随从去森林打猎；突然，浓云密布，沉重的空气弥漫四周，遮蔽了天空和太阳；深沉的黑夜与巨大的黑暗笼罩了森林。由于每个猎人都为自己的安全担忧，他们各自四散奔逃。国王独自徘徊时，想起了基督，正如人们在危险时刻惯常所做的那样。他决定，若能从眼前的困境中得救，他将在天主面前行走并敬拜祂。就在这些念头萦绕他心头的瞬间，黑暗消散，天空晴朗，阳光射入森林，国王安全地走了出来。他将所发生的事告诉王后，派人请来女俘虏，命令她教导他应当如何敬拜基督。在她给予了女人所当说当做的教导之后，国王召集臣民，向他们清楚宣告了赐予他和王后的神圣仁慈，并虽未受训，却向民众宣讲了基督的教义。全体国民被说服接受基督信仰，男子因国王的陈述而信服，女子因王后和女俘虏的陈述而信服。很快，在全体国民的一致同意下，他们热忱地准备建造一座教堂。外墙建成后，运来了机械以竖起柱子并固定在柱座上。据说，当第一和第二根柱子被这样扶正后，第三根柱子却遇到了极大困难，无论技艺还是体力都无济于事，尽管有许多人帮忙拉拽。傍晚时分，女俘虏独自留在那地方，她整夜在那里祈求天主使柱子容易立起，尤其是因为其他人都因失败而沮丧地离开了；因为柱子只升起了一半，悬在那里，一端深嵌在基础中无法向下移动。天主愿意借此，以及先前的奇迹，使伊比利亚人对天主更加坚定。清晨，当他们来到教堂时，看见了奇妙的景象，仿佛梦境一般。前一天还无法移动的柱子，现在竟直立着，并高出其底座一点。所有在场的人都惊叹不已，一致承认唯独基督是真天主。在他们观看时，柱子悄然自发地滑下，犹如被机械调整一般落在底座上。其余的柱子便轻易立起，伊比利亚人更加迅速地完成了建筑。教堂如此快速建成后，伊比利亚人按照女俘虏的建议，派遣使者到君士坦丁皇帝那里，提出结盟和约的建议，并请求派遣司铎到他们民族。使者到达后，报告了所发生的事，以及整个民族如何虔诚地敬拜基督。罗马皇帝对使团感到喜悦，同意了所提出的每一项请求后，遣返了使者。伊比利亚人就这样认识了基督，直到今日他们仍虔诚地敬拜祂。\par

% structure: p0021 source=h2 target=subsection reason=relative-heading-rank; base=h2

\subsection{第八章：亚美尼亚人和波斯人如何皈依基督信仰。}

\Emph{Subsequently}基督信仰传到了邻近的部落，并大大地传播开来。我听说\Place{亚美尼亚}人是最先接纳基督信仰的。据说当时的君主\Person{提里达特斯}因在他家中发生的一个奇妙的天主记号而成了基督徒；他命令传令官向所有臣民宣布，要采纳同样的宗教。我认为\Place{波斯}人的归化是由于他们与\Place{奥斯若恩}人和\Place{亚美尼亚}人的来往；因为他们很可能与这些神圣的人交谈，并体验了他们的德行。\par

% structure: p0023 source=h2 target=subsection reason=relative-heading-rank; base=h2

\subsection{第九章 \Place{波斯}王\Person{沙普尔}被激动反对基督徒。\Person{西默盎}，\Place{波斯}主教，和\Person{乌斯塔扎内斯}，宦官，遭受殉道之苦。}

当时间推移，基督徒的人数增加，并开始形成教会，任命司铎和执事时，作为祭司部族从起初历代以来守护波斯宗教的\OriginalTerm{贤士}{Magi}大大被激怒反对他们。因嫉妒而在某种程度上天然反对基督信仰的犹太人也同样被触怒。因此他们向在位的君王\Person{沙普尔}控告当时担任\Place{波斯}王城\Place{塞琉西亚}和\Place{泰西封}总主教的\Person{西默盎}，指控他是罗马\OriginalTerm{凯撒}{Cæsar}的朋友，并将波斯事务通报给他。\Person{沙普尔}相信了这些控告，起初以重税压迫基督徒，尽管他知道他们大多数人是自愿选择贫困的。他将征收税赋的事托付给残忍的人，希望由于缺乏生活必需品以及收税者的暴行，他们会被迫背弃自己的宗教；因为这是他的目的。后来，他命令将司铎和敬礼天主的领导人用剑杀死。圣堂被拆毁，器皿存入国库，\Person{西默盎}被逮捕，作为王国和波斯宗教的叛徒。这样，贤士在犹太人的合作下，迅速摧毁了祈祷之所。\Person{西默盎}被捕时，被铁链捆绑，带到国王面前。在那里，此人显示了他的卓越和勇气；因为当\Person{沙普尔}命令将他带去受酷刑时，他不畏惧，也不愿俯伏朝拜。国王大为恼怒，质问他为何不像从前那样俯伏朝拜。\Person{西默盎}回答说：\Quote{从前我不是被捆绑着带来，为要我背弃天主的真理，所以我那时不反对向王权表示惯常的尊敬；但现在我这样做就不合适了，因为我站在这里是为捍卫虔诚和我们信仰。}他说完后，国王命令他朝拜太阳，并许诺要赐给他礼物，尊荣他；另一方面，威胁说若不顾从，就要降灾给他和全体基督徒。国王见他既未被威吓吓倒，也未被许诺动摇，而\Person{西默盎}坚定不移，拒绝朝拜太阳或背弃自己的宗教，就命令暂时把他捆绑起来，大概以为他会改变主意。\par

当\Person{西默盎}被押往监狱时，年迈的太监\Person{乌斯塔扎内斯}——\Person{沙普尔}的养父兼宫廷总管——正坐在宫门口，起身向他致敬。\Person{西默盎}以高声傲慢的语调责备地禁止他，转过脸去，径直走了过去。原来这太监从前是基督徒，但最近屈服于权势，崇拜了太阳。此举深深触动了太监，他放声痛哭，脱下所披的白衣，像哀悼者一样穿上黑衣，然后坐在宫殿前，哭泣哀叹说：\Quote{我有祸了！我既然否认了天主，还有什么不能临到我呢？因此\Person{西默盎}，我素昔的密友，认为不值得与我说话，竟转身离我而去。}\Person{沙普尔}听闻此事后，召来太监，询问他悲伤的原因，问他家中是否遭了灾祸。\Person{乌斯塔扎内斯}回答说：\Quote{大王，家中并无灾祸；但若遭其他任何祸患，都比这强，也更容易忍受。如今我哀哭，是因为我还活着，而本该早已死去；我还能看见太阳——这太阳，我并非自愿，而是为取悦你才宣称崇拜的。因此，我理应死去：我既是基督的背叛者，又是你的欺骗者。}接着他指着天地的主宰起誓，决不再动摇自己的信念。\Person{沙普尔}对太监奇妙的转变大为惊愕，却更加恼怒基督徒，仿佛他们是靠邪术实现此事。不过，他饶了老人一命，并竭尽全力，时而温和、时而严厉，想使他回心转意。但见他徒劳无功，而\Person{乌斯塔扎内斯}坚持宣称自己绝不会愚昧到崇拜受造物而不崇拜造物主，他便怒火中烧，下令用剑斩下太监的首级。刽子手上前执刑时，\Person{乌斯塔扎内斯}请求他们稍候，以便向国王禀告一些事。于是他召来一位最忠诚的太监，命他对\Person{沙普尔}说：\Quote{大王，我从幼年至今，一向对你家忠心耿耿，伺候你父王与你，尽忠尽职。我无需证人证实我的话，这些事实确凿。为了我往日屡次乐意服侍你的一切事，求你赐我这赏报：不要让不明真相的人以为我因背叛国政或犯有其他罪行而受此刑；而要由传令官宣告并公布：\Person{乌斯塔扎内斯}斩首，并非因为他在宫中犯有任何奸恶，而是因为他是基督徒，并因不肯听从王命否认自己的天主。}太监传达此言后，\Person{沙普尔}依\Person{乌斯塔扎内斯}所请，命传令官按所愿宣告；因为国王心想，若有人想到他连年迈的养父和珍爱的家仆都牺牲了，必定不会放过其他基督徒，这样就能轻易劝阻他人信奉基督教。但\Person{乌斯塔扎内斯}坚信，正如他因胆怯而同意崇拜太阳曾使许多基督徒恐惧，如今藉着公开详述受苦的原因，许多人得知他为宗教而死，必能受启迪，从而效法他的刚毅。\par

% structure: p0026 source=h2 target=subsection reason=relative-heading-rank; base=h2

\subsection{第十章 \Person{沙普尔}在\Place{波斯}杀害基督徒。}

\Person{西默盎}在狱中得知\Person{乌斯塔扎内斯}以如此可敬的方式结束生命后，便为他的缘故感谢天主。次日，恰逢一周的第六天，也是纪念救主受难的前夕——即复活节前夕的年度纪念日——国王下令将\Person{西默盎}斩首；因为他再次从狱中被带到宫中，在教义问题上与\Person{沙普尔}进行了极其高贵的辩论，并表明决意不崇拜国王或太阳。同日，另有百名囚犯被下令处死。\Person{西默盎}目睹他们受刑，最后他自己也被处决。这些殉道者中包括主教、长老及其他不同品级的圣职人员。当他们被带往刑场时，贤士的首领走近他们，问他们是否愿意服从国王的宗教、崇拜太阳以保全性命。因无人肯从，他们便被带至刑场，刽子手开始屠杀这些殉道者。\Person{西默盎}站在待戮者身旁，劝勉他们恒心忍耐，论证死亡、复活与虔敬之事，并从圣经指出，像这样的死才是真正的生命；而活着却因惧怕否认天主，则无异于真死。他又告诉他们，即使无人杀害他们，死亡也终必临到；因为死亡是我们出生的自然结果。此世之后的事是永恒的，并非人人相同；而是如同用规矩丈量一般，必须对此生的历程作出精确交代。凡行善者，必得永生的赏报，逃避行恶者所受的惩罚。他还说，最伟大、最幸福的行为莫过于为天主而死。\Person{西默盎}正这样论说，如同家仆一般指点他们如何赴战，各人听罢，便勇敢地走向杀戮。刽子手处决了百人之后，\Person{西默盎}本人也被杀害；与他一同被囚的两位本教会年迈长老——\Person{阿贝德查拉斯}和\Person{阿纳尼亚斯}——也与他一同殉道。\par

% structure: p0028 source=h2 target=subsection reason=relative-heading-rank; base=h2

\subsection{第十一章 \Person{沙普尔}工艺总监\Person{普西克斯}。}

王工监督\Person{Pusices}当时在场，看见\Person{Anannias}在准备受刑时颤抖，就对他说：\Quote{老人家，暂且闭目，鼓起勇气，因为你很快就要看见基督的光了。}他说完这话，立即被捕，带到王前。他坦白承认自己是基督徒，并大胆地向王谈论自己的意见和那些殉道者，因此被判处一种异常残酷的死刑，因为对王如此直言是不合法的。刽子手刺穿他颈部的肌肉，将他的舌头拔了出来。因某些人的指控，他那已献身于圣洁童贞生活的女儿同时被审讯并处决。次年，在纪念基督苦难的那一天，当准备庆祝祂从死者中复活的节日时，\Person{Sapor}王在波斯颁发了一道极其残酷的敕令，判处所有承认自己是基督徒的人死刑。据说当时有更多的基督徒死于刀下；因为\OriginalTerm{玛日僧}{Magi}在城乡中仔细搜寻那些藏匿起来的人；还有许多基督徒自愿投案，免得因沉默而显得是否认基督。在这些被毫不留情地牺牲的基督徒中，有许多王宫中的侍从被杀，其中包括一个特别受王宠爱的阉人\Person{Azades}。\Person{Sapor}听到他的死讯，悲痛欲绝，便停止了普遍屠杀基督徒的行动，并下令只杀宗教导师。\par

% structure: p0030 source=h2 target=subsection reason=relative-heading-rank; base=h2

\subsection{第十二章   \Person{Tarbula}，\Person{Symeon}的姐妹，及其殉道。}

大约同一时期，王后患了一种病，主教\Person{Symeon}的姐妹\Person{Tarbula}，一位圣洁的童贞女，与她的女仆（过着同样生活）一同被捕，还有\Person{Tarbula}的一位姐妹，她在丈夫去世后誓守寡居，过着类似的生活。她们被捕的原因是犹太人的指控，他们报告说，这些人因对\Person{Symeon}之死感到愤怒，用法术伤害了王后。由于病人容易相信最令人厌恶的指控，王后相信了这指控，尤其因为它来自犹太人，因为她接受了他们的观点，并遵守犹太礼仪生活，她极其相信他们的诚实和对她的忠诚。\OriginalTerm{玛日僧}{Magi}抓住了\Person{Tarbula}和她的同伴，判处她们死刑；他们将她们锯成两半，然后固定在柱子上，让王后从柱子中间走过，作为驱除疾病的方法。据说这个\Person{Tarbula}容貌美丽，体态非常端庄，有一个玛日僧深深迷恋她，秘密派人提出苟合的建议，并承诺如果她同意，就救她和她的同伴。但她不听他的淫荡之言，轻蔑地对待玛日僧，斥责他的情欲。她宁愿勇敢地死去，也不愿失去童贞。\par

按照我们上文提到的\Person{Sapor}王的敕令，基督徒不应被不加区别地屠杀，而只应杀死司铎和教义的导师，\OriginalTerm{玛日僧}{Magi}和\OriginalTerm{大玛日僧}{Arch-Magi}走遍整个波斯地区，刻意虐待主教和长老。他们尤其到\Place{Adiabene}地区（波斯领土的一部分）搜寻，因为那里已完全基督化了。\par

% structure: p0033 source=h2 target=subsection reason=relative-heading-rank; base=h2

\subsection{第十三章  圣\Person{Acepsimas}及其同伴的殉道。}

大约在这时期，他们逮捕了\Person{阿塞普西玛斯}主教和他的许多神职人员。他们商议之后，只满足于追捕首领；其余的人被没收财产后释放。然而，一位名叫\Person{雅各伯}的司铎自愿跟随\Person{阿塞普西玛斯}，获得麻葛许可与他同囚，并勇敢地服侍这位老人，尽力减轻他的不幸，为他包扎伤口；因为在他被捕后不久，麻葛曾残忍地用生牛皮鞭殴打他，强迫他崇拜太阳；他拒绝后，他们又把他关进锁链。两位司铎，名叫\Person{艾塔拉斯}和\Person{雅各伯}，以及两位执事，名叫\Person{阿扎达内斯}和\Person{阿布迪耶苏斯}，被麻葛极其羞辱地鞭打后，被迫因信仰关在狱中。过了很长时间，大麻葛向国王禀报了这些人的情况，请求惩处；并获准自由处置他们，除非他们同意崇拜太阳。大麻葛向囚犯们宣布了\Person{沙普尔}的这一决定。他们公开回答说，绝不会背叛\Person{基督}的事业，也不会崇拜太阳；于是他毫不留情地折磨他们。\Person{阿塞普西玛斯}英勇宣认信德，坚持到底，直到死亡结束他的痛苦。一些被波斯人扣为人质的亚美尼亚人秘密偷走了他的尸体并埋葬了。其他囚犯虽然同样遭受鞭打，却奇迹般地活了下来，由于他们不肯改变主意，又被戴上锁链。其中\Person{艾塔拉斯}在被打时被拉伸，胳膊因极度扭曲而从肩膀脱臼；他双手像死了一样垂着摆动，别人不得不把食物送到他嘴里。在这统治下，无数司铎、执事、隐修士、圣洁贞女以及其他服务于教会并为其教义献身的人，殉道而终。以下是我所能查明的主教的名字：\Person{巴尔巴西梅斯}、\Person{保禄}、\Person{加迪阿贝斯}、\Person{萨比努斯}、\Person{马雷阿斯}、\Person{莫丘斯}、\Person{若望}、\Person{霍尔米斯达斯}、\Person{帕帕斯}、\Person{雅各伯}、\Person{罗马斯}、\Person{马雷斯}、\Person{阿加斯}、\Person{博赫雷斯}、\Person{阿布达斯}、\Person{阿布迪耶苏斯}、\Person{若望}、\Person{阿布拉明斯}、\Person{阿格德拉}、\Person{萨波雷斯}、\Person{依撒格}、\Person{道萨斯}。后者被波斯人俘虏，从名叫\Place{扎布代乌斯}的地方带来。他大约在这时期为教义而死；乡村主教\Person{马雷阿布德斯}以及他约二百五十名也被波斯人俘虏的神职人员，与他一同殉难。\par

% structure: p0035 source=h2 target=subsection reason=relative-heading-rank; base=h2

\subsection{第十四章 \Person{米勒斯}主教殉道及其品行。\Person{沙普尔}治下\Place{波斯}一万六千名杰出男士殉道，另有无数无名者。}

大约在这时期，\Person{米勒斯}殉道。他最初为波斯人服军役，但后来放弃此业，选择了修道生活。据说他被祝圣为一座\Place{波斯}城市的主教，遭受了各种苦难，受了许多伤和折磨；他努力使居民皈依基督信仰，却未能成功，就诅咒那城，然后离去。不久，一些主要市民冒犯了国王，一支带着三百头象的军队被派去攻打他们；城市被彻底摧毁，土地被犁过并撒了种。\Person{米勒斯}只带着他的行囊，里面装着神圣的福音经书，前往\Place{耶路撒冷}祈祷；然后从那里去了\Place{埃及}，为见隐修士们。我们听说他所完成的非凡而值得称颂的事迹，由叙利亚人证实，他们写下了他的言行录。至于我自己，我认为关于他和\Person{沙普尔}统治期间在\Place{波斯}殉道的其他殉道者，我说得已经够了；因为要详细叙述每个人的情况，诸如他们的名字、国籍、殉道的方式以及所受的酷刑，是困难的；因为人数众多，波斯人热衷于采用这类方法，甚至残忍至极。我简要地说，经查明的这一时期殉道的男女，统计为一万六千人；此外还有不计其数的群众，因此，对于波斯人、叙利亚人和\Place{埃德萨}居民来说，列举他们的名字似乎很困难，他们曾为此投入大量心血。\par

% structure: p0037 source=h2 target=subsection reason=relative-heading-rank; base=h2

\subsection{第十五章 \Person{君士坦丁}致信\Person{沙普尔}，请求停止迫害基督徒。}

罗马皇帝\Person{君士坦丁}听说基督徒在波斯遭受的苦难，非常愤怒，难以忍受。他极其渴望帮助他们，却不知如何实现这一目标。此时，波斯国王的一些使节抵达他的宫廷。他满足了他们的请求并打发他们走后，认为这是为波斯基督徒向沙普尔进言的好机会，便写信给沙普尔，承认如果他能善待那些钦佩他治下基督徒教导的人，那将是一个极大且永远难以言喻的恩惠。他说：\Quote{他们的宗教中没有任何可指责之处；他们仅以不流血的祈祷向天主呈上恳求，因为天主不喜悦流血，只喜悦纯洁的灵魂献身于德行和宗教；因此，相信这些事的人是值得称赞的。}随后，皇帝向沙普尔保证，如果他宽待基督徒，天主会对他仁慈，并以瓦莱里安和他自己为例证。他自己因信靠基督，并藉着神圣眷顾的帮助，从西海的岸边出发，使整个罗马世界归顺，平息了许多对外邦人和篡位者的战争；但他从未求助于祭献或占卜，而仅在军队前使用十字架的标记，并献上不流血、无玷污的祈祷来争取胜利。瓦莱里安的统治在他停止迫害教会时本是兴盛的；但后来他开始迫害基督徒，便因神圣的报复被交到波斯人手中，被他们俘虏并残忍地处死。\par

\Person{君士坦丁}以这样的口吻写信给沙普尔，敦促他善待这一宗教；因为皇帝对所有地区的基督徒，无论是罗马人还是外邦人，都怀有警醒的关怀。\par

% structure: p0040 source=h2 target=subsection reason=relative-heading-rank; base=h2

\subsection{第十六章。\Person{欧塞比乌}和\Person{特奥格尼斯}在尼西亚公会议上曾赞同阿里乌的著作，后恢复原职。}

尼西亚公会议后不久，阿里乌被从流放地召回；但禁止进入亚历山大的禁令仍未撤销。他如何力争获准返回埃及，将在适当的地方叙述。不久之后，尼科美底亚的主教欧塞比乌和尼西亚的主教特奥格尼斯重新获得了他们的教会，驱逐了接替他们被祝圣的安菲翁和克雷斯特斯。他们的复职归功于他们向主教们提交的一份文件，其中包含一份撤回声明：\Quote{尽管我们未经审判就被你们的虔敬判罪，我们仍认为对你们虔敬的判决保持沉默是合适的。但若继续沉默，而沉默被视为对诽谤者所言之真实性的证明，那将是荒谬的，所以我们现向你们声明：我们也同意这一信仰，在仔细审查了\InnerQuote{同质}一词的含义后，我们完全致力于维护和平，我们从未追求任何异端。我们为教会的安全提出了我们想到的论点，并已完全信服，也完全说服了那些本应被我们说服的人，我们签署了信经；但我们没有签署绝罚令，不是因为我们反对信经，而是因为我们不相信被指控者是我们所听到的那种人；我们从他那里收到的信件以及他在我们面前提出的论点，充分使我们相信他并非那样的人。但愿神圣的会议相信，我们并非意在反对，而是赞同你们所准确定义的各点，并且通过这份文件，我们证明我们同意这些点：这不是因为我们厌倦了流放，而是因为我们希望消除一切异端的嫌疑；因为如果你们现在屈尊接纳我们到你们面前，你们会发现我们在所有问题上与你们意见一致，服从你们的决定，那么你们的虔敬就可以仁慈地对待那个在这些问题上被指控的人，并把他召回。如果受审方已被召回并为自己辩护免于指控，那么我们若以沉默证实诽谤我们散布的谣言，那是荒谬的。因此，我们恳求你们，以你们合乎虔敬、爱基督之心，致函我们最蒙天主爱戴的皇帝，递交我们的请愿书，并尽快就你们对我们的意见给予忠告。}欧塞比乌和特奥格尼斯正是通过这种方式，在改变态度后，重新回到了他们的教会。\par

% structure: p0042 source=h2 target=subsection reason=relative-heading-rank; base=h2

\subsection{第十七章。亚历山大主教\Person{亚历山大}去世后，根据他的建议，\Person{阿塔纳修斯}登上主教宝座；关于他的青年时代：他如何是自学成才的司铎，并受到大\Person{安东尼}的爱戴。}

大约在此时，\Place{亚历山大里亚}主教\Person{亚历山大}即将离世时，遵从天主旨意（我深信如此）指引对他的投票，留下\Person{阿塔纳修斯}作为他的继任者。据说\Person{阿塔纳修斯}起初试图逃避这荣誉而逃亡，但他虽不情愿，后来却被\Person{亚历山大}强制接受了主教之职。\Place{叙利亚}人\Person{阿波利纳留斯}以如下言辞作证：\Quote{在所有这些事中，不虔敬激起了许多骚乱，但其最初的效果却临到了此人的有福导师，他当时作为助手在场，并像儿子对待父亲一样行事。后来这位圣洁者本人也经历了同样的遭遇，因为当他被任命为主教继任者时，他逃亡以逃避荣誉；但藉着天主的帮助，他在藏身之处被发现，天主曾以神圣显现预言给他的有福前任，继任将归于他。因为当\Person{亚历山大}临死时，他呼唤当时不在场的\Person{阿塔纳修斯}。一位同名者恰好在场，听到他这样呼唤便回答了他；但\Person{亚历山大}对他沉默，因为他并未召唤此人。他又呼唤，正如常发生的那样，在场者保持沉默，于是不在场者被显明。此外，有福的\Person{亚历山大}预言般地喊道：\InnerQuote{啊，\Person{阿塔纳修斯}，你以为可以逃脱，但你逃脱不了}，意思是\Person{阿塔纳修斯}将被召入冲突。}这就是\Person{阿波利纳留斯}关于\Person{阿塔纳修斯}的记述。\par

亚略派者断言，\Person{亚历山大}去世后，该主教和\Person{美利提乌斯}各自的追随者彼此共融，来自\Place{底比斯}和\Place{埃及}其他地区的五十四位主教聚集，共同宣誓通过共同投票选出一个能有益地管理\Place{亚历山大里亚}教会的人；但七位主教违反誓言，并违背所有人的意见，秘密祝圣了\Person{阿塔纳修斯}；因此许多百姓和\Place{埃及}神职人员与他断绝了共融。就我而言，我深信\Person{阿塔纳修斯}是由天主上智继承大司祭之职的；因为他口才出众、智慧聪颖，能够应对阴谋，而当时的时代极其需要这样的人。他在行使教会职务中表现出极大的才能，并适合司祭之职，而且可以说是自幼自学成才。据说他年轻时发生过以下事件。\Place{亚历山大里亚}人习惯每年隆重纪念一位名叫\Person{伯多禄}的殉道主教。当时管理教会的\Person{亚历山大}参与了这个庆典，在完成了敬礼后，他留在现场等待他预期共进早餐的客人。在此期间，他偶然向海望去，看见一些儿童在海滩上玩耍，并模仿主教和教会的礼仪自娱。起初他认为这种模仿是无害的，并乐于观看；但当他们触及不可言传之事时，他感到不安，并将此事告诉了神职人员的首领。儿童们被召集起来，询问他们玩的是什么游戏，以及在这娱乐中他们做了什么、说了什么。起初他们因害怕而否认；但当\Person{亚历山大}以酷刑威胁时，他们承认\Person{阿塔纳修斯}是他们的主教和首领，许多尚未受洗的儿童已被他施洗。\Person{亚历山大}仔细询问了他们游戏中的司祭习惯说什么或做什么，以及他们回答或学到了什么。发现教会的精确礼仪已被他们准确遵循后，他与周围的司祭商议，决定没有必要给那些因纯朴而被认为堪得天主恩宠的人重新施洗。因此他只为他们举行了那些只允许被祝圣者施行的入门礼。然后他以天主为证人，将\Person{阿塔纳修斯}和其他扮演司铎和执事的儿童交给他们的亲属，以便他们为教会和为所模仿的领袖身份被培养。不久之后，他接纳\Person{阿塔纳修斯}作为他的餐桌伴侣和秘书。\Person{阿塔纳修斯}受过良好教育，通晓文法和修辞，且当他成年后，在就任主教之前，已向与他交谈的人证明他是一位智者和有才智的人。但当\Person{亚历山大}去世后，继任之责落在他身上时，他的声望大大增加，并由他个人的德行和隐修士伟大的\Person{安当}的见证所维持。这位隐修士应他请求前往见他，访问城市，陪同他去教堂，并在关于天主性的观点上与他一致。伟大的\Person{安当}向他显示了无限的友谊，并回避了他敌人和反对者的交往。\par

% structure: p0045 source=h2 target=subsection reason=relative-heading-rank; base=h2

\subsection{第十八章　亚略派者和美利提乌斯派者给\Person{阿塔纳修斯}带来声望；关于\Person{欧瑟伯}及其要求\Person{阿塔纳修斯}接纳亚略共融之事；关于术语\OriginalTerm{同质}{Consubstantial}；\Person{欧瑟伯·庞斐禄}和\Place{安提约基雅}的\Person{欧斯达提乌斯}主教特别引起了骚动。}

然而，\Person{阿塔纳修斯}的声望尤其因亚略派者和美利提乌斯派者而增长；尽管他们不断策划，却似乎从未能正确抓住并囚禁他。首先，\Person{欧瑟伯}写信敦促他接纳亚略派者共融，并隐含威胁说，若他拒绝如此行，将虐待他。但\Person{阿塔纳修斯}不屈服于他的陈述，坚持认为那些在真理上创新并发明异端、且已被尼西亚会议谴责的人不应被接纳进入教会，于是\Person{欧瑟伯}设法使皇帝支持亚略，从而促成了他的返回。我将稍后陈述所有这些事是如何发生的。\par

此时，主教们之间又发生了一场吵闹的争论，关于\Quote{\OriginalTerm{同质}{consubstantial}}一词的确切含义。一些人认为这个词不能被接受而不亵渎；它暗示天主圣子不存在，并包含了\OriginalTerm{蒙塔努斯}{Montanus}和\OriginalTerm{撒伯里乌}{Sabellius}的错谬。另一方面，那些辩护这个词的人则认为他们的对手是希腊人（或外邦人），并认为他们的观点导致多神论。\OriginalTerm{欧瑟比}{Eusebius}，别号\OriginalTerm{庞斐禄}{Pamphilus}，和\Place{安提约基雅}主教\OriginalTerm{尤斯塔修斯}{Eustathius}在这场争论中带头。两人都承认天主圣子是\OriginalTerm{位格存在}{hypostatically}，然而他们却如同误解了彼此而争论。\OriginalTerm{尤斯塔修斯}{Eustathius}控告\OriginalTerm{欧瑟比}{Eusebius}篡改了\Place{尼西亚}大公会议所批准的道理，而\OriginalTerm{欧瑟比}{Eusebius}则宣称他赞同所有尼西亚道理，并指责\OriginalTerm{尤斯塔修斯}{Eustathius}坚持\OriginalTerm{撒伯里乌}{Sabellius}的异端。\par

% structure: p0048 source=h2 target=subsection reason=relative-heading-rank; base=h2

\subsection{第十九章　安提约基雅会议；\OriginalTerm{尤斯塔修斯}{Eustathius}被不公正地免职；\OriginalTerm{厄弗罗尼}{Euphronius}接掌主教座；君士坦丁大帝致函会议和\OriginalTerm{欧瑟比·庞斐禄}{Eusebius Pamphilus}，后者拒绝\Place{安提约基雅}主教职。}

在\Place{安提约基雅}召开了一次会议，\OriginalTerm{尤斯塔修斯}{Eustathius}被剥夺了该城教会的职务。最普遍的看法是，他被免职仅仅因为他坚持\Place{尼西亚}大公会议的信德，并因为他控告\OriginalTerm{欧瑟比}{Eusebius}、\Place{提洛}主教\OriginalTerm{保利诺}{Paulinus}和\Place{西托波利}主教\OriginalTerm{帕特罗斐禄}{Patrophilus}（东方司铎赞同他们的观点）偏爱\OriginalTerm{亚略}{Arius}的异端。然而，他被免职的借口是他用不圣洁的行为玷污了司祭职。他的免职在\Place{安提约基雅}引起了如此大的骚动，以至于人民几乎要拿起武器，全城陷入混乱。这极大地损害了他在皇帝眼中的形象；因为当皇帝了解到所发生的事，并得知该教会的人分为两派时，他非常愤怒，并怀疑\OriginalTerm{尤斯塔修斯}{Eustathius}是骚乱的制造者。然而，皇帝派了他宫中一位杰出的官员，授予全权，在不诉诸暴力或伤害的情况下安抚民众，结束动乱。\par

那些免除了\OriginalTerm{尤斯塔修斯}{Eustathius}职务并因此聚集在\Place{安提约基雅}的人，以为如果他们能成功地将一个与他们意见一致、为皇帝所知、并以学识和口才著称的人安放在\Place{安提约基雅}教会之上，他们的观点就会被普遍接受，并能获得其余人的服从，于是他们将目光投向了\OriginalTerm{欧瑟比·庞斐禄}{Eusebius Pamphilus}担任该主教座。他们就此事写信给皇帝，称这一安排将极受欢迎。事实上，所有敌视\OriginalTerm{尤斯塔修斯}{Eustathius}的圣职人员和教友都曾寻找他。然而，\OriginalTerm{欧瑟比}{Eusebius}写信给皇帝，拒绝这一尊位。皇帝赞许地同意了这一拒绝；因为有一条教会法律禁止将主教从一个教区调任到另一个教区。他写信给人民和\OriginalTerm{欧瑟比}{Eusebius}，赞同他的判断，并称他有福，因为他堪当不仅为一座城，而是为整个世界担任主教。皇帝还写信给\Place{安提约基雅}教会的人民，关于同心合意，告诉他们不应渴望其他地区的主教，正如他们不应贪图他人的财产。此外，他还向私下开会的会议发出了另一封信，同样赞扬了\OriginalTerm{欧瑟比}{Eusebius}拒绝主教职，如同他在写给\OriginalTerm{欧瑟比}{Eusebius}的信中所做的那样；并且由于确信\Place{卡帕多细雅}的长老\OriginalTerm{厄弗罗尼}{Euphronius}和\OriginalTerm{亚勒图撒的乔治}{George of Arethusa}是信理上得到认可的人，他命令主教们从中选出一位，或选择任何一位似乎配得上这一荣誉的人，并为\Place{安提约基雅}教会祝圣一位首领。收到皇帝这些信后，\OriginalTerm{厄弗罗尼}{Euphronius}被祝圣了；我听说\OriginalTerm{尤斯塔修斯}{Eustathius}平静地承受了这不公正的诽谤，认为这样更好，因为他除了美德和优秀品质外，还因其优美的口才而受到公正的钦佩，正如他流传下来的作品所证明的那样，这些作品在用词的选择、表达的韵味、情感的温和、叙述的优雅和魅力方面都受到高度评价。\par

% structure: p0051 source=h2 target=subsection reason=relative-heading-rank; base=h2

\subsection{第二十章　关于继\OriginalTerm{玛加略}{Macarius}担任\Place{耶路撒冷}主教座的\OriginalTerm{马克西模}{Maximus}。}

约在这个时期，继\OriginalTerm{西尔维斯特}{Silvester}之后担任主教、在位时间很短的\OriginalTerm{马尔谷}{Mark}去世了，\OriginalTerm{儒略}{Julius}被提升到\Place{罗马}的主教座。\OriginalTerm{马克西模}{Maximus}继\OriginalTerm{玛加略}{Macarius}之后担任\Place{耶路撒冷}主教。据说\OriginalTerm{玛加略}{Macarius}曾祝圣他为\Place{达斯颇里}教会的主教，但\Place{耶路撒冷}教会的成员坚持要他留在他们中间。因为由于他是一位\OriginalTerm{宣信者}{confessor}，并且在其他方面都很优秀，在人民的心目中，他早已被秘密选为\OriginalTerm{玛加略}{Macarius}死后的主教候选人。为了避免触怒人民并引发骚乱，他们为\Place{达斯颇里}另选了一位主教，而\OriginalTerm{马克西模}{Maximus}则留在\Place{耶路撒冷}，与\OriginalTerm{玛加略}{Macarius}共同行使司祭职务；在\OriginalTerm{玛加略}{Macarius}去世后，他治理了该教会。然而，那些确切了解这些情况的人都知道，\OriginalTerm{玛加略}{Macarius}同意人民想要留任\OriginalTerm{马克西模}{Maximus}的愿望；因为据说他对祝圣\OriginalTerm{马克西模}{Maximus}感到后悔，认为他理应被保留为自己的继承人，因为他持有关于天主的正确观点以及他的宣认，这使他深受人民爱戴。他也担心，在他死后，\OriginalTerm{欧瑟比}{Eusebius}和\OriginalTerm{帕特罗斐禄}{Patrophilus}（他们接受了亚略主义）的追随者会利用这个机会将他们中一个持同样观点的人安放在他的主教座上；因为即使\OriginalTerm{玛加略}{Macarius}在世时，他们也试图引入一些革新，但由于他们将要与他分离，他们为此保持沉默。\par

% structure: p0053 source=h2 target=subsection reason=relative-heading-rank; base=h2

\subsection{第二十一章　\OriginalTerm{梅利提安派}{Melitians}与\OriginalTerm{亚略派}{Arians}意见一致；\OriginalTerm{欧瑟比}{Eusebius}和\OriginalTerm{德奥格尼}{Theognis}试图重新点燃\OriginalTerm{亚略}{Arius}的病症。}

与此同时，最初在埃及人中掀起的争端仍无法平息。亚略异端已被尼西亚大公会议明确谴责，而\Person{梅利提乌斯}的追随者则根据上述规定被接纳共融。当\Person{亚历山大}返回\Place{埃及}，\Person{梅利提乌斯}将那些他非法篡夺治理权的教会交还给他，然后回到了\Place{吕库斯}。不久之后，他感到自己大限将至，违背尼西亚大公会议的法令，指定他最亲密的朋友之一\Person{若望}为继承人，从而在教会中产生了新的不和起因。当亚略派察觉到梅利提乌斯派在引入革新时，他们也骚扰教会。因为，正如在类似骚乱中常见的那样，有人赞赏\Person{亚略}的观点，而另一些人则主张应由\Person{梅利提乌斯}所祝圣的人治理教会。这两派异端此前一直相互敌对，但察觉到天主教会的司铎受到众人追随，他们便心生嫉妒，结盟联合，对\Place{亚历山大城}的圣职人员表现出共同的敌意。他们的攻防策略长期协同进行，以至于久而久之，在\Place{埃及}梅利提乌斯派普遍被称为亚略派，尽管他们仅在教会领导权问题上持不同意见，而亚略派则在关于天主的观点上与\Person{亚略}相同。尽管他们各自否认对方的信条，却为了在敌对的共融中达成隐秘的一致而伪装，违背自己的观点；同时每个人都期望能轻易地在自己所愿之事上取胜。然而，从这一时期起，梅利提乌斯派在讨论了那些议题之后，接受了\Person{亚略}的教义，并在关于天主的问题上持有与\Person{亚略}相同的观点。这重新点燃了关于\Person{亚略}的最初争论，部分平信徒和圣职人员脱离了与其他人的共融。关于\Person{亚略}教义的争论在其他城市再次燃起，尤其是在\Place{比提尼亚}、\Place{赫勒斯滂}和\Place{君士坦丁堡}城。简而言之，据说\Place{尼苛米底亚}的主教\Person{尤塞比乌斯}和\Place{尼西亚}的主教\Person{狄奥格尼斯}贿赂了皇帝委托保管尼西亚大公会议文件的书记官，抹去了他们的签名，并公开试图教导说圣子不应被视为\OriginalTerm{与圣父同性同体}{consubstantial with the Father}。\Person{尤塞比乌斯}被指控这些违规行为并带到皇帝面前，他极其大胆地回应，展示了他衣服的一部分。他说：\Quote{如果这件长袍在我面前被撕成碎片，我不能肯定那些碎片都是同一种质料。}皇帝对这些争端深感悲伤，因为他原以为这类问题已由尼西亚大公会议最终决断，但事与愿违，他见它们再度激荡。他尤其遗憾\Person{尤塞比乌斯}和\Person{狄奥格尼斯}接纳了某些\Place{亚历山大城}的人进入共融，尽管主教会议曾因他们持有异端意见建议他们悔改，而且他本人也曾谴责他们，将之流放离乡，视之为骚乱的煽动者。有些人断言，正是由于上述原因皇帝愤怒地将\Person{尤塞比乌斯}和\Person{狄奥格尼斯}流放；但正如我已说过的，我的信息来自那些深谙此事之人。\par

% structure: p0055 source=h2 target=subsection reason=relative-heading-rank; base=h2

\subsection{第二十二章 亚略派与梅利提乌斯派反对\Person{圣亚大纳削}的徒劳阴谋。}

降在\Person{亚大纳削}身上的种种灾祸，主要是由\Person{欧瑟伯}和\Person{德奥尼削}引起的。由于他们在皇帝面前享有极大的言论自由和影响力，便竭力争取让\Person{亚略}返回\Place{亚历山大里亚}——他们与亚略和睦友好——同时将反对他们的人驱逐出教会。他们在\Person{君士坦丁}面前控告亚大纳削，说他是所有扰乱教会的骚乱和麻烦的肇事者，并指他排斥那些渴望加入教会的人；他们声称，只要将他一人撤除，和睦就会恢复。许多与\Person{若望}一同的主教和圣职人员证实了对他的指控，他们勤勉地求得觐见皇帝的机会；他们假装非常正统，并将所有流血、捆绑、不公的打击、创伤和教堂的焚烧都归咎于亚大纳削及其教派的主教们。但当亚大纳削向皇帝证明若望追随者祝圣的非法性、他们对尼西亚大公会议法令的革新、他们信仰的不纯正，以及对那些持守天主正确见解者的侮辱时，\Person{君士坦丁}不知该相信谁。由于双方互相指控，每一方都频繁激起许多控告，而他热切渴望恢复人民的同心合意，便写信给亚大纳削，说任何人都不应被排除在外。如果此事被泄露到最后，他将不顾后果地派一个人，将亚大纳削驱逐出\Place{亚历山大里亚}城。如果有人想看到皇帝这封信，他在这里会找到有关此事的段落：\Quote{既然你现在知晓了我的旨意，即你应当向所有渴望进入教会的人提供不受阻碍的进入。因为我若听说任何愿意加入教会的人被你禁止或妨碍，我将立即派遣一位官员，按照我的命令将你撤职，并将你转移到别处。}然而，亚大纳削写信给皇帝，使他确信亚略派不应被天主教会接纳共融；\Person{欧瑟伯}看出只要亚大纳削反对，他的计划就永远无法实施，便决定不惜一切手段除掉他。但他找不到足够的借口来实现这一计划，便许诺梅利提安派，如果他们提出控告，他就会让皇帝和当权者关心他们的利益。于是，第一项指控来了：说亚大纳削对埃及人征收亚麻短衫税，并且这种贡赋是从控告者那里强征的。亚大纳削教会的司铎\Person{阿丕斯}和\Person{玛加略}当时恰好在宫中，清楚地证明了这项持续指控是虚假的。在奉命为这项罪行答辩时，亚大纳削又被指控阴谋反对皇帝，并为此目的将一箱金子送给了一个名叫\Person{非路门}的人。皇帝识破了控告者的诽谤，送亚大纳削回家，并写信给\Place{亚历山大里亚}的人民，证明他们的主教非常节制、信仰正确；他乐意接见亚大纳削，并认出他是天主的人；而且，由于嫉妒是他被指控的唯一原因，他反而比控告者更显光辉；皇帝听闻亚略派和梅利提安派宗派分子在\Place{埃及}挑起了分裂，在同一封信中，他劝勉民众要仰望天主，留意祂的审判，彼此和好，全力追捕那些阴谋破坏他们同心合意的人；皇帝就这样写信给人民，劝勉他们全体同心合意，努力防止教会内的分裂。\par

% structure: p0057 source=h2 target=subsection reason=relative-heading-rank; base=h2

\subsection{第二十三章 关于圣\Person{亚大纳削}与\Person{阿尔色尼}手部事件的诽谤}

梅利提安派初次尝试失败后，又秘密捏造了对\Person{阿塔纳修}的其他指控。一方面，他们控告他打碎了一只圣爵，另一方面，说他杀了一个名叫\Person{阿尔塞尼乌斯}的人，并砍下他的手臂用于魔法目的。据说这位\Person{阿尔塞尼乌斯}原是神职人员，但因犯下某种罪行，害怕被主教定罪和惩罚，逃到一个隐蔽处。\Person{阿塔纳修}的敌人利用此事策划了最严重的攻击。他们竭力搜寻\Person{阿尔塞尼乌斯}，找到了他；他们对他大献殷勤，承诺确保他得到一切善意和安全，并秘密将他带到隐修院司铎\Person{帕特林斯}那里，\Person{帕特林斯}是他们的同伙之一，与他们的利益一致。这样仔细藏匿他之后，他们便在市场和公共集会中大肆散布说他已被\Person{阿塔纳修}杀死的谣言。他们还贿赂一位名叫\Person{若望}的隐修士来证实这个指控。这个恶名传遍四方，甚至传到了皇帝的耳中，\Person{阿塔纳修}担心在受这些假谣言影响而偏见的法官面前难以辩护，便采取了与对手类似的手段。他尽一切努力防止真理被他们的攻击所遮蔽；但因\Person{阿尔塞尼乌斯}没有出现，民众无法被说服。因此，他考虑到除非能证明据说已死的\Person{阿尔塞尼乌斯}还活着，否则对他的怀疑无法消除，于是派一位最值得信任的执事去寻找他。执事前往\Place{特贝}，从一些隐修士的声明中查明了\Person{阿尔塞尼乌斯}的住处。当他来到\Person{帕特林斯}那里——\Person{阿尔塞尼乌斯}曾被藏匿在\Person{帕特林斯}处——却发现\Person{阿尔塞尼乌斯}不在那里；因为一听说执事到来，他就被转移到了\Place{下埃及}。执事逮捕了\Person{帕特林斯}，并把他带到\Place{亚历山大}，还有他的同伙之一\Person{厄里亚}，据说\Person{厄里亚}是把\Person{阿尔塞尼乌斯}转移到别处的人。执事将两人交给埃及军队统帅，他们供认\Person{阿尔塞尼乌斯}还活着，曾秘密藏在他们家中，现在住在\Place{埃及}。\Person{阿塔纳修}确保所有这些事实都被报告给\Person{君士坦丁}。皇帝回信给他，希望他专心履行司铎职务，在人民中维持秩序和虔敬，不要因梅利提安派的阴谋而不安，显然嫉妒是流传对他的虚假指控和教会骚乱的唯一原因。皇帝还说，今后他不再相信这些报告；除非诽谤者保持平安，否则他一定会以国法的严厉对待他们，让正义得以伸张，因为他们不仅不正义地密谋反对无辜者，还可耻地滥用了教会的良善秩序与虔敬。这就是皇帝给\Person{阿塔纳修}的信的主旨；他还命令在公众面前高声宣读，以使所有人都知晓他的意图。梅利提安派受到这些威胁而惊恐，暂时安静了一段时间，因为他们焦虑地看待统治者的威胁。\Place{埃及}各地的教会享有深切的平安，在这位伟大司铎的主持下，因众多异教徒和其他异端者的皈依而日益增长。\par

% structure: p0059 source=h2 target=subsection reason=relative-heading-rank; base=h2

\subsection{第24章 有些印度民族当时通过两个俘虏——\Person{福鲁门提乌斯}和\Person{埃德修斯}——接受了基督信仰。}

我们听说，大约在此时，一些最遥远的民族，即我们称之为印度人的，他们未曾听闻\Person{巴尔多禄茂}的宣讲，却通过\Person{弗鲁孟提}分享了我们的教义；\Person{弗鲁孟提}成了他们当中神圣学问的司铎和教师。但为了使我们能通过发生在\Place{印度}的奇迹知道，基督信仰的教义应被接受为一个并非来自人的体系（尽管它对于一些人而言像是一系列奇迹），有必要叙述\Person{弗鲁孟提}被祝圣的理由。情况如下：\Place{希腊}最著名的哲学家们曾探索未知的城市和地区。\Person{苏格拉底}的朋友\Person{柏拉图}曾在\Place{埃及}居留一段时间，以熟悉他们的风俗习惯。同样地，他航行到\Place{西西里}去观赏那里的火山口，从那里如同从泉源中自发地涌出火流，常常泛滥，像河流一样冲毁邻近地区，以致至今许多田地仍显烧焦，不能播种或栽种树木，就如他们所叙述的\Place{索多玛}之地一样。这些火山口也曾被\Person{恩培多克勒}探索，他在\Place{希腊}人中以哲学闻名，并用英雄体诗歌阐述了他的知识。他出发去调查这次火喷发，要么是因为他认为这样的死亡方式优于其他任何方式，要么是因为（更真实地说）他也许不知道为何要以这种方式结束生命，他跳进火中而丧命。\Place{科斯}的\Person{德谟克利特}探索了许多城市、气候和民族，他关于自己说，他一生中有八十年是在异国旅行中度过的。除了这些哲学家，成千上万的\Place{希腊}智者，无论古代还是现代的，都致力于这种旅行。与之效仿，\Place{腓尼基}的\Place{提洛}的哲学家\Person{梅洛皮乌}远行至\Place{印度}。据说他带着两个青年，名叫\Person{弗鲁孟提}和\Person{厄德西乌}；他们是他的亲属；他指导他们的修辞训练，并以自由教育培养他们。在尽可能探索了\Place{印度}之后，他启程回家，并登上了一艘即将驶往\Place{埃及}的船。碰巧由于缺水或其他必需品，船只不得不在某个港口停靠，\Place{印度}人冲上船，杀死了所有的人，包括\Person{梅洛皮乌}。这些\Place{印度}人刚刚解除了与罗马人的同盟；他们因怜悯这两个男孩年轻，将他们作为活俘虏带走，并带到了他们的国王面前。国王任命较年轻的一个做他的司酒者；较年长的\Person{弗鲁孟提}，他任命他管理他的家，并让他做他的财宝管理者；因为他看出他聪明且非常能干。这些青年在多年间有用且忠实地服侍国王，当国王感到他的末日临近时，他的儿子和妻子还在世，他赏赐了仆人的良好意愿以自由，并允许他们去他们想去的地方。他们急于回到\Place{提洛}，他们的亲属居住在那里；但国王的儿子年幼，他的母亲恳求他们再留一段时间，管理公共事务，直到她的儿子达到成年。他们听从了她的恳求，并指导了\Place{印度}王国和政府的政务。\Person{弗鲁孟提}出于某种神圣的冲动（也许是因为天主自发地感动了他），询问\Place{印度}是否有任何基督徒，或者罗马商人中是否有航船到那里的。成功地找到了他询问的对象后，他召唤他们到面前，以爱与友善对待他们，召集他们祈祷，集会按照罗马的惯例进行；当他建造了祈祷之所后，他鼓励他们不断尊崇天主。\par

当国王的儿子达到成年时，\Person{弗鲁孟提}和\Person{厄德西乌}恳求他和王后，并不容易地说服统治者让他们离开，于是作为朋友告别后，他们作为罗马属民返回。\Person{厄德西乌}去了\Place{提洛}看望亲属，不久以后被晋升为司铎。\Person{弗鲁孟提}却没有返回\Place{腓尼基}，而是去了\Place{亚历山大}；因为对他而言，爱国和孝道都从属于宗教热忱。他与\Place{亚历山大}教会的首领\Person{亚大纳削}商谈，向他描述了\Place{印度}的情况，以及为居住在该国的基督徒任命一位主教的必要性。\Person{亚大纳削}召集了本地司铎，与他们商议此事；他祝圣\Person{弗鲁孟提}为\Place{印度}主教，因为他特别适合并且擅长在那群人中间做大量的服务，他是第一个在他们中间显扬基督徒之名的人，并且播下了参与教义的种子。因此，\Person{弗鲁孟提}回到\Place{印度}，据传说，他如此卓越地履行了司铎职能，以致成为普遍钦佩的对象，并被尊崇宛若一位宗徒。天主极大地尊荣他，使他能行许多奇妙的医治，并施行征兆和奇迹。这就是印度司铎制度的起源。\par

% structure: p0062 source=h2 target=subsection reason=relative-heading-rank; base=h2

\subsection{第二十五章 提洛会议；圣亚大纳削被非法罢黜。}

\Person{阿塔纳修斯}的仇敌所设的阴谋，又使他陷入新的麻烦，激起了皇帝对他的仇恨，并招来了大批控告者。皇帝被他们的纠缠所厌烦，便在\Place{巴勒斯坦}的\Place{凯撒勒雅}召开了一次教务会议。\Person{阿塔纳修斯}被召前往；但他惧怕该城主教\Person{优西比乌}、\Place{尼科米底亚}主教\Person{优西比乌}及其同党的诡计，拒绝出席，并坚持拒绝达三十个月之久，尽管不断被催促出席。然而在那段时期结束时，他受到更紧迫的催逼，便前往\Place{提洛}，那里聚集了许多东方主教，他们命令他接受控告者对他的指控。\Person{若望}一派的主教\Person{卡利尼科斯}和一位名叫\Person{伊苏里亚斯}的人，控告他打碎了一只神秘爵杯，推倒了一把主教座；并多次使\Person{伊苏里亚斯}——尽管他是一位长老——被锁上镣铐；又假意向\Place{埃及}总督\Person{希吉努斯}告发，说他曾向皇帝雕像投掷石块，致使\Person{伊苏里亚斯}被投入监狱；又控告他罢免了\Place{培琉喜阿姆}天主教会的主教\Person{卡利尼科斯}，并说除非\Person{卡利尼科斯}能消除有关他打碎神秘爵杯的某些嫌疑，否则将禁止他与教会共融；又控告他将\Place{培琉喜阿姆}教会交给已被罢免的长老\Person{马尔谷}管理；并将\Person{卡利尼科斯}置于军兵看守之下，并对他施以司法酷刑。\Person{若望}一派的主教\Person{厄乌普略}、\Person{帕科弥乌}、\Person{依撒格}、\Person{阿基拉斯}和\Person{赫尔梅翁}，控告他施加殴打。他们一致坚持说，他是借某些人的伪证而获得主教尊位的，因为曾有法令规定，凡无法洗清所受指控之罪的人，不得领受圣秩。他们更进一步指控说，他们被他欺骗后，便与他断绝了共融关系，而他不但没有消除他们的疑虑，反而暴力对待他们，将他们投入监狱。\newline{}\par

此外，\Person{阿尔塞尼乌}的事件又被重新提起；正如在这种精心策划的阴谋中常发生的那样，甚至许多原本被认为是他朋友的人也出人意料地以控告者的身份出现。接着宣读了一份文件，其中载有民众的怨言，说\Place{亚历山大里亚}的人民因他的缘故无法继续前往教堂集会。\Person{阿塔纳修斯}被催促为自己辩护，他多次出庭；成功地反驳了其中一些指控，并请求对其他指控给予调查的宽限。他极为困惑，因为他思及控告者得到了审判官的宠爱，又思及\Person{阿里乌}派和\Person{梅利提乌斯}派中有许多证人出庭反对他，还思及那些指控虽已被推翻，告密者却仍受到纵容。尤其是在关于\Person{阿尔塞尼乌}的起诉中，他被指控砍下\Person{阿尔塞尼乌}的手臂用于巫术；以及在关于某位女子的起诉中，他被指控送给她不洁的礼物，并在夜间强奸了这名不愿从命的女子。这两项起诉后来都被证明是荒谬的，充满了虚假的刺探。当那女子在主教面前作证时，站在\Person{阿塔纳修斯}身旁的\Place{亚历山大里亚}长老\Person{弟茂德}，按照他秘密商定的计划走近她，对她说：\Quote{那么，妇人，难道是我玷污了你的贞洁吗？}她回答说：\Quote{难道不是你吗？}并提到了她遭受强迫的地点和相关情况。他也把\Person{阿尔塞尼乌}领到他们中间，向审判官展示了他的双手，并要求他们让控告者解释他们曾展示的那只手臂。因为事情发生：\Person{阿尔塞尼乌}，或是受了神力的驱使，或是如人们所说，被\Person{阿塔纳修斯}的计划隐藏起来，当他听说那主教因他而面临危险时，便在审判的前一天夜间逃走，并抵达了\Place{提洛}。但这些指控如此被迅速驳回，以致无需辩护，在案卷中也没有提及第一项指控；我认为最可能的原因是，整件事被认为太不体面、太荒谬，故而不值得记录。至于第二项指控，控告者试图为自己辩护说，\Person{阿塔纳修斯}属下的一位名叫\Person{普路西安}的主教，曾奉其上级之命，烧毁了\Person{阿尔塞尼乌}的房屋，把他绑在柱子上，用皮鞭抽打他，然后把他锁在一间小屋里。他们还进一步说，\Person{阿尔塞尼乌}从窗户逃出了小屋，当人们搜寻他时，他暂时躲藏起来；由于他没有露面，他们自然以为他死了；他作为一个人和宣信者所获得的声誉，使他受到\Person{若望}一派主教的喜爱；于是他们寻找他，并为他向地方长官申诉。\newline{}\par

\Person{亚大纳削}一想到这些事，便充满忧虑，开始怀疑敌人正暗中策划置他于死地。经过数次会议，当主教会议乱作一团，控诉者与法庭周围众多人群高声喊叫，说亚大纳削应该被罢黜，因为他是个巫术师和暴徒，根本不配司铎之职时，被皇帝指派维护会议秩序的官员，被迫秘密地让被告离开审判厅；因为他们害怕在骚乱中他们会成为杀他的人，这在混乱中常有发生。发现自己在\Place{提洛}已无法安全停留，且面对众多敌视他的审判官，毫无希望获得公正对待，他便逃往\Place{君士坦丁堡}。主教会议在他缺席期间谴责了他，罢黜了他的主教职，并禁止他居住在\Place{亚历山大里亚}，说恐怕他会引起骚乱和叛乱。\Person{若望}及其所有支持者恢复了共融，仿佛他们曾不公正地受屈，每人恢复了原有神职品级。随后主教们向皇帝报告了他们的所作所为，并致信各地主教，吩咐他们不可与亚大纳削共融，不可写信给他或收他的信，因为他已因他们所调查的罪行而被定罪，且因其逃跑，在那些未审理的指控中也有罪。他们还在信中声明，他们不得不对他做出如此判决，因为前一年皇帝命令他去\Place{凯撒勒雅}出席东方主教会议时，他违命不从，让主教们久等，并藐视统治者的命令。他们还指出，当主教们在\Place{提洛}集会时，他带着大批随从到该城，目的是在会议中引起骚乱和混乱；他在那里时而拒绝回答对他的指控，时而单独辱骂主教们；当被传唤时，有时不服从，有时不屑于受审。他们在同一封信中详细说明，他明显有打碎神秘爵杯的罪行，这一事实由\Place{尼凯亚}主教\Person{泰奥格尼斯}、\Place{卡尔西顿}主教\Person{马里斯}、\Place{赫拉克利亚}主教\Person{德奥多罗}、\Person{瓦伦提诺}和\Person{乌尔萨基乌斯}，以及\Person{马其顿纽斯}所证实，后者曾被派往\Place{埃及}的村庄，据称爵杯即在彼处打破，以查明真相。主教们就这样依次详述了对亚大纳削的每一项指控，其手法正如诡辩家们为了加强诽谤效果所采用的。然而，许多当时在场的司铎都看出指控的不公。据传，宣信者\Person{帕夫努提乌斯}当时也在会议中，他起身拉住\Place{耶路撒冷}主教\Person{马克西穆斯}的手，要领他离开，仿佛这些为敬虔而眼睛被挖出的宣信者，不应参与恶人的集会。\par

% structure: p0066 source=h2 target=subsection reason=relative-heading-rank; base=h2

\subsection{第二十六章 \Person{君士坦丁大帝}在\Place{耶路撒冷}\Place{哥耳哥达}建造圣殿及其奉献礼。}

被称为\Quote{大殉道堂}的圣殿，建在\Place{耶路撒冷}髑髅地，约在\Person{君士坦丁}在位第三十年完成；一位名为\Person{马里亚诺斯}的官员，是皇帝的速记员，来到\Place{提洛}，向会议递交了皇帝的信函，命令他们速往\Place{耶路撒冷}，以祝圣该圣殿。虽然此事先前已决定，但皇帝认为有必要首先调停在\Place{提洛}聚集的主教们之间的争议，清除一切不和与悲伤，然后才去祝圣圣殿。因为对于这样的节日，司铎们应当同心合意。当主教们到达\Place{耶路撒冷}时，圣殿遂被祝圣，同样还有皇帝送来的众多装饰品和礼物，至今仍保存在圣殿中；其奢华与壮丽令人一见便惊叹不已。自那时起，\Place{耶路撒冷}教会隆重举行祝圣周年纪念；节庆持续八天，并施行圣洗圣事；天下各地的人均在此期间涌向\Place{耶路撒冷}，朝拜圣地。\par

% structure: p0068 source=h2 target=subsection reason=relative-heading-rank; base=h2

\subsection{第二十七章 关于那位劝谏\Person{君士坦丁}召回\Person{阿黎乌}和\Person{欧佐伊乌斯}脱离流放的司铎；关于其可能为虔诚信仰的论著，以及\Person{阿黎乌}如何再次被在\Place{耶路撒冷}聚集的主教会议所接纳。}

那些接受\Person{阿黎乌}观点的主教们找到了一个有利时机，藉着热心争取在\Place{耶路撒冷}城召开一次会议，恢复了\Person{阿黎乌}和\Person{欧佐伊乌斯}的共融。他们以下列方式实现了他们的计划：——\par

某位司铎非常倾慕\Person{阿黎乌}的学说，与皇帝的妹妹关系密切。起初他隐藏自己的观点；但他经常拜访，并逐渐与\Person{君士坦丁}的妹妹——名为\Person{君士坦提亚}——更加熟悉，便鼓起勇气向她陈述，\Person{阿黎乌}是因为\Place{亚历山大里亚}教会主教\Person{亚历山大}的嫉妒和私仇而被不公正地流放，并被逐出教会。他说\Person{亚历山大}的嫉妒是因民众对\Person{阿黎乌}的尊敬所激起的。\par

\Person{君士坦提娅}相信这些陈述是真的，但在她有生之年并未采取任何步骤更改\Place{尼西亚}的法令。她患上一种可能致死的疾病，便恳求前来看望她的皇帝兄弟，作为最后的恩赐，答应她将要请求的事：即接纳上述长老为亲密朋友，并信赖他，视其为对天主性持有正确见解的人。\newline{} 她说：\Quote{就我而言，我正濒临死亡，不再关心此生的事务；现在我唯一的恐惧，是担心你招致天主的义怒，遭受任何灾祸或失去帝国，因为你被诱导错误地把正义良善的人判以永久流放。} 从那时起，皇帝开始宠信这位长老，允许他自由地与他交谈，并就他姊妹所嘱咐的同样主题进行讨论，并认为有必要重新审查\Person{亚略}的案件；皇帝作出这一决定，可能是受了那些攻击可信性的影响，或是为了满足他姊妹的愿望。不久，他便从流放地召回\Person{亚略}，并要求他提交一份关于天主性信仰的书面声明。\Person{亚略}避免使用他先前发明的新术语，而用简单且为圣经所承认的措辞写成了另一份声明；他宣誓说，他持守这份声明中所阐述的教义，他 \OriginalTerm{由衷}{ex animo} 感到这些陈述，并且除此之外别无其他想法。声明如下：\Person{亚略}和\Person{欧佐伊}，长老，致我们最虔诚、最蒙天主宠爱的\Person{君士坦丁}皇帝。\par

遵照您所钟爱的虔敬所命令的，君主皇帝，我们在此提交我们自己信仰的书面声明，并在天主面前声明，我们以及所有与我们同在的人，都相信此处所陈述的。\par

我们相信唯一的天主，全能的\Person{圣父}，并相信祂的圣子主\Person{耶稣基督}，祂在万世之前由祂所生，是天主圣言，万有藉祂而造，无论是天上的还是地上的；祂降来取了肉身，受难并复活，升了天，从那里祂将再次来审判生者与死者。\par

\Quote{我们相信\Person{圣神}，相信肉身的复活，相信来世的生命，相信天国，并相信一个遍及普世的天主的大公教会。我们是从福音领受这信仰的，主在其中对祂的门徒说：\BibleQuote{去教导万民，因\Person{圣父}及\Person{圣子}及\Person{圣神}之名给他们授洗。} 如果我们不如此相信，如果我们不真正接受关于\Person{圣父}、\Person{圣子}和\Person{圣神}的教义，正如整个大公教会和圣经所教导的，并且我们在每一点上如此相信，愿天主在现今和将来的日子作我们的审判者。因此，我们恳求您的虔敬，我们最蒙天主宠爱的皇帝，既然我们被列于神职人员的行列，并持守教会和圣经的信仰与思想，就通过您的缔造和平与虔诚的虔敬，使我们公开地与我们的母亲教会和好；好使无益的问题和争论被抛开，我们与教会和平共处，我们全体共同为您和平与虔诚的帝国以及您的全家献上惯常的祈祷。}\par

许多人认为这一信仰声明是狡猾的编造，表面措辞不同，实则支持\Person{亚略}的教义；其中表述模糊，可作多种解释。皇帝以为\Person{亚略}和\Person{欧佐伊}与\Place{尼西亚}会议的主教们持相同意见，对此事感到欣喜。然而，他并未试图在没有那些按教会法律在这些事务上是权威之人的判断和认可的情况下，恢复他们的共融。因此，他将他们送到当时聚集在\Place{耶路撒冷}的主教们那里，并写信要求他们审查\Person{亚略}和\Person{欧佐伊}提交的信仰声明，并影响会议，以便无论他们发现其教义正统，且其敌人的嫉妒是定罪的唯一原因，还是他们没有理由责备那些定罪者而改变了心意，都能给予他们仁爱的决定。那些长久以来热心于此的人，借着皇帝的信件为掩护，抓住了机会，接纳他进入共融。他们立即写信给皇帝本人、\Place{亚历山大里亚}教会，以及埃及、\Place{忒拜}和利比亚的主教和神职人员，恳切劝勉他们接纳\Person{亚略}和\Person{欧佐伊}进入共融，因为皇帝在他的一封书信中证明了他们信仰的正确，并且皇帝的决定已得到会议的投票确认。\par

这些就是\Place{耶路撒冷}会议热切讨论的主题。\par

% structure: p0077 source=h2 target=subsection reason=relative-heading-rank; base=h2

\subsection{第二十八章 \Person{君士坦丁}皇帝致\Place{提洛}主教会议的书信，以及圣\Person{亚大纳修}因\Person{亚略}派的阴谋而被放逐。}

\Person{亚大纳修}从\Place{提洛}逃出后，前往\Place{君士坦丁堡}，来到\Person{君士坦丁}皇帝面前，当着那些定他罪的主教们的面，控告他所遭受的，并恳求皇帝准许将\Place{提洛}会议的决议提交皇帝审查。\Person{君士坦丁}认为这一请求合理，便以下列措辞写信给聚集在\Place{提洛}的主教们：——\par

\Quote{朕不知道你们的会议在混乱和激烈中通过什么法令；但似乎由于某种扰乱性的无序，通过了不符合真理的法规，你们不断彼此激怒显然阻止了你们考虑什么是中悦天主的。但这是天主眷顾的工作，要驱散从这场争端中产生的恶事，并向朕清楚地表明，你们的判断是否因私人好恶而误导。因此朕命令你们全都毫不迟延地来到朕的虔敬面前，以便朕能够准确了解你们的事务。朕将向你们解释朕以此语气写信给你们的缘由，你们将从下文知道朕为何通过此文件将你们召到朕面前。当朕骑马返回那座以朕之名命名的城市（朕视之为朕极为繁荣的国家）时，\Person{阿塔纳修斯}主教与几个随从出乎意料地出现在大路中间，朕见到他感到极为惊讶。洞察万有的天主为朕作证，起初朕不知道他是谁，但朕的一些侍从查明此事以及他所遭受的不公，才告诉了朕必要的信息。朕当时没有接见他。但他坚持请求面谈；虽然朕拒绝了他，并几乎命令将他从朕面前带走，他更大胆地告诉朕，他请求朕的恩惠只是让朕将你们召来，以便他能在你们面前抱怨他所遭受的不必要的苦难。既然这个请求似乎合理且及时，朕认为以这种语气写信给你们是恰当的，并命令所有参加\Place{提尔}会议的人尽快赶到朕的仁爱朝廷，以便你们能以行动向朕证明你们决定的纯洁与坚定，而朕你们不能否认是天主真诚的仆人。藉着朕事奉天主的热情，和平已普世建立，天主的名在迄今不知真理的蛮族中受真诚赞美；并且明显的是，不知道真理的人就不认识天主。尽管如此，如上所述，蛮族藉着朕的媒介已学会真诚认识并敬拜天主；因为他们看见天主的保护处处时时临于朕身；他们因敬畏朕的权能而更深地敬畏天主。但朕必须宣布忍耐的奥迹（因朕不愿说朕遵守它们），朕，朕说，不应做任何可能导致纷争或仇恨，或者直截了当地说，导致人类毁灭的事。那么，如朕所说，尽快来到朕这里，并确信朕将尽一切努力以如此方式保存天主法律中所有特别无误的部分，使任何过失或异端都无法捏造；而那些以圣名伪装、试图引入变异和亵渎的法律仇敌，已被公开驱散、彻底粉碎、完全压制。}\par

皇帝的这封信激起了某些主教们的恐惧，以至于他们启程回家。但\Place{尼科美底亚}主教\Person{优西比乌}和他的同党来到皇帝面前，陈述说\Place{提尔}会议没有制定任何不基于正义反对\Person{阿塔纳修斯}的法令。他们提出\Person{特奥格尼斯}、\Person{马里斯}、\Person{特奥多尔}、\Person{瓦伦斯}和\Person{乌尔萨基乌斯}作为证人，作证说他打碎了神秘之杯，并在说了许多其他诽谤之后，他们的指控占了上风。皇帝要么相信他们的陈述是真实的，要么认为如果\Person{阿塔纳修斯}被撤职，主教们之间就会恢复一致，于是将他流放到\Place{高卢}西部的\Place{特雷维斯}；因此他被押送到那里。\par

% structure: p0081 source=h2 target=subsection reason=relative-heading-rank; base=h2

\subsection{第二十九章 \Person{亚历山大}，\Place{君士坦丁堡}主教；他拒绝接纳\Person{阿里乌}进入共融；\Person{阿里乌}在寻求自然排泄时爆裂。}

耶路撒冷会议之后，亚略去了埃及，但由于未能获准与亚历山大教会共融，便返回了君士坦丁堡。所有赞同他观点的人，以及依附尼苛默底亚主教优西比乌的人，都狡猾地聚集在该城，意图召开一次会议。当时负责管理君士坦丁堡教区的亚历山大尽力阻止这次会议召开。但他的努力受阻，于是公开拒绝与亚略订立任何盟约，声称推翻他们自己以及从普天下几乎所有地区聚集在尼西亚的那些主教们的投票，既不正义，也不符合教会法规。当优西比乌的党羽发现他们的论据对亚历山大夫生效力时，便转而使用侮辱，并威胁说，除非他在指定日期接纳亚略共融，否则他将被逐出教会，另选一位愿意与亚略共融的人接替他。然后他们分手了：优西比乌的党羽等待他们实施威胁的指定时间，而亚历山大则祈祷，愿优西比乌的话不至于成为现实。他主要的恐惧源于皇帝已被说服让步。在指定日期的前一天，他匍匐在祭台前，整夜向天主祈祷，愿他的敌人无法实施反对他的阴谋。下午较晚的时候，亚略突然感到胃痛，被迫前往为此类紧急情况设立的公共处所。他进去一段时间未出来，一些在外等候的人进去，发现他已死去，仍坐在马桶上。他的死讯传开后，众人对此事的看法不一。有些人相信他是在那一刻死于突发心脏病，或因事情如其所愿而感到兴奋过度所致；另一些人则认为这种死法是因他的不敬而受的惩罚。持他观点的人认为，他的死是邪术造成的。在此不妨引述亚历山大主教亚大纳削对此事的陈述。他的叙述如下：\par

% structure: p0083 source=h2 target=subsection reason=relative-heading-rank; base=h2

\subsection{第三十章 伟大的亚大纳削关于亚略之死的记述}

\Quote{亚略，异端的创始人，优西比乌的同伙，应优西比乌党羽的请求，被带到最可敬的君士坦丁奥古斯都面前，奉命以书面形式阐述他的信仰。他极其狡猾地起草了这份文件，像魔鬼一样，用圣经的简单话语掩盖他不敬的论断。最可敬的君士坦丁对他说：\InnerQuote{如果你心中没有其他想法，就为真理作证；因为如果你发假誓，主必惩罚你。}这可怜的人发誓说，他除了文件中现在列出的观点外，从未持有或设想任何其他观点，即使他曾经断言过别的；此后不久他出去，审判便临到他：他向前弯腰，肚腹崩裂。所有人的共同结局是死亡。我们不应仅仅因为一个人死了就责备他，即使他是敌人，因为我们不知道我们是否能活到晚上。但亚略的结局如此奇特，似乎值得一提。优西比乌的党羽威胁要让他重返教会，君士坦丁堡主教亚历山大反对他们的意图；亚略信赖优西比乌的权力和威胁；因为那天是安息日，他期望第二天被重新接纳。争论激烈；优西比乌的党羽大声威胁，而亚历山大则祈祷。主是审判者，并宣布自己反对不义之人。日落前不久，亚略因自然的迫切需求，被迫进入为此类紧急情况设立的处所，在这里他同时失去了共融恢复和生命。最可敬的君士坦丁听到此事时大为震惊，视之为发假誓的证明。于是每个人都知道，优西比乌的威胁绝对徒劳，亚略的期望也落空了。同时也显明，亚略的疯狂无论在现世还是在首生者的教会中，都不能与救主共融。那么，岂不令人惊奇吗？居然还有人试图为被主定罪的人开脱，并捍卫那个主证明为不配共融的异端，因为主不允许其创始人进入教会。我们得到可靠消息，亚略的死法就是这样。据说此后很长一段时间，无人使用他死时坐的那个马桶。那些因自然需要（如同在人群中常有的情况）而不得不去公共处所的人，进去时互相提醒避开那个马桶，该处后来被回避，因为亚略在那里受到了他不敬的惩罚。后来，一个信奉亚略异端的有钱有势的人，从公众那里买下了那块地方，并在该处建了一所房子，以使此事被遗忘，不再有亚略之死的永久纪念。}\par

% structure: p0085 source=h2 target=subsection reason=relative-heading-rank; base=h2

\subsection{第三十一章 亚略死后亚历山大事变。君士坦丁大帝致该教会信函}

亚略之死并未终止他所引发的教义争论。那些持守他主张的人并未停止阴谋反对持相反意见的人。亚历山大里亚人民大声抱怨亚大纳削被放逐，并为他归来献上恳求；著名的隐修士安当多次写信给皇帝，恳求他不要相信梅利提安派的暗示，而要拒绝他们的控告，视之为诽谤；然而皇帝并未被这些理由说服，写信给亚历山大里亚人，指责他们愚昧和混乱的行为。他命令圣职人员和圣洁的贞女保持安静，宣布他不会改变心意，也不会召回亚大纳削，他说，他视其为煽动者，被教会公正地判罪。他回复安当，说他不应忽视会议的法令；他说，即使少数主教出于恶意或取悦他人的动机，也很难相信如此众多明智而卓越的主教会被这些动机驱使；他补充说，亚大纳削傲慢自大，是分裂与煽动的原因。亚大纳削的敌人尤其指控他这些罪行，因为他们知道皇帝特别厌恶这些行为。当他听说教会分裂为两派，一派支持亚大纳削，另一派支持若望时，他愤慨至极，并将若望本人放逐。这位若望继承梅利提乌斯，并与那些持同样见解的人一起，被提洛会议重新纳入共融并恢复圣职。他的放逐违背了亚大纳削敌人的意愿，但此事仍然发生，提洛会议的法令并未使若望受益，因为皇帝对于任何涉嫌煽动基督徒子民分裂或骚乱的人，都不接受任何恳求或请求。\par

% structure: p0087 source=h2 target=subsection reason=relative-heading-rank; base=h2

\subsection{第三十二章 君士坦丁颁布法律反对一切异端，禁止民众在任何地方聚会敬拜，除非在天主教会内；因此大多数异端消失。支持尼科美底亚的欧瑟伯的亚略派人，诡诈地企图抹去\Quote{同质}一词。}

尽管许多人在辩论中热心支持亚略的教义，但尚未形成一个可以称为亚略派的独特派别；因为除了诺瓦提安派、所谓弗里吉亚派、华伦提尼派、马西昂派、保利安派以及少数持守早已发明的异端者外，所有人都作为一个教会聚集，彼此共融。然而皇帝颁布法律，要拆除他们自己的祈祷之所；他们应在教会中聚会，不得在私宅或公共场所聚会。他认为在天主教会中共融更好，并劝告他们聚集在她的围墙内。藉着这条法律，我相信几乎所有的异端都消失了。在先前的皇帝统治时期，所有敬拜基督的人，无论他们在意见上有多大差异，都受到外邦人同样的对待，并遭受同样残酷的迫害。这些共同的灾难，他们同样都要承受，阻止了他们对自己之间意见分歧进行深入探究；因此各派成员各自聚会很容易，并且通过彼此不断商议，无论人数多么稀少，他们都不会分裂。但在该法律通过后，他们既不能在公开场合聚集，因为被禁止；也不能秘密举行集会，因为他们受到所在城市的主教和圣职人员的监视。因此，这些分裂教派的大多数因害怕后果而加入天主教会。那些持守原初见解的人，在死时没有留下门徒传播其异端，因为他们既不能聚集到同一地方，也无法安全地教导同见解的人。由于异端教义的荒谬，或由于创造和教导这些异端者的完全无知，每个异端的追随者从一开始就人数很少。只有诺瓦提安派，他们拥有优秀的领袖，并且关于天主性的意见与天主教会相同，从开始就人数众多，并且一直如此，未受该法律多大损害；我相信皇帝乐意对他们放宽法令的严厉，因为他只想让臣民感到恐惧，并无意迫害他们。当时君士坦丁堡这一异端的主教阿刻西乌斯，因他的圣善生活而深受皇帝尊敬；很可能正是为了他的缘故，他所治理的教会得到了保护。弗里吉亚派在所有罗马省份（除弗里吉亚及其邻近地区外）受到与其他异端者相同的对待，因为在这里他们自蒙丹时代起就大量存在，直到今日仍然如此。\par

大约此时，尼科美底亚主教欧瑟伯和尼刻亚主教特奥格尼斯的追随者开始在著作中对尼西亚会议所颁布的信经进行革新。他们不敢公开拒绝圣子与圣父同质的说法，因为皇帝坚持这一说法；但他们提出另一份文件，并告知东方主教，他们接受尼西亚教义的字句，但加以言辞解释。由于这一声明和反思，先前的争论重新陷入新的讨论，似乎已经平息的问题再次被搅动起来。\par

% structure: p0090 source=h2 target=subsection reason=relative-heading-rank; base=h2

\subsection{第三十三章 安其拉主教马塞鲁；他的异端与废黜。}

同一时期，\Place{迦拉达}地区\Place{安奇拉}的主教\Person{玛策禄}，被在\Place{君士坦丁堡}集会的众主教免职并逐出教会，因他引进了一些新学说，教导天主圣子的存在始于祂由\Person{玛利亚}诞生之时，且祂的王国将有终结；此外，他还撰写了一份书面文件，在其中阐述了这些观点。其后，一位极富口才与学识的人\Person{巴西略}，被授予\Place{迦拉达}教区的主教职。他们也致信邻近地区的教会，要求他们寻找\Person{玛策禄}所写之书的抄本并销毁，并引领那些被发现接受其观点的人回归正途。他们说该书篇幅过长，无法在信函中全文抄录，但插入了一些引文以证明他们所谴责的学说确实在该书中得到倡导。然而，有些人坚持认为\Person{玛策禄}只是提出了一些问题，这些问题被\Person{欧瑟伯}的追随者曲解，并作为实际的信仰宣示呈报给了皇帝。\Person{欧瑟伯}及其党羽对\Person{玛策禄}极为恼怒，因为他没有同意\Place{腓尼基}会议所提出的定义，也没有同意\Place{耶路撒冷}为支持\Person{亚略}所作的规定；并且他同样拒绝出席大殉道殿的奉献礼，以避免与他们共融。在致皇帝的信中，他们大篇幅强调了后一件事，并将其作为指控提出，声称拒绝出席他在\Place{耶路撒冷}所建圣殿的奉献礼，是对他个人的侮辱。促使\Person{玛策禄}撰写这部著作的原因是，\Place{卡帕多细雅}人\Person{阿斯德里}，一位诡辩家，写了一篇为亚略主义辩护的论文，并在各城市、向众主教以及他出席的几次会议上朗读。\Person{玛策禄}着手反驳他的论点，并在从事此事时，或有意或无意地陷入了\Person{萨莫萨塔的保禄}的见解。然而，后来他在\Place{撒尔德}会议中证明了自己并不持有此类思想后，恢复了他的主教职务。\par

% structure: p0092 source=h2 target=subsection reason=relative-heading-rank; base=h2

\subsection{第三十四章 君士坦丁大帝驾崩；他在领受圣洗圣事后逝世，葬于宗徒圣殿。}

皇帝已将帝国分给他的儿子们，他们称为凯撒。他把西部区域赐给\Person{君士坦丁}和\Person{君士坦斯}，东部赐给\Person{君士坦提乌斯}。当他不适，需要借助温泉治疗时，为此他前往\Place{彼提尼雅}的\Place{厄肋诺波里}城。然而，他的病情加重，他前往\Place{尼科美底亚}，并在该城的一个郊区领受了神圣的圣洗圣事。仪式后，他满心喜乐，感谢天主。然后，他按照先前的分配，确认了帝国在他儿子们之间的划分，并赐予\Place{旧罗马}和以他自己命名的城市某些特权。他将遗嘱交给那位不断赞扬\Person{亚略}的司铎手中，此人是由他妹妹\Person{君士坦提亚}临终时推荐为德性高尚之人的；他附加誓言命令该司铎在\Person{君士坦提乌斯}回来时交给他，因为\Person{君士坦提乌斯}和其他凯撒都不在垂死的父亲身边。安排好这一切后，\Person{君士坦丁}只活了没几天；他在六十五岁、在位三十一年时逝世。他是基督宗教的有力保护者，也是第一位开始热心于教会并赐予她崇高恩惠的皇帝。他在一切事业上都比任何其他君主更为成功；因为我确信，他没有谋划过任何不依靠天主的事。他在对抗哥特人和萨尔马提人的战争中，以及实际上在他所有的军事行动中都取得了胜利；他如此轻松地按照自己的意志改变了政府形式，创建了另一个元老院和另一个以自己名字命名的帝国城市。他攻击外教宗教，并在短时间内颠覆了它，尽管它已在君主和人民中盛行多时。\par

\Person{君士坦丁}死后，他的遗体被放入金棺，运往\Place{君士坦丁堡}，停放在宫中某个平台上；宫中的人对他保持了与他生前相同的尊荣和礼仪。\Person{君士坦提乌斯}当时在东方，闻知父亲去世后，急忙赶回\Place{君士坦丁堡}，以最隆重的仪式安葬了王的遗体，并将之安放在逝者生前下令在\Place{宗徒教堂}修建的墓穴中。从此以后，将此后基督宗教皇帝的遗体安葬在同一墓穴成为惯例；同样，主教们也埋葬于此，因为圣统的尊严不仅在荣誉上与皇权相等，且在神圣场所中甚至占据优势。\par

\SourceRecord{Translated by Chester D. Hartranft.；Nicene and Post-Nicene Fathers, Second Series；Vol. 2.；Edited by Philip Schaff and Henry Wace.；Buffalo, NY: Christian Literature Publishing Co.,；1890.；<http://www.newadvent.org/fathers/26022.htm>.}

% structure: p0001 source=h1 target=section reason=page-title

\section{教会史（第三卷）}

% structure: p0002 source=h2 target=subsection reason=relative-heading-rank; base=h2

\subsection{第一章　\Person{君士坦丁大帝}驾崩后，\Person{欧瑟伯}与\Person{德奥格尼斯}的追随者攻击尼西亚信德}

我们已经看到\Person{君士坦丁}在位期间教会中所发生的事件。他去世后，在\Place{尼西亚}所颁布的教义受到重新审查。虽然这教义并未获得普遍赞同，但在\Person{君士坦丁}生前，无人敢公开拒绝。然而他一去世，许多人便放弃了这个意见，尤其是那些先前被怀疑有背叛之心的人。其中，\Place{比提尼亚}省的\Person{欧瑟伯}与\Person{德奥格尼斯}主教竭尽全力，试图使\Person{亚略}的主张占上风。他们相信，如果能阻止\Person{亚大纳削}从流放地返回，并将\Place{埃及}各教会交给一位与他们意见相同的主教管理，这个目的便容易实现。他们找到了一位得力的帮手，就是那位曾从\Person{君士坦丁}那里获得召回\Person{亚略}的长老。此人因将\Person{君士坦丁}的遗嘱交付给\Person{君士坦提乌斯}皇帝而深受器重；由于受到信任，他大胆把握时机，最终成为皇帝妻子及后宫有权势的太监们的亲信。此时，\Person{欧瑟伯}被委任管理皇宫事务，他狂热拥护\Person{亚略}主义，便诱使皇后及许多宫廷人员接受同样的观点。因此，关于教义的争论再次在私下和公开场合盛行，辱骂与仇恨重新燃起。这种局面正合\Person{德奥格尼斯}及其同党的意。\par

% structure: p0004 source=h2 target=subsection reason=relative-heading-rank; base=h2

\subsection{第二章　伟大的\Person{亚大纳削}从\Place{罗马}返回；\Person{君士坦丁大帝}之子恺撒\Person{君士坦丁}的书信；\Person{亚略}派对抗\Person{亚大纳削}的新阴谋；\Place{贝罗亚}的\Person{阿卡西乌}；\Person{君士坦斯}与\Person{君士坦丁}之间的战争}

此时，\Person{亚大纳削}从\Place{高卢}返回\Place{亚历山大里亚}。据说\Person{君士坦丁大帝}本打算召回他，并在遗嘱中甚至下了这样的命令。但他因去世未能实行这一意图，而他的儿子——与他同名、当时正在西\Place{高卢}指挥军队的\Person{君士坦丁}——召回了\Person{亚大纳削}，并就此事致信\Place{亚历山大里亚}的民众。我找到这封信从拉丁文译成希腊文的抄本，将按原样插入如下：\par

恺撒\Person{君士坦丁}，致\Place{亚历山大里亚}城中天主教会的民众。\par

\Quote{我相信你们不会不知道，可敬法律的诠释者亚大纳削，由于他嗜血成性、充满敌意的仇敌的残酷行为持续不断，危及他的神圣人身，因而被暂时送到高卢，以免因这些卑鄙对手的乖戾而遭受不可挽回的绝境；为使这危险落空，他从紧逼他的人的口中被救出，并奉命住在我近旁，使得在他所居住的城市里，他能充足地获得一切必需品；但他的德行如此卓越非凡，因为他确信天主的助佑，以至于他蔑视命运的一切严酷重担。我主、我父——蒙福的纪念——君士坦丁·奥古斯都，原有意恢复这位主教于其原位，并尤其以此将他归还给你们深爱的虔诚；然而，由于他被人的命运所阻，未及实现其意图便去世了，我作为他的继承人，决意执行这位神圣纪念的皇帝的计划。亚大纳削将告诉你们，当他见到你们的面容时，他是如何受到我的极大尊敬。我对他如此行事并不奇怪，因为你们自己渴望的形象和如此高贵之人的面貌，促使并推动我采取这一步骤。愿天主的眷顾守护你们，我亲爱的弟兄们。}\par

由于皇帝的这封信，\Person{亚大纳削}回家，重新执掌\Place{埃及}各教会。那些依附\Person{亚略}教义的人惊慌失措，无法保持平静；他们不断煽动骚乱，并对他采取其他阴谋诡计。\Person{欧瑟伯}的同党在皇帝面前控告他，说他是好乱之人，并违背教会法律、未经主教同意就推翻了流放令。我将在适当的地方叙述他们如何通过阴谋使\Person{亚大纳削}再次被逐出\Place{亚历山大里亚}。\par

被称为\Person{旁菲鲁斯}的\Person{欧瑟伯}大约此时去世，\Person{阿卡西乌}继任\Place{巴勒斯坦}的\Place{凯撒利亚}的主教职位。他热切模仿\Person{欧瑟伯}，因为他曾受教于\Person{欧瑟伯}的圣经教导；他头脑敏捷，辞令文雅，留下了许多值得称道的著作。不久之后，皇帝\Person{君士坦丁}在\Place{阿奎莱亚}向他的兄弟\Person{君士坦斯}宣战，并被自己的将领所杀。罗马帝国在幸存的两兄弟之间瓜分：西方归于\Person{君士坦斯}，东方归于\Person{君士坦提乌斯}。\par

% structure: p0010 source=h2 target=subsection reason=relative-heading-rank; base=h2

\subsection{第三章　\Place{君士坦丁堡}主教\Person{保禄}，以及\Quote{敌对圣神者}\Person{马其顿尼}}

\Person{亚历山大}在这时去世，保罗继任了君士坦丁堡的主教之职。亚略派和马其顿尼派的追随者断言，他是擅自占据此位，并未征求尼苛米底亚主教欧瑟伯或色雷斯地区赫拉克肋亚主教德奥若多的意见；按规定，祝圣之权属于最近的主教们。然而，许多人根据他所继任的亚历山大的证词主张，他是由当时聚集在君士坦丁堡的主教们祝圣的。因为当亚历山大九十八岁、已精力充沛地执行主教职务二十三年、濒临死亡时，他的神职人员问他希望将教会托付给谁。\Quote{如果，}他回答说，\Quote{你们寻求一位通晓神圣事物、善于教导的人，就要保罗；但若你们想要一位熟悉公共事务和统治者会议的人，马其顿尼更好。}马其顿尼派自己承认亚历山大给出了这一证词；但他们说保罗更擅长事务处理和演说艺术；而他们为马其顿尼强调他生命的见证；他们指责保罗耽于柔弱和品行不端。然而，从他们自己的承认来看，保罗是一位有口才的人，并且在教导教会方面卓有成效。事实证明，他无力对抗生活中的意外事件，或与当权者交往；因为他从未成功挫败敌人的阴谋，不像那些善于处理事务的人。尽管他深受人民爱戴，但他因那些当时拒绝尼西亚所盛行教义的人的背叛而遭受严重苦难。首先，他被驱逐出君士坦丁堡教会，似乎对他有一些行为不当的指控。然后他被判处流放，最后，据说他成为敌人阴谋的牺牲品，被勒死。但这些后来的事件发生在之后时期。\par

% structure: p0012 source=h2 target=subsection reason=relative-heading-rank; base=h2

\subsection{第四章 在保罗的祝圣上激起了一场骚动}

保罗的祝圣在君士坦丁堡教会中引起了巨大骚动。在亚历山大活着时，亚略派人士并未十分公开行动；因为人民由于留心听他而得到了良好的治理，并尊崇神圣事物，尤其相信发生在亚略身上的意外事件——他们认为他遭遇了那样的死亡——是天主义怒的惩罚，是由亚历山大的诅咒带来的。然而，这位主教去世后，人民分裂为两派，关于教义的争论和斗争公开进行。亚略的追随者希望祝圣马其顿尼，而那些坚持圣子与圣父同性同体的人则希望保罗作他们的主教；后一派获胜。保罗祝圣后，当时恰巧不在家的皇帝返回君士坦丁堡，对所发生的事表现出极大的不悦，就好像主教职被授予了一个不配的人。通过保罗敌人的阴谋，召开了一次会议，他被逐出教会。会议将君士坦丁堡教会交给了尼苛米底亚主教欧瑟伯。\par

% structure: p0014 source=h2 target=subsection reason=relative-heading-rank; base=h2

\subsection{第五章 安提约基雅的部分会议；它废黜了亚大纳削；它任命了额我略；它的两份信仰声明；同意它们的人}

这些事发生不久之后，皇帝前往叙利亚城市安提约基雅。那里已经建成了一座教堂，其规模与宏伟超群。君士坦丁在世时便已开始建造，此建筑刚由其子君士坦提乌完工，欧瑟伯的党羽——他们向来对此事热心——便认为这是召集公会议的有利时机。因此，他们与来自各地、持有相同见解的人一同在安提约基雅会集；与会主教约有九十七位。他们声称的目的是为新竣工的教堂举行祝圣典礼，但实际上别无所图，只为废除尼西亚大公会议的法令，其后来的行动充分证明了这一点。当时安提约基雅教会由普拉塞图斯治理，他是优弗罗纽的继任者。君士坦丁大帝去世约在此时之前五年。当全体主教在君士坦提乌皇帝面前聚集时，多数人表示极大愤慨，严厉指控亚大纳削藐视他们所颁布的司铎法规，在未先获得公会议认可的情况下占据了亚历山大主教区。他们还指控他是因他返回而煽动的叛乱中数人死亡的原因，并称许多人因此被捕并移交司法审判。通过这些指控，他们设法使亚大纳削名誉扫地，并决议任命额我略治理亚历山大教会。随后他们转而讨论教义问题，对尼西亚大公会议的法令未加挑剔。他们向各城的主教发送公函，在其中宣称：他们本人身为主教，并未追随亚略。\Quote{因为，}他们说，\Quote{我们怎能追随他呢？他不过是一位长老，而我们的地位在他之上。}既然他们是其信仰的检验者，便乐意接纳了他；他们相信自始由传统传承的信仰。他们在信函末尾进一步解释这一点，但未提及圣父或圣子的实体，也未使用\Quote{同体}一词。事实上，他们采用了如此模糊的措辞，以致无论是亚略派还是尼西亚大公会议法令的追随者，都无法质疑他们措辞的安排，仿佛他们不熟悉圣经。他们有意避免任何被一方或另一方拒绝的表达形式，仅使用普遍接受的用语。他们明认圣子与圣父同在、是独生子、是天主、在万有之先存在，并取得肉身，完成了圣父的旨意。他们承认这些以及类似的真理，但未描述圣子与圣父同永恒或同体的教义，也未描述相反的教义。此后，看来他们改变了关于这一信式的心意，又发布了另一份信式，我认为它与尼西亚大公会议的十分接近——除非其中含有我不明显的隐秘含义。尽管他们——不知出于何种动机——未明说圣子是同体的，但他们承认祂是不变的、祂的神性不容改变、祂是圣父实体、旨意、德能与荣耀的完美肖像、祂是一切受造物的首生者。他们宣称他们已找到这信式，且完全由在尼科美底亚殉道的路济安所写；路济安是一位备受赞誉、极其精确通晓圣经的人。我不知道这说法是否属实，抑或他们仅为借殉道者的尊荣加重自己文件的分量而提出此说。参与这次公会议的不仅有欧瑟伯（他在保禄被逐后，从尼科美底亚调任君士坦丁堡宗座），还有阿卡西乌（欧瑟伯·庞斐禄的继任者）、西托波利斯主教帕特罗斐禄、赫拉克勒亚（旧名佩林图斯）主教德奥多禄、革尔马尼西亚主教欧多克修（他后来在马其顿纽之后治理君士坦丁堡教会），以及被选担任亚历山大教会的额我略。普遍承认所有这些主教持有相同见解，例如卡帕多细亚凯撒勒雅主教狄亚尼、叙利亚劳迪基雅主教乔治，以及许多其他担任都主教区和重要教会的主教。\par

% structure: p0016 source=h2 target=subsection reason=relative-heading-rank; base=h2

\subsection{第六章　别号厄梅森的欧瑟伯；额我略接受亚历山大教会；亚大纳削赴罗马寻求庇护。}

\Person{欧瑟伯}，别号厄梅森的欧瑟伯也出席了公会议。他出身于奥斯若恩的厄德萨城贵族家庭。按本国习俗，他自幼学习圣经，后来在则时旅居其故乡的教师指导下学习了希腊学问。随后，他在欧瑟伯·庞斐禄和西托波利斯主教帕特罗斐禄的引导下，更深入地通晓了圣经。当欧斯达提因齐禄的指控被罢黜时，他去了安提约基雅，与继任者优弗罗纽亲密同居。为逃避被授予司铎尊位，他逃往亚历山大，并常出入哲学家的学园。在熟悉了他们的学规后，他返回安提约基雅，与优弗罗纽的继任者普拉塞图斯同住。在该城举行公会议期间，君士坦丁堡主教欧瑟伯恳请他接受亚历山大教区；因为人们认为，藉其圣德的盛名和卓越的口才，他很容易在埃及人心中取代亚大纳削。然而，他拒绝受任，理由是这样做只会招致亚历山大人的立即仇恨，他们除了亚大纳削不愿接受任何其他主教。因此，额我略被任命为亚历山大教会的主教，而欧瑟伯被任命为厄梅萨教会的主教。\par

他在那里遭受了一场叛乱；因为民众控告他从事那种被称为占星术的天文学，他被迫逃亡以寻求安全，便前往\Place{劳狄刻雅}，与那座城市的主教\Person{乔治}同住，\Person{乔治}是他的密友。后来他陪同这位主教前往\Place{安提约基雅}，并获准从主教\Person{普拉凯图斯}和\Person{纳尔奇苏斯}那里返回\Place{厄梅萨}。他深受\Person{君士坦提乌斯}皇帝的敬重，并随同皇帝出征\Place{波斯}。据说天主藉着他施行了奇迹，正如\Place{劳狄刻雅}的\Person{乔治}所见证的，他曾记载了关于他的这些和其他事迹。\par

然而，尽管他拥有如此众多的高尚品质，他还是无法逃脱那些因目睹他人德行而恼怒之人的嫉妒。他承受了接纳\Person{撒伯里乌斯}学说的指责。但在当时，他与在\Place{安提约基雅}召集的主教们一同投票。据说\Place{耶路撒冷}主教\Person{马克西穆斯}故意回避了这次会议，因为他后悔在无意中同意了撤职\Person{阿塔纳修}的决议。罗马教廷的管事者，以及来自\Place{意大利}东部或\Place{罗马}以外地区的代表都没有出席\Place{安提约基雅}会议。在同一时期，法兰克人侵扰了\Place{高卢}西部；东方各省，特别是\Place{安提约基雅}，在会议之后遭受了剧烈的地震。会议之后，\Person{格列高利}率领大批士兵前往\Place{亚历山大里亚}，这些士兵受命确保他平安无扰地进入城市；阿里乌派信徒也支持他，他们急切希望驱逐\Person{阿塔纳修}。\Person{阿塔纳修}担心人民因他而受苦，便在夜间将众人聚集在教堂里，当士兵前来占领教堂时，祈祷结束后，他首先命人唱了一首圣咏。在吟唱这首圣咏期间，士兵们留在外面，安静地等待其结束，与此同时\Person{阿塔纳修}从歌者中间穿过，秘密逃脱，逃往\Place{罗马}。\Person{格列高利}就这样占据了\Place{亚历山大里亚}的主教席位。人民的愤怒被激起，他们焚烧了以他们先前一位主教\Person{狄约尼削}命名的教堂。\par

% structure: p0020 source=h2 target=subsection reason=relative-heading-rank; base=h2

\subsection{第七章　\Place{罗马}与\Place{君士坦丁堡}的大祭司；\Person{欧瑟伯}之后\Person{保禄}的复职；军队统帅\Person{赫尔默根纳}的被杀；\Person{君士坦提乌斯}从\Place{安提约基雅}前来，撤去\Person{保禄}的职务，并对该城发怒；他让\Person{马塞多尼乌斯}处于悬而未决的状态，然后返回\Place{安提约基雅}。}

就这样，那些支持各种异端反对真理之人的阴谋得以成功执行；他们罢免了那些在整个东方坚决维护尼西亚大公会议教义的主教。这些异端者占据了最重要的主教座，如埃及的\Place{亚历山大里亚}、叙利亚的\Place{安提约基雅}，以及\Place{赫勒斯滂}的帝都，并使所有被说服的主教都服从于他们。\Place{罗马}教会的首领和西方的所有司铎都视这些行为为个人侮辱；因为他们从一开始就赞同尼西亚与会者投票作出的所有决定，至今也没有改变这种想法。\Person{阿塔纳修}到达后，他们亲切地接待了他，并在他们中间支持他的事业。被这种干涉激怒，\Person{欧瑟伯}写信给\Person{尤利乌斯}，敦促他自任为审判官，审查\Place{提尔}会议针对\Person{阿塔纳修}所颁布的法令。但在\Person{尤利乌斯}尚未明确意见之前，实际上在\Place{安提约基雅}会议后不久，\Person{欧瑟伯}就去世了。此事一发生，那些坚持尼西亚大公会议教义的\Place{君士坦丁堡}市民便将\Person{保禄}带入教堂。与此同时，反对派的群众也抓住机会聚集在另一座教堂，其中包括\Place{尼西亚}主教\Person{特奥格尼斯}、\Place{赫拉克勒亚}主教\Person{特奥多雷}以及同一派别的其他在场者，他们祝圣了\Person{马塞多尼乌斯}为\Place{君士坦丁堡}主教。这引发了城中频繁的骚乱，完全呈现出战争的面貌，因为人民互相攻击，许多人丧生。城市充满了骚动，以致当时在\Place{安提约基雅}的皇帝听闻所发生之事后，勃然大怒，颁布了驱逐\Person{保禄}的法令。骑兵统帅\Person{赫尔默根纳}试图执行皇帝的这道谕令；因为他被派往\Place{色雷斯}，途中经过\Place{君士坦丁堡}，他想凭借军队武力将\Person{保禄}从教堂中逐出。但人民非但没有屈服，反而公开抵抗他；当士兵为了执行所接到的命令而更加暴力时，叛乱者闯入\Person{赫尔默根纳}的住所，放火焚烧，杀死了他，并用绳子系住他的尸体，拖过全城。皇帝一接到这个消息，就骑马赶往\Place{君士坦丁堡}，以惩罚人民。但当他看到人民流着泪、恳求着前来迎接他时，他便宽恕了他们。他将城中每年从\Place{埃及}贡赋中由他父亲\Person{君士坦丁}从国库拨给他们的粮食减少了一半，大概是认为奢侈和过度使民众游手好闲、倾向于叛乱。他将怒气转向\Person{保禄}，下令将他逐出城市。他对\Person{马塞多尼乌斯}也极为不悦，因为他是将军和其他人被杀的原因，而且他是在未经皇帝批准的情况下受祝圣的。然而，他返回了\Place{安提约基雅}，既没有确认也没有解除\Person{马塞多尼乌斯}的祝圣。与此同时，阿里乌派教义的热心者罢免了\Person{格列高利}，因为他对其教义的支持漠不关心，而且因他进城时城市所遭受的灾难，尤其是教堂被焚，引起了\Place{亚历山大里亚}人的反感。他们选举了\Place{卡帕多细亚}人\Person{乔治}接替他；这位新主教因其活力和对阿里乌派教义的热心而受到赞赏。\par

% structure: p0022 source=h2 target=subsection reason=relative-heading-rank; base=h2

\subsection{第八章 东方主教们到达罗马；罗马主教犹利关于他们的信函；藉犹利的信函，保禄与亚大纳削重获各自的教座；东方总司铎致犹利信函的内容。}

\Person{亚大纳削}离开亚历山大时，已经逃往\Place{罗马}。君士坦丁堡主教\Person{保禄}、安基拉主教\Person{玛尔塞禄}和加沙主教\Person{阿斯克莱帕}也同时前往那里。\Person{阿斯克莱帕}反对亚略派，因此被罢免，曾被一些异端者指控推倒了一座祭台；\Person{金提安}被任命接替他治理加沙教会。还有哈德良堡主教\Person{路济}，因另一项指控被免去所管理的教会，此时也居住在\Place{罗马}。罗马主教在了解了每个人的控诉，并发现他们关于尼西亚信理持相同见解后，便接纳他们进入共融，因为他们同样拥有正统信仰；而且由于对他教座尊严的关怀涵盖一切，他使他们全都恢复了各自的教会。他写信给东方的主教们，责备他们不公正地审判了这些主教，以及因放弃尼西亚教义而骚扰教会。他召集他们中的少数人在指定日期到他面前，以便向他说明他们所判定的理由，并威胁说，除非他们停止创新，否则不再容忍他们。这就是他信函的主旨。\Person{亚大纳削}和\Person{保禄}恢复了各自的教座，并将犹利的信函转交给了东方的主教们。主教们几乎无法容忍这样的文件，他们聚集在\Place{安提约基雅}，草拟了一份答复犹利的信函，措辞优美，法律技巧精湛，却充满了相当的讽刺，并沉溺于最强烈的威胁。在这封书信中，他们承认罗马教会理应受到普世的尊敬，因为它既是宗徒的学校，又从一开始就成为虔敬的母城，尽管教义的引入者是从东方定居在那里的。他们补充说，在荣誉方面不应将第二位置归于他们，因为他们教会在规模或数量上并不占优势；因为他们以德行和决断力超越罗马人。他们指责犹利接纳亚大纳削的追随者进入共融，并表达了对他的愤慨，因为他侮辱了他们的主教会议并废除了他们的法令；他们攻击他的行为是不公正的，与教会权利不一致。在这些谴责和对这些冤情的抗议之后，他们继续陈述说，如果犹利承认被他们驱逐的主教的罢免，以及承认他们任命代替者，他们就应许和平与共融；但除非他同意这些条件，他们将公开宣布反对他。他们补充说，在他们之前管理东方教会的司铎们，对罗马教会罢免\Person{诺瓦替安}一事并未提出反对。他们在信中丝毫没有提及他们偏离尼西亚大公会议教义的任何行为，只是说他们有各种理由来为他们所采取的行动辩护，并且认为当时没有必要对他们的行为进行任何辩护，因为他们被怀疑在各方面都违反了正义。\par

% structure: p0024 source=h2 target=subsection reason=relative-heading-rank; base=h2

\subsection{第九章 保禄与亚大纳削被逐；马其顿尼受付托治理君士坦丁堡教会。}

在如此写信给犹利之后，东方的主教们向皇帝君士坦提乌斯控告了他们所罢免的人。于是，当时在\Place{安提约基雅}的皇帝写信给君士坦丁堡总督\Person{斐理伯}，命令他将教会交付给\Person{马其顿尼}，并将\Person{保禄}逐出城市。总督害怕民众的骚动，在皇帝的命令被泄露之前，他前往名为\Place{采克西普斯}的公共浴场，那是一座显眼而庞大的建筑，并召来了\Person{保禄}，仿佛他想与他商谈一些公共利益的事务。\Person{保禄}一到，总督便向他展示了皇帝的谕令。\Person{保禄}依照命令，被秘密地通过毗邻浴场的宫殿带往海边，登上一艘船，被送往\Place{得撒洛尼}，据说他的祖先原籍那里。他被严格禁止接近东方地区，但不禁止前往\Place{伊利里亚}和更远的省份。\par

在离开法庭后，\Person{斐理伯}在\Person{马其顿尼}的陪同下前往教堂。民众在此期间已无数地聚集起来，迅速挤满了教堂；他们分裂成的两派，即亚略异端的支持者和\Person{保禄}的追随者，都努力夺取教堂。当总督和\Person{马其顿尼}到达教堂大门时，士兵们试图强行推开人群，以便为这些显要人物让路；但由于人群如此拥挤，他们既无法后退——因为他们已紧密地挤到最远的地方——也无法让路；士兵们以为人群不愿退让，就用剑杀死了许多人，还有许多人被踩死。皇帝的谕令就这样实现了，\Person{马其顿尼}接受了教会，而\Person{保禄}意外地被逐出了君士坦丁堡的教会。\par

与此同时，\Person{亚大纳削}已逃走并躲藏起来，害怕皇帝君士坦提乌斯的威胁，因为他曾威胁要处死他；因为异端者使皇帝相信他是一个煽动叛乱的人，并且他在返回主教职时导致了几人的死亡。但皇帝的愤怒主要是由于这样的陈述而被激起的：即\Person{亚大纳削}变卖了皇帝君士坦丁赐予亚历山大穷人的粮食，并占用了价款。\par

% structure: p0028 source=h2 target=subsection reason=relative-heading-rank; base=h2

\subsection{第十章 罗马主教致书东方主教维护亚大纳削；东方主教遣使至罗马，与罗马主教一同调查对东方主教的指控；\Person{君士坦斯}凯撒遣散该使团。}

埃及的主教们提交了一份书面声明，声称这些指控是虚假的，\Person{儒略}得知\Person{亚大纳削}在埃及远非安全，便将他召至他自己的城市。同时，他回复了在\Place{安提约基雅}召集的主教们的信件——恰好那时他收到了他们的书信——指责他们秘密引入违背\Place{尼西亚}大公会议信理的革新，并违反了教会的法律，因为忽略了邀请他参加他们的主教会议；他声称有一项司铎法规宣布：凡违反罗马主教判决而制定的法令均属无效。他还责备他们在\Place{提洛}和\Place{玛雷奥提斯}针对\Person{亚大纳削}的所有诉讼中偏离了正义，并指出，因\Person{阿尔塞尼乌斯}之手事件的诬告，前一个城市所颁布的法令已被废除；又因\Person{亚大纳削}缺席，后一个城市的法令也被废除。最后，他谴责了他们书信的傲慢风格。\par

基于所有这些理由，\Person{儒略}承担起为\Person{亚大纳削}和\Person{保禄}辩护的责任；后者不久前抵达意大利，曾对这些灾难沉痛哀叹。当\Person{儒略}发现他写给东方那些身负司铎尊位者的信函毫无效果时，便将此事禀报给皇帝\Person{君士坦斯}。于是，\Person{君士坦斯}致信其弟\Person{君士坦提乌斯}，请求他派遣几位东方主教，以便为他们的罢黜敕令说明理由。为此选出了三位主教：即\Place{基里基雅}伊勒诺波利斯的主教\Person{纳尔奇苏斯}、\Place{色雷斯}赫拉克利亚的主教\Person{狄奥多}，以及\Place{叙利亚}阿勒图萨的主教\Person{玛尔谷}。他们抵达意大利后，竭力为其行为辩护，并试图使皇帝相信东方主教会议所作判决是公正的。当被要求陈述他们的信仰宣认时，他们隐瞒了在\Place{安提约基雅}起草的信经，而提交了另一份书面的信仰宣认，同样与\Place{尼西亚}所认可的教义相悖。\Person{君士坦斯}看出他们不公正地陷害了\Person{保禄}和\Person{亚大纳削}，将他们逐出共融，并非如控告所声称的因其品行的问题，而仅仅是由于教义上的分歧；因此，他遣散了使团，对他们前来的陈述完全不予采信。\par

% structure: p0031 source=h2 target=subsection reason=relative-heading-rank; base=h2

\subsection{第十一章 长信经与\Place{撒尔底迦}会议颁布的法令。\Person{罗马}主教\Person{儒略}与西班牙主教\Person{奥西乌}因与\Person{亚大纳削}等人共融而被东方主教废黜。}

三年后，东方主教们向西方主教们送交了一份信经，因其措辞和思想的冗长超过以往任何信条，被称为\OriginalTerm{长信经}{\GreekText{μακρόστιχος} \GreekText{ἔκθεσις}}。在这份信经中，他们未提及天主的本体，但绝罚了那些主张圣子出于\Quote{无有}、或祂属于另一个\OriginalTerm{位格}{hypostasis}而非出于天主、或曾有一段时间或一个时代祂不存在的人。\Person{欧多克西}（当时仍是\Place{革尔马尼西亚}的主教）、\Person{玛尔提里}和\Person{马其顿尼}携带了这份文件，但西方司铎并未接受它；因为他们宣称对\Place{尼西亚}确定的教义已完全满意，认为完全不必过于探究这些细节。\par

\Person{君士坦斯}皇帝请求其弟恢复\Person{亚大纳削}派主教的职位，但因敌对异端的影响而未果；此外，\Person{亚大纳削}派与\Person{保禄}派恳请\Person{君士坦斯}召开一次主教会议，以应对废除正统教义的阴谋，于是两位皇帝一致同意东西方主教在指定日期会集于\Place{依利里亚}的\Place{撒尔底迦}城。东方主教们起初聚集于\Place{色雷斯}的\Place{斐理伯波里}，致书已集会于\Place{撒尔底迦}的西方主教们，声称除非他们将\Person{亚大纳削}派逐出会议并断绝共融——因为后者已被罢黜——否则不予加入。他们随后前往\Place{撒尔底迦}，但声明若不将已被他们罢黜的人排除于教会之外，他们便不进入教堂。西方主教们答复说，他们从未驱逐过那些人，如今也不愿让步，尤其因为罗马主教\Person{犹利乌}在调查此案后并未定他们的罪，况且他们也在场，愿意为自己辩护并再次驳斥所受的指控。但这些声明毫无成效；由于他们为调解分歧所定的期限已过，双方最终就这些问题互致信函，反而加深了原有的敌意。他们分别集会，做出了相反的裁决：东方主教们确认了此前对\Person{亚大纳削}、\Person{保禄}、\Person{马尔塞鲁}及\Person{阿斯克勒帕斯}所颁布的判决，并罢黜了罗马主教\Person{犹利乌}，因为他首先接纳了被他们定罪的人进入共融；宣信者\Person{何西乌}也遭罢黜，部分由于同一原因，部分由于他是\Place{安提约基雅}教会首领\Person{保利诺}与\Person{欧斯塔提乌}的朋友。\Place{特雷维斯}主教\Person{玛克西穆}被罢黜，因为他首倡接纳\Person{保禄}进入共融，并促成其返回君士坦丁堡，且拒绝与前往高卢的东方主教共融。此外，他们还罢黜了\Place{撒尔底迦}主教\Person{普罗托根}与\Person{高登提乌}：前者因偏袒\Person{马尔塞鲁}——尽管此前已定其罪——后者因行事异于前任\Person{基里亚库}，并支持了许多当时已被他们罢黜的人。颁布这些判决后，他们通告各地区主教，不得与被罢黜者共融，不得致信他们，亦不得接收他们的来信。他们还命令众人相信其附于信函后的信条中所陈述的关于天主的道理；该信条未提及\OriginalTerm{同质}{consubstantial}一词，但绝罚了那些说有三神、或说基督不是神、或说父、子、圣神是同一者、或说子是非受生者、或说曾有某一时间或世代子不存在的人。\par

% structure: p0034 source=h2 target=subsection reason=relative-heading-rank; base=h2

\subsection{第十二章  \Person{犹利乌}派与\Person{何西乌}派的主教另开会议，罢黜东方大司祭，并制定信条。}

\Person{何西乌}的追随者在此期间聚集在一起，宣告他们无罪：\Person{亚大纳削}，因为那些在\Place{提洛}集会的人对他进行了不义的阴谋；\Person{马塞鲁斯}，因为他并未持有被指控的意见；\Person{阿斯克勒帕斯}，因为他已由\Person{优西比乌·庞非勒}和许多其他审判官的投票重新被确立在他的教区；对此，他通过审判记录证明了其真实性；最后，\Person{路济乌斯}，因为他的指控者已经逃逸。他们写信给每位被宣告无罪者的教区，命令他们接受并承认自己的主教。他们声明，\Person{额我略}并非由他们任命为\Place{亚历山大里亚}主教；\Person{巴西略}亦非\Place{安基拉}主教；\Person{昆提亚努斯}亦非\Place{迦萨}主教；他们未曾接纳这些人共融，甚至不视他们为基督徒。他们罢免了以下主教：\Place{色雷斯}主教\Person{德奥多罗}、\Place{伊雷诺波利}主教\Person{纳尔奇苏斯}、\Place{凯撒勒雅（巴勒斯坦）}主教\Person{阿卡西乌}、\Place{厄弗所}主教\Person{梅诺芬托}、\Place{莫西亚的锡吉杜努姆}主教\Person{乌尔撒西乌}、\Place{潘诺尼亚的穆尔西亚}主教\Person{瓦伦斯}、\Place{劳狄刻雅}主教\Person{乔治}，尽管后者未与东方主教一同出席主教会议。他们将上述人物逐出司铎职务和共融，因为他们将圣子与圣父的本体分离，并接纳了那些因持守\OriginalTerm{亚略异端}{Arian heresy}而先前被罢免的人，且将他们擢升到天主服务的最高职位。在因这些偏差而开除他们并宣告他们为\OriginalTerm{天主教会}{Catholic Church}之外的人之后，他们致信各国主教，命令他们确认这些法令，并在教义问题上与他们同心一意。他们还撰写了另一份信仰文件，比\OriginalTerm{尼西亚}{Nicæa}文件更为详尽，尽管相同的思想被谨慎保存，且该文件的措辞变化甚微。在\Place{撒狄卡}集会的西方主教中位居首位的\Person{何西乌}和\Person{普罗托格内斯}，或许害怕被怀疑对\OriginalTerm{尼西亚公会议}{Nicene council}的教义有任何创新，写信给\Person{犹利}，证明他们坚定持守这些教义，但因清晰需要，他们不得不扩充相同的思想，以免亚略派利用文件的简短，将不谙辩证法的人引入荒谬。当每一方都完成了我所叙述的事情之后，会议解散，成员各自返回家乡。这次主教会议是在\Person{鲁菲努斯}和\Person{优西比乌}担任执政官期间，约在\Person{君士坦丁}去世后十一年召开的。约有三百位西方城市的主教，以及超过七十六位东方主教出席，其中包括\Person{伊斯基里翁}，他是由\Person{亚大纳削}的敌人任命为\Place{马勒奥提斯}主教的。\par

% structure: p0036 source=h2 target=subsection reason=relative-heading-rank; base=h2

\subsection{Chapter 13. 主教会议之后，东方与西方分离；西方高尚地持守尼西亚公会议的信仰，而东方则在此信条上处处为纷争所扰乱。}

这次主教会议之后，东方教会与西方教会停止了同信仰者之间通常存在的交往，并拒绝彼此共融。西方基督徒与所有直到\Place{色雷斯}的人分离；东方基督徒则与直到\Place{伊利里亚}的人分离。可以想见，教会这种分裂状态伴随着异议和诽谤。尽管他们此前在教义问题上已有分歧，但恶事尚未达到严重程度，因为他们仍保持共融，且常有相同感受。整个西方教会完全以教父的教义为准则，远离一切关于教义的争论和咬文嚼字。尽管\Person{奥克森提乌}（已出任\Place{米兰}主教）以及\Place{潘诺尼亚}主教\Person{瓦伦斯}和\Person{乌尔撒西乌}曾试图将帝国那一部分引入亚略教义，但罗马教座的主席和其他司铎已提前慎重防范，铲除了这一烦扰异端的种子。至于东方教会，尽管自\Place{安提约基雅}会议以来它就被纷争撕裂，且已公开与尼西亚信仰形式分歧，但我认为，多数人的意见实际上是统一的，他们宣认圣子与圣父同体。然而，有一些人喜爱争论，反对\Quote{同质}一词——因为最初反对这个词的人，据我推断，并如大多数人所常见，认为显得被征服是耻辱。其他人则因频繁辩论这些主题而最终确信了关于天主的真理，此后始终坚定持守。还有一些人，知道不应发生争论，便倾向于双方都满意的一方，或是由于友谊的影响，或是被各种原因所左右——这些原因常诱使人接受本该拒绝的东西，并在需要彻底确信的情形下不敢行动。另有许多人，认为在言辞上耗费时间是荒谬的，便安静地采纳\Place{尼西亚}公会议所教导的见解。\Person{保禄}（\Place{君士坦丁堡}主教）、\Person{亚大纳削}（\Place{亚历山大里亚}主教）、全体隐修士、仍然在世的\Person{大圣安当}及其门徒，以及埃及和罗马帝国其他地区的无数人，在整个东方其他地区坚定而公开地持守尼西亚公会议的教义。既然我提到隐修士，我将简要叙述在\Person{君士坦提乌斯}统治期间活跃的那些圣人。\par

% structure: p0038 source=h2 target=subsection reason=relative-heading-rank; base=h2

\subsection{Chapter 14. 关于此时在\Place{埃及}活跃的圣人，即\Person{安当}、两位\Person{玛加略}、\Person{赫拉克利乌}、\Person{克罗尼乌}、\Person{帕夫努提乌}、\Person{普图巴斯图斯}、\Person{阿尔西修斯}、\Person{塞拉皮翁}、\Person{皮图里翁}、\Person{帕科米乌}、\Person{阿波罗尼乌}、\Person{阿努夫}、\Person{希拉里翁}，以及许多其他圣人的名录。}

我将从埃及以及两位名叫玛加略的人开始叙述。他们是斯凯提斯和邻近山野的著名首领：一位是埃及本地人，另一位被称为\OriginalTerm{波利提库斯}{Politicus}，因为他是一位市民，出身于亚历山大。二人都以神圣的知识和哲学如此奇妙地富有恩宠，以至于魔鬼都对他们感到恐惧；他们还行了许多非凡的奇事和神奇的治愈。据传说，那位埃及人曾使一个死人复活，为要说服一个异端者相信死人复活的真理。他活了大约九十岁，其中六十年是在旷野中度过的。当他年轻时开始学习哲学，进步神速，以至于隐修士们给他起了个绰号叫\Quote{\Emph{老孩子},}。他在四十岁时被祝圣为司铎。另一位玛加略在晚年成为司铎；他精通各种苦修操练，有些是他自己发明的，而他从其他苦修者那里听到的，他都以各种形式成功付诸实践，以至于皮肤完全干枯，胡须不再生长。\Person{潘波}、\Person{赫拉克利德}、\Person{克罗尼乌}、\Person{帕夫努提乌}、\Person{普图巴斯图}、\Person{阿尔西修斯}、\Person{大塞拉皮昂}、\Person{皮图里翁}（他住在特拜附近）以及\Person{帕科米乌}（即被称为塔本纳修士的隐修会的创始人）都在同一地方和时期兴盛。这一修会的衣着和管理方式在某些方面与其他隐修士不同。然而，其成员致力于德行，轻视尘世事物，激发灵魂向往天上的默观，并预备它欣然离开肉身。他们穿着皮衣，似乎是出于对厄里亚的纪念，因为我认为他们以为这样就能常在记忆中保存先知的德行，并像他一样能够勇敢地抵抗情欲的诱惑，被同样的热忱所影响，并因盼望同样的赏报而受到激励去实践节制。据说这些埃及隐修士的特殊服装与他们的哲学的一些秘密有关，并非没有充分理由就与其他人不同。他们穿无袖的束腰外衣，为要教导手不应预备去做傲慢的恶事。他们头上戴一种称为风帽的遮盖物，为表明他们应当像被乳养、戴着同样形状的帽子的婴儿一样，以同样的纯真与纯洁生活。他们的腰带和一种在他们腰部、肩膀和手臂上的围巾，提醒他们应当随时准备为天主服务和工作。我知道人们为他们的服装特色提出了其他理由，但我所说的似乎已经足够。据说\Person{帕科米乌}最初独自在一个洞穴里修炼哲学，但一位神圣天使向他显现，命令他召集一些年轻隐修士，与他们一起生活，因为他独自修炼哲学已很有成就；现在天使要赐给他律法，让他训练他们，而他要作为团体的领袖占有并造福许多人。然后一块板子被赐给他，这块板子至今仍被小心保存。上面刻有诫命，规定他必须允许每个人按照自己的能力吃喝、工作、禁食；那些食量大的要从事艰苦的劳动，而苦修者则被分配较轻松的工作；他受命建造许多小室，每室住三个隐修士，他们要在公共食堂沉默进食，围坐在桌子旁，脸上蒙着面纱，使他们看不到彼此，也看不到桌子及其上的东西以外的任何事物；他们不得与陌生人一同进食，但旅客除外，要对他们款待；那些希望与他们一起生活的人，须先经过三年的考验，在此期间要从事最艰苦的工作，然后才能分享他们的团体生活。他们要穿皮衣，戴饰有紫色钉子的羊毛头巾，以及亚麻束腰外衣和腰带。他们要穿着束腰外衣和皮衣睡觉，躺在特制的长椅上，椅子两侧封闭，每张椅子能容纳卧具材料。在每周的第一天和最后一天，他们要走近祭台领受神圣奥迹的共融，然后解开腰带，脱下皮袍。他们每天祈祷十二次，晚上同样多次，夜间也要献上同样次数的祈祷。在第九时辰，他们要祈祷三次，在将进食时，每祈祷前要唱一首圣咏。整个团体要分成二十四个班，每班由希腊字母中的一个来标记，使得每班都有一个与其行为和习惯等级相称的称呼。例如，较简单的被称为\OriginalTerm{伊欧塔}{Iota}，而弯曲者被称为\OriginalTerm{泽塔}{Zeta}或\OriginalTerm{克西}{Xi}，其他字母的名称则根据该秩序的目的最相称地对应于字母的形状来选取。\par

这些就是\Person{帕科米乌}管理自己门徒的律法。他是一位爱人且为天主所爱的人，以致能预知未来，并常蒙准与神圣天使交往。他居住在特拜依德的\Place{塔本纳}，因此塔本纳修士的名称一直延续至今。通过采纳这些规则进行管理，他们变得非常著名，并随着时间的推移增长到七千人之多。但与\Person{帕科米乌}同住的塔本纳岛上的团体约有一千三百人；其余的人住在特拜依德和埃及其他地方。他们都遵守同一生活规则，并拥有一切公有。他们把塔本纳岛上建立的团体视为母亲，其统治者视为父亲和君王。\par

大约同一时期，\Person{阿颇罗尼}因宣信奉行隐修哲学而闻名。据说他从十五岁起就在旷野中致力于哲学，当他到了四十岁时，按他所领受的一项神圣命令，前往有人居住的地区居住。他在\Place{泰拜依德}也有一个团体。他深得天主喜爱，并被赋予施行奇迹治愈和非凡事迹的能力。他恪尽职守，以极大的仁慈和善意教导他人哲学。他的祈祷如此蒙悦纳，以至于他向天主所求的没有一样被拒绝，但他如此明智，总是提出审慎的、天主乐于恩赐的请求。\par

我相信，神圣的\Person{阿努夫}生活于这一时期。我获知，自迫害时期他首次宣示皈依基督教以来，他从未说过谎，也不贪求世间之物。他向天主的一切祈祷和恳求都得蒙应允，并由一位圣天使教导他各种美德。不过，我们关于埃及隐修士所言已足矣。\par

同一种哲学在埃及被习得后，约在这一时期也在\Place{巴勒斯坦}得到培育，神圣的\Person{依拉利雍}因而声名大噪。他出生于\Place{塔巴塔}，一个靠近\Place{迦萨}镇、朝向南方的小村庄，紧邻一条流入大海的溪流，当地人给溪流取了与村庄相同的名字。当他正在\Place{亚历山大里亚}学习文法时，他前往旷野拜见伟大的隐修士\Person{安当}，并在他陪伴下学会了同样的哲学。在那里逗留短暂时光后，他返回本国，因为由于涌向安当的人群，他无法如自己所愿那样安静。发现父母已故后，他将遗产分给兄弟和穷人，未为自己保留任何东西，便去住在靠近海边、距本村约二十斯塔迪翁的一片旷野中。他的斗室是一座非常小的房子，由砖块、木屑和碎瓦片建成，其宽度、高度和长度使得无人能在其中不低头站立，或不解脚躺着；因为他在一切事上力求使自己习惯于艰苦并克制奢侈安逸。在我们所认识的人中，无人能与他那无夸耀而受称赞的节制相匹敌。他抵挡饥渴寒暑及其他身体和灵魂的苦难。他行为恳切，言谈庄重，对圣经有良好的记忆和准确的理解。他如此为天主所爱，以至于至今许多受苦和附魔之人在他的坟墓上得治愈。值得注意的是，他最初被安葬在\Place{塞浦路斯}岛上，但他的遗骸现在安放在\Place{巴勒斯坦}；因为恰好他在塞浦路斯居住期间去世，当地人以极大的荣誉和尊敬安葬了他。但\Person{赫西加斯}，他最著名的门徒之一，盗走了遗体，运到巴勒斯坦，安放在自己的隐修院中。从此，居民们每年举行盛大而辉煌的庆典；因为在巴勒斯坦，人们有习俗将这种荣誉赐予那些以善行著称的人，例如\Person{奥勒留}、\Person{安特多尼}、\Person{阿历克西雍}（\Place{贝塔加通}人）和\Person{阿拉菲雍}（\Place{阿萨莱亚}人），他们在\Person{君士坦提乌}统治期间，虔敬而勇敢地实践哲学，并通过他们的个人德行，使仍处于异教迷信下的城邑和村庄中的信仰大为增长。\par

大约同一时期，\Person{犹利安}在\Place{厄德萨}附近实践哲学；他尝试了一种极其严厉和不涉形体的生活方式，以至于他似乎只有骨头和皮而没有肉。这段历史的陈述归功于叙利亚作家\Person{厄弗冷}，他写了犹利安生平的故事。天主亲自证实了人们对他的高度评价；因为天主赐予他驱逐魔鬼和治愈各种疾病的能力，无需借助药物，而仅仅通过祈祷。\par

除上述外，还有许多其他教会哲学者在\Place{厄德萨}和\Place{阿米达}地区以及被称为\Place{高加留}的山附近兴盛；其中有\Person{达尼尔}和\Person{西默盎}。但我现在不再谈叙利亚隐修士；以后若蒙天主允许，我将更全面地描述他们。\par

据说治理亚美尼亚的\Place{塞巴斯特}教会的\Person{欧斯塔修}在\Place{亚美尼亚}、\Place{帕夫拉戈尼亚}和\Place{本都}建立了一个隐修士团体，并成为严格纪律的创始人，涉及应食用或避免何种肉类、应穿着何种服装、以及应采纳何种习俗和精确的行为准则。有些人断言他是通常归功于\Place{卡帕多细亚}的\Person{巴西略}的苦修论著的作者。据说他的极度精确使他在某些方面陷于全然违反教会法律的过度行为。然而，许多人为他辩护，将责任推给一些门徒，他们谴责婚姻，拒绝与已婚者一起向天主祈祷，轻视已婚的司铎，在主日守斋，在私人房屋中举行集会，谴责富人完全没有天国之份，鄙视那些食用动物肉类的人。他们不保留惯常的束腰外衣和斯托拉作为服装，而是采用奇异而不寻常的装束，并做出许多其他革新。许多妇女受他们迷惑离开了丈夫；但由于无法坚持节欲，她们陷入了奸淫。另一些妇女以宗教为借口剃去头发，穿着男装，行为不当。\Place{帕夫拉戈尼亚}的首府\Place{甘格拉}附近的主教们聚集起来，宣布所有接受这些意见的人应与天主教会隔离，除非他们按照主教会议的决议弃绝上述每一种习俗。据说从那时起，\Person{欧斯塔修}将自己的服装换为斯托拉，并像其他司铎一样穿着旅行，从而证明他引入并实行这些新奇之事并非出于固执，而是为了虔诚的苦修。他以言论和生活的纯洁著称。说实话，他并非雄辩，也从未学习过雄辩术；但他有卓越的见识和很高的说服力，以至于他劝服了几个生活在淫乱中的男女开始节制而认真的生活。据说有一男一女按照教会的习俗献身于贞洁生活，被指控同居。他努力使他们停止交往；发现劝诫对他们无效时，他深深叹息说，有一次，一个合法结婚的妇女听他讲论节欲的好处，深受感动，自愿地停止与丈夫的合法同房；而另一方面，他说服能力的不足则被上述两人坚持非法行为的事实所证明。这些就是在上述地区开创隐修纪律的人。\par

虽然色雷斯人、伊利里亚人和其他欧洲民族在隐修团体方面仍缺乏经验，但他们并非完全没有献身于哲学的人。其中最著名的是\Person{玛尔定}，他出身于\Place{潘诺尼亚}的\Place{萨博利亚}的贵族家庭。他原本是一位著名的战士和军队统帅，但认为事奉天主是更光荣的职业，便接受了哲学的生活，首先住在\Place{伊利里亚}。在那里他热心地维护正统教义，反对亚略派主教的攻击，在被人们图谋陷害并屡次殴打后，被驱逐出境。于是他前往\Place{米兰}，独自居住。然而，由于该地区的主教\Person{奥克森修}（他并未坚定持守尼西亞信德）的阴谋，他很快被迫离开隐居之地，前往一个名为\Place{加莱纳里亚}的岛屿，在那里停留了一段时间，以植物根茎为生。\Place{加莱纳里亚}是\Place{第勒尼安海}中一个无人居住的小岛。\Person{玛尔定}后来被任命为\Place{塔拉西纳}（\Place{都尔}）教会的主教。他富有神恩，以至于使一个死人复活，并行了其他像宗徒们所做的一样奇妙的征兆。我们听说\Person{依拉略}，一个在生活和言谈上属神的人，大约在同一时期和同一地区生活；像\Person{玛尔定}一样，他因热心维护信德而被迫离开住所。\par

我现在已叙述了关于那些在虔诚和教会礼仪中实践哲学的人所能查知的内容。大约在同一时期，还有许多人在教会中因他们伟大的口才而著称，其中最卓越的是：在\Place{厄梅萨}执行司铎职务的\Person{欧瑟伯}；\Place{博斯特拉}的主教\Person{弟铎}；\Place{特穆伊斯}的主教\Person{塞拉皮翁}；\Place{安齐拉}的主教\Person{巴西略}；\Place{革尔马尼基}的主教\Person{尤多克修}；\Place{凯撒勒雅}的主教\Person{阿卡基}；以及掌管\Place{耶路撒冷}教区的\Person{济利禄}。他们所受教育的证明在于他们所撰写并留下的书籍，以及许多值得记述的事情。\par

% structure: p0049 source=h2 target=subsection reason=relative-heading-rank; base=h2

\subsection{第十五章　盲人\Person{狄狄摩}与异端者\Person{阿埃齐}}

\OriginalTerm{狄迪莫斯}{Didymus}，一位教会作家兼\OriginalTerm{亚历山大里亚}{Alexandria}圣学院院长，大约在同一时期闻名。他通晓每一门科学，熟悉诗歌与修辞学、天文学与几何学、算术以及各种哲学理论。他凭借自己的心智努力，借助听觉，获得了所有这些知识，因为他起初学习基础时就失明了。当他进入青年时期，表现出对言语和修养的强烈渴望，为此他常拜访这些学科的教师，但仅靠听觉学习，他进步如此迅速，以至很快理解了数学中困难的定理。据说他通过刻有字母的刻板用手触摸学会了字母；并通过注意力和记忆力的力量，以及仔细聆听声音，使自己熟悉了音节和单词。他的情况非常特殊，许多人专程到亚历山大里亚，目的就是听他讲论，或者至少是见他一面。他捍卫\OriginalTerm{尼西亚大公会议}{Nicæan council}教义的坚定态度，令\OriginalTerm{亚略派}{Arians}极为不悦。他通过劝服而非推理的力量，轻易使听众信服，并使每个人成为模糊点的判断者。他深受\OriginalTerm{天主教会}{Catholic Church}成员的追捧，获得埃及各种隐修会士和\Person{大安当}的称赞。\par

相传当\Person{安当}离开旷野前往\Place{亚历山大里亚}，为\Person{亚大纳削}的教义作证时，他对\Person{狄迪莫斯}说：\Quote{啊，\Person{狄迪莫斯}，被剥夺了老鼠、小鼠和最卑贱动物所拥有的视觉器官，并不是一件严酷的事，也不值得为此悲伤；但拥有如同天使的眼睛，籍此能够敏锐地默观神圣者，并准确地看到真知识，却是极大的福分。}在意大利及其属地，我已经提到过的\Person{欧瑟伯}和\Person{依拉略}，因在运用本国语言上的力量而杰出，据说他们关于信德和反对异端者的论著，备受赞扬地流传开来。据传，\Person{路济弗尔}是那以他名字命名的异端的创始人，并活跃于这一时期。\Person{阿埃齐乌}在异端者中也备受推崇；他是一位辩证家，擅长三段论，精于辩论，勤奋研究此类形式，但缺乏技巧。他如此大胆地推论天主的本性，以至许多人称他为\Quote{无神论者}。据说他原本是\Place{安提约基雅}的一位医生，因经常参加教会聚会，思考\WorkTitle{圣经}，结识了当时身为凯撒的\Person{加卢斯}；加卢斯非常尊崇宗教，珍爱其教授者。似乎由于阿埃齐乌通过这些辩论赢得了凯撒的尊重，他便更加专心致志于这些追求，以期在皇帝面前得到进展。据说他通晓亚里士多德哲学，并常去\Place{亚历山大里亚}教授此哲学的学校。\par

除了以上提到的个人外，教会中还有许多人有能力教导人民，并就\WorkTitle{圣经}的教义进行推理。若试图列举他们全部，将是一项过于艰巨的任务。我对上述异端的领袖或狂热者给予称赞，请不要感到奇怪。我钦佩他们的口才和言谈的说服力。我把他们的教义留给有权评判的人去判断。因为我不是被派来记录此类事务的，这也与历史不相称；我只叙述实际发生的事件，不添加自己的补充。关于当时在教育和言谈方面最为杰出、并使用罗马和希腊语言的人，我已在上文尽可能多地列举了据我所知的人。\par

% structure: p0053 source=h2 target=subsection reason=relative-heading-rank; base=h2

\subsection{第十六章：关于\Person{圣厄弗冷}}

叙利亚人\Person{厄弗冷}享有最高的荣誉，是\OriginalTerm{天主教会}{Catholic Church}最伟大的装饰。他出生于\Place{尼西比斯}，或其家族来自邻近地区。他献身于隐修哲学；尽管没有接受过教育，却出乎意料地精通叙利亚的学问和语言，以至他能够轻松理解最深邃的哲学定理。他的写作风格充满华丽的演讲、思想的丰富与适度，超越了最受认可的希腊作家。如果这些作家的作品被翻译成叙利亚语或其他语言，并舍弃了希腊语的美感，它们将保留不了多少原来的优雅和价值。厄弗冷的作品则没有这个缺点：它们在他生前就被翻译成希腊语，而且翻译至今仍在进行，但它们仍保留了许多原有的力量，以至于他的作品在阅读希腊文时，与阅读叙利亚文时一样令人钦佩。后来成为\Place{卡帕多细亚}首都主教的\Person{巴西略}是\Person{厄弗冷}的极大仰慕者，对他的博学感到惊讶。\Person{巴西略}被公认为他那个时代最有口才的人，我认为他的评价是对\Person{厄弗冷}价值的更有力的见证，超过了任何赞颂他的作品。据说他写了三十万行诗句，并且有许多弟子热切地追随他的教义。他最著名的弟子是\Person{阿巴斯}、\Person{则诺比乌斯}、\Person{亚巴郎}、\Person{玛拉斯}和\Person{西默盎}，叙利亚人和所有追求准确学问的人都极为称赞他们。\Person{保拉纳斯}和\Person{亚兰纳德}因其完美的演讲而受到赞扬，尽管据说他们偏离了纯正教义。\par

我不知道，在奥斯若恩（\Place{Osroëne}）曾有过一些非常博学之人，例如巴尔德撒纳斯（\Person{Bardasanes}），他创立了以其名字命名的异端，以及他的儿子哈尔摩尼乌斯（\Person{Harmonius}）。据说后者深谙希腊学识，并首次将本地语言纳入格律和音乐法则；他将这些诗句交给唱诗班，直到现在叙利亚人常常歌唱的，并非哈尔摩尼乌斯的原诗，而是同样的旋律。因为哈尔摩尼乌斯并未完全摆脱他父亲的错误，并对希腊哲学家所教导的灵魂、身体的生成与毁灭、以及重生持有各种观点，他将其中一些情感引入他所创作的抒情歌曲中。\newline{}当厄弗冷（\Person{Ephraim}）察觉到叙利亚人因文辞的优雅和旋律的节奏而着迷时，他开始担心，唯恐他们吸入同样的意见；因此，尽管他不懂希腊学术，他仍致力于理解哈尔摩尼乌斯的格律，并依照教会的教义创作了类似的诗篇，还创作了圣歌和赞美无欲之人的颂歌。从那时起，叙利亚人按照哈尔摩尼乌斯所确立的颂歌法则歌唱厄弗冷的颂歌。仅此著作就足以证明厄弗冷的天赋。他以其所行的善事和所坚守的严格纪律而闻名。他特别喜爱安静。他如此严肃并小心避免惹人毁谤，以至于连看妇女也避开。\newline{}据说曾有一位生活放纵的妇女，或是想试探他，或是受雇于他人，某次设法与他正面相遇，并定睛注视着他；他斥责她，命令她低头看地。\Quote{我为何要听从你的命令？}那妇人回答说，\InnerQuote{我不是从地里生的，而是从你生的。你若低头看你所从出的地，反倒更为合理；而我看着你，因为我从你而生。}厄弗冷被这妇人惊住，便将整个事件记录在一本书中，大多数叙利亚人将其视为他最好的作品之一。\newline{}也有人说，虽然他天生易怒，但自从他度隐修生活以来，从未对任何人表现出愤怒。有一次，他按惯例禁食数日后，他的侍从在递送食物时，失手打翻了盛食物的盘子。厄弗冷见他羞愧恐惧，便对他说：\Quote{鼓起勇气；我们到食物那里去，既然食物不来找我们。}他随即坐在碎片旁边，吃完了晚饭。\newline{}我下面要讲的事足以表明他完全不受虚荣之爱的羁绊。他被选为某城的主教，人们试图带他去祝圣。他一意识到意图，便跑到市场上，行为像疯子一样，蹒跚而行，拖着衣服，当众进食。那些来带他去作主教的人看到他这副模样，以为他疯了，便离开了；他趁这机会逃走了，躲藏起来，直到另一个人被祝圣代替他。\newline{}关于厄弗冷，我说这些就足够了，尽管他的同胞还讲述了许多关于他的其他轶事。然而，他在临终前不久的一次行为在我看来如此值得纪念，我将在此记录。当时厄德萨（\Place{Edessa}）城遭受严重饥荒，他离开了修道的斗室，责备富人任凭穷人饿死在他们周围，而不将他们的多余之物分给穷人；他通过他的哲学向他们指出，他们如此小心积攒的财富将变成对他们的谴责，并导致灵魂的毁灭，而灵魂比一切财富、身体和其他一切价值都更有价值；他证明他们因自己的行为而轻视了自己的灵魂。富人们敬畏此人和他的话语，回答说：\Quote{我们并非有意囤积财富，而是我们不知道可以托付谁来分配我们的财物，因为人人都贪图利益，会辜负所托。}厄弗冷问：\Quote{你们认为我怎样？}他们承认认为他是一位能干、卓越、善良的人，值得信赖，正如他的名声所确认的那样。于是他主动承担分配赈济的工作。他一收到钱，就在公共廊柱下安置了大约三百张床；在此他照料因饥荒而患病受苦的人，无论是外乡人还是周边地区的本地人。饥荒过后，他回到先前居住的斗室；数日后，他便去世了。他最终未获得高于执事的圣职，但他的美德并不逊色于那些被授予司铎之位、因善行的交谈和学识而受人钦佩的人。\newline{}我已经讲述了厄弗冷美德的一些事迹。若要完全描述他的品格以及与他在同一时期投身于哲学生活和生涯的其他杰出人物的性格，需要一位比我更有经验的手笔；有些事情甚至需要像他本人那样的作家。由于语言能力薄弱，以及对这些人及其事迹的无知，我无力胜任。有些人隐居于旷野；另一些人则生活在城市交往中，竭力保持平凡的外表，似乎与常人无异，操练德行，隐藏对自己的真实评价，以避免他人的赞扬。因为他们专注于未来赏报的交换，只让天主作他们思想的见证，不关心外在的荣耀。\par

% structure: p0056 source=h2 target=subsection reason=relative-heading-rank; base=h2

\subsection{第十七章 该时期的动态与皇帝及总司铎共同努力下基督信仰教义的进展。}

那些在此时期管理教会的人，以个人的品行著称，且可想见，他们所治理的百姓热切地依附于对基督的敬拜。因着司铎们及教会哲学家的热忱、德行和奇事，宗教每日进展；这些哲人吸引了异教徒的注意，并引导他们放弃迷信。当时在位皇帝在保护教会方面，与其父一样热忱，且赐予圣职人员、其子女及奴隶尊荣与免税。他们确认了其父所颁布的法律，并强制执行新法，禁止献祭、敬拜偶像或任何其他异教习俗。他们下令关闭一切神庙，无论位于城内或乡村。其中一些神庙被赠予教会，当所需用地或建筑材料时。他们以最大的关怀对待祈祷之所，那些因年久失修而破损的得以修缮，另一些则以极其壮丽的风格从根基重建。\Place{厄梅撒}的教堂最值得一观，且以其美丽著称。严令禁止犹太人购买任何不属于他们自己异端的奴隶。若他们违反此法，奴隶将被没收归公；但若他们对其施行犹太人的割礼，则刑罚为死刑及全部财产充公。因皇帝们渴望以各种方式促进基督信仰的传播，他们认为有必要阻止犹太人使那些祖辈信奉另一种宗教的人改宗，而那些怀有希望公开信奉基督信仰的人则被小心地为教会保留；因为基督宗教正是从异教群众中增长的。\par

% structure: p0058 source=h2 target=subsection reason=relative-heading-rank; base=h2

\subsection{第十八章 论\Person{君士坦丁}诸子所持守的教义。\Quote{同质}（homoousios）与\Quote{同似}（homoiousios）二词的区别。由此\Person{君士坦提乌斯}很快放弃了正确信仰。}

皇帝们从一开始就保存了他们父亲关于教义的观点；因为他们都偏爱尼西亚信条。\Person{君士坦斯}直到去世都坚持这些见解；\Person{君士坦提乌斯}有段时间也持类似看法；然而，当\Quote{同质}（\OriginalTerm{同质}{homoousios}）一词受到诽谤时，他放弃了先前的观点，但他并未完全停止承认圣子与圣父是同似的。\Person{优西比乌}的追随者及东部其他主教（他们以其言论和生命受人钦佩）——如我们所知——在\Quote{同质}（\OriginalTerm{同质}{homoousios}）与\Quote{同似}（他们以\Quote{\Emph{homoiousios}}一词来指称）之间做了区分。他们说，\Quote{同质}（homoousios）一词适当地属于有形体，如人和其他动物、树木和植物，其参与和起源在类似的事物中；而\Quote{同似}（\OriginalTerm{同似}{homoiousios}）一词则专属于无形体，如天主和天使，对其中每一个，都按其自身特有的本质形成概念。\Person{君士坦提乌斯}被这一区分所欺骗；虽然我确信他持守与其父及兄弟相同的教义，但他却采用了措辞上的变化，不再使用\Quote{同质}一词，而改用\Quote{同似}一词。我们前面提到的教师们主张，在使用术语时有必要如此精确，否则我们将有危险把无形体者想象为形体。然而，许多人视这一区分为荒谬的，\Quote{因为，}他们说，\Quote{由心智所领会的事物，只能以来自可见事物的名称来指称；只要在观念上没有错误，使用词语并无危险。}\par

% structure: p0060 source=h2 target=subsection reason=relative-heading-rank; base=h2

\subsection{第十九章 有关\Quote{同质}一词的更多细节。\Place{阿里米努姆}会议，其召开的方式、源泉和原因。}

君士坦提乌斯皇帝之所以被诱导采纳\OriginalTerm{相似同体}{homoiousios}一词，并不令人惊讶，因为许多顺从尼西亚大公会议教义的司铎都已接受此词。许多人互换使用这两个词来表达相同的含义。因此，在我看来，亚略派严重背离了真理，因为他们声称在尼西亚大公会议之后，许多司铎（其中包括优西比乌斯和狄奥格尼斯）拒绝承认圣子与圣父\OriginalTerm{同性同体}{homoousios}，而君士坦丁为此非常愤怒，以致将他们判处流放。他们说，后来通过一个梦或来自天主的异象，他的妹妹获知这些主教持有正统教义并遭受了不公正的对待；于是皇帝召回他们，并询问他们为何背离尼西亚教义，因为他们曾参与在那次会议上制定的关于信德的文件；他们回答说，他们并非出于确信而同意那些教义，而是出于恐惧：如果当时的争论持续下去，皇帝（他当时刚开始接受基督教，尚未受洗）可能会被驱使回到异教（这似乎是可能的），并迫害教会。他们断言，君士坦丁对这种辩护感到满意，并决定召集另一次大公会议；但由于死亡阻止了他实施这一计划，任务便落在了他的长子君士坦提乌斯身上。他对他表示，除非他能建立整个帝国的统一敬礼，否则拥有皇权对他毫无益处；而君士坦提乌斯，据说在父亲的鼓动下，在阿里米努姆召开了一次大公会议。这个说法很容易被看出是粗劣的捏造，因为会议是在希帕提乌斯和优西比乌斯担任执政官期间召开的，且是在君士坦提乌斯在其父去世后继承帝位二十二年之后。在二十二年的间隔期间，举行了多次会议，其中就\OriginalTerm{同性同体}{homoousios}和\OriginalTerm{相似同体}{homoiousios}这两个词进行了辩论。似乎没有人敢于否认圣子与圣父有相似的本体，直到埃提乌斯提出相反的意见，激怒了皇帝，以至于为了制止异端的蔓延，他命令司铎们在阿里米努姆和塞琉西亚聚集。因此，这次大会的真正原因并非君士坦丁的命令，而是埃提乌斯所挑起的问题。这一点将在我们之后所述中变得更加明显。\par

% structure: p0062 source=h2 target=subsection reason=relative-heading-rank; base=h2

\subsection{第二十章  亚他那修再次被君士坦提乌斯的信函恢复职位，并接受其教座。安提约基雅的大司铎。君士坦提乌斯向亚他那修提出的问题。圣咏中对天主的赞美。}

当君士坦斯得知萨迪卡所议决之事时，他写信给他兄弟，请求他将亚他那修和保禄的追随者恢复到他们自己的教会。由于君士坦提乌斯似乎犹豫不决，他再次写信，并以战争威胁他，除非他同意接纳这些主教。君士坦提乌斯与东方主教们商议此事后，认定因这事引发内战是愚蠢的。因此，他将亚他那修从意大利召回，并派出公共马车护送他返乡，还写了几封信请求他尽快返回。亚他那修当时住在阿奎莱亚，收到君士坦提乌斯的信件后，便前往罗马向尤利乌斯和他的朋友们告别。尤利乌斯以极大的友好表示与他道别，并给他一封致亚历山大神职人员和民众的信，信中称他是一个了不起的人，因他所经历的众多磨难而值得称颂，并祝贺亚历山大教会迎来了这样一位优秀的司铎的归来，劝勉他们跟随他的教导。\par

他随即前往叙利亚的安提约基雅，当时皇帝正驻跸于此。莱昂提乌斯治理该地区的教会；因为依斯塔提乌斯逃亡后，持异端思想的那些人夺取了安提约基雅的主教座。他们首先任命的主教是欧弗洛尼乌斯；继他之后是普拉凯图斯；随后是斯德望。后者因不配此尊位而被罢免，莱昂提乌斯获得了主教职。阿塔纳修斯视之为异端者而回避他，并与那些被称为\Quote{依斯塔提乌斯派}、在一所私宅中聚会的人共融。由于他发现君士坦提乌斯态度友善、乐于倾听，似乎皇帝会将他自己的教会归还给他，君士坦提乌斯在敌对异端首领的煽动下答复如下：\Quote{我准备履行我召回你时所承诺的一切；但你理当同样赐我一恩惠，即：你将你治下众多教会中的一座让给那些不愿与你共融的人。} 阿塔纳修斯回答说：\Quote{陛下，服从你的命令是极其公正和必要的，我决不违抗；但在安提约基雅城，有许多人避免我们与异端者之间的共融，我请求同样获准一座教会，使我们能在那里安全地聚会。} 由于阿塔纳修斯的请求在皇帝看来合理，异端者认为保持沉默更为明智；因为他们思量，他们的特殊见解在亚历山大里亚永远无法立足，因为阿塔纳修斯既能留住持相同见解的人，也能引导持相反见解的人；而且，如果他们让出安提约基雅的一座教会，人数众多的依斯塔提乌斯派就会聚集，届时可能尝试革新，因为他们能毫无风险地留住他们所掌握的人。此外，异端者察觉到，尽管教会的治理在他们手中，但全体圣职人员和百姓并非都遵从他们的教义。当他们按习惯分成两队咏唱赞美天主的圣诗时，每首颂歌结束时，各人宣告自己独特的见解。有些人赞美\Quote{圣父\Emph{与}圣子}，视二者在光荣中同等；另一些人则荣耀\Quote{圣父\Emph{藉着}圣子}，以此插入介词表示他们认为圣子低于圣父。当这些事情发生时，敌对派系的主教莱昂提乌斯，当时掌管安提约基雅主教座，不敢禁止按照尼西亚大公会议传统咏唱赞美天主的圣诗，因为他惧怕激起民众的骚乱。然而，据传他曾举手摸到自己已完全白发的头说：\Quote{当这雪融化时，将会有许多泥泞。} 他以此意指，他死后，不同的咏诗方式将引发大骚乱，他的继任者不会像他一样对民众表现同样的体恤。\par

% structure: p0065 source=h2 target=subsection reason=relative-heading-rank; base=h2

\subsection{第二十一章 君士坦提乌斯为阿塔纳修斯写给埃及人的信。耶路撒冷会议。}

皇帝在送阿塔纳修斯回埃及后，为他写信给该国的各位主教和司铎，以及亚历山大里亚教会的百姓；他证明其品行纯全、道德高尚，并劝勉他们心思一致，在他的指导下共同祈祷、事奉天主。他补充说，若有任何心怀恶意之人挑起骚乱，将按律法为此等罪行受罚。他还命令，他先前颁布反对阿塔纳修斯及其共融者的法令应从公共登记册中抹去，他的圣职人员应恢复先前享有的豁免权；将这些谕令发送给埃及和利比亚的总督。\par

阿塔纳修斯一抵达埃及，就撤换了他所知效忠于亚流主义的人，将教会治理和尼西亚大公会议的宣认交给他认可的人，并劝勉他们坚定持守。当时据说，他在游历其他国家时，若偶然访问亚流派管辖的教会，也实行同样的改变。他肯定被指控胆敢在他无权祝圣的城市施行祝圣礼。但由于他成功返回，尽管敌人不愿，且似乎不易受到怀疑，因他受皇帝君士坦斯的友谊尊敬，他比以前更受敬重。许多先前与他为敌的主教，尤其是巴勒斯坦的主教们，接纳他进入共融。当他当时访问后者时，他们友善地接待他。他们在耶路撒冷召开会议，马克西穆斯及其他人为他写了以下信件。\par

% structure: p0068 source=h2 target=subsection reason=relative-heading-rank; base=h2

\subsection{第二十二章 耶路撒冷会议为阿塔纳修斯所写的信函。}

神圣的耶路撒冷会议致埃及、利比亚和亚历山大里亚的司铎、执事及百姓，我们亲爱的、最珍视的弟兄：在主内安好。\par

\Quote{我们永远，亲爱的，无法为天主、万有的创造者，因祂现在所成就的奇妙工作，特别因祂藉着恢复你们的牧者与主、我们的同工亚大纳削而赐予你们各教会的恩惠，而献上相称的感谢。有谁曾希望亲眼见到这事，而你们如今在事实上正在实现呢？但诚然，你们的祈祷已被眷顾祂教会的宇宙的天主垂听，祂已垂念你们的眼泪和抱怨，并因此俯允了你们的请求。因为你们曾如无牧之羊般被分散、被撕裂。因此，那位从天上眷顾你们、关心祂自己羊群的真牧者，已将你们所渴望的那位归还给你们。看，我们为教会的和平行一切事，并受像你们这样的爱所感。因此，我们接纳并拥抱你们的牧者，并通过他与你们共融，我们发出这篇致辞和我们的感恩祈祷，好使你们知晓我们如何藉着爱的纽带与他及你们联合。你们理应为最蒙天主所爱的皇帝们的虔诚祈祷，他们既已认识到你们对他的渴望和他的纯洁，便决意以一切尊荣将他归还给你们。那么，伸出双手迎接他，并热心为那赐予你们这些恩惠的天主而向他献上必要的感恩祈祷；愿你们永远与天主同乐，并在基督耶稣我们的主内荣耀主，愿荣耀藉着祂归于圣父，直到永远。阿们。}\par

% structure: p0071 source=h2 target=subsection reason=relative-heading-rank; base=h2

\subsection{第23章　属于亚略派系别的瓦伦斯和乌尔撒基乌斯向罗马主教承认他们对亚大纳削提出了虚假指控。}

这是巴勒斯坦召开的主教会议所写的信。不久之后，亚大纳削满意地看到提洛会议对他所作判决的不公被公开承认。瓦伦斯和乌尔撒基乌斯曾与特奥格尼斯及其追随者一同被派往玛勒奥提斯，就如我们先前所提，为调查伊斯基里翁指控亚大纳削打碎的那只圣爵；他们向罗马主教尤利乌斯写了以下撤回书：\par

乌尔撒基乌斯和瓦伦斯致最可福的尤利乌斯主教大人。\par

由于我们先前，如同众所周知，多次在书信中针对主教亚大纳削提出各种指控，并且虽然我们一直受阁下有关此事之信函的催促，但我们公开指控却未能为我们的控告给出理由；因此，我们现在在众位司铎、我们的弟兄面前，向阁下承认：你们所听到关于上述亚大纳削的一切都是完全虚假、虚构的，并且全然违背他的本性。为此，我们欣然与他共融，尤其因为您的虔诚，按照您内在的良善之爱，已原谅了我们的错误。此外，我们向您声明：如果东方的诸位主教，甚至亚大纳削本人，任何时候恶意地召唤我们受审，我们绝不脱离您对此案的判断和处置。我们现在并永远憎恨，正如我们先前在米兰提交的备忘录中所做的那样，异端者亚略及其追随者，他们声称有一段时期圣子不存在，并说基督是从无中来的，并且否认基督在万世之前就是天主和天主子。我们再次以亲手所写声明，我们将永远谴责上述亚略异端及其始作俑者。\par

\Quote{我，乌尔撒基乌斯，以我自己的签名签署这认罪书。瓦伦斯亦同样。}\par

这就是他们写给尤利乌斯的认罪书。也有必要附上他们给亚大纳削的信，内容如下：\par

% structure: p0077 source=h2 target=subsection reason=relative-heading-rank; base=h2

\subsection{第24章　瓦伦斯和乌尔撒基乌斯致伟大的亚大纳削的调解信。其他东方主教恢复各自教座。马其顿尼再次被逐，保禄继任教座。}

主教乌尔撒基乌斯和瓦伦斯致我们在主内的弟兄亚大纳削。\par

\Quote{我们趁我们的弟兄和同工司铎穆塞乌斯前往您尊贵的处所之时，亲爱的弟兄，通过他向您从阿奎莱亚致以最丰盛的问候，并希望我们的信函发现您健康。如果您回复此信，您将给我们极大的鼓励。要知道，我们与您和平共融，并处于教会共融之中。}\par

因此亚大纳削在如此情况下从西方返回埃及。保禄、玛尔塞鲁斯、阿斯克莱帕斯和路济乌斯，因皇帝谕令而从流放中归来，接受了他们各自的教座。保禄一返回君士坦丁堡，马其顿尼便退隐，并私下举行教会聚会。在安基拉，当巴西略从那边的教会中被免职、玛尔塞鲁斯复职时，发生了大骚乱。其他主教则毫无困难地复职。\par

\SourceRecord{Translated by Chester D. Hartranft.；Nicene and Post-Nicene Fathers, Second Series；Vol. 2.；Edited by Philip Schaff and Henry Wace.；Buffalo, NY: Christian Literature Publishing Co.,；1890.；<http://www.newadvent.org/fathers/26023.htm>.}

% structure: p0001 source=h1 target=section reason=page-title

\section{\WorkTitle{教会史（卷四）}}

% structure: p0002 source=h2 target=subsection reason=relative-heading-rank; base=h2

\subsection{第一章 \Person{君士坦斯·凯撒}之死。在\Place{罗马}发生的事件。}

在\Place{萨迪卡}会议四年后，\Person{君士坦斯}在\Place{西高卢}被杀。策划谋杀他的\Person{玛格嫩提乌斯}将\Person{君士坦斯}的整个政府置于自己的控制之下。与此同时，\Person{维特拉尼奥}在\Place{锡尔米乌姆}被\Place{伊利里亚}军队拥立为皇帝。前任皇帝之妹的儿子\Person{内波提安}召集了一群角斗士，争夺皇权，而古\Place{罗马}在这些祸患中受害最深。然而，\Person{内波提安}被\Person{玛格嫩提乌斯}的士兵处死。\Person{君士坦提乌斯}发现自己是帝国唯一的主人，被宣告为唯一统治者，并迅速废黜了暴君。与此同时，\Person{亚大纳削}抵达\Place{亚历山大里亚}，准备召集\Place{埃及}主教会议，并确认了在\Place{萨迪卡}和\Place{巴勒斯坦}为他通过的法令。\par

% structure: p0004 source=h2 target=subsection reason=relative-heading-rank; base=h2

\subsection{第二章 \Person{君士坦提乌斯}再次驱逐\Person{亚大纳削}，并流放代表同质论者。\Place{君士坦丁堡}主教\Person{保禄}之死。\Person{马其顿尼}：他第二次篡夺该教座及其恶行。}

皇帝被异端者的诽谤所欺骗，改变了主意，违抗\Place{萨迪卡}会议的法令，流放了他先前恢复的诸位主教。\Person{马塞鲁}再次被废黜，\Person{巴西尔}重新获得了\Place{安基拉}的主教职。\Person{路济}被投入监狱，死在那里。\Person{保禄}被判永久流放，被解往\Place{亚美尼亚}的\Place{库库苏姆}，并在那里去世。然而，我始终无法查明他是否死于自然原因。仍有传闻说他是被\Person{马其顿尼}的追随者勒死的。他一被流放，\Person{马其顿尼}就夺取了他的教会的治理权；并且，在他在\Place{君士坦丁堡}所并入的几个隐修士修会的协助下，以及通过与他多位邻省主教的联盟，据说开始了一场针对持有\Person{保禄}意见之人的迫害。他首先将他们逐出教会，然后强迫他们与自己共融。许多人在争斗中受伤而死；一些人被剥夺了财产；一些人被剥夺了公民权；还有的人被用铁器在额头上烙印，以标记为可耻之人。皇帝听到这些事后很不高兴，并将这些事的罪责归于\Person{马其顿尼}及其追随者。\par

% structure: p0006 source=h2 target=subsection reason=relative-heading-rank; base=h2

\subsection{第三章 神圣公证者的殉道。}

迫害加剧，导致了流血事件。\Person{玛尔提里乌斯}和\Person{马尔西安}是被杀者中的两位。他们曾住在\Person{保禄}家中，被\Person{马其顿尼}交给总督，指控他们谋杀了\Person{赫尔莫根尼斯}，并挑起了先前反抗他的叛乱。\Person{玛尔提里乌斯}是副执事，\Person{马尔西安}是圣咏歌唱者和圣经诵读员。他们的坟墓很有名，位于\Place{君士坦丁堡}城前，作为殉道者的纪念；它被安置在一所祈祷室内，这祈祷室由\Person{若翰}开始建造，由\Person{西西尼乌斯}完成；这两人后来都主持了\Place{君士坦丁堡}教会。因为那些曾被不配地判定为无分于殉道荣誉的人，却被天主所尊荣，因为那处曾斩首被处死者、此前因鬼魂而无人敢近的地方，如今被洁净了，受魔鬼影响的人脱离了疾病，并且在那坟墓旁行了许多其他著名的奇迹。这就是关于\Person{玛尔提里乌斯}和\Person{马尔西安}应当陈述的细节。如果我所述似乎难以令人置信，则容易向那些更准确地了解情况的人寻求进一步信息；关于他们，也许有比我详细记录的更奇妙得多的故事。\par

% structure: p0008 source=h2 target=subsection reason=relative-heading-rank; base=h2

\subsection{第四章 \Person{君士坦提乌斯}在\Place{锡尔米乌姆}的征战，与关于\Person{维特拉尼奥}和\Person{玛格嫩提乌斯}的细节。\Person{加卢斯}接受\OriginalTerm{凯撒}{Cæsar}称号，并被派往\Place{东方}。}

在\Person{亚大纳削}被驱逐（发生在大约此时）之后，\Person{乔治}在整个\Place{埃及}迫害所有拒绝顺从自己意见的人。皇帝进军\Place{伊利里亚}，进入\Place{锡尔米乌姆}，\Person{维特拉尼奥}已按约到达那里。先前拥立他为皇帝的士兵突然改变主意，向\Person{君士坦提乌斯}欢呼为唯一君主和\OriginalTerm{奥古斯都}{Augustus}，因为皇帝及其支持者都竭力促成此举。\Person{维特拉尼奥}意识到自己被出卖，便作为恳求者投向\Person{君士坦提乌斯}脚前。\Person{君士坦提乌斯}确实怜悯他，但剥夺了他的帝国装饰和紫袍，迫使他回归平民生活，从国库中慷慨供应他的需要，并告诉他，对老人而言，远离帝国忧虑、过宁静生活更为适宜。在安排好对\Person{维特拉尼奥}的这些处理之后，\Person{君士坦提乌斯}派遣一支大军进入\Place{意大利}对抗\Person{玛格嫩提乌斯}。然后他将\OriginalTerm{凯撒}{Cæsar}称号授予他的表弟\Person{加卢斯}，并派他进入\Place{叙利亚}保卫\Place{东方}各行省。\par

% structure: p0010 source=h2 target=subsection reason=relative-heading-rank; base=h2

\subsection{第五章 \Person{西里尔}在\Person{马克西穆}之后主持司铎职务，而最大形状的十字架，光辉胜过太阳，再次出现在天空，并可见数日。}

在玛息莫之后治理耶路撒冷教会的济利禄时代，天上出现了十字架的标记。它光芒璀璨，并非像彗星那样光线分散，而是聚集了大量光线，看似稠密却又透明。其长度从哥耳哥达到橄榄山约有十五斯塔狄，宽度与长度成比例。如此非凡的现象引起普遍的恐惧。男女老幼离开家、市场或各自的工作，跑到教堂，在那里一起向基督咏唱赞美诗，并自愿宣认对天主的信仰。这消息立即惊动了我们的整个帝国，其传播迅速；因为按照惯例，有来自世界各地（可谓如此）的旅客住在耶路撒冷祈祷或参观名胜，他们目睹了这一标记，并将事实告知家乡的朋友。皇帝也得知了此事，部分通过当时流传的众多报告，部分通过主教济利禄的信件。据说这一征兆应验了圣经中包含的古代预言。它导致许多外邦人和犹太人皈依基督信仰。\par

% structure: p0012 source=h2 target=subsection reason=relative-heading-rank; base=h2

\subsection{第六章　西尔米乌姆主教福提诺及其异端，以及为此召开的西尔米乌姆公会议。三份信仰信条。这位空谈者被安西拉的巴西略驳斥。福提诺被罢免后，虽受邀请，却拒绝和解。}

大约此时，治理西尔米乌姆教会的福提诺向当时驻在该城的皇帝陈述了他早先发起的一种异端。他天生口才流利且有说服力，使许多人接受了他的想法。他承认有一位全能的天主，万物由祂的圣言创造，但不承认圣子的受生与存在是在万世之前；相反，他声称基督的存在源自玛利亚。这一观点一经传开，立即激起了东西方主教的愤慨，他们共同认为这是对各自信仰的革新，因为无论是维护尼西亚公会议教义的人，还是支持亚流派主张的人，都同样反对它。皇帝也厌恶这异端，并在其驻地西尔米乌姆召开了一次公会议。东方主教中，治理亚历山大教会的乔治、安西拉主教巴西略、阿勒图萨主教马尔谷出席了会议；西方主教中有穆尔萨主教华伦斯和精修者奥西乌。后者曾参加尼西亚公会议，此次却是勉强参与；此前不久，他因亚利安派的阴谋被判流放；亚利安派因相信若能说服或强迫这位普受钦佩和尊敬的人（如奥西乌）加入他们，就能加强其势力，便强求皇帝下令召他参加西尔米乌姆公会议。公会议召开的时间是在色尔吉与尼格里尼安任执政官期满后的第二年；这一年因暴君叛乱，东西方均无执政官。福提诺因被指控支持撒伯流和萨莫萨塔的保禄的错误而被该会议罢免。会议随后起草了三份额外的信仰信条，其中一份用希腊文写成，其余两份用拉丁文写成。但它们彼此不一致，也与先前任何教义说明在措辞或含义上不一致。希腊文信条并未说圣子与圣父同质或类似本质，但声明凡主张圣子无开始、或主张圣子来自圣父本质的扩展、或主张圣子与圣父结合而不服从圣父的人，均被开除教籍。在一份罗马信条中，禁止说圣子与圣父在神性本质（罗马人称本质）上同质或类似本质，因为这样的说法并不见于圣经，且超出人的理解与知识。它说，必须承认圣父在尊荣、尊严、神性以及由父之名所提示的关系上高于圣子；必须承认圣子如同一切受造物一样服从圣父，圣父无开始，圣子的受生除了圣父以外无人可知。据说，这份信条完成后，主教们意识到其中包含的错误，试图将其撤回并加以修正；皇帝威胁要惩罚那些保留或隐藏任何抄本的人。但一经公布，便无法完全压制了。\par

第三份信条与其它信条含义相同。它禁止使用\Quote{本质}一词，原因在于拉丁文所用术语，而希腊文术语因教父们用得太简单，且对许多无知民众成为绊脚石，因为它在圣经中找不到根据，\Quote{我们认为完全有理由弃绝使用它；并且我们命令所有提及神性时使用该词的做法一概省略，因为圣经中从未说圣父、圣子、圣神是同质的，其中写到\InnerQuote{位格}一词。但我们依照圣经说，圣子与圣父相似。}\par

关于信仰，这是在皇帝面前作出的决定。奥西乌起初拒绝同意。然而，强迫被采用；据说他年事已高，在挨打之下屈服了，并给予了同意和签名。\par

\Person{福提诺}被撤职后，主教会议认为，如果能设法说服他改变观点，或许不无裨益。但当主教敦促他，并承诺若他放弃自己的教义，赞同他们的信式，便恢复他的主教职位时，他并不顺从，反而向他们提出辩论挑战。为此，在指定的日子，主教们与皇帝委任主持集会的审判官一同聚集；这些审判官在口才与地位上皆属宫廷之首。\Place{安基拉}主教\Person{巴西略}被选来开始与\Person{福提诺}的论战。双方提问与答辩纷繁，并由速记员即刻记录，故论战持续良久；但最终胜利归于\Person{巴西略}。\Person{福提诺}被定罪并流放，但他并未因此停止扩张自己的教义。他用希腊文和拉丁文撰写并出版了许多著作，在其中力图证明除他本人观点外，其他一切见解均为谬误。关于\Person{福提诺}及其名称所属异端，我所要说的现已尽述。\par

% structure: p0017 source=h2 target=subsection reason=relative-heading-rank; base=h2

\subsection{第七章　暴君\Person{玛格嫩提乌斯}与背教者\Person{西尔瓦努斯}之死。犹太人在巴勒斯坦的骚乱。凯撒\Person{加卢斯}因涉嫌革命而被杀。}

与此同时，\Person{玛格嫩提乌斯}控制了古罗马城，并处死了许多元老与平民。他听说\Person{君士坦提乌斯}军队逼近，便退入高卢；双方在此多次交战，互有胜负。最终，\Person{玛格嫩提乌斯}战败，逃往\Place{穆尔萨}——那是高卢的一处要塞；他眼见士兵因战败而士气低落，便立于高处，试图重振其勇气。但士兵们虽照常以皇帝欢呼声迎接\Person{玛格嫩提乌斯}，并准备在其公开露面时呐喊，却暗中不约而同地改呼\Person{君士坦提乌斯}为皇帝，以取代\Person{玛格嫩提乌斯}。\Person{玛格嫩提乌斯}由此断定，天主并未注定他执掌帝国权柄，便力图从要塞撤往远方。但他遭\Person{君士坦提乌斯}军队追击，在名为\Place{色琉古山}之地被赶上；他独自逃脱，奔向\Place{卢格敦}。抵达后，他杀了自己的母亲和已被他立为凯撒的兄弟，最后自杀。不久，他的另一个兄弟\Person{德森提乌斯}也自尽。然而，公众骚乱并未平息；因为不久之后，\Person{西尔瓦努斯}在高卢自立为至尊，但随即被\Person{君士坦提乌斯}的将领处死。\par

\Place{狄奥凯撒里亚}的犹太人也侵扰了巴勒斯坦及邻近地区；他们拿起武器，企图摆脱罗马枷锁。凯撒\Person{加卢斯}当时在\Place{安提约基雅}，闻讯后派兵攻打，击败了他们，并摧毁了\Place{狄奥凯撒里亚}。\Person{加卢斯}因胜利而陶醉，不能承受自己的成功，竟图谋最高权力，杀死了财务官\Person{玛格努斯}与东方总督\Person{多米提安}，因他们向皇帝报告了他的叛乱。\Person{君士坦提乌斯}震怒，召他觐见。\Person{加卢斯}不敢抗命，遂启程。然而，他到达\Place{埃拉沃纳岛}时，被皇帝下令处死；此事发生在他执政第三年，\Person{君士坦提乌斯}执政第七年。\par

% structure: p0020 source=h2 target=subsection reason=relative-heading-rank; base=h2

\subsection{第八章　\Person{君士坦提乌斯}抵达\Place{罗马}。在\Place{意大利}召开会议。记述\Person{亚大纳削}大帝因亚略派阴谋所遭遇之事。}

僭主们死后，\Person{君士坦提乌斯}预料和平即将恢复，动乱即将止息，便离开\Place{西尔米乌姆}，返回\Place{古罗马}，享受战胜僭主后的凯旋殊荣。他还打算在\Place{意大利}召开一次会议，若有可能，使东方和西方的主教们在教义上达成一致。大约此时，\Person{尤利乌斯}去世了，他治理\Place{罗马}教会二十五年；\Person{利伯略}继任。那些反对\Place{尼西亚}会议教义的人认为这是一个有利的时机，可以诽谤他们已废黜的主教，并设法将他们赶出教会，指斥他们为错误教义的支持者和公共安宁的扰乱者；并且指控他们在\Person{君士坦斯}在世时，曾试图挑起皇帝之间的误解；的确，正如我们上文所述，\Person{君士坦斯}曾以战争威胁其兄弟，除非他同意接纳正统主教。他们的努力主要针对\Person{亚他那修}，对他怀有如此强烈的憎恶，以致即使他在\Person{君士坦斯}的保护下，并享有\Person{君士坦提乌斯}的友谊时，他们也未能掩饰其敌意。\Person{纳尔奇苏斯}（\Place{西利西亚}主教）、\Person{狄奥多}（\Place{色雷斯}主教）、\Person{欧根尼}（\Place{尼西亚}主教）、\Person{帕特罗菲鲁斯}（\Place{西托波利斯}主教）、\Person{默诺凡特}（\Place{厄弗所}主教）及其他三十位主教在\Place{安提约基雅}集会，致信各地所有主教，声称\Person{亚他那修}违反教会规章返回其主教职位，在任何会议上都未为自己辩护，只得到自己一派一些人的支持；他们劝诫众人不要与他共融，也不要写信给他，而要与已被祝圣接替他的\Person{乔治}共融。\Person{亚他那修}对这些举动嗤之以鼻；但他即将经历比以往更大的考验。\Person{玛格嫩提乌斯}一死，\Person{君士坦提乌斯}成为罗马帝国唯一的统治者，便全力以赴，诱使西方主教们承认：子与父是\OriginalTerm{相似本体}{like substance}。然而，在实施这一计划时，他首先并未诉诸强迫，而是试图通过劝说，使其他主教赞同东方主教们反对\Person{亚他那修}的决议；因为他想，如果能在这点上使他们意见一致，他就能轻易地妥善处理与宗教有关的事务。\par

% structure: p0022 source=h2 target=subsection reason=relative-heading-rank; base=h2

\subsection{第九章 米兰会议。亚他那修的逃亡。}

皇帝极急于在\Place{米兰}召开会议，但只有少数东方主教前往；有些人似乎以生病为由推辞出席；另一些人则因路途遥远和艰难而不来。然而，会议中约有三百多位西方主教。东方主教坚持认为\Person{亚他那修}应被判流放，并被逐出\Place{亚历山大里亚}；其他主教或因惧怕、或因欺骗、或因无知，同意了这项措施。只有\Person{狄约尼削}（\Place{阿尔巴}主教，\Place{意大利}首府）、\Person{欧瑟伯}（\Place{利古里亚}地区\Place{韦尔切利}主教）、\Person{保利诺}（\Place{特里尔}主教）、\Person{罗德努斯}和\Person{路济弗尔}反对这一决定；他们宣称不应以如此轻微的借口谴责\Person{亚他那修}，并且邪恶不会因他的谴责而终止；相反，那些支持关于天主性的正统教义的人会立即遭到迫害。他们指出，整个措施是皇帝与\OriginalTerm{亚略派}{Arians}共同策划的阴谋，旨在压制\OriginalTerm{尼西亚信仰}{Nicene faith}。他们的勇敢招致了立即流放的敕令，\Person{依拉略}也与他们一同被放逐。结果极其清楚地表明，\Place{米兰}会议是为了何种目的而召开的。因为不久后在\Place{阿里米诺}和\Place{塞琉西亚}召开的会议，显然是为了改变\Place{尼西亚}会议所确立的教义，正如我即将说明的。\par

\Person{阿塔纳修}得知宫廷中有人对他设下阴谋，便认为谨慎起见不要亲自去见皇帝，因为他知道那会危及生命，也认为这样做无济于事。然而，他从埃及主教中挑选了五位，其中包括\Place{图米斯}的主教\Person{塞拉皮翁}——一位以超凡圣德和雄辩能力著称的主教——派遣他们连同教会中的三位司铎前往当时在西方的皇帝那里。他们奉命尽可能设法安抚皇帝；如有必要，则反驳敌对派的诽谤；并采取他们自认为对教会和\Person{阿塔纳修}最有利的措施。他们起航后不久，\Person{阿塔纳修}收到皇帝的一些信件，召他入宫。\Person{阿塔纳修}和全体教会之民对此命令深感不安；因为他们认为，无论顺从还是违抗一位持异端信念的皇帝，都无法获得安全。然而，他决定留在\Place{亚历山大}，而信使未达成任何事便离开了该城。次年夏天，另一位皇帝的使者与各省总督一同抵达，受命催促\Person{阿塔纳修}离开该城，并对神职人员采取敌对行动。然而，当他看到教会之民充满勇气并准备拿起武器时，也未能完成使命便离开了该城。不久之后，驻扎在\Place{埃及}和\Place{利比亚}的所谓罗马军团军队开进了\Place{亚历山大}。据报\Person{阿塔纳修}藏在名为\Quote{特敖纳斯}的教堂内，军队指挥官和皇帝再次委托处理此事的\Person{希拉里}下令撞开教堂大门，从而闯入；但他们搜遍各处，并未在墙内找到\Person{阿塔纳修}。据说他是藉天主眷顾逃脱了这次和许多其他危险；天主先前已揭示此事；当他刚出去时，士兵们就占据了教堂大门，几乎要抓住他了。\par

% structure: p0025 source=h2 target=subsection reason=relative-heading-rank; base=h2

\subsection{第十章。\Person{亚略}派人对\Person{阿塔纳修}的种种阴谋，以及他藉天主眷顾从各种危险中逃脱。\Person{乔治}在驱逐\Person{阿塔纳修}后在\Place{埃及}所行的恶事。}

毫无疑问，\Person{阿塔纳修}深受天主喜爱，并蒙受预见未来的恩赐。还可以举出比我们刚才叙述的更奇妙的事实，以证明他深谙未来。在\Person{君士坦斯}生前，皇帝\Person{君士坦提乌斯}曾一度决心虐待这位圣人；但\Person{阿塔纳修}逃走，并藏匿在某位熟人处。他长期住在一个地下无光、曾用作储水池的居所里。除了一个看似忠心并服侍他的女仆外，无人知道他的藏身之处。然而，由于异端者急切地想活捉\Person{阿塔纳修}，似乎最终通过礼物或承诺腐蚀了那个仆人。但天主预先警告\Person{阿塔纳修}她的背叛，他便逃离了那个地方。女仆因对主人做虚假证言而受了惩罚，而主人则逃离了本国；因为对异端者来说，接待或藏匿\Person{阿塔纳修}并非小罪，而是被视为违抗皇帝明确命令的行为，是危害帝国之罪，并受到民事法庭的惩处。我听说\Person{阿塔纳修}另一次以类似方式得救。他又一次因同样原因被迫逃命；他沿尼罗河上行，打算退到\Place{埃及}更远的地区，但敌人得知他的意图并追踪他。因天主预先警告他将被追踪，他便告知同船之人，命令他们返回\Place{亚历山大}。当他顺流而下时，他的谋害者划船经过。他安全抵达\Place{亚历山大}，并成功藏身于城中众多相似的房屋中。他成功躲过这些和其他许多危险，导致外邦人和异端者指控他行巫术。据传，有一次他穿城而过时，一只乌鸦呱呱叫，一些在场的异端者嘲笑地问他乌鸦在说什么。他微笑着回答说：\Quote{它发出\OriginalTerm{明天}{cras}的声音，在拉丁语中意为\InnerQuote{明天}，它借此向你们宣布，明天对你们将不吉利；因为这意味着罗马皇帝明天将禁止你们庆祝你们的节日。}尽管\Person{阿塔纳修}的这个预言看似荒谬，却应验了；因为次日皇帝的命令被转发给总督，命令禁止外邦人在神庙中聚集举行惯常仪式，也不得庆祝他们的节日；于是外邦人所保留的最隆重、最盛大的节日被废除了。我所说的足以表明这位圣人拥有预言的恩赐。\par

\Person{阿塔纳修斯}以我们所描述的方式逃脱了那些试图逮捕他的人之后，他的神职人员和民众一度继续拥有教堂；但最终，\Place{埃及}总督和军队指挥官强行驱逐了所有支持\Person{阿塔纳修斯}观点的人，以便将教堂的管理权交给那些支持\Person{乔治}的人，当时\Person{乔治}即将抵达。不久后他到达该城，教堂被置于他的管辖之下。他以武力而非司铎的温和方式进行统治；他试图在民众心中制造恐怖，对\Person{阿塔纳修斯}的追随者进行残酷迫害，并囚禁和残害了许多男女，因此被视作暴君。由于这些原因，他招致了普遍的憎恨；民众对他的行为极为愤怒，冲进教堂，差点将他撕碎；在这种极端的危险中，他艰难逃脱，逃往皇帝那里。那些支持\Person{阿塔纳修斯}观点的人随后占据了教堂。但他们并未长久控制这些教堂；因为驻扎在\Place{埃及}的军队指挥官回来，将教堂归还给了\Person{乔治}的党羽。后来，一位帝国文书（属于公证员阶层）被派来惩罚骚乱的领导者，他拷打和鞭笞了许多市民。当\Person{乔治}稍后返回时，他比以往更加可怕，似乎也引起了比以前更大的厌恶，因为他煽动皇帝犯下许多恶行；此外，\Place{埃及}的隐修士们公开宣称他背信弃义、傲慢自大。这些隐修士的意见总是被民众采纳，他们的见证被普遍接受，因为他们以德行和生活哲学著称。\par

% structure: p0028 source=h2 target=subsection reason=relative-heading-rank; base=h2

\subsection{第十一章　\Place{罗马}主教\Person{利贝里乌斯}及其被\Person{君士坦提乌斯}流放的原因。他的继任者\Person{斐理克斯}。}

虽然我所记载的这些事件并非在\Person{君士坦斯}去世后同一时期发生在\Person{阿塔纳修斯}和\Place{亚历山大}教会身上，但我认为为了更清晰起见，有必要将这些事件全部合在一起叙述。\Place{米兰}会议毫无结果地解散了，皇帝将所有反对\Person{阿塔纳修斯}敌人计划的人判以流放。由于\Person{君士坦提乌斯}希望在教会内建立统一的教义，并使司铎们保持一致的意见，他计划召集所有宗教的主教召开一次西部会议。他意识到实施这一计划的困难，因为有些主教需要穿越广阔的陆地和海洋，但他并未完全失去成功的希望。当这个计划萦绕在他心头，他尚未准备凯旋进入\Place{罗马}时，他派人召来了\Place{罗马}主教\Person{利贝里乌斯}，并试图劝说他与随行司铎（其中包括\Person{欧多克修斯}）保持一致的意见。然而，\Person{利贝里乌斯}拒绝服从，并声明他决不会在这一点上让步，于是皇帝将他流放到\Place{色雷斯}的\Place{贝罗亚}。据称，\Person{利贝里乌斯}被流放的另一个借口是，他不愿退出与\Person{阿塔纳修斯}的共融，而是勇敢地对抗皇帝，因为皇帝坚持认为\Person{阿塔纳修斯}损害了教会，导致他两个哥哥中较年长者的死亡，并在\Person{君士坦斯}和他自己之间播下了仇恨的种子。由于皇帝重新恢复了各次会议（尤其是\Place{提尔}会议）所通过的所有针对\Person{阿塔纳修斯}的法令，\Person{利贝里乌斯}告诉他，不应重视那些出于仇恨、偏袒或恐惧而发布的法令。他要求各地区的主教都应签署\Place{尼西亚}会议所编纂的信经，并且那些因坚持该信德而被流放的主教应被召回。他建议，在这些问题得到纠正之后，所有主教都应自费，不配备公共交通工具或金钱，以免显得成为负担和破坏者，前往\Place{亚历山大}，对真相进行准确的检验，因为在那座城市比在其他地方更容易进行，因为受害者和施加伤害者以及指控的辩驳者都居住在那里。然后，他展示了\Person{瓦伦斯}和\Person{乌尔萨西乌斯}写给他在\Place{罗马}宗座的前任\Person{尤利乌斯}的信，他们在信中请求他的宽恕，并承认在\Place{马雷奥提斯}对\Person{阿塔纳修斯}提出的指控是虚假的；他恳求皇帝不要在\Person{阿塔纳修斯}缺席的情况下定他的罪，也不要相信那些显然是经由他敌人的阴谋而获得的法令。关于所谓的对他两个兄弟造成的伤害，他恳求皇帝不要借司铎之手报复，因为司铎是由天主设立，不是为了执行报复，而是为了圣化，并施行公正和仁慈的行为。\par

皇帝看出\Person{利贝里乌斯}不打算服从他的命令，便命令将他押送到\Place{色雷斯}，除非他在两天内改变主意。\Quote{对于我，皇帝啊，}\Person{利贝里乌斯}回答说，\Quote{无需考虑；我的决心早已形成并确定，我已准备好前往流放之地。}据说，当他被押往流放地时，皇帝送给他五百金币；但他拒绝接受，并对送钱来的使者说：\Quote{去，告诉送这金子的人，把它给他周围的谄媚者和伪君子，因为他们贪得无厌，使他们陷入永无止境的匮乏，永远无法缓解。基督，在各方面都与他的圣父相似，供应我们食物和一切美物。}\par

\Person{利贝里乌斯}因上述理由被从\Place{罗马}教会废黜后，其治理权转交给该地圣职人员中的一位执事\Person{费利克斯}。据说\Person{费利克斯}始终持守尼西亚信德，在宗教事务上行为无可指摘。惟一指控他的事是：在其被祝圣之前，他曾与异端者共融。当皇帝进入\Place{罗马}时，民众高声要求\Person{利贝里乌斯}，恳求其返回；皇帝与同他在一起的诸位主教商议后答复说，如果\Person{利贝里乌斯}同意接受与宫廷司铎所持相同的意见，他便召回\Person{利贝里乌斯}并归还给民众。\par

% structure: p0032 source=h2 target=subsection reason=relative-heading-rank; base=h2

\subsection{第十二章 叙利亚人\Person{艾狄乌斯}和\Place{安提约基雅}\Person{莱昂提乌斯}的继任者\Person{欧多克修斯}。论\Quote{同体}一词。}

约在此时期，\Person{艾狄乌斯}提出了他关于天主性的奇特见解。他当时是\Place{安提约基雅}教会的执事，由\Person{莱昂提乌斯}祝圣。他与\Person{亚略}一样，主张圣子是受造物，是从无中受造的，且与圣父不相似。由于他极其好辩，在神学主题上断言极为大胆，且倾向于使用极为微妙的论证方式，连那些与他持相同意见者都视他为异端者。当他因此被异端者开除教籍时，他假装拒绝与他们共融，因为他们不义地接纳了\Person{亚略}进入共融，而\Person{亚略}曾向\Person{君士坦丁}皇帝发誓说他持守\Place{尼西亚}大公会议的教义。这就是关于\Person{艾狄乌斯}的记述。\par

当皇帝在西部时，\Place{安提约基雅}主教\Person{莱昂提乌斯}去世的消息传来。\Person{欧多克修斯}请求皇帝允许他返回\Place{叙利亚}，以便管理该教会的事务。获准后，他急速前往\Place{安提约基雅}，在未获得\Place{劳迪塞亚}主教\Person{乔治}、\Place{阿勒图萨}主教\Person{马尔谷}、其他\Place{叙利亚}主教以及任何有权祝圣的主教的同意下，自行就任该城主教。据传他的行动得到皇帝及宫廷宦官们的赞同，这些人像\Person{欧多克修斯}一样，支持\Person{艾狄乌斯}的教义，相信圣子与圣父不相似。当\Person{欧多克修斯}占据\Place{安提约基雅}教会后，他敢于公开支持这一异端。他在\Place{安提约基雅}召集所有与他持相同意见的人，其中包括\Place{提洛}主教\Person{阿卡基乌斯}，并拒绝了\Quote{相似实体}和\Quote{同体}这两个词，借口说它们已被西部主教们谴责。因为\Person{奥西乌斯}与那里的一些司铎，为了遏止\Person{瓦伦斯}、\Person{乌尔萨基乌斯}和\Person{革尔马尼乌斯}挑起的争论，确实曾在\Place{西尔米乌姆}（虽是被迫，据传如此）同意避免使用\Quote{同体}和\Quote{相似实体}两词，理由是这样词并未出现在圣经中，且超出人的理解。他们寄了一封信给主教们，似乎这些主教支持\Person{奥西乌斯}在这点上所写的，并向\Person{瓦伦斯}、\Person{乌尔萨基乌斯}和\Person{革尔马尼乌斯}表达感谢，因为他们给了西部主教们正确观点的动力。\par

% structure: p0035 source=h2 target=subsection reason=relative-heading-rank; base=h2

\subsection{第十三章 \Person{欧多克修斯}在\Person{乔治}（\Place{劳迪塞亚}主教）信件中被指责的革新。\Place{安基拉}会议派往\Person{君士坦提乌斯}的代表团。}

\Person{欧多克修斯}引入这些新教义后，\Place{安提约基雅}教会中许多反对者被开除教籍。\Place{劳迪塞亚}主教\Person{乔治}给了他们一封信，让他们带给那些应\Person{巴西略}之邀从邻近城镇来到\Place{加拉提雅}的\Place{安基拉}参加他建造的一座教堂祝圣典礼的主教们。这封信如下：——\par

\Person{乔治}致他最尊敬的诸位主上\Person{马其顿尼乌斯}、\Person{巴西略}、\Person{塞克罗皮乌斯}和\Person{尤金纽斯}，在主内问安。\par

\Quote{几乎整座城池都因\Person{艾狄乌斯}的船难而受苦。这个恶人的门徒——你们曾鄙视他们——被\Person{欧多克修斯}所鼓励，并被他提拔到圣职岗位，而\Person{艾狄乌斯}本人也被提升到最高荣誉。因此，请前去援助这座大城，免得因它的船难而全世界都被淹没。你们要聚集起来，征求其他主教的签名，好使\Person{艾狄乌斯}被逐出\Place{安提约基雅}教会，并使那些已被\Person{欧多克修斯}预先操纵进入圣职人员名册的门徒被去除。如果\Person{欧多克修斯}坚持与\Person{艾狄乌斯}一同断言圣子与圣父不相似，并偏袒支持这一教义的人而非拒绝者，那么\Place{安提约基雅}城就失于你们了。}这就是\Person{乔治}书信的内容。\par

聚集在\Place{安基拉}的主教们从\Person{欧多克修斯}在\Place{安提约基雅}的法令中清楚地看出，他图谋在教义上引入革新；他们将此事告知皇帝，并恳求他确认在\Place{萨尔底卡}、\Place{西尔米乌姆}及其他会议所确立的教义，特别是圣子与圣父相似的教义。为了向皇帝提出这个请求，他们派出由以下主教组成的代表团：\Place{安基拉}主教\Person{巴西略}、\Place{塞巴斯特}主教\Person{优斯塔修斯}、\Place{库齐库斯}主教\Person{厄琉修斯}，以及内廷司铎\Person{莱昂提乌斯}。他们抵达皇宫时，发现\Place{安提约基雅}的一位司铎\Person{阿斯发利乌斯}——他是\Person{艾狄乌斯}异端的狂热分子——即将离开，他已办完他此行的事务，并已从皇帝那里获得一封信。然而，当\Person{君士坦提乌斯}收到来自\Place{安基拉}代表团关于该异端的情报后，他谴责了\Person{欧多克修斯}及其追随者，撤销了托付给\Person{阿斯发利乌斯}的信，并写了以下这封信：——\par

% structure: p0040 source=h2 target=subsection reason=relative-heading-rank; base=h2

\subsection{第十四章 \Person{君士坦提乌斯}皇帝反对\Person{欧多克修斯}及其同党的信件。}

\Emph{君士坦提乌斯·奥古斯都}凯旋者，致\Place{安提约基雅}的圣教会。\par

\Quote{欧多克修斯来时不具我们的权柄；无人应假设他有权柄，因我们远非以恩惠看待此类人。若他们在此类交易中对他人施以欺诈，便证明他们将在更高之事上精炼真理。因为那些为权力之故，绕城而行，如游荡者从一城跳至另一城，窥探每一角落，为贪求更多之人，他们自愿放弃什么呢？据说这些人中有某些江湖骗子和诡辩家，其名几乎不可容忍，其行为邪恶且最不敬。你们皆知我所指之人；因你们完全熟悉阿埃蒂乌斯的教义及其所培养的异端。他与随从专务败坏百姓之业；这些聪明之徒竟敢公开声称我们赞同他们的祝圣。他们如是散布谣言，仿效多言者之方式；但这不真实，且远离真理。请回忆我们首次宣认信仰时所使用的言辞；因我们承认我们的救主是天主子，与圣父相似实体。但这些人敢于将关于天主性的任何想象随意提出，离无神论不远；且他们竭力在他人中传播其见解。我们确信其不法行径将自食其果。同时，将其逐出会议与普通商议已足够；因我现在不提及将来必须降临于他们的惩罚，除非他们停止疯狂。他们聚拢最邪恶之人，仿佛奉旨而行，并挑选异端领袖为神职人员，从而贬低可敬的圣秩，如同他们可任意妄为，这祸患何其大！谁能容忍这些人以不敬充满城市，在遥远地区隐匿污秽，且仅以伤害义人为乐？这是何等邪恶的合一，跛行着登上神圣宝座！如今是那些汲取真理者进入光明之时，先前因恐惧受抑，现愿摆脱因袭之人，请步入中央；因这些恶人的诡计已被彻底驳斥，无法发明任何伎俩使其脱离不敬之行。善人之责在于持守祖传信德，并可说拓展之，而不必操心他事。我恳切劝勉那些虽刚逃离此异端悬崖之人，赞同那些聪明于神圣学问的主教们为更善所正确定定的法令。}\par

因此我们看见，通常称为\OriginalTerm{不相似派}{Anomian}的异端，在这个时期几乎就要占主导地位了。\par

% structure: p0044 source=h2 target=subsection reason=relative-heading-rank; base=h2

\subsection{第15章　君士坦提乌斯皇帝前往西尔米乌姆，召回利贝里乌斯，并将其恢复至罗马教会；他也命令菲利克斯协助利贝里乌斯执行司铎职务。}

这些事件过后不久，皇帝从罗马返回西尔米乌姆；在接见西方主教的代表团后，他从贝罗亚召回利贝里乌斯。君士坦提乌斯在东方主教的代表及营中其他司祭面前，催促利贝里乌斯承认圣子与圣父非同一实体。此举受到巴西略、尤斯塔提乌斯和尤西比乌斯的煽动，他们对皇帝影响甚大。他们汇编了一份文件，将反对萨莫萨塔的保禄及西尔米乌姆主教福提诺的法令集合在一起；并附上一份在安提约基雅举行教堂奉献礼时制定的信条格式，仿佛有人以\OriginalTerm{同质}{consubstantial}一词为借口，企图建立自己的异端。利贝里乌斯、阿塔纳修斯、亚历山大、塞维里亚努斯及非洲司祭克雷森斯被说服同意此文件，同样还有乌尔萨西乌斯、西尔米乌姆主教革尔曼尼乌斯、穆尔萨主教瓦伦斯及在场所有东方主教。他们部分赞成了利贝里乌斯所拟的一份信仰宣认，其中他声明，凡主张圣子在实体及一切其他方面与圣父不相似者，均予绝罚。因为当安提约基雅的欧多克修斯及其同党（支持阿埃蒂乌斯异端）收到奥西乌斯的信时，他们散布谣言，说利贝里乌斯已放弃\OriginalTerm{同质}{consubstantial}一词，并承认圣子与圣父不相似。西方主教们制定这些法令后，皇帝允许利贝里乌斯返回罗马。当时在西尔米乌姆集会的主教们致信治理罗马教会的菲利克斯及其他主教，要求他们接纳利贝里乌斯。他们指示二人应共享宝座，以和谐之心共同履行司祭职责；且无论菲利克斯的祝圣或利贝里乌斯的放逐中可能有何违法之事，均予遗忘。罗马人民视利贝里乌斯为极优秀之人，因他反对皇帝时所表现之勇气而高度敬重他，以致为他发起骚乱，甚至流血。菲利克斯仅存活不久；利贝里乌斯便独自执掌教会。此事无疑是天主所定，使伯多禄的座位不致因两位主教同居而受辱；因如此安排是分裂之兆，且违反教会法。\par

% structure: p0046 source=h2 target=subsection reason=relative-heading-rank; base=h2

\subsection{第16章　皇帝因阿埃蒂乌斯异端与安提约基雅的革新，意图在尼科美底亚召开大公会议；但因该城发生地震及其他许多事务干扰，会议先于尼西亚，后于阿里米努姆和塞琉西亚召开。关于精修者阿尔萨西乌斯的记述。}

这些就是发生在\Place{西尔米乌姆}的事件。在此期间，由于害怕触怒皇帝，\Place{东方}教会和\Place{西方}教会似乎已经在同一教义的宣认上联合了。皇帝已决定在\Place{尼西亚}召开一次公会议，以审议在\Place{安提约基雅}引入的革新，以及阿埃蒂乌斯异端。然而，由于\Person{巴西略}及其同党反对在该城举行公会议，因为此前那里曾就教义问题发生过争执，于是决定在\Place{彼提尼雅}的\Place{尼科美底亚}举行公会议；诏令已发出，召集各国最聪明、最有口才的主教们，在指定日期准时到达，以便国家所有司铎都有幸参与这次主教会议，并出席其决议。这些主教中已有许多人起程，这时传来\Place{尼科美底亚}遭遇灾难的消息，说天主已将整座城市震为平地。由于城市毁灭的传说到处流传并愈演愈烈，主教们便停止了行程；因为正如这类情况中常见的那样，远方的人所听到的传闻远比实际发生的情况严重得多。据传，\Place{尼西亚}、\Place{佩林图斯}及其邻近城市，甚至\Place{君士坦丁堡}，都遭受了同样的灾难。正统主教们对此感到极为悲痛；因为宗教的敌人趁着一座宏伟教堂倒塌之机，向皇帝进言说，许多主教、男人、女人和儿童逃入教堂希望得救，结果全部丧生。这一报道并不真实。地震发生在白天第二时辰，那时教堂里并没有集会。唯一遇难的主教是\Place{尼科美底亚}主教\Person{凯克罗皮乌斯}和一位来自\Place{博斯普鲁斯}的主教，而且致命意外发生时他们都在教堂外面。城市在瞬间被震动，所以人们即使有逃亡的愿望，也没有能力寻求安全；在遭遇危险的第一刻，他们要么得救，要么就当场死亡。\par

据说这场灾难是由\Person{阿尔萨西乌斯}预言的。他是\Place{波斯}人，原是一名士兵，负责照料皇帝的狮子；但在\Person{李锡尼}统治期间，他成了著名的明认信仰者，并离开了军队。随后他来到\Place{尼科美底亚}的城堡，在城墙内过着隐修哲学的生活。在这里，他从天上得到了一个异象，奉命立即离开该城，以便从即将发生的灾难中得救。他极其急切地跑到教堂，恳求圣职人员向天主献上祈祷，求祂平息愤怒。但他发现他们不仅不相信他，反而把他当作笑柄，认为他在揭露未曾预料到的苦难，于是他便回到自己的塔楼，俯伏在地祈祷。就在这一刻，地震发生了，许多人丧生。那些幸免于难的人逃往乡间和旷野。正如在一座繁荣的大城市中常发生的那样，每户人家的火盆和灭火器、浴场的炉灶以及所有在工艺中使用火的人的熔炉中都燃着火；当房屋框架倒塌时，火焰被杂物围困，当然还有干燥的木材混杂其中，其中许多是含油的——这助长了迅速蔓延的大火，无限制地滋养着火焰；火焰四处蔓延，吞噬一切周围的物质，使整座城市，可以说，成为一片火海。由于无法进入房屋，那些从地震中逃生的人涌向城堡。\Person{阿尔萨西乌斯}被发现死在未倒塌的塔楼里，俯伏在地，保持着他开始祈祷时的姿势。据说他曾恳求天主让他死去，因为他宁愿死，也不愿目睹一座他初次认识基督并实践隐修哲学的城市遭受毁灭。既然我已被引导去谈论这位善人，那么不妨提及天主赐予他驱魔和净化受魔鬼困扰之人的能力。一次，一个附魔的人手持一把光剑跑过市场。人们纷纷逃离，全城一片混乱。\Person{阿尔萨西乌斯}迎上前去，呼求基督的名号，就在那名号下，魔鬼被驱逐，那人恢复了理智。除了上述事迹外，\Person{阿尔萨西乌斯}还行了其他许多超越人能力和技艺的事。有一条龙，或其他某种爬虫，盘踞在路边的一个洞穴里，用它的气息杀死过往行人。\Person{阿尔萨西乌斯}来到那里，开始祈祷，那蛇自动从洞里爬出，用头撞地，杀死了自己。所有这些细节，我是从那些听见过见过\Person{阿尔萨西乌斯}的人叙述的人那里得知的。\par

主教们因获悉尼科美底亚发生灾祸的消息而不敢继续行程，有的等待皇帝进一步命令，有的则在书信中表明自己关于信仰的看法。皇帝迟疑不决，不知应采取何种措施，便写信咨询巴西略是否应召开会议。巴西略在回信中赞扬了皇帝的虔诚，并以圣经中的事例安慰他因尼科美底亚遭毁而带来的伤痛；他劝勉皇帝为宗教起见应尽快召开会议；不要放弃这一表现宗教热忱的机遇，也不要把为此已聚集并启程前来的司铎们遣散，直至处理完某些事务。他还建议会议改在尼凯亚而非尼科美底亚举行，以便争议之点能在最初被质疑的原地得到最终裁决。巴西略写此信时相信皇帝会乐于接受此提议，因为皇帝本人最初也曾建议在尼凯亚举行会议。皇帝收到巴西略的信后，命令在夏初时主教们（除身体有疾者外）聚集于尼凯亚；后者则应派遣司铎与执事陈述意见，并就教义争议共同商议，对一切分歧作出相同决定。他规定从西方教会选出十名代表，从东方教会也选出十名代表，负责审议可能颁布的法令，判断其是否符合圣经，并总体监督会议事务。经过进一步商议，皇帝颁布法令，在决定会议地点并收到通知前往之前，主教们应留居原地或各自教会。随后他写信给巴西略，指示他致函东方主教们，询问他们认为应在何处举行会议，以便在初春公开宣布；因为皇帝认为因该省近期发生地震，在尼凯亚召开会议并不明智。巴西略致信各省主教，敦促他们共同商议，尽快决定举行会议最适宜的地点，并在信前附上皇帝信函的副本。正如在类似情况下常见的那样，主教们对此问题意见分歧，巴西略便前往皇帝当时所在的西尔米乌姆。他在该城遇到几位因私事前来的主教，其中包括阿雷图萨的马可，以及被任命主持亚历山大里亚教会的乔治。最终，瓦伦斯及其追随者决定会议在伊苏里亚的城市塞琉西亚举行；当时瓦伦斯正居留于西尔米乌姆，因他们支持亚诺米派的异端，便敦促在宫廷的主教们签署一份已拟定的信仰陈述书，其中未提及\Quote{实体}一词。然而，当召开会议的准备积极进行时，欧多克修斯、阿卡西乌斯、乌尔萨西乌斯和瓦伦斯等人考虑到许多主教仍忠于尼凯亚信经，另一些则倾向于安提约基亚教堂奉献礼所拟定的陈述书，但双方都保留使用\Quote{实体}一词，并主张圣子在各方面与圣父相似；他们意识到若两方聚集一处，必将轻易谴责埃提乌斯的学说与其各自信经相悖，便设法使西方主教们聚集于阿里米努姆，而东方主教们则聚集于伊苏里亚的塞琉西亚。他们想，说服少数人比说服多数人容易，因此认为分别处理或许能引导双方赞成自己的观点，或至少能成功说服一方，以免其异端遭受普遍谴责。他们通过一位名叫优西比乌的太监——他是皇宫总管，与欧多克修斯交好并持相同教义，且许多掌权者都试图讨好这位优西比乌——实现了这一目的。\par

% structure: p0050 source=h2 target=subsection reason=relative-heading-rank; base=h2

\subsection{第十七章 阿里米努姆会议的进展。}

皇帝被说服，认为将全体主教召集到同一地点举行会议，对公家而言耗费巨大，对主教们而言路途遥远，皆非适宜。于是他写信给当时在阿里米努姆的主教们，以及当时在塞琉西亚的主教们，指示他们着手调查关于信仰的争议问题，然后关注耶路撒冷主教济利禄及其他主教对不公正的免职和放逐令提出的申诉，并审查针对其他主教所作的各种判决的合法性。事实上，当时有多项控告针对不同主教悬而未决。埃及人控告乔治掠夺与暴行。最后，皇帝命令从每个会议中选出十名代表到他那里，报告各自会议的进展。\par

按照诏书，主教们在指定的城市聚集。里米尼公会议首先开始议事；约有四百名成员组成。那些最敌视亚大纳削的人认为没有什么进一步可针对他制定的法令。当他们开始探讨教义问题时，瓦伦斯和乌尔萨基乌，在革尔美纽、欧克森提乌、卡尤和德莫菲卢的支持下，走到集会中央，要求压制所有先前编纂的信经格式，只保留他们不久前在西尔米乌姆用拉丁文提出的格式。在这个格式中，根据圣经，教导说圣子与圣父相似；但完全没有提到天主的本质。他们宣称这个格式已获得皇帝批准，公会议有责任采纳它，而不是过于审慎地征求每位成员的个人意见，以免如果这些术语被交付辩论和精确论证，就会产生争议和分裂。他们还补充说，这能使那些不太善于言辞的人更正确地理解天主，胜过于引入新奇术语，这些术语近乎争辩式的戏法。通过这些陈述，他们旨在谴责使用\Quote{\OriginalTerm{同质}{consubstantial}}一词，因为他们说这个词在圣经中找不到，并且对大众来说晦涩难懂；他们希望用\Quote{圣子与圣父在一切事上相似}这一表述来替代，这一表述有圣经支持。在宣读了包含上述陈述的格式之后，许多主教对他们说，不应制定任何新的信经格式，先前编纂的格式已足够应付一切目的，他们聚会正是为了阻止所有革新。这些主教随后敦促那些编纂并宣读格式的人公开谴责亚流派教义，认为它是引起各地教会动荡的根源。乌尔萨基乌、瓦伦斯、革尔美纽、欧克森提乌、德莫菲卢和卡尤对此抗议，公会议命令宣读其他异端的陈述，以及尼西亚所颁布的陈述；以便那些支持各种异端的格式被谴责，而那些符合尼西亚教义的格式获得认可；以便不再有进一步争论的理由，也不必再召开公会议，而能做出有效的决定。他们指出，编制这么多格式是荒谬的，好像他们刚刚开始认识信仰，好像他们想轻视教会的古老传统，这些传统曾由他们自己及其前任治理教会，其中许多人曾作出美好的宣认，并获得了殉道的冠冕。这些就是主教们提出的论据，以证明不应尝试任何革新。由于瓦伦斯和乌尔萨基乌及其同党拒绝被这些论据说服，而是坚持倡导采纳他们自己的格式，他们被罢免，并决定拒绝他们的格式。有人指出，该格式开头部分声称是在君士坦提乌面前、在\Quote{永恒的奥古斯都}在场时、在优西比乌和希帕提乌执政期间于西尔米乌姆编制的，这是荒谬的。亚大纳削在给他朋友的信中评论了同样的事情，他说称君士坦提乌为永恒皇帝，却不敢承认天主子是永恒的，这是可笑的；他还嘲讽了该格式标注的日期，好像是要谴责前时代的信仰以及在此之前已接受信仰的人。\par

在里米尼事件之后，瓦伦斯和乌尔萨基乌及其追随者因被罢免而恼怒，急忙赶到皇帝那里。\par

% structure: p0054 source=h2 target=subsection reason=relative-heading-rank; base=h2

\subsection{第十八章 从里米尼公会议写给君士坦提乌皇帝的信}

公会议选出二十名主教，派他们作为使团前往皇帝处，携带以下信件，该信已从拉丁文译为希腊文：——\par

我们相信，是天主之命令，也是陛下之虔敬之安排，使我们从西方各城聚集到\Place{阿里米农}，旨在宣明天主教会之信德，并揭发那些反对它而提出异端者。经过长期的考察，我们得出结论：最好持守那自古延续、由先知、福音作者、我们主耶稣基督的宗徒们所宣讲的信德，祂是您帝国的守护者、您力量的保护者，并坚持到底，直到终点。若对\Place{尼西亚}主教们在最杰出的皇帝、您的父亲\Person{君士坦丁}之同意下所如此正确公正地提出的教义作任何改变，将是荒谬且不合法的；其教导和思想已传出并在普世的听闻和默想中被宣讲；它是亚略异端的对抗者和毁灭者；通过它，不仅那偏离信德者，而且所有其他异端都被摧毁。对这些教义有所增减是极其危险的；其中任何一条都不可作丝毫改动，否则就会给敌人机会为所欲为。\Person{乌尔撒基乌}和\Person{瓦伦斯}，在被怀疑参与和商讨亚略教义之后，曾被割断与我们的共融。为了恢复共融，他们承认了错误并获得了宽恕，正如他们自己的书信所证明的，通过这些书信他们得以豁免并获得赦免。宽恕谕令颁布之时，是在\Place{米兰}会议期间，当时罗马教会的司铎也在场。\par

既然我们知道在\Place{尼西亚}所制定的信德格式是在堪受记念的\Person{君士坦丁}面前，以最大的谨慎和精确编纂而成的，他一生持守之，并在他的圣洗和当他在天堂享受应得的平安时，我们认为企图更改它，并忽视如此众多的圣洁宣信者、殉道者以及这教义的文章和作者，他们对此付出很多思考，并延续了天主教会的古老法令，那是荒谬的。天主通过我们的主耶稣基督，将他们的信德知识传至您所生活的时代，因着祂您统治并治理世界。这些可悲的、依我们看来值得哀叹的人，又以不法的鲁莽宣布自己为不虔敬观点的传令者，并企图推翻整个真理体系。因为根据您的谕令，会议被召集，这些人暴露了他们自己欺骗的观点；他们试图用一种诡计和扰乱引入革新，并找到了一些同伴，他们为此邪恶交易捕获了这些人：即\Person{革尔玛尼乌}、\Person{奥克森提乌}和\Person{加尤}，他们引起了争辩和不和。这些人的教导虽然是一致的，却超出了所有亵渎的范围。当他们意识到自己终究不属于同一异端，而且在他们邪恶建议的任何一点上想法都不一致时，他们转向我们的信经，以便它看起来像是另一份文件。时间虽短，但足以驳斥他们的意见。为了教会事务不因他们而失事，并约束那将一切抛来抛去的动乱和骚动，看来稳妥的做法是持守古代不可动摇的定界，并将上述人士从与我们共融中驱逐。为此，我们派遣了我们的重新受命的代表到您仁爱座前，并让他们带信申明会议的感情。这些代表尤其受我们委任，要持守奠基者们正确定义的真理，并向陛下圣洁教导\Person{瓦伦斯}和\Person{乌尔撒基乌}声称的谬误，他们说对正义真理作些改变会在教会中产生和平。因为那些破坏和平的人如何能产生和平？他们更可能将争辩和骚乱引入其他城市和罗马教会。因此我们恳求陛下以和蔼的听闻和温和的面容考虑我们的代表，并不要让死者因任何新奇改变而受辱。我们祈求您允许我们留在我们从祖先接受的定界和法令中，我们要说，祖先们以准备好的心智、智慧并借着圣神完成了他们的工作。因为这些革新不仅使信徒陷入不信，也欺骗不信者走向不成熟。我们还恳求您命令那些目前离开自己教会的主教——其中一些受年老体弱之苦，另一些受贫穷匮乏之苦——提供他们返回家乡的途径，以便教会不再长期失去他们的牧职。\par

再者，我们恳求您，不要从前次决定中取走任何东西，也不要增加任何东西；让一切保持不变，正如从您父亲的虔敬直到现在所保存的那样；以便我们将来不再疲劳，被迫成为我们自己堂区的外人，而是主教和人民能在和平中共处，并能够致力于为您的个人救恩、帝国及和平祈祷和恳求，愿天主恩赐您不间断地享有这一切。\par

\Quote{我们的代表将向您展示主教们的签名和名字，其中一些将从圣经中为陛下圣洁提供教导。}\par

% structure: p0060 source=h2 target=subsection reason=relative-heading-rank; base=h2

\subsection{第十九章 关于会议的使节和皇帝的书信；\Person{乌尔撒基乌}和\Person{瓦伦斯}的拥护者后来与所发布的书信的一致；总主教的流放。关于\Place{尼西亚}会议，以及会议为何在\Place{阿里米农}举行。}

我们现在已经抄录了阿利米努姆公会议的来信。\Person{乌尔撒基乌斯}和\Person{瓦伦斯}及其同党，抢在公会议代表抵达之前，向皇帝展示了他们已宣读的文件，并诽谤公会议。皇帝因这份信式是在\Place{西尔米乌姆}他面前拟定的而对其被否决感到不悦，因此他厚待\Person{乌尔撒基乌斯}和\Person{瓦伦斯}；相反，他对代表们表现出极大的轻蔑，甚至迟迟不予接见。然而，他最终致函主教会议，告知他们他因被迫征讨蛮族而无法与代表们商议；因此他命令代表们留在\Place{阿德里亚诺波尔}直到他返回，以便在处理完公务后，他能心无旁骛地听取并审查代表们的陈述；他说：\Quote{因为研究神圣之事，理当带着不受其他俗务搅扰的心思。}这就是他信函的大意。\par

主教们回复说，他们绝不能背离他们已形成的决定，正如他们此前书面声明并由代表们所宣告的那样；他们恳求他善待他们，接见他们的代表，并阅读他们的信函。他们告诉他，如此众多的教会失去主教，这在他眼中必定是痛心之事；如果合他心意，他们希望在冬天之前返回各自的教会。写了这封充满恳求和哀告的信后，主教们等待了一段时间的答复；但未获任何回应，他们后来便返回了自己的城市。\par

我上述内容清楚证明，在\Place{阿利米努姆}召开的主教们确认了古时在\Place{尼西亚}所颁布的谕令。现在让我们看看他们最终是如何同意由\Person{瓦伦斯}、\Person{乌尔撒基乌斯}及其追随者所编纂的信仰信式的。关于此事，我得到了各种不同的说法。有人说，皇帝因主教们未经他许可就离开\Place{阿利米努姆}而恼怒，于是放任\Person{瓦伦斯}及其同党按自己的意愿治理西方教会，颁布他们自己的信式，将拒绝签署者从教会中驱逐，并另立他人接替其位。据说，\Person{瓦伦斯}利用这一权力，强迫一些主教签署信式，并将许多拒绝服从者逐出教会，首当其冲的是罗马主教\Person{利贝里乌斯}。还有人断言，当\Person{瓦伦斯}及其同党在意大利如此行事之后，他们决定以同样方式处置东方教会。这些迫害者穿越色雷斯时，据说他们在该省的城市\Place{尼西亚}停留。他们在那里召开了一次会议，宣读了\Place{阿利米努姆}的信式，并将其译为希腊语，通过声称该信式已获大公会议批准，使其在\Place{尼西亚}得到采纳；然后他们狡猾地称之为尼西亚信式，企图借助名称的相似性欺骗单纯之人，使之误以为是尼西亚公会议所颁布的古信式。这是某些人所提供的说法。另一些人则说，在\Place{阿利米努姆}公会议召开的主教们因被滞留该城而疲惫不堪，因为皇帝既未回复他们的信函，也未准许他们返回各自的教会；此时，持对立异端的人向他们指出，全世界的司祭不应为了一个词而起纷争，只需承认圣子相似圣父即可平息所有争议；因为东方的主教们永远不会罢休，除非\OriginalTerm{本体}{substance}一词被弃绝。据说，通过这些游说，公会议成员终于被说服，同意了\Person{乌尔撒基乌斯}竭力向他们强加的信式。\Person{乌尔撒基乌斯}及其同党担心公会议派往皇帝的代表会宣告西方主教们起初所表现的坚定态度，并揭露弃绝\OriginalTerm{同质}{consubstantial}一词的真正原因，于是以无法获得公共交通工具、道路不便旅行为借口，将这些代表滞留在色雷斯的\Place{尼西亚}整个冬天；据说，他们还诱使代表们将他们已接受的信式从拉丁文译成希腊文，并发送给东方主教们。通过这种方式，他们预见到信式会按照他们的意图产生效果而不被察觉欺诈；因为无人能证明\Place{阿利米努姆}公会议成员并非自愿弃绝\OriginalTerm{本体}{substance}一词，以顺从反对使用该词的东方主教们。但这显然是一个虚假的说法；因为公会议全体成员，除少数人外，坚决主张圣子在本质上相似于圣父，他们之间唯一的分歧在于，有些人说圣子与圣父\OriginalTerm{同质}{consubstantial}，而另一些人则断言他与圣父\OriginalTerm{相似实体}{like substance}。有些人以某种形式陈述此事，另一些人则以不同形式陈述。\par

% structure: p0064 source=h2 target=subsection reason=relative-heading-rank; base=h2

\subsection{第二十章　东方教会发生的事件：\Person{马拉松尼乌斯}、\Place{基齐库斯}的\Person{埃勒乌西乌斯}和\Person{马其顿尼乌斯}驱逐坚持\OriginalTerm{同质}{Consubstantial}一词的人。关于诺瓦提安教会：一座教堂如何被迁移；诺瓦提安派与正统派共融。}

当我在上面述说的事件在\Place{意大利}发生时，\Place{东方}，甚至在塞琉西亚会议召开之前，已是重大骚乱的舞台。\Person{阿卡西乌}和\Person{帕特罗斐禄}的追随者驱逐了\Person{玛克西穆}，将\Place{耶路撒冷}的教会转交给\Person{济利禄}。\Person{马其顿尼}侵扰\Place{君士坦丁堡}和邻近的城市；他得到\Person{厄勒乌西乌}和\Person{马拉托尼乌}的协助。后者原是他自己教会的执事，热衷于管理男女隐修士居所的穷人事务，\Person{马其顿尼}擢升他为\Place{尼科美底亚}的主教。\Person{厄勒乌西乌}，他并非毫无名望，原隶属宫廷军事服务，被祝圣为\Place{基齐库斯}的主教。据说\Person{厄勒乌西乌}和\Person{马拉托尼乌}在行为上都是好人，但他们热心迫害那些主张圣子与圣父同体的人，尽管他们不像\Person{马其顿尼}那样明显地残忍，\Person{马其顿尼}不仅驱逐那些拒绝与他共融的人，还囚禁一些人，并将其他人拖到法庭。在许多情况下，他强迫不情愿的人共融。他抓住未受洗的儿童和妇女，为他们施洗，并在不同地方摧毁了许多教堂，借口皇帝已命令拆毁所有承认圣子与圣父同体的祈祷之所。\par

在这一借口下，位于\Place{君士坦丁堡}的诺瓦替派教会，坐落于该城称为佩拉尔古斯的部分，被摧毁了。据说这些异端者与天主教会成员共同做了一件勇敢的事，他们达成一致。当那些受命摧毁这座教会的人即将开始拆毁工作时，诺瓦替派聚集在一起；一些人拆下材料，另一些人将其运送到该城一个名为西凯的郊区。他们迅速完成了这一任务；因为男女老少都参与了，通过将他们的劳动奉献给天主，他们异常振奋。由于这种热忱，教会很快得以重建，并因此而得名\Quote{阿纳斯塔西亚}。\Person{君士坦提乌斯}去世后，他的继任者\Person{尤利安}将诺瓦替派先前拥有的土地赐予他们，并允许他们重建教会。民众振奋地利用了这一许可，从西凯运回了原建筑的同样材料。但这发生在比我们目前回顾的更晚的时期。在这一时期，诺瓦替派教会与天主教会几乎达成联合；因为他们对天主性的看法相同，并遭受共同的迫害，两教会的成员聚集在一起共同祈祷。那时天主教没有祈祷之所，因为亚略派从他们手中夺走了。而且，由于两教会成员之间的频繁交往，他们似乎认为彼此之间的差异是徒然的，并决定彼此共融。我想，本来肯定已经实现了和解，若非群众的愿望被少数人的诽谤所挫败，这些人声称有一项古老的法律禁止教会的联合。\par

% structure: p0067 source=h2 target=subsection reason=relative-heading-rank; base=h2

\subsection{第二十一章 \Person{马其顿尼}在\Place{曼提尼乌姆}的行动。他试图移动\Person{君士坦丁大帝}灵柩时被免职。\Person{尤利安}被立为凯撒。}

大约在同一时间，\Person{厄勒乌西乌}完全摧毁了\Place{基齐库斯}的诺瓦替派教会。\Place{帕夫拉戈尼亚}其他地区的居民，特别是\Place{曼提尼乌姆}的居民，遭受了类似的迫害。\Person{马其顿尼}得知这些人中的大多数是\Person{诺瓦图}的追随者，并且教会权力本身不足以驱逐他们，便说服皇帝派遣四个步兵队攻击他们。因为他以为不习惯武器的人，一见士兵出现就会吓住，顺从他的意见。但结果相反，因为\Place{曼提尼乌姆}的人民拿起镰刀、斧头和任何手边可得的武器，向军队进发。一场激烈的冲突发生了，许多\Place{帕夫拉戈尼亚}人倒下，但几乎所有的士兵都被杀。\Person{马其顿尼}的许多朋友责备他造成了如此巨大的灾难，皇帝也不悦，对他的好感比以前少了。另一事件更加强了敌意。\Person{马其顿尼}计划移动\Person{君士坦丁大帝}的灵柩，因为隐藏它的建筑正在倒塌。人民对此意见分歧：一些人赞同这一计划，另一些人反对，认为这是不虔诚的，类似于挖坟。那些坚持尼西亚教义的人持后一种意见，坚持不应侮辱\Person{君士坦丁}的遗体，因为这位皇帝持有与他们相同的教义。而且，我想，他们急于反对\Person{马其顿尼}的计划。然而，\Person{马其顿尼}不再拖延，命人将灵柩运到同一座教堂，即安放殉道者\Person{阿卡西乌}陵墓的教堂。人民分成两派，一派赞成，一派谴责这一行为，在同一教堂内互相冲撞，导致如此多的杀戮，以至于祈祷之所和毗邻的地方充满了鲜血和尸体。皇帝（当时在西方）听闻此事后深为愤怒；他指责\Person{马其顿尼}是侮辱他父亲和屠杀人民的起因。\par

皇帝本已决定前往东方，继续他的行程；他授予他的堂弟\Person{尤利安}\Quote{凯撒}的称号，并派他去\Place{高卢西部}。\par

% structure: p0070 source=h2 target=subsection reason=relative-heading-rank; base=h2

\subsection{第二十二章 \Place{塞琉西亚}会议。}

大约在同一时期，东方的主教们约一百六十人，在伊苏里亚的塞琉西亚城集会。此时正值尤西比乌与希帕提乌执政。持宫廷显赫军职的莱奥纳斯奉君士坦提乌斯之命赶赴此会议，以便在他的监督下进行教义宣认。该省军事长官劳里基乌斯在场准备所需一切，因为皇帝诏书命令他提供此服务。在本次会议的首场会议中，几位主教缺席，其中有斯基托波利斯的帕特罗菲卢斯主教、君士坦丁堡的玛塞多尼乌斯主教及安基拉的巴西略主教。他们以各种借口为其缺席辩解：帕特罗菲卢斯以眼疾为由，玛塞多尼乌斯自称身体不适；但人们怀疑他们因害怕受到多项指控而缺席。由于其他主教拒绝在他们缺席期间着手调查争议问题，莱奥纳斯命令他们立即开始审查被激起的疑问。于是，一些人认为必须从教义主题的讨论开始，而另一些人则主张应首先调查其中那些被指控者的行为，如同对耶路撒冷主教济利禄、塞巴斯特主教尤斯塔提乌斯及其他人的情况那样。皇帝诏书的模棱两可——时而规定这样做，时而那样做——引发了这场争执。由此产生的争论变得如此激烈，以致他们之间的一切团结都被破坏，分裂为两派。然而，希望从教义审查开始之人的意见占了上风。当他们着手研究术语时，一些人希望摒弃使用\OriginalTerm{本体}{substance}一词，并援引不久前玛尔谷在西尔米乌姆编纂、且得到宫廷主教（其中包含安基拉主教巴西略）接纳的信仰宣言的权威。另许多人则渴望采纳安提约基雅教堂奉献礼上起草的信仰宣言。前一阵营包括尤多克西乌斯、阿卡基乌斯、帕特罗菲卢斯、亚历山大主教乔治、提尔主教乌拉尼乌斯及其他三十二位主教。后一阵营由叙利亚劳迪塞亚主教乔治、基齐库斯主教埃勒乌西乌斯、帕弗拉戈尼亚庞培城主教索弗洛尼乌斯支持；多数人赞同他们。人们有理由怀疑阿卡基乌斯及其同党缺席的原因在于他们与上述主教意见分歧，以及他们希望逃避对其提出的某些指控的调查；因为，尽管他们先前曾书面承认该子与圣父完全相似且\OriginalTerm{同质}{consubstantial}，但如今他们完全回避了先前的声明。经过长时间的辩论和争执，塔尔索主教西尔瓦努斯以高声且不容置疑的语调宣称，不应引入任何新的信仰宣言，而只应采纳在安提约基雅所批准的，且只应这一宣言有效。由于这一提议与阿卡基乌斯的追随者相悖，他们便退出会场；其他主教则宣读了安提约基雅宣言。次日，这些主教在教堂集合，关闭大门，私下确认了这一宣言。阿卡基乌斯谴责这一做法，并将他所主张的宣言私下呈递给莱奥纳斯和劳里基乌斯。三天后，同一批主教再次集会，此前缺席的玛塞多尼乌斯和巴西略也加入其中。阿卡基乌斯及其同党声明，除非被罢免和受指控者离开集会，否则他们将不参与会议的任何进程。他们的要求得以满足，因为对立派的主教们决心不给他任何借口解散会议——这显然是阿卡基乌斯的目的，以阻止即将对埃提乌斯异端及其自身与同党所受指控的审查。当全体成员集合时，莱奥纳斯表示他持有阿卡基乌斯同党交给他的一份文件；这是他们的信仰宣言并附有引言。其他主教对此一无所知，因为与阿卡基乌斯持相同观点的莱奥纳斯乐意为整个事件保密。当这份文件被宣读时，整个会场充满骚动；因为其中某些声明的要旨是：尽管皇帝禁止在信仰宣言中引入任何圣经中未出现的术语，但被罢免的主教们从各省被带至集会，连同其他非法被祝圣者，导致会议陷入混乱，一些成员受辱，另一些人被剥夺发言权。文件还补充道，阿卡基乌斯及其同党并不拒绝在安提约基雅编纂的宣言，尽管在该城集会的人们正是为了针对当时出现的困难而草拟了它；但因\Quote{同质}和\Quote{类似质}两词使某些人苦恼，且近来有人宣称子与父不同质，故有必要因此摒弃圣经中未出现的\Quote{同质}和\Quote{类似质}二词，谴责\Quote{不同质}一词，并明确承认子与父相似；因为祂正如宗徒保禄在某处所说：\BibleQuote{不可见的天主的肖像}。这些引言之后附有一份宣言，它既不符合尼西亚信经，也不符合安提约基雅信经，措辞如此巧妙，以致亚流和埃提乌斯的追随者若如此宣认信仰，也不会显得有误。此宣言省略了尼西亚会议在谴责亚流主义教义时所使用的措辞，对安提约基雅会议关于子之神性不变、以及祂是父之本体、旨意、能力及荣耀不变肖像的声明只字未提，仅仅简略表达了对圣父、圣子及圣神的信仰；在对几位从未卷入任何一方教义争论的个人施以一些粗俗谩骂后，凡持有与本宣言不同意见者均被宣告为与天主教会无关。莱奥纳斯提交的文件内容即如此，该文件已由阿卡基乌斯及其赞同者签字。宣读后，帕弗拉戈尼亚主教索弗洛尼乌斯喊道：\Quote{若我们每天将个人的意见当作信仰声明来接受，我们将无法达到真理的精确性。}阿卡基乌斯反驳说，编纂新宣言并未被禁止，因为尼西亚信经本身也曾被修改过一次或多次，埃勒乌西乌斯如此回应：\Quote{但本次会议并非为学习已知之事或接受安提约基雅与会者已批准的宣言之外的任何其他宣言而召开；并且，无论生死，我们都会坚持这一宣言。}争论至此，他们转而进行另一项调查，问阿卡基乌斯同党他们认为子在哪方面与父相似。他们回答说子仅在意志上相似，而非在本体上；对方于是坚持祂在本体上相似，并通过阿卡基乌斯先前所写的一部著作证明他曾持此观点。阿卡基乌斯回答说他不应被自己的著作所审判；争论激昂地持续了一段时间，基齐库斯主教埃勒乌西乌斯如此发言：\Quote{无论玛尔谷或巴西略有任何过犯，无论他们或阿卡基乌斯的追随者彼此之间有何指控，对会议而言并不重要；会议也没有责任审查他们的宣言是否值得推荐；维持已由九十七位司铎在安提约基雅确认的宣言就够了；若有人想引入其中未包含的任何教义，就应被视为与宗教和教会无关。}与他意见相同的人为他的发言鼓掌；然后集会散会。次日，阿卡基乌斯和乔治的同党拒绝参加会议；莱奥纳斯既已公开宣布与他们意见一致，也不顾一切恳求拒绝出席。被派去请他的代表们在阿卡基乌斯同党的宅邸中找到他；他以会议中分歧过大为托辞拒绝邀请，并说他奉皇帝之命仅在成员一致时才参加会议。许多时间就这样消耗了；其他主教屡次恳请阿卡基乌斯同党参加集会；但他们有时要求在莱奥纳斯宅邸举行特别会议，有时又声称他们受皇帝之命审判被指控者；因为他们既不接受其他主教通过的信仰宣言，又不清理自己所受的指控；他们也拒绝审查已被他们罢免的济利禄的案件；也无人强迫他们这样做。然而，会议最终罢免了亚历山大主教乔治、凯撒勒雅主教阿卡基乌斯、提尔主教乌拉尼乌斯、斯基托波利斯主教帕特罗菲卢斯、安提约基雅主教尤多克西乌斯以及另外数位主教。许多人也被开除共融，直到他们能洗清所控罪行。这些事项以书面形式传达至每位神职人员的堂区。安提约基雅司铎阿德里安被祝圣为该教会主教，代替尤多克西乌斯；但阿卡基乌斯的同党逮捕了他，将他交给莱奥纳斯和劳里基乌斯。他们把他交给士兵看管，但后来将他流放。\par

我们现在已对塞琉西亚会议的终结作了简短叙述。欲知更详尽信息者，须查阅会议记录，该记录由在场的速记员抄录。\par

% structure: p0073 source=h2 target=subsection reason=relative-heading-rank; base=h2

\subsection{第23章 阿卡奇与艾提乌：以及两个会议——阿里米努姆会议与塞琉西亚会议——的代表如何被皇帝引导接受同一教义。}

\Emph{紧接}上述事件之后，阿卡奇的支持者便前往皇帝处；但其他主教各自返回家乡。被一致推举为皇帝代表的十位主教抵达宫廷时，遇到了阿里米努姆会议的十位代表，以及阿卡奇的同党。后者已争取到宫廷中的重要人物，并借助他们的影响赢得了皇帝的青睐。据报，这些改宗者中有一些早先就已赞同阿卡奇的观点；有些则被教会的财富所收买；另一些则被呈现在他们面前的论点之微妙以及劝说者的威严所诱骗。事实上，阿卡奇并非等闲之辈；他天生具有卓越的智力和口才，在实现其计划时从不缺乏技巧或手法。他主领一座著名的教会，并以尤西比乌·潘菲鲁为老师而自豪，他继承了这位老师的主教职位，并且因他的著作的声誉和传承而比任何人都更显赫。凭借所有这些优势，他无论从事何事都轻而易举地成功。\par

此时在君士坦丁堡共有二十位代表，每个会议各十位，此外还有许多主教因各种原因来到该城。皇帝在前往西方之前任命为君士坦丁堡首席长官的霍诺拉图斯，奉命在大公会议监督的参预下，审查有关艾提乌及其异端的报告。君士坦提乌斯偕同一些统治者最终承担了这一案件的调查；由于证明艾提乌引入了与信仰根本对立的教义，皇帝和其他审判官对他的亵渎言论感到愤慨。据说阿卡奇派起初假装对这一异端无知，其目的在于引诱皇帝及其周围的人去审查它；因为他们以为艾提乌的口才不可抗拒；他必定能说服听众；他的异端将征服不情愿的人。然而，当结果证明他们的期望落空时，他们便要求阿里米努姆会议所接受的信仰表述应得到塞琉西亚会议代表的批准。当后者抗议他们决不会放弃使用\Quote{本质}一词时，阿卡奇派向他们宣誓：他们并不认为圣子在本质上与圣父不相似；相反，他们准备谴责这一意见为异端。他们补充说，他们更珍视西方主教在阿里米努姆编纂的信仰表述，因为\Quote{本质}一词出人意料地被删除了；他们说，如果接受这一表述，便不会再提及\Quote{本质}或\Quote{同质}这两个词，而许多西方司铎因对尼西亚公会议的尊敬而特别执着于后一词汇。\par

正是由于这些原因，皇帝才批准了该信仰表述；当他回忆起在阿里米努姆聚集的主教人数众多时，又想到说\Quote{圣子与圣父相似}或\Quote{与圣父同质}并无错误；并且进一步考虑到，如果用其他等价且无可争议的表述（例如\Quote{相似}一词）来替代圣经中未出现的术语，不会导致意义上的差别，他便决定批准该信仰表述。抱持着这样的看法，他命令主教们接受这一信仰表述。次日，准备举行宣告他为执政官的盛大仪式，按照罗马习俗，该仪式在正月之初举行，皇帝与主教们共度了那一整天及翌夜的部分时间，最终说服了塞琉西亚会议的代表接受从阿里米努姆送来的信仰表述。\par

% structure: p0077 source=h2 target=subsection reason=relative-heading-rank; base=h2

\subsection{第24章 阿里米努姆会议的信仰表述获阿卡奇派认可。被革职的大司祭名单及其被定罪的原因。}

阿卡西的支持者一度留在君士坦丁堡，并邀请了几位卑提尼亚的主教前往，其中有加采东主教马里斯和哥特人主教乌尔菲拉斯。这些主教聚集在一起，约有五十人，他们批准了在阿里米努姆大公会议上宣读的信经格式，并补充规定：永远不再使用\Quote{实体}和\Quote{位格}这两个词来指称天主。他们还宣布，过去制定的所有其他信经格式，以及将来可能编纂的任何信经格式，都应受到谴责。然后，他们罢免了埃提乌斯的执事之职，因为他撰写了一些充满争执和一种与教会圣召相悖的虚妄知识的著作；他在写作和辩论中使用了一些不敬的言辞；并且他是教会内动乱和叛乱的起因。许多人声称，他们并非心甘情愿地罢免他，而只是为了消除皇帝对他们的怀疑——他们曾被指控持有埃提乌斯的观点。\newline{}那些持这些观点的人利用了皇帝因上述原因对马其顿尼所怀的怨恨，因此他们罢免了马其顿尼，同样也罢免了基齐库斯主教埃莱乌西乌斯、安基拉主教巴西尔、撒尔迪斯主教赫奥塔西乌斯和佩尔加穆斯主教德拉孔提乌斯。虽然他们在教义上与这些主教存在分歧，但在罢免他们时，并未指责他们的信仰，而是提出了与所有人共通的指控：他们扰乱了和平，违反了教会的法律。他们特别指出：当长老狄奥吉内斯从亚历山大里亚前往安基拉时，巴西尔没收了他的文件并殴打了他；他们还指控巴西尔未经审判就将许多来自安提约基雅、幼发拉底河沿岸、西里西亚、加拉提亚和亚细亚的神职人员交给各省的统治者，流放并施以残酷的惩罚，以至于许多人被戴上锁链，被迫贿赂押送他们的士兵以免受虐待。他们补充说：有一次，皇帝命令将埃提乌斯及其一些追随者带到塞克罗皮乌斯面前，以便他们回答针对他们的各种指控，巴西尔却建议负责执行这一命令的人根据自己的判断行事。他们说，巴西尔写信给叙利亚总督和省长赫莫根尼斯，指明谁应被流放以及流放地点；当皇帝召回流放者时，他不同意他们返回，反而反对统治者和神职人员的意愿。他们还指控巴西尔煽动西尔米乌姆的神职人员反对日耳曼尼乌斯；尽管他书面宣称他与日耳曼尼乌斯、瓦伦斯和乌尔萨西乌斯保持了共融，但他却将他们作为罪犯置于非洲主教们的法庭面前；当被问及此事时，他否认并发了假誓；后来被证实后，他又试图用诡辩来为他的假誓辩护。他们补充说，他是伊利里亚、意大利、非洲和罗马教会中纷争和叛乱的起因；他将一名女仆投入监狱，强迫她作伪证陷害女主人；他为一名行为放荡、与女子非法同居的人施洗，并提拔他为执事；他未将一名导致数人死亡的江湖医生逐出教会；他和一些神职人员曾在圣台前发誓互不控告。他们说，这是神职人员的首领为了保护自己免受原告的指控而采取的策略。简而言之，这就是他们罢免巴西尔所列举的理由。他们说，罢免欧斯塔提乌斯是因为他担任长老时曾被自己的父亲、卡帕多细亚凯撒利亚教会主教尤拉里乌斯定罪并剥夺了祈祷中的共融；还因为他被本都地区内凯撒利亚城的一次会议绝罚，并被君士坦丁堡主教尤西比乌斯因未能忠实履行某些职责而罢免。他还被在甘格瑞召集的人因教导、行为和思想违背健全教义而被剥夺了主教职位。他被安提约基雅会议判定犯有伪证罪。他还试图推翻在梅利提纳召开的会议的决议；尽管他犯有许多罪行，却厚颜无耻地渴望成为他人的审判者，并斥责他们为异端。他们罢免埃莱乌西乌斯，因为他轻率地提拔了一个名叫赫拉克利乌斯的提尔人为执事；此人曾在提尔担任赫拉克勒斯神的祭司，被控告并审判犯有巫术罪，因此退居基齐库斯并假装皈依基督教；此外，埃莱乌西乌斯在得知这些情况后并未将他逐出教会。他还未经调查就按立了一些来到基齐库斯的人，而这些人已被参与此次会议的加采东主教马里斯定罪。赫奥塔西乌斯被罢免，是因为他未经吕底亚主教的同意就被按立为撒尔迪斯主教。他们罢免了佩尔加穆斯主教德拉孔提乌斯，因为他先前在加拉提亚已担任过另一个主教职位，并且据称他在两次任命中均非法按立。在这些事件之后，举行了第二次会议，罢免了塔尔苏斯主教西尔瓦努斯、帕夫拉戈尼亚庞培波利斯主教索弗洛尼乌斯、萨塔拉主教埃尔皮狄乌斯以及伊苏里亚塞琉西亚主教内奥纳斯。他们为罢免西尔瓦努斯所提出的理由是，他在塞琉西亚和君士坦丁堡领导了一个愚蠢的派别；此外，他任命了塞奥菲卢斯为卡斯塔巴拉教会的主席，而塞奥菲卢斯先前已被巴勒斯坦的主教们按立为厄琉特罗波利斯主教，并曾发誓未经他们许可绝不接受任何其他主教职位。索弗洛尼乌斯因贪婪而被罢免，因为他出售了部分献给教会的祭品以牟利；此外，在接到一次和二次传唤后，他才勉强出庭，而且不是回答对他的指控，而是向其他法官上诉。内奥纳斯被罢免，是因为他在按立安提约基雅主教安尼亚努斯时使用了暴力，并试图在自己的教会中进行按立；他还将一些先前从政且完全不懂圣经及教会法典的人按立为主教，这些人按立后更愿意享受自己的财产而不愿履行主教职责，并书面声明他们宁愿管理自己的产业也不愿在没有财产的情况下履行主教职务。埃尔皮狄乌斯被罢免，是因为他参与了巴西尔的恶行并引起了极大的混乱；还因为他违反梅利提纳会议的决议，恢复了一个名叫尤西比乌斯的人在长老团中原有的品级，而此人曾因将一位因违反协议和誓言而被绝罚的妇女奈克塔里亚按立为女执事而被罢免；授予她这一尊荣显然违反了教会的法律。\par

% structure: p0079 source=h2 target=subsection reason=relative-heading-rank; base=h2

\subsection{第25章 耶路撒冷主教\Person{济利禄}被罢黜的原因。主教们之间的相互纷争。\Person{梅利提乌}由亚略派人士祝圣，并取代\Person{欧斯塔修}占据\Place{塞巴斯特}主教职。}

除了上述主教外，\Person{济利禄}，耶路撒冷主教，也被罢黜，因为他接纳了\Person{欧斯塔修}和\Person{厄尔丕迪乌}进入共融，而此二人曾反对在\Place{梅利提纳}召集的会议所颁布的法令，\Person{济利禄}本人也在该会议中；又因为他还在\Person{巴西尔}和\Place{劳狄刻亚}主教\Person{乔治}被罢黜后，在\Place{巴勒斯坦}接纳他们进入共融。\Person{济利禄}最初就任耶路撒冷主教时，曾与\Place{凯撒勒雅}主教\Person{阿卡奇乌}就其都主教权利发生争执，他主张自己的主教座堂为宗徒之所。这场争执激起了两位主教之间的敌意，他们互相指责对方关于天主性的教义不纯正。事实上，二人此前已受怀疑：\Person{阿卡奇乌}偏袒\Person{亚略}的异端，而\Person{济利禄}则支持那些主张圣子在本体上与圣父相似的人。\Person{阿卡奇乌}既对\Person{济利禄}怀有敌意，又得到省内同自己意见一致的主教们的支持，便借以下借口图谋罢黜\Person{济利禄}。\Place{耶路撒冷}和邻近地区曾一度遭遇饥荒，大批穷人向身为主教的\Person{济利禄}恳求必需的食物。因他无钱购买所需用品，便为此变卖了圣堂的帷幔和圣器。据说有一人认出他在祭台上奉献的供品成为了一名女演员服装的一部分，便特意查问它从何处购得；查明是一名商人卖给了女演员，而主教卖给了商人。\Person{阿卡奇乌}正是借这个借口罢黜了\Person{济利禄}。\par

经过查问，我发现事实如下。据说亚加奇派随后将所有上述已被罢黜的主教逐出\Place{君士坦丁堡}。他们自己派别中有十位主教拒绝在这些罢免法令上签字，被与其他主教分开，并被禁止履行职务或管理其教会，直到他们同意签名。规定除非他们在六个月内遵守，并同意会议的所有法令，否则他们将被罢黜，并应召集各省主教选举其他人取代他们。这些决议和行为之后，信件被送到所有主教和圣职人员，要求遵守并执行其法令。\par

结果，不久之后，欧多克修派的人士被到处安插。\Person{欧多克修}本人占据了\Person{马其顿尼}的主教职位；\Person{亚大纳西}被安置在\Person{巴西尔}的教会；后来成为以自己名字命名的异端首领的\Person{欧诺米乌}占据了\Person{厄勒乌西}的教座；\Person{梅利提乌}被任命到\Place{塞巴斯特}的教会，取代\Person{欧斯塔修}。\par

% structure: p0083 source=h2 target=subsection reason=relative-heading-rank; base=h2

\subsection{第26章 \Place{君士坦丁堡}主教\Person{马其顿尼}之死。\Person{欧多克修}在教导中所说的话。\Person{欧多克修}与\Person{阿卡奇乌}竭力寻求废除在\Place{尼加亚}和\Place{阿里米农}所颁布的信德宣言；由此在教会中引发的纷扰。}

\Person{马其顿尼}被逐出\Place{君士坦丁堡}教会后，退隐到该城市郊，在那里去世。\Person{欧多克修}在\Person{君士坦提乌斯}执政第十年、被称为凯撒的\Person{犹利安}执政第三年，占据了他的教会。据说，在被称为\Quote{索非亚}的大教堂奉献典礼上，他站起来教导民众，以以下命题开始他的讲道：\Quote{圣父是不敬虔的，圣子是敬虔的}；这些话在民众中引起了极大骚动，他于是补充说：\Quote{安静；圣父是不敬虔的，因为他不崇拜任何人；圣子是敬虔的，因为他崇拜圣父。}这个解释使听众哄堂大笑。\Person{欧多克修}和\Person{阿卡奇乌}共同竭尽全力试图使\Place{尼加亚}会议的法令被人遗忘。他们将\Place{阿里米农}宣读的信条加上他们自己的各种解释补充，送往帝国各省，并从皇帝那里获得一项法令，流放所有拒绝签署的人。但这项在他们看来似乎很容易执行的事业，却成了最大灾祸的开端，因为它激起了全帝国的骚动，使各地教会遭受了比在异教皇帝统治下所受的迫害更为严重的迫害。因为即使这次迫害没有像先前那样造成身体上的折磨，但就它的可耻性质而言，在所有正确思考的人看来更为严重；因为迫害者和被迫害者都属于教会；而这一点更加可耻，即同一宗教的人以教会法律禁止对敌人和陌生人表现出的残酷程度对待自己的同伴。\par

% structure: p0085 source=h2 target=subsection reason=relative-heading-rank; base=h2

\subsection{第27章 \Person{马其顿尼}在被拒于教座之后，亵渎圣神。通过\Person{马拉松尼乌}及其他人的工具，他的异端得到传播。}

创新的精神是自我标榜的，因此它越来越深入，并带着日益增长的自负滑向更大的新奇事物，在轻视教父们的同时，它制定了自已的法律，也不尊重古人对天主的教义，而是不断想出奇怪的教义，并不安地将新奇叠加于新奇，正如如今的事件所表明的那样。因为当\Person{马其顿纽斯}从\Place{君士坦丁堡}教会中被废黜后，他放弃了\Person{阿卡基乌斯}和\Person{尤多克修斯}的教义。他开始教导说，圣子是天主，祂在各方面且在本质上与圣父相似。但他断言，圣神并不享有同等的尊荣，并称祂为仆人和侍者，并将那些可以无误地应用于圣天使的言辞应用于祂。这一教义被\Person{埃勒乌修斯}、\Person{尤斯塔修斯}以及所有其他在\Place{君士坦丁堡}被废黜的主教，即相反异端的追随者所接受。他们的榜样很快被\Place{君士坦丁堡}、\Place{比提尼亚}、\Place{色雷斯}、\Place{赫勒斯滂}及邻近省份的民众中不小的一部分所效仿。因为他们的生活方式影响力不小，而民众对此特别关注。他们表现出极大的庄重举止，其纪律如同隐修士一般；他们的言谈朴素，且具有说服力的风格。据说所有这些品质都集中在\Person{马拉松纽斯}身上。他最初在军队中担任公职，受命于总督。在任职期间积攒了一些钱财后，他离开了军事职业，并负责管理救济病患和贫困者的机构。后来，在\Place{塞巴斯特}主教\Person{尤斯塔修斯}的建议下，他选择了苦修的生活方式，并在\Place{君士坦丁堡}建立了一所隐修院，存留至今。他为支持上述异端带来了如此多的热忱和如此多的个人财富，以至于马其顿派被许多人称为马拉松派，在我看来并非没有理由；因为似乎正是他本人及其机构，使得该异端在\Place{君士坦丁堡}未被完全消灭。实际上，在\Person{马其顿纽斯}被废黜后，马其顿派直到\Person{阿卡狄乌斯}统治时期既没有教堂也没有主教。\par

亚略派人士将所有持不同意见的人逐出教堂并严加迫害，剥夺了他们所有这些特权。要列举这一时期被逐出各自城巿的司铎的名字绝非易事；因为我相信帝国没有一个省份免于这样的灾祸。\par

% structure: p0088 source=h2 target=subsection reason=relative-heading-rank; base=h2

\subsection{第二十八章。亚略派人士以为神圣的\Person{梅利修斯}支持他们的观点，便将他从\Place{塞巴斯特}迁至\Place{安提约基雅}。当他大胆地宣认正统教义时，他们感到困惑，并在废黜他之后，将\Person{欧佐伊乌斯}置于该主教座。\Person{梅利修斯}形成了自己的教会；但那些坚持同体的人却回避他，因为他是被亚略派人士祝圣的。}

在\Person{尤多克修斯}获得\Place{君士坦丁堡}教会治理权的时期，有许多人觊觎\Place{安提约基雅}的主教座；正如在这种情况下常见的那样，争斗和骚乱分裂了该教会的圣职人员和民众。\par

各方都急于将教会的治理委托给一位同己派别的主教；因为关于教义的无穷争论在他们中盛行，他们甚至在咏唱圣咏的方式上也无法达成一致；如前所述，每人皆按自己独特的信条咏唱圣咏。安提约基雅的教会既处于此种状态，欧多克修斯派认为将塞巴斯特主教默肋提乌斯托付给该城主教职位是妥当的，因他拥有伟大而具说服力的口才，卓越的生活，且如他们所想，一切都与己见相同。他们相信他的声誉会吸引安提约基雅及邻近城市的居民皈依他们的异端，尤其是那些称为欧斯塔修派的人，他们始终坚定持守尼西亚信理。然而他们的期望完全落空了。据说他初到安提约基雅时，一大群由亚流派及与保利努斯共融的人组成的民众蜂拥而至。有些人想见此人，因他尚未到来时名声已大；另一些人则急于听他所说的话，并查明其见解的性质；因为传闻已流传开来，后来证明属实，即他持守尼西亚会议所定的教义。在他的初次讲道中，他仅限于教导民众我们所称的伦理；但后来他公开宣告圣子与圣父同性同体。据说听到这些话时，教会总执事——当时是那里的神职人员之一——伸出手捂住传道者的口；但他却更清楚地用手指而非言语表达自己的见解。他只向民众伸出三根手指，然后合拢，只留下一根手指伸出，以手势表达了他不能言语说出的话。当总执事窘迫中抓住他的手时，他释放了口；舌头自由了，默肋提乌斯更大声更清楚地宣告他的见解，劝诫听众持守尼西亚会议的信条，并向听众作证，持其他见解的人偏离了真理。当他坚持用言语或手势表达同一见解，而总执事封住他口时，双方发生了争斗，不亚于全能搏击；欧斯塔修派的追随者大声欢呼、喜乐跳跃，而亚流派则沮丧不已。欧多克修斯及其同党因这番讲道而极为愤怒，用阴谋诡计将默肋提乌斯逐出安提约基雅。但不久之后，他们又召他回来，因为他们以为他放弃了先前的见解而采取了他们的见解。然而，他持守尼西亚信理的坚定与不可改变很快显明，他被逐出教会，且奉皇帝命令被流放；安提约基雅主教之位授予欧佐伊乌斯，他先前曾与亚流一同被流放。默肋提乌斯的跟随者与亚流派分离，自行举行集会，因为那些从一开始就主张圣子与圣父同性同体的人拒绝接纳他们共融，因为默肋提乌斯是由亚流主教祝圣的，且他的跟随者是由亚流司铎施洗的。为此他们虽持相同见解却仍分离。\par

皇帝得知波斯将要发生叛乱，便前往安提约基雅。\par

% structure: p0092 source=h2 target=subsection reason=relative-heading-rank; base=h2

\subsection{第二十九章　阿卡西乌斯派并不安宁，仍力图废除\Quote{同性同体}一词，并确认亚流异端。}

阿卡西乌斯派不能保持平静；于是他们与少数其他人在安提约基雅聚集，谴责他们自己先前颁布的谕令。他们决定从阿里米努姆和君士坦丁堡已宣读的信仰公式中删除\Quote{相似}一词，并声称在各方面，在本性及旨意上，圣子与圣父不相似，且他从无中生出，正如亚流从起初所教导的。埃提乌斯的追随者也加入他们，埃提乌斯是继亚流之后第一个公开冒险宣认这些意见的人；因此埃提乌斯被称为无神论者，他的赞成者被称为阿诺米乌派和埃克苏康提派。\par

当持守尼西亚信理的人质问阿卡西乌斯派：既然在他们的信仰公式中明认圣子是\Quote{出自天主的天主}，他们怎能说圣子与圣父不相似且从无中生出时，他们回答说，宗徒保禄曾宣称\Quote{一切都是出于天主}\BibleRef{格前 11:12}，而圣子包含在\Quote{一切}之中；并说正是在此意义上，且依照圣经，他们公式中的表述必须被理解。他们所用的正是这样的含糊其辞与诡辩。最终，他们发现无法提出任何有效论证来在那些就此问题追问他们的人面前为自己辩护，于是在君士坦丁堡的信仰公式被第二次宣读后，他们便退出集会，回到自己的城。\par

% structure: p0095 source=h2 target=subsection reason=relative-heading-rank; base=h2

\subsection{第三十章　安提约基雅主教乔治与耶路撒冷大司祭。三位大司祭先后继济利禄之位；济利禄恢复耶路撒冷主教之位。}

在此期间，\Person{亚大纳削}不得不躲藏起来，\Person{乔治}返回\Place{亚历山大里亚}，开始残酷地迫害外教人和与他意见不同的基督徒。他强迫双方按他指定的方式敬拜，凡有反对者，便用强制手段迫使他们服从。统治者憎恨他，因为他藐视他们，并向官员发号施令；民众也因他的暴政而厌恶他，因为他的权力超过其他所有人。外教人对他比基督徒更为憎恶，因为他禁止他们献祭和庆祝祖先的节日；还因为他曾带领埃及总督和武装士兵进城，洗劫了他们的偶像、还愿物和神庙装饰品。这实际上是他死亡的原因，关于此事我将详述。\par

在\Person{济利禄}被废黜后，\Person{厄勒纽斯}获得了耶路撒冷的教会；他由\Person{赫拉克利乌斯}继任，\Person{赫拉克利乌斯}又由\Person{希拉里乌斯}继任；因为我们从传统中得知，在那个时期，这些人管理了那里的教会，直到\Person{德奥多西}在位时，\Person{济利禄}再次恢复了他的教座。\par

\SourceRecord{Translated by Chester D. Hartranft.；Nicene and Post-Nicene Fathers, Second Series；Vol. 2.；Edited by Philip Schaff and Henry Wace.；Buffalo, NY: Christian Literature Publishing Co.,；1890.；<http://www.newadvent.org/fathers/26024.htm>.}

% structure: p0001 source=h1 target=section reason=page-title

\section{\WorkTitle{教会史（卷五）}}

% structure: p0002 source=h2 target=subsection reason=relative-heading-rank; base=h2

\subsection{第一章 背教者\Person{尤利安}的背教。\Person{君士坦提乌斯}皇帝之死。}

以上是东方教会所发生的事。然而，在此期间，凯撒\Person{尤利安}攻击并征服了居住在莱茵河畔的蛮族；他杀了许多人，又俘虏了其他人。由于这场胜利极大地增加了他的声望，也因他的节制与温和使他深受军队爱戴，他们便拥立他为奥古斯都。他非但没有就此事向\Person{君士坦提乌斯}表示歉意，反而撤换了由\Person{君士坦提乌斯}选任的官员，并积极散发信件，其中\Person{君士坦提乌斯}曾请求蛮族进入罗马领土，帮助他对抗\Person{马格嫩提乌斯}。接着他突然改变了信仰，尽管他先前曾宣认过基督教，却宣布自己为大祭司，频繁出入异教神庙，献祭，并邀请他的臣民采纳那种敬拜方式。\par

由于预期波斯人会入侵罗马领土，且\Person{君士坦提乌斯}为此已前往\Place{叙利亚}，\Person{尤利安}认为他可以在不交战的情况下夺得\Place{伊利里库姆}；于是他启程前往该行省，借口要向\Person{君士坦提乌斯}解释他未经其同意便接受了皇权象征之事。据说，当他到达\Place{伊利里亚}边境时，葡萄藤上结满了青绿的葡萄，尽管收成季节已过，昴星团也已沉落；并且他的随从受到大气中露水的溅洒，每一滴露水都印有十字圣号。他及其许多随从视这些不合时令的葡萄为吉兆；而露水只是偶然落在衣服上形成了那样的形状。\par

另一些人则认为，在两个征兆中，青绿的葡萄意味着皇帝将早逝，其统治将极为短暂；而第二个征兆，即由露滴形成的十字架，表明基督教信仰来自天上，所有人都应当接受十字圣号。就我而言，我深信那些视这两个现象为\Person{尤利安}不祥之兆的人并未弄错；时间的推移证明了他们看法的准确性。\par

当\Person{君士坦提乌斯}听说\Person{尤利安}率军向他进攻时，他放弃了原本针对波斯人的征讨，启程前往\Place{君士坦丁堡}；但他行至\Place{莫普苏克雷奈}时去世，该地靠近\Place{托鲁斯山脉}，位于\Place{基利家}与\Place{卡帕多细亚}之间。\par

他死时年四十五岁，与父亲\Person{君士坦丁}共同执政十三年，并在那位皇帝去世后又独自统治了二十五年。\par

\Person{君士坦提乌斯}去世后不久，早已控制\Place{色雷斯}的\Person{尤利安}进入\Place{君士坦丁堡}并被拥立为皇帝。异教徒声称，占卜者和魔鬼在\Person{君士坦提乌斯}前往\Place{加拉太}之前已经预言了他的死亡以及局势的变动，并曾建议他进行这次远征。若非\Person{尤利安}的生命在不久后便告终结，且他只是如同梦境一般品尝了皇权，这或许可以被视为真实的预言。但我认为，相信他在听闻\Person{君士坦提乌斯}之死的预言、并被警告自己将命丧波斯人之手后，依然投身于明显的死亡——这于世人不曾给他带来比缺乏谋略、统兵无方更高的声誉——而且他若活着，恐怕会使得罗马领土的大部分落入波斯人的轭下，这是荒谬的。不过，我插入此议论，只是担心因遗漏而被指责。我让每个人自行判断。\par

% structure: p0009 source=h2 target=subsection reason=relative-heading-rank; base=h2

\subsection{第二章 \Person{尤利安}的生平、教育及训练，及其登基为帝。}

\Emph{立刻}，在\Person{君士坦提乌斯}死后，教会中升起了对迫害的恐惧，基督徒因预料这场灾难而感到的痛苦，超过了实际发生时的痛苦。这种情绪的产生，是由于长期没有经历这样的危险而变得不习惯，也由于对暴君们曾加诸其祖先的酷刑的回忆，以及他们深知皇帝对其教义的仇恨。据说他公开背弃对基督的信仰如此彻底，以至于他通过异教徒称为弃绝性的献祭与赎罪，以及牲畜的血，来净化自己脱离我们的圣洗。从那时起，他专注于占卜和举行异教仪式，无论公开还是私下。据说有一天，他在察验牺牲的内脏时，看到其中有一个冠冕环绕的十字架。这一景象令参与仪式的人惊恐不已，因为他们判断这表示宗教的力量，以及基督教教义的永恒性；因为环绕它的冠冕是胜利的象征，又因其连续性，圆环无所始亦无所终，在任何方向上都没有界限。首席占卜者命令\Person{尤利安}放心，因为在他看来牺牲是吉兆，而且它们包围着基督教教义的符号，实际上正将其挤入其中，使它无法随意扩展蔓延，因为它被圆周所限制。\par

我也曾听说，有一天尤利安下到一个最著名而可怖的\OriginalTerm{内殿}{adytum}中，或是为了参与某些\OriginalTerm{入教仪式}{initiation}，或是为了询问\OriginalTerm{神谕}{oracle}；借着为此目的而设的机关或法术，极其可怖的\OriginalTerm{幽灵}{specters}突然出现在他面前，以致在惊慌恐惧中，他忘记了在场的人——因为他转信新宗教时已是成年人——无意中又恢复了早年的习惯，以基督的标记画了十字，正如基督徒在遭遇不测危险时通常所做的那样。幽灵立刻消失了，它们的计谋也被挫败。主持仪式者起初对此感到惊讶，但得知鬼魔逃遁的原因后，便宣称此举是\OriginalTerm{亵圣}{profanation}；他劝勉皇帝要勇敢，绝不可在言行上求助基督教的一切，然后又带他重新参加了入教仪式。皇帝对这类事情的狂热使基督徒们颇为忧伤，并使他们极度焦虑，尤其他本人从前曾是基督徒。他生于虔敬的双亲，在襁褓中便按教会惯例受了圣洗圣事，在圣经的知识中长大，由主教和教会的士人养育。他与加卢斯是君士坦提乌斯的儿子，这位君士坦提乌斯与皇帝君士坦丁是同父的兄弟，也是达耳马提乌斯的兄弟。达耳马提乌斯有一个同名的儿子，曾被立为凯撒，在君士坦丁死后被兵士所杀。加卢斯和尤利安当时是孤儿，本也难逃此命运，但加卢斯因当时患病而被饶恕，人们以为他不久便会自然死亡；尤利安则因年纪尚幼，年仅八岁而得以幸免。在这次奇妙的保全之后，两兄弟被安置在卡帕多细亚一个名为马塞隆的宫殿中居住；这座皇家宅邸位于阿尔革斯山附近，离凯撒勒雅不远；它有一座宏伟的宫殿，配有浴场、花园和长流的泉源。他们在那里接受与其出身高贵相应的培养和教育；语言教师和圣经诠释者教导他们科学和适龄的体育锻炼，以致他们被登记在圣职人员之列，并向民众朗读教会书籍。他们的习惯和行为并未显露出任何背离虔敬之处。他们尊敬圣职人员和其他善良的人以及热心于教义的人；他们经常去教堂，并适当地向殉道者的坟墓致敬。\par

据说他们曾着手将殉道者圣玛玛斯的坟墓安置在一座大建筑物中，并彼此分工；当他们试图在荣誉的竞争中相互超越时，发生了一件事，其惊人之处，若不是有许多仍在我们中间的人从目击者那里听到，实在是完全难以置信。\par

加卢斯所负责的那部分建筑物进展迅速，遂其所愿；但尤利安所负责的部分中，一部分倒塌了；另一部分从地面拱起；第三部分刚一接触地基便无法保持竖直，反而向后翻倒，仿佛下方有某种抗拒而强大的力量在推挤它。\par

众人普遍视此为异兆。然而人们从这现象中并未得出任何结论，直到后来发生的事才显明了它的含义。但自那一刻起，已有少数人怀疑尤利安宗教的真实性，并猜测他不过是因怕得罪当时仍是基督徒的皇帝而表面上持守虔敬，实则因不便吐露而隐藏了自己的真实想法。据称，他最初是因与占卜者交往而秘密地背离了他祖先的宗教；因为当君士坦提乌斯对两兄弟的怒气平息后，加卢斯前往亚细亚，定居在以弗所，他的大部分产业都在那里；而尤利安则前往君士坦丁堡，常去学校，在那里他出色的天赋和在各门学问上敏捷的造诣并未被隐藏。他身着平民的装束出现在公共场合，并有许多同伴；但因他与皇帝有亲属关系，又有能力处理事务，且被预期会成为皇帝，关于此事的议论甚多，正如在人口众多与都城中所常有的那样，他奉命退隐至尼科美底亚。\par

在那里他结识了以弗所的哲学家马克西穆斯，这位哲学家教授他哲学，并激发他对基督宗教的仇恨，还向他保证那广为流传的关于他的预言是真的。尤利安正如许多人在预料到严峻处境时通常会做的那样，被这些有利的希望所软化，视马克西穆斯为友。这些事传到了君士坦提乌斯的耳中，尤利安感到害怕，于是剃了发，外表上采取了隐修士的生活方式，而暗地里却持守着另一种宗教。\par

当他到达成年时，他更容易被迷惑，但仍对这些倾向感到焦虑；他钦佩那预言未来的技艺（如果真有这种技艺的话），认为掌握它是必要的；于是他进行了一些基督徒所不允许的实验。从那时起，追随此技艺的人便成了他的朋友。抱着这种想法，他从尼科美底亚来到亚细亚，在那里与从事此类活动的人交往，变得更加热衷于占卜。\par

当他的兄弟加卢斯已被立为凯撒，却因被控叛乱而被处死后，君士坦提乌斯也怀疑尤利安怀有对帝国的野心，因此将他置于卫兵的看管之下。\par

\Person{君士坦提乌斯}的妻子\Person{尤西比亚}为\Person{尤利安}获准退隐\Place{雅典}；他遂在那里定居，借口参加外教的训练和学校，但据传言，他与占卜者商谈自己的未来前景。\Person{君士坦提乌斯}召回他，并宣布他为凯撒，许诺将他的姊妹\Person{君士坦提娅}许配给他，并派他去\Place{高卢}；因为先前\Person{君士坦提乌斯}为对付\Person{马格嫩提乌斯}而雇佣的蛮族，发现自己的服务未被需要，已瓜分了那个地区。由于\Person{尤利安}非常年轻，一同派去了负责谨慎事务的将领；但这些将领沉溺于享乐，他作为凯撒亲临战场，并负责战争。他坚定了士兵们的战斗精神，并以其他方式激励他们冒险；他还下令，凡杀死一名蛮族者，应给予固定赏金。在赢得士兵们的爱戴后，他写信给\Person{君士坦提乌斯}，告知将领们的轻率；当另一名将领被派来后，他攻击蛮族并取得了胜利。蛮族派遣使节求和，并出示了\Person{君士坦提乌斯}请求他们进入罗马领土的信函。他故意拖延遣返使节；他出其不意地攻击了许多敌人并征服了他们。\par

有人说\Person{君士坦提乌斯}出于敌意的预谋，将这战役托付给他；但在我看来这似乎并不可能。因为，既然任命他为凯撒全凭\Person{君士坦提乌斯}一人，他为何要赐予他这个称号？为何将他的姊妹嫁给他，或听取他对无效将领的抱怨，并派去一位称职的将领以完成战争，如果他不对\Person{尤利安}怀有好意？\par

但我推测，他赐予\Person{尤利安}凯撒称号是因为他对\Person{尤利安}怀有好意；但在\Person{尤利安}未经他同意而被宣布为皇帝后，他通过\Place{莱茵河}上的蛮族谋害他；我认为这要么出于恐惧，害怕\Person{尤利安}会报复他和他兄弟\Person{加卢斯}在青年时期所受的虐待，要么如常理般，出于对他获得同样荣誉的嫉妒。但关于此事，存在着许多不同的意见。\par

% structure: p0021 source=h2 target=subsection reason=relative-heading-rank; base=h2

\subsection{第三章  \Person{尤利安}在稳居帝位后，开始暗中煽动对基督宗教的反对，并巧妙地引入异教。}

当\Person{尤利安}发现自己成为帝国的唯一主宰时，他下令在整个东方重新开放所有异教神庙；那些被废弃的应予修缮；那些已沦为废墟的应予重建，并恢复祭台。他为此拨出了大量资金；他恢复了城市中的古代风俗和祖先礼仪，以及献祭的实践。\par

他亲自奠祭并公开献祭；向热心举行这些仪式的人赐予荣誉；恢复入教者、司祭、圣职解说者和偶像仆役的旧有特权；并确认了之前皇帝为他们制定的法律；他免除了他们的义务和其他负担，正如他们先前的权利；他向神庙守护者恢复了已被废除的供应，并命令他们戒除肉类，并根据外教的说法，戒除一切与宣称要过洁净生活的人相宜之物。\par

他还下令，应将\Place{尼罗河}水位计、符号和先前的祖先牌位保存在\Place{塞拉皮斯神庙}中，而非按照\Person{君士坦丁}所制定的规定存放在教堂里。他频繁写信给那些他知道异教兴盛的城市居民，敦促他们请求他们可能渴望的礼物。相反，他对基督徒公开表示厌恶，拒绝以他的临在尊荣他们，或接见那些被派去报告冤屈的代表。\par

当\Place{尼西比斯}的居民向他请求援助以抵御即将入侵罗马领土的波斯人时，他拒绝帮助他们，因为他们完全基督信仰化，既不愿重新开放他们的神庙，也不前往圣地；他威胁说，除非他听到他们已返回异教，否则他不会帮助他们，不会接见他们的使节，也不会进入他们的城市。\par

他还指控\Place{巴勒斯坦}的\Place{君士坦提亚}居民依恋基督宗教，使他们的城市成为\Place{加沙}的附属。\Place{君士坦提亚}，正如我们之前所说，原先称为\Place{马朱马}，被用作\Place{加沙}船只的港口；但\Person{君士坦丁}听说其居民多数是基督徒后，将其提升为城市，并以自己儿子的名字命名，并给予独立的政府形式；因为他认为它不应隶属于一个沉迷异教仪式的城市\Place{加沙}。\Person{尤利安}登基后，\Place{加沙}居民起诉\Place{君士坦提亚}居民。皇帝亲自担任法官，并判决\Place{加沙}胜诉，命令\Place{君士坦提亚}成为该城市的附属，尽管它距离\Place{加沙}有二十斯塔迪亚之遥。\par

他被废弃了其原名，从此被称为\Place{加沙}的沿海地区。他们现在有相同的市府官员、军事长官和公共规章。然而，在教会事务方面，他们仍被视为两个城市。他们各有自己的主教和自己的圣职界；他们各自庆祝纪念本城殉道者的节日，并纪念先后治理他们的司铎；相邻田地的边界仍被保留，这些田地划分了属于主教的祭台。\par

在我们记忆所及之内，曾发生这样一件事：加萨主教在\Place{马余玛}教会的主持者去世后，试图将该城的圣职人员并入自己的管辖之下；他提出的理由是：一个城由两位主教治理是不合规定的。马余玛居民反对这一计划，省议会审理了此项争端，并另立了一位主教。省议会决定：那些因虔诚而被视为配得上城市荣誉的人，完全不应该因一位异教皇帝基于不同理由的决定，而被剥夺教会圣职和地位所赋予的特权。\par

但这些事件发生在目前所讨论的时期之后。\par

% structure: p0030 source=h2 target=subsection reason=relative-heading-rank; base=h2

\subsection{第四章  尤利安加害\Place{凯撒勒雅}居民。\Place{卡尔西顿}主教\Person{马里斯}的勇敢忠信。}

大约在同一时期，皇帝将\Place{卡帕多细雅}的大而富庶的都市——位于\Place{阿尔革乌斯山}附近的\Place{凯撒勒雅}——从城市名单中删除，甚至剥夺了它在\Person{克劳狄凯撒}统治期间获得的\Quote{凯撒勒雅}之名（原名\Place{马匝卡}）。他长久以来对这座城市的居民怀有极深的厌恶，因为他们热诚信奉基督信仰，并曾摧毁了供奉祖先\Person{阿波罗}的神庙和该城守护神\Person{朱庇特}的神庙。唯一尚存于城中的\Person{福尔图娜}神庙，在他即位后被基督徒推倒；闻知此事后，他痛恨整座城市，几乎无法容忍。他还责备少数异教徒，说他们本应赶赴神庙，若有需要，应甘愿为\Person{福尔图娜}殉道。他下令严加搜查并没收属于\Place{凯撒勒雅}城及其郊区的所有教会财产和金钱；由此搜得约三百磅黄金，被送入国库。他又命令所有圣职人员编入省总督麾下的军队，这在罗马人中被视为最艰苦且最不体面的服役。\par

他下令对基督徒居民进行人口普查，包括妇女和儿童，并像对村庄课税那样向他们征收重税。\par

他进一步威胁说，除非他们的神庙迅速重建，否则他的愤怒不会平息，反而要临到这座城市，直至没有一个\OriginalTerm{加里肋亚人}{Galileans}存活（这是他惯于用来嘲弄基督徒的名称）。毫无疑问，若不是死亡迅速介入，他的恐吓必会完全实现。\par

他起初对待基督徒比先前迫害者更为人道，这并非出于对基督徒的同情，而是因为他发现异教并未从酷刑中获得任何好处，反而基督信仰特别增长，并因那些为捍卫信德而死者之坚毅而更加受尊崇。\par

纯粹是由于嫉妒他们的光荣，他没有像以前的迫害者那样用火与剑对付他们、虐待他们的身体，或把他们扔进海里、活埋以逼其改变信仰，而是诉诸辩论和劝说，试图以此使他们皈依异教；他期望通过放弃一切暴力手段并表现出意外的仁慈，更易达到目的。据说有一次，当他在君士坦丁堡的\Person{福尔图娜}神庙献祭时，\Place{卡尔西顿}主教\Person{马里斯}出现在他面前，公开斥责他为不敬神的人、无神论者和背教者。尤利安无以回敬，只能指责他的失明（因年迈视力衰退，由一个孩童牵引）。尤利安按照惯常亵渎基督的方式，随后嘲弄道：\Quote{你的天主，那位\OriginalTerm{加里肋亚人}{Galilean}，不会治愈你。}\Person{马里斯}回答说：\Quote{我为我的失明感谢天主，因为它使我不能看见一个背离我们宗教的人。}尤利安没有回答就走了，因为他认为通过对基督徒表现出个人的、意外的忍耐和温和，会使异教更加进展。\par

% structure: p0036 source=h2 target=subsection reason=relative-heading-rank; base=h2

\subsection{第五章  尤利安恢复基督徒自由，以便在教会中制造更多麻烦。他所策划的虐待基督徒的方式。}

出于这些动机，尤利安召回所有在\Person{君士坦提乌斯}统治时期因宗教信仰而被流放的基督徒，并归还他们被依法没收的财产。他吩咐人们不要对基督徒行任何不义，不要侮辱他们，也不要强迫他们违背意愿献祭。他命令，如果他们自愿走近祭坛，必须首先平息异教徒认为能避祸的邪魔之怒，并按常规的赎罪方式洁净自己。然而，他剥夺了\Person{君士坦丁}赐予圣职人员的豁免权、荣誉和供给；废除了为他们制定的法律，并恢复了他们的法定义务。他甚至强迫那些因贫穷而被列入圣职人员名单的贞女和寡妇，退还公家拨给他们的供给。因为\Person{君士坦丁}在调整教会世俗事务时，从每个城市的税收中拨出足够的部分，用于维持各地的圣职人员；为确保这一安排稳定，他制定了一项法律，从尤利安去世后一直有效至今。据说这些行为极其残酷苛刻，从收款人向被勒索者出具的收据可以看出，这些收据旨在表明根据\Person{君士坦丁}法律收到的财产已被退还。\par

但是，无论如何，皇帝的敌意并未因宗教而减弱。他怀着对信仰的刻骨仇恨，抓住每一个机会摧毁教会。他剥夺了教会的财产、奉献物和圣器，并判处那些在\Person{君士坦丁}和\Person{君士坦提乌斯}统治期间拆毁神庙的人重建神庙，或支付重建费用。基于此，由于他们无力支付款项，也由于对圣金的追查，许多司铎、圣职人员和其他基督徒遭到残酷的折磨并被投入监狱。\par

从以上所述可以得出结论：尽管\Person{尤利安}流的血比之前教会的迫害者少，并且他为折磨身体设计的刑罚更少，然而在其他方面他却更为严厉；因为他似乎在各个方面都对教会施加祸害，除了他召回了曾被\Person{君士坦提乌斯}皇帝判处流放的司铎；但据说他发布这道命令是为了他们，并非出于仁慈，而是为了通过他们之间的争竞，使教会陷入兄弟纷争，并使其丧失自身权利，或者因为他想诋毁\Person{君士坦提乌斯}；因为他以为，通过偏袒与他意见相同的异教徒，以及对那些为基督受苦者施以怜悯，作为不公正对待，他可以使已故君主几乎为所有臣民所憎恨。他将宦官逐出宫廷，因为已故皇帝对他们好感。他判处宫廷总管\Person{优西比乌}死刑，因为他怀疑是他的建议导致了其兄\Person{加卢斯}被杀。他将欧诺米异端的领袖\Person{阿埃提乌斯}从\Person{君士坦提乌斯}流放他的地区召回，此人因与\Person{加卢斯}的亲密关系而另受怀疑；\Person{尤利安}给他写了充满仁慈的信，并为他提供公家交通工具。出于类似原因，他判处\Place{基齐库斯}的主教\Person{厄勒乌西乌斯}以最重的刑罚，令其两个月内自费重建他在\Person{君士坦提乌斯}任内拆毁的诺瓦替安派教堂。还可以发现许多其他出于对前任的仇恨而做的事情，要么他亲自完成，要么允许他人完成。\par

% structure: p0040 source=h2 target=subsection reason=relative-heading-rank; base=h2

\subsection{第六章 \Person{亚大纳削}在一位聪慧美貌的贞女家中藏匿七年后，再次公开出现并进入\Place{亚历山大里亚}的教堂。}

在此期间，长期藏匿的\Person{亚大纳削}听闻\Person{君士坦提乌斯}的死讯，夜间出现在\Place{亚历山大里亚}的教堂中。他的突然出现引起了极大的惊异。他曾逃脱\Place{埃及}总督的逮捕，该总督奉皇帝之命并应\Person{乔治}之友的请求，如前所述，计划逮捕他；他藏匿在\Place{亚历山大里亚}一位圣洁贞女的家中。据说她拥有非凡的美丽，以至于见到她的人都视她为自然之奇观；而那些拥有节操和谨慎的男人都远离她，以免因猜疑而受到指责。她正值青春芳华，极其谦逊和谨慎，这些品质即使在天性未赐予美貌的情况下，也足以使身体显得精致优美。因为并不是像有些人断言的那样：\Quote{身体是什么样，灵魂也是什么样。}相反，身体的状态是由灵魂的活动所塑造的；任何积极活动的人，只要他如此积极活动，就会表现出那种性质。\par

我相信，所有仔细研究过这个问题的人都承认这真理。据说\Person{亚大纳削}是奉天主的启示逃到这位圣洁贞女的家中，天主以此方式拯救他。\par

当我思考由此产生的结果时，我无法怀疑这一切事件都是由天主引导的；这样，如果任何人企图因\Person{亚大纳削}而骚扰他的亲属，并迫使他们发誓，他们也不至于痛苦。在一位如此可爱的贞女家中藏匿一位司铎，没有任何值得怀疑之处。然而，她有勇气接待他，并通过她的谨慎保全了他的性命。她是他最忠实的守护者和勤勉的仆人；因为她为他洗脚，为他送食物，而且独自服侍其他一切自然在严格需求中所要求的事情；他需要的书籍，她通过他人照料；在提供这些服务的长久时间里，\Place{亚历山大里亚}的居民无人知晓。\par

% structure: p0044 source=h2 target=subsection reason=relative-heading-rank; base=h2

\subsection{第七章 \Place{亚历山大里亚}主教\Person{乔治}的暴死与凯旋。密特拉神庙中某些事件的结果。\Person{尤利安}关于此严重事件的信件。}

在\Person{亚大纳削}这样被保全并突然出现在教堂之后，没有人知道他来自何处。然而，\Place{亚历山大里亚}的民众为他的归来而欢喜，并将他的教堂归还给他。\par

亚略派人士既从教会中被逐出，只得在私宅中聚会，并立\Person{路济乌斯}代替\Person{乔治}为其异端的主教。\Person{乔治}此前已被杀害；因为当官员向公众宣布\Person{君士坦提乌斯}驾崩、\Person{尤利安}独揽大权时，\Place{亚历山大}的外教人起来叛乱。他们用呼喊和辱骂攻击\Person{乔治}，仿佛要立刻杀死他。这些仓促攻击的抵抗者随后将他投入监狱；不久之后，他们在清晨冲进监狱，杀死了他，将尸体扔到骆驼上，并在白天百般侮辱后，于夜间焚烧了它。我并非不知亚略异端者声称\Person{乔治}遭受了\Person{阿塔纳修斯}追随者的残酷对待；但在我看来，这些行为的实施者更可能是外教人；因为他们比任何其他人更有理由恨他，尤其因为他侮辱了他们的偶像和庙宇，并且还禁止他们献祭或举行祖传仪式。此外，他在宫廷中获得的影响力加深了人们对他的仇恨；正如民众对掌权者常有的态度，他们认为他不可忍受。\par

在名为\Place{密特里乌姆}的地方也发生了一场灾祸；那里原本是一片荒野，\Person{君士坦提乌斯}曾将其赐予\Place{亚历山大}的教会。当\Person{乔治}清理地基以建造祈祷之所时，发现了一个内殿。里面发现了一些偶像和某些用于启蒙或成全的器物，在观者看来显得可笑而怪异。基督徒们将它们公开展出，并举行游行以刺激外教人；但外教人聚集成群，手持刀剑、石头和随手拿到的武器，冲向并攻击基督徒。他们杀死了许多基督徒，并为了嘲笑其宗教，将另一些人钉在十字架上，还使许多人受伤。\par

这导致基督徒已开始的工作被放弃，而外教人一听说\Person{尤利安}即位为帝，就杀害了\Person{乔治}。这一事实为皇帝本人所承认，若非受真理所迫，他不会承认；因为我想，他宁愿让基督徒（无论他们是谁）而不是外教人成为杀害\Person{乔治}的凶手；但真相无法隐藏。这在他就此写给\Place{亚历山大}居民的信中显而易见，信中表达了他严厉的看法。在这封信中，他只谴责而略过了惩罚；因为他说他敬畏他们的守护神\Person{塞拉比斯}、他们的建立者\Person{亚历山大}，以及他的叔叔\Person{尤利安}，后者曾任\Place{埃及}和\Place{亚历山大}总督。这叔叔如此偏爱外教人、极端憎恨基督教，以至于违背皇帝意愿，将基督徒迫害至死。\par

% structure: p0049 source=h2 target=subsection reason=relative-heading-rank; base=h2

\subsection{第八章 关于\Place{安提约基雅}圣器保管者\Person{狄奥多尔}。叛徒的叔叔\Person{尤利安}如何因这些圣器而成为蠕虫的猎物。}

据说当皇帝的叔叔\Person{尤利安}一心要搬走\Place{安提约基雅}教会的许多贵重奉献物，并将其放入皇家宝库，同时关闭祈祷之所时，所有神职人员都逃走了。唯有一位名叫\Person{狄奥多里托}的长老没有离开城市；\Person{尤利安}抓住他，作为宝物保管者，且能提供有关宝物的情况，便极其残暴地虐待他；最后，在每一种酷刑下他都勇敢地回应，并因信仰宣认而备受称许之后，\Person{尤利安}下令用剑将他处死。当\Person{尤利安}掠夺了圣器后，他将它们扔在地上并开始嘲弄；在尽情亵渎\Person{基督}之后，他坐在圣器上并加重了其侮辱行为。立时，他的生殖器和直肠腐烂；它们的肉变成腐臭，化为蛆虫。这病超出了医生的技能。然而，出于对皇帝的敬畏和惧怕，他们尝试各种药物，并宰杀最昂贵、最肥的禽鸟，将脂肪敷在腐烂的部分，希望蛆虫因此被引到表面，但毫无效果；因为蛆虫深埋，钻入活肉，不停啃咬直至结束他的生命。看来这场灾祸是神怒的惩罚，因为皇家宝库保管者和其他宫廷高官曾嘲弄教会，都因神怒而死于异常可怕的方式。\par

% structure: p0051 source=h2 target=subsection reason=relative-heading-rank; base=h2

\subsection{第九章 \Person{圣尤西比乌}、\Person{圣内斯塔布斯}和\Person{圣泽诺}在\Place{迦萨}城的殉道。}

我既已在我的历史中讲述至此，并记述了\Person{革奥尔革}及\Person{特奥多里托}之死，我以为宜当叙述关于三位兄弟——\Person{欧瑟比}、\Person{奈斯塔布}和\Person{泽农}——之死的若干详情。\Place{加沙}居民因激怒于他们，便将其从藏身的家中拖出，投入监牢，加以殴打。然后他们聚集在剧场，大声呼喊指控他们，声称他们在其神庙中犯有亵渎之罪，并利用过去的机会伤害和侮辱异教。通过这些呼喊，以及互相怂恿杀害兄弟，他们变得狂怒；当彼此煽动后，如同暴乱的群众惯常所为，他们冲向监牢。他们极其残忍地对待这些人；有时脸朝地，有时背朝地，受害者被拖拽，路面将其撞碎。我听说甚至妇女放下纺锤，用织布梭子刺他们，市场上的厨师从台子上抓起沸腾冒着热气的锅，将热水浇在受害者身上，或用铁签刺穿他们。当他们剥去其肉、砸碎其头颅，以致脑浆流到地上后，他们的尸体被拖出城外，扔在通常用作兽类尸体堆置的地方；然后点燃大火，焚烧尸体；未被火焚尽的骨头与骆驼和驴的骨头混在一起，以免轻易被找到。但这些骨头并未隐藏太久；因为一位基督徒妇女，虽是\Place{加沙}居民，但非本地人，在天主的指引下夜间收集了骨头。她将其放入陶罐，交给他们的表兄弟\Person{泽农}保管，因为天主在梦中如此指示她，并告诉了那位妇女这人的住处：她尚未见到他时，他便被显示给她，因为她先前不认识\Person{泽农}；当迫害近期激烈时，他一直隐藏。他几乎被\Place{加沙}人抓住并处死；但他趁着人们专注于杀害他表兄弟的时候逃脱，逃往\Place{安特东}，一个距\Place{加沙}约二十斯塔迪翁的沿海城市，同样倾向于异教且沉迷于偶像崇拜。当该城居民发现他是基督徒时，便用棍棒狠狠打他的背，并把他赶出城。他于是逃往\Place{加沙}港口躲藏；那位妇女在那里找到他，将遗骨交给他。他小心地将遗骨保存在家中，直到\Person{特奥多西}皇帝当政时，他被祝圣为主教；他在城外建造了一座祈祷之所，在那里设立祭台，并将殉道者的遗骨放在宣信者\Person{奈斯托尔}的身旁。\Person{奈斯托尔}与他的表兄弟们关系密切，与他们一起被\Place{加沙}人抓住，被监禁并鞭打。但那些拖着他穿过城市的人被他的俊美容貌所打动；因怜悯，他们还未等他完全断气就将他扔出城外。有些人发现了他，将他抬到\Person{泽农}家中，他在伤口包扎时咽了气。当\Place{加沙}居民开始反思其罪行的严重性时，他们战栗不安，生怕皇帝会惩罚他们。\par

曾有传言说皇帝极为愤怒，并决定惩罚十人团；但此传言是虚假的，除了罪犯们的恐惧和自责外毫无根据。\Person{尤利安}非但没有像在\Place{亚历山大城}因\Person{革奥尔革}之死而对那里的人那样表现出愤怒，甚至没有写信斥责\Place{加沙}人。相反，他罢免了该省总督，并视其为可疑，声称仅因仁慈才未将其处死。对总督的控罪是：他曾逮捕了一些据报最先引发骚乱和谋杀的\Place{加沙}居民，并将他们监禁起来，等待依法审判。\Quote{他有什么权利，}皇帝问道，\Quote{只因公民们对几个加里肋亚人报复他们及其神祇所受的伤害，就逮捕他们？}据说此事属实。\par

% structure: p0054 source=h2 target=subsection reason=relative-heading-rank; base=h2

\subsection{第十章 论圣\Person{依拉里昂}与\Place{赫利奥波利斯}被猪毁灭的贞女。\Place{阿雷图萨}主教\Person{玛尔克}的奇特殉道。}

同一时期，\Place{加沙}居民搜寻隐修士\Person{依拉里昂}；但他已逃往\Place{西西里}。他在那里于旷野和山中拾柴，并肩负进城售卖，以此获得维持身体的食物。但当他被一位他曾驱魔的贵人认出后，他便退到\Place{达尔马提亚}；在那里，他藉天主的德能行了许多奇迹，并通过祈祷遏止了海潮，使海浪返回其边界；他再次离开，因为与赞美他的人同住并不令他喜悦；但他变换住所时，希望不被注意，并通过频繁迁徙摆脱关于他的名声。最终他航行前往\Place{塞浦路斯}岛，但在\Place{帕福斯}靠岸；应\Place{塞浦路斯}主教的恳求，他喜爱那里的生活，并在一个名为\Place{哈尔布里斯}的场所实践哲学。\par

在那里他仅藉逃跑逃脱了殉道；因为他服从了神圣诫命，该诫命命我们不可自陷于迫害；但若落入迫害者之手，则要以我们自身的刚毅战胜压迫者的暴力。\par

我所描述的对基督徒施加这种暴行的不仅限于加沙和亚历山大的居民。黎巴嫩山附近的赫利奥波利斯以及叙利亚的阿雷图萨的居民似乎在残暴程度上更胜一筹。前者犯下了一种几乎难以置信的野蛮行径，若非有目睹者的证言佐证。他们剥去从未被众人瞩目的圣童贞女的衣服，使其赤身裸体，作为公开的\OriginalTerm{奇观}{spectacle}和侮辱的对象。在诸多其他折磨之后，他们最终剃光她们的头发，剖开她们的腹部，将通常喂给猪的食物藏入她们的内脏；由于猪无法分辨，却因对惯常食物的需求所驱使，它们也撕裂了人肉。\par

我相信，赫利奥波利斯的居民对圣童贞女犯下这种野蛮行径，是因为禁止了古代习俗，即女子在与未婚夫结婚前，须先与任何偶然相遇者行淫。这一习俗被君士坦丁颁布的法律所禁止，他在摧毁赫利奥波利斯的维纳斯神庙，并在其废墟上建立一座教堂之后。\par

阿雷图萨的主教马尔谷，一位因白发与生活而可敬的老人，被该城居民处以极其残酷的死刑。这些居民长期对他怀有敌意，因为在君士坦丁统治时期，他更以精神而非说服的方式将异教徒提升至基督教，并拆毁了一座最神圣壮丽的神庙。尤利安即位后，他看到民众对主教群情激愤；一道敕令发布，命令主教要么支付重建费用，要么重建神庙。他思考到前者不可能，后者对基督徒、尤其是对司铎而言是非法的，于是起初逃离了城市。然而，当他听说许多人因他而受苦，有些人被拖到法庭，另一些人受酷刑时，他返回并自愿承受民众可能对他施加的任何惩罚。全体民众非但没有更加钦佩他，认为他表现出合乎哲学家的行为，反而认为他出于对他们的蔑视，于是冲向他，将他拖过街道，挤压、拉扯和击打他们碰到的任何肢体。各性别、各年龄段的人们都带着敏捷和狂怒参与这场残暴的行为。他的耳朵被细绳割下；经常上学的男孩们戏弄他，将他抛向空中并翻滚，将他向前送，又接住他，毫不留情地用他们的笔刺他。当他全身布满伤口，却仍然呼吸时，他们用蜂蜜和某种混合物涂抹他，将他放入一个用灯心草编织的鱼篮中，升高到一个高处。据说当他在这个姿势中，黄蜂和蜜蜂落在他身上吞噬他的肉时，他对阿雷图萨的居民说，他被高举在他们之上，可以俯视他们，这让他想起了他们在未来生活中将存在的差异。还相传那位地方长官，虽是异教徒，却品行高尚，至今在该国仍受纪念，他钦佩马尔谷的自制，并大胆地斥责皇帝允许自己被一位暴露于无数酷刑的老人击败；他补充说，这样的行为使皇帝蒙羞，而被迫害的人的名声却同时变得光辉。就这样，那位有福的人以如此坚定的刚毅承受了阿雷图萨居民对他施加的所有折磨，甚至连异教徒也称赞他。\par

% structure: p0060 source=h2 target=subsection reason=relative-heading-rank; base=h2

\subsection{第十一章．论马其顿纽斯、特奥杜卢斯、格拉提安、布西里斯、巴西尔和尤普西基乌斯，他们在那个时期殉道。}

大约同一时期，马其顿纽斯、特奥杜卢斯和塔提安（他们是弗里吉亚人）勇敢地忍受了殉道。弗里吉亚城米索斯的一座神庙，在关闭多年后，被该省总督重新开放，这些殉道者夜间进入其中，摧毁了神像。当其他人被逮捕并即将因该行为受罚时，他们坦承自己是行动者。他们本可通过向偶像献祭而逃脱进一步的惩罚；但总督无法说服他们接受这些条件的赦免。他的劝说无效，便以多种形式虐待他们，最终将他们放在一个点着火的大烤架上。当他们被烧时，他们对总督阿马库斯（因为那是他的名字）说：\Quote{如果你想要煮熟的肉，下令将我们的身体翻转到另一面烤火，以免在你看来我们半生不熟。}就这样，这些人高贵地承受，并在刑罚中献出了生命。\par

据说布西里斯也在伽拉提亚城安基拉因他光辉而最英勇的信仰宣认而闻名。他属于被称为恩克拉提特派的异端；该省总督逮捕了他，并因他嘲笑异教徒而意图虐待他。他公开将他带到刑讯室，命令将他吊起。布西里斯将双手举到头上，使两侧暴露，并告诉总督，行刑者把他吊到刑具上再放下来是多余的，因为他已经准备好，无需如此就能接受所需的各种酷刑。总督对此提议感到惊讶；但他的惊异因随后的事而增加，因为布西里斯坚守不动，双手高举，按照总督的指示，他的两侧被钩子撕扯，承受打击。紧接着，布西里斯被投入监狱，但不久后因尤利安死讯宣布而被释放。他活到狄奥多西统治时期，放弃了他先前的异端，加入了天主教会。\par

据说大约在这时期，安基拉教会司铎巴西略和卡帕多细亚凯撒勒雅的贵族欧普绪基——他刚娶妻，仍为新郎——以殉道结束了生命。我相信欧普绪基是因拆毁命运神庙而被判刑的，正如我前述，此事激怒了皇帝对凯撒勒雅全体居民。事实上，所有参与此事者都已被判刑，有的处死，有的流放。巴西略长久以来在捍卫信仰方面表现出极大的热忱，并在君士坦提乌斯在位期间反对亚略派；因此，欧多克西的同党禁止他举行公开集会。但朱利安登基后，他四处奔走，公开而坦率地劝勉基督徒坚持自己的教义，并戒绝以异教祭献和奠祭玷污自己。他敦促他们将皇帝可能赐予的荣誉视为无物，因为这类荣誉转瞬即逝，且导致永久的耻辱。他的热忱已使他成为异教徒猜疑和憎恨的对象。一日他恰巧经过，看见他们献祭。他深深叹息，发出祈祷，愿没有任何基督徒陷入类似谬误。他当场被拿住，解送至行省总督处。他受了许多酷刑，在勇敢忍受这痛苦中，他获得了殉道的冠冕。\par

即使这些暴行是违反皇帝意愿而犯下的，但它们仍足以证明，他的统治期间有不少不凡的殉道者。\par

为清晰起见，我将所有这些事件集中叙述，尽管殉道实际上发生在不同时期。\par

% structure: p0066 source=h2 target=subsection reason=relative-heading-rank; base=h2

\subsection{第十二章　论西方主教路济弗尔和欧瑟伯。欧瑟伯与亚大纳削及其他主教在亚历山大里亚召开会议，通过定义圣神与圣父圣子同质来确认尼西亚信仰；他们关于本质与位格的法令。}

亚大纳削返回后，撒丁尼亚卡利亚里主教路济弗尔和意大利利古里亚城韦尔切利主教欧瑟伯从上底拜依斯返回。他们曾被君士坦提乌斯判处终身流放该地。为规整和全面系统化教会事务，欧瑟伯来到亚历山大里亚，并与亚大纳削共同召开会议，以确认尼西亚教义。\par

路济弗尔派一位执事随同欧瑟伯代他出席会议，而他自己去了安提约基雅，探望该教会在纷扰中的状况。\par

当时亚略派在欧佐依领导下，以及梅利提乌斯的追随者——如上所述，他们甚至与赞同自己见解的人也意见不合——挑起了一个分裂。由于梅利提乌斯当时尚未从流放地返回，路济弗尔祝圣了保利诺为主教。\par

与此同时，许多城市的主教已与亚大纳削和欧瑟伯在亚历山大里亚集会，并确认了尼西亚教义。他们宣认圣神与圣父圣子同体，并使用了\Quote{天主圣三}一词。\par

他们宣告，天主圣言所摄取的人性不仅包括完全的身体，也包括完全的灵魂，正如古代教会哲学家所教导的。因教会对\Quote{本质}与\Quote{位格}这两个术语有争议，且关于这些词的争辩频繁，他们明智地（我认为）规定，今后在涉及天主的开端处不应使用这些术语，除非为了驳斥撒伯里乌的见解；以免因术语匮乏，同一事物似乎有三个名称；而应让人藉各自独特的术语以三重方式理解每个事物。\par

这些是在亚历山大里亚集会的主教们所通过的法令。亚大纳削在会议中宣读了他为自辩所写的关于自己逃亡的文件。\par

% structure: p0073 source=h2 target=subsection reason=relative-heading-rank; base=h2

\subsection{第十三章　论安提约基雅的司祭长保利诺和梅利提乌斯；欧瑟伯与路济弗尔如何彼此对立；欧瑟伯与希拉留辩护尼西亚信仰。}

会议结束后，欧瑟伯前往安提约基雅，发现百姓中纷争不已。那些依附梅利提乌斯的人不愿加入保利诺，而是各自分开举行集会。欧瑟伯对事态深感痛心；因为祝圣本应在全体人民一致同意下进行；然而出于对路济弗尔的尊重，他没有公开表达不满。\par

他拒绝与任何一方共融，但承诺通过会议调解双方的怨愤。正当他努力恢复和谐与一致时，梅利提乌斯从流放地返回，发现持他见解的人已脱离另一方，于是他在城外与他们举行集会。与此同时，保利诺在城内召集他自己的会众；因为他的温和、圣善的生活和高龄，竟赢得了亚略派首领欧佐依的尊重，不仅未被逐出城，反而被分配了一座教堂供他自己使用。欧瑟伯发现他为恢复和谐所做的一切努力受挫后，离开了安提约基雅。路济弗尔自觉受了他伤害，因为欧瑟伯拒绝批准保利诺的祝圣；他不悦之下，与欧瑟伯断绝了共融。仿佛纯粹出于争辩的欲望，路济弗尔开始谴责亚历山大里亚会议的决议；这样他似乎创生了以他命名的异端——路济弗尔派。\par

接受他主张的人脱离了教会；但尽管他对事态的发展深感懊恼，然而因他曾派一位执事随同欧瑟伯代他出席，他仍服从亚历山大里亚会议的决议，并与天主教会教义保持一致。大约在这时期，他返回撒丁尼亚。\par

与此同时，\newline{}\Person{Eusebius} 遍历东方各省，使那些背弃信仰者回归，并教导他们必须相信什么。经过\Place{Illyria}后，他去了\Place{意大利}，在那里遇到了\Place{Aquitania}的\Place{Poictiers}主教\Person{Hilarius}。\Person{Hilarius}在\Person{Eusebius}之前就已从流放中返回，并教导意大利人和高卢人应当接受什么教义、拒绝什么教义；他用拉丁语雄辩地表达自己，据说写了许多杰出的著作反驳\OriginalTerm{亚略派的教义}{Arian dogmas}。这样，\Person{Hilarius}和\Person{Eusebius}在西方地区维护了\OriginalTerm{尼西亚大公会议}{Nicæan council}的教义。\par

% structure: p0078 source=h2 target=subsection reason=relative-heading-rank; base=h2

\subsection{第十四章。马其顿派与亚略派就\Person{阿卡西乌}一事发生争论。}

此时，\Person{Macedonius}的追随者——其中包括\Person{Eleusius}、\Person{Eustathius}和\Person{Sophronius}，他们现在开始公开被称为\OriginalTerm{马其顿人}{Macedonians}，构成了一个独立的宗派——在\Person{Constantius}去世后采取了大胆的措施，召集了那些曾在\Place{Seleucia}集会且与自己见解相同的人，并举行了几次公会议。他们谴责了\OriginalTerm{阿卡西乌派}{partisans of Acacius}以及在\Place{Ariminum}所确立的信仰，同时确认了在\Place{安提约基雅}提出、后来又经\Place{Seleucia}批准的教义。\par

当被问及他们与阿卡西乌派争论的原因时——尽管他们原与阿卡西乌派见解相同，也曾共融——他们通过\Place{Paphlagonia}主教\Person{Sophronius}之口回答说：西方基督徒坚持使用\OriginalTerm{同本体}{consubstantial}一词，而东方的\Person{Aëtius}追随者则维护\OriginalTerm{本体不相似}{dissimilarity as to substance}的教义；前者通过使用\Quote{同本体}一词，将圣父和圣子的不同位格不规则地编织成一体，而后者则描绘了圣父与圣子本性之间的关系中存在着过大的差异；但他们自己保持了两种极端之间的中庸之道，避免了两种错误，虔诚地维护圣子在\OriginalTerm{本体}{hypostasis}上与圣父相似。马其顿派正是通过这样的申辩来为自己辩护。\par

% structure: p0081 source=h2 target=subsection reason=relative-heading-rank; base=h2

\subsection{第十五章。\Person{阿塔纳修}再次被放逐；关于\Place{基齐库斯}主教\Person{埃勒乌西乌}和\Place{博斯特拉}主教\Person{提多}；提及作者的祖先。}

皇帝得知\Person{Athanasius}在\Place{Alexandria}的教堂里举行集会，勇敢地教导百姓，并使许多异教徒归信基督教后，便命令他以最严厉的刑罚离开\Place{Alexandria}。强制执行这一敕令的借口是：\Person{Athanasius}在受到\Person{Constantius}放逐后，未经在位皇帝的许可便重新担任主教之职；因为\Person{Julian}宣称，他从未打算让被\Person{Constantius}放逐的主教们恢复教会职务，而只是让他们返回故乡。在宣布立即离开的命令时，\Person{Athanasius}对围在他身边哭泣的基督徒群众说：\Quote{你们要鼓起勇气；这不过是一片云，很快就会消散。}说完这些话，他告别了；然后他把教会的管理托付给最热心的朋友，离开了\Place{Alexandria}。\par

大约在同一时期，\Place{Cyzicus}的居民派遣使团到皇帝那里，向他陈述一些私事，特别是恳求恢复异教神庙。他赞赏他们的先见之明，并答应满足他们所有的请求。他驱逐了该城的主教\Person{Eleusius}，因为他拆除了一些神庙，用侮辱的方式亵渎了神圣区域，为寡妇提供住所，为圣洁的贞女建造房屋，并引导异教徒放弃他们祖传的礼仪。\par

皇帝禁止一些陪伴他的外国基督徒进入\Place{Cyzicus}城，似乎是出于担心他们会与城内的基督徒一起因宗教问题煽动叛乱。有许多人与他们聚集在一起，也持有与城内基督徒相同的宗教观点，他们从事羊毛制品生产以供应国家，并铸造钱币。他们人数众多，分为两个人口众多的阶层；他们得到先前皇帝的许可，可以带着妻子和财产居住在\Place{Cyzicus}，条件是每年向国库提供士兵的服装和新铸造的钱币。\par

尽管\Person{尤利安}渴望用一切手段推进异教，但他仍认为对拒绝献祭的人施以武力或报复是极为不智的。此外，每个城市都有如此多的基督徒，以至于统治者甚至很难数清他们。他甚至没有禁止他们聚集敬拜，因为他知道当自由意志受到质疑时，强制是完全无用的。他将圣职人员和教会首领逐出所有城市，以终止这些集会，并正确地说，若没有他们，民众的聚会将确实解散；如果确实没有召集教会的人，也没有人教导或分施奥秘，那么宗教本身迟早会落入遗忘。他为此举提出的借口是，圣职人员是民众中叛乱的领导者。以此为由，他将\Person{厄肋乌西乌}及其朋友逐出\Place{基齐库斯}，尽管该城连一点叛乱的迹象或征兆都没有。他还公开呼吁\Place{博斯特拉}市民驱逐他们的主教\Person{蒂托}。看来皇帝曾威胁要弹劾\Person{蒂托}和其他圣职人员，称他们是民众中可能发生的任何叛乱的制造者；\Person{蒂托}于是写信告诉他，虽然基督徒在人数上接近异教徒，但根据他的劝勉，他们愿意保持安静，不太可能起来叛乱。\Person{尤利安}为了避免引起\Place{博斯特拉}居民对\Person{蒂托}的敌意，在致他们的信中表示，他们的主教对他们进行了诽谤，说他们不叛乱是出于他的劝勉而不是出于他们自己的意愿；\Person{尤利安}劝他们将\Person{蒂托}作为公敌逐出城市。\par

看来在其他地方，基督徒也遭受了类似的不公；有时出于皇帝的命令，有时由于民众的愤怒和冲动。这些行为的责任可以公正地归咎于统治者；因为他没有使违法者受到法律的强制，反而出于对基督宗教的仇恨，只用言语斥责肇事者，而他的行动却促使他们继续如此。因此，虽然基督徒并未受到皇帝的绝对迫害，但他们不得不从一个城市逃往另一个城市，从一个村庄逃往另一个村庄。我的祖父和许多祖先被迫这样逃亡。我的祖父出身异教，和他的家人以及\Person{阿拉菲昂}的家人一起，是\Place{贝特利亚}最先接受基督教的人。\Place{贝特利亚}是\Place{加沙}附近一个人口众多的城镇，那里有一些因古老和建筑精美而深受当地人尊敬的庙宇。其中最著名的是万神庙，建在一座人工高地上，俯瞰整个城镇。据推测，该地方因这座庙宇而得名，这座庙宇最初的名字是叙利亚语，后来译为希腊语，用一个表示\Quote{所有神明居所}的词来表达。\par

据说上述家庭是通过隐修士\Person{希拉里翁}的媒介而归化的。\Person{阿拉菲昂}似乎被魔鬼附身；异教徒和犹太人都无法用任何咒语和法术解救他脱离此苦；但\Person{希拉里翁}只凭呼求基督之名就驱逐了魔鬼，\Person{阿拉菲昂}和他的全家立刻接受了基督教。\par

我的祖父拥有卓越的自然禀赋，他成功地将其应用于解释圣经；他在通识方面有所成就，并且通晓算术。他深受\Place{阿斯卡隆}、\Place{加沙}以及周围地区基督徒的爱戴；并因其诠释圣经疑难之处的恩赐而被视为宗教必需。对于另一个家庭的美德，无人能以适当的言辞表达。在该地区建立的首批圣堂和隐修院是由该家庭的成员创建，并凭其权势和对旅客及穷人的善举而维持。一些属于该家庭的好人甚至在我们这个时代仍然兴盛；在我青年时，我曾见过其中几位，但他们当时已非常年迈。在叙述历史过程中，我将有机会更多地谈到他们。\par

% structure: p0089 source=h2 target=subsection reason=relative-heading-rank; base=h2

\subsection{第十六章　\Person{尤利安}建立异教并废除我们习俗的努力。他写给异教大司祭的书信。}

皇帝深感悲伤，因为他发现所有确保异教优势的努力完全无效，并看到基督教声名鹊起；因为虽然庙宇大门敞开，虽然献祭，古老节庆在所有城市恢复，但他远未满足；因为他清楚预见，一旦他影响力撤除，整个局面将迅速改变。他特别懊恼地发现许多异教祭司的妻子、儿女和仆人皈依了基督教。考虑到基督教的一个主要支柱是其信奉者的生活和行为，他决定将基督教的秩序和纪律引入异教庙宇，设立各种等级和职务，任命教师和经师来教导异教教义和劝诫，并命令在特定日子和时刻祈祷。此外，他还决心为那些渴望过哲学隐居生活的男女建立隐修院，并为接济陌生人和穷人以及其他慈善目的建立医院。他希望在异教徒中引入基督教对自愿和非自愿过犯的补赎制度；但他最钦佩并希望在异教徒中建立的教会纪律是主教们给外出旅行者的推荐信，其中在任何地方和任何情况下，他们将旅行者推荐给其他主教的好客和善意。\Person{朱利安}就这样努力将基督教的习俗嫁接于异教。但如果我所说似乎不可信，我不必远寻证据来证实我的断言；因为我可以拿出皇帝本人关于此事的信。他写道如下：——\par

致\Place{加拉提亚}大祭司\Person{阿尔萨修斯}。异教尚未达到所期望的繁荣程度，这是由于其信徒的行为。然而，对诸神的崇拜以最宏大最壮丽的规模进行，远超我们的祈祷和希望；愿我们的\Person{阿德拉斯泰亚}对这些话仁慈，因为没有人敢期望在如此短的时间内看到如此广泛和惊人的变化。但我们是否满足于已经实现的？难道我们不应该考虑无神论的进展主要归因于基督徒对陌生人表现出的仁爱、对死者的尊敬以及他们生活中假设的欺骗性的庄重吗？我们每个人都应该勤勉尽责；我指的不只是你一人，因为那还不够，而是指加拉提亚的所有祭司。\par

你必须要么使他们感到羞耻，要么尝试说服的力量，要么剥夺他们的祭司职务，如果他们不与他们的妻子、儿女和仆人一同侍奉诸神，或者如果他们支持加里肋亚人的仆人、儿子或妻子不敬地对待诸神并偏爱无神论而非虔诚。然后劝诫祭司不要经常去戏院，不要在酒馆饮酒，不要从事任何贸易或任何邪恶技艺。\par

尊重那些听从你告诫的人，驱逐那些无视的人。在每个城市建立客栈，以便邻国和远方的陌生人能根据各自需要受益于我们的仁爱。\par

我已提供必要开支的资金，并向整个加拉提亚发出指示，你每年应获得三万蒲式耳谷物和六万桶酒，其中五分之一用于支持侍奉祭司的穷人；其余分配给陌生人和我们自己的穷人。因为，既然犹太人中没有需要的人，甚至不虔敬的加里肋亚人不仅供应自己一方有需要的人，也供应与我们同一立场的人，如果我们让自己的人民遭受贫困，那将是可耻的。\par

教导异教徒合作从事这项慈善工作，让异教城镇的初熟之果献给诸神。\par

使异教徒习惯于这种慷慨，向他们展示这种行为如何得到远古实践的认可；因为\Person{荷马}描述\Person{尤迈奥斯}说——\par

\Quote{我的客人！我若以蔑视对待陌生人，\newline{}便是冒犯，\newline{}即使比你自己更穷的人到来；\newline{}因为所有穷人和所有陌生人都是\Person{宙斯}所顾念的。}\par

让我们不让他人在善行上超过我们；让我们不以暴力羞辱自己，反而让我们在敬拜诸神上居首位。如果我听说你按我的指示行事，我将充满喜乐。不要常到总督府中拜访他们，而要常写信给他们。当他们进入城市时，不要让任何祭司去迎接他们；当他们去诸神庙宇时，不要让祭司陪伴他们超过门廊；在这种情况下，也不要有士兵在他们前面行进；但让愿意跟随的人跟随他们。因为一旦他们进入神圣界限，他们便只是普通个人；因为在那里，如你所知，你的职责是按照神圣法令主持。谦卑遵守此法的人表明他们拥有真宗教；而藐视此法的人则骄傲虚荣。\par

我准备援助\Place{佩西努斯}的居民，只要他们向众神之母求恩；但如果他们忽视这一职责，他们将招致我最大的不悦。\par

\Quote{我本人若在此接待，\newline{}然后护送离开，\newline{}一个被诸神憎恶的人，如同这些，便是犯罪。}\par

\Quote{因此，说服他们，如果他们渴望我的帮助，他们必须向众神之母献上恳求。}\par

% structure: p0102 source=h2 target=subsection reason=relative-heading-rank; base=h2

\subsection{第十七章。为了不被视为暴君，\Person{朱利安}狡猾地对付基督徒。废除十字架记号。他迫使士兵献祭，尽管他们不愿意。}

当尤利安\Person{尤利安}以所述方式行事和写作时，他期望通过这种手段轻易地诱使臣民改变宗教信仰。尽管他热切希望废除基督宗教，但他显然羞于采用暴力手段，以免被认为暴虐。然而，他使用了一切可能设想的办法来引导臣民回归异教；他尤其对士兵施加压力，时而单独训话，时而通过军官传达。为了使他们在一切事上习惯敬拜诸神，他恢复了罗马军队军旗的古老形式——正如我们已说过的，君士坦丁曾奉天主之命将其改为十字圣号。尤利安还下令在官方画像中将自己与诸神并列：要么是朱庇特从天而降，将皇权的象征（王冠或紫袍）授予他，要么是玛尔斯或墨丘利凝视着他，仿佛在赞叹他的口才与军事才能。他将诸神的画像与自己的画像并列，目的是让人民在尊崇他的借口下暗中敬拜诸神；他滥用了古老习俗，试图向臣民隐藏其意图。他认为，若他们在这一点上服从，就会在其他一切事上更乐意服从；但若他们胆敢拒绝服从，他就有理由惩罚他们，视其为破坏罗马传统、冒犯皇帝与国家之人。只有极少数人（法律对他们加以制裁）看穿了他的计谋，拒绝向他的画像行传统敬意；但多数人因无知或单纯，仍遵循旧制，轻率地向他的画像行礼。皇帝从这一诡计中获益甚微；但他并未停止改变宗教的努力。\par

尤利安接下来的计谋比前者更少诡诈、更为暴力；许多宫廷士兵的刚毅因此受到考验。当规定的向军队发放金钱的日子到来时（这日子通常落在罗马某节日的周年，例如皇帝诞辰或某皇城奠基日），尤利安\Person{尤利安}考虑到士兵天性轻率而单纯，且贪财，因此认为这是一个引诱他们敬拜诸神的良机。于是，当每个士兵上前领取金钱时，他被命令献祭——火和乳香已预先按罗马古老习俗放置在皇帝附近。有些士兵有勇气拒绝献祭，也拒绝接受黄金；另一些士兵则因习惯于遵守法律和习俗而照做，不认为自己在犯罪。还有一些士兵被黄金的光泽迷惑，或因恐惧和眼前的考验而被迫顺从异教礼仪，让自己陷入本应逃避的诱惑。\par

据说，有些无知而陷入此罪的人正坐在桌旁互相饮酒时，其中一人碰巧在杯间提起了基督的名。另一位宾客立刻喊道：\Quote{你真奇怪，刚才你还为了皇帝的赏赐而抛洒乳香否认了他，现在却呼求基督\OriginalTerm{基督}{Christ}。}听到这话，他们突然意识到自己所犯的罪；他们起身离席，冲到街上，大声哭喊，向所有人见证他们过去是、将来仍是基督徒，他们是在无意中、仅用手而非出于判断的同意而献了乳香。然后他们来到皇帝面前，退还他的金子，勇敢地请他收回自己的赏赐，并恳求处死他们，宣称无论为了基督，他们因手所犯的罪而在身体其他部分遭受何种折磨，他们绝不会背弃自己的信仰。\par

无论皇帝对他们如何不悦，他都没有杀死他们，以免他们获得殉道者的荣耀；因此他只剥夺了他们的军职，将他们逐出宫廷。\par

% structure: p0107 source=h2 target=subsection reason=relative-heading-rank; base=h2

\subsection{第十八章. 他禁止基督徒进入市场和审判席，并剥夺其参与希腊教育的权利。巴西略\Person{巴西略}、神学家额我略\Person{额我略}和阿波利纳里\Person{阿波利纳里}对这项法令的抵制。他们迅速将\WorkTitle{圣经}译成希腊表达方式。阿波利纳里\Person{阿波利纳里}和纳齐盎的额我略\Person{额我略}在这件事上比巴西略\Person{巴西略}做得更多：一人以修辞风格，另一人以史诗风格并模仿每一位诗人。}

尤利安\Person{尤利安}对所有基督徒都怀有如上所述的情感，每当有机会时他便表现出来。那些拒绝向诸神献祭的人，即使在其他方面无可指摘，也被剥夺了公民权利，以及参与集会和广场活动的特权；他不允许他们担任法官或行政长官，也不允许他们分享官职。\par

他禁止基督徒的子女上公立学校，也不准他们学习希腊诗人和作家的著作。他对叙利亚人\Person{阿波利纳里}——一位博学多识、精通语文的人——以及卡帕多西亚人\Person{巴西略}和\Person{额我略}（当时最著名的演说家），还有其他博学善辩之士怀有极大的敌意，其中一些人拥护尼西亚信理，另一些人则支持亚略的教义。他禁止基督徒父母的子女接受希腊学问的唯一动机，是他认为这些学问有助于获得辩论和说服的能力。因此，\Person{阿波利纳里}运用其渊博学识和才智，创作了一部关于希伯来古代史直到撒乌耳为王的英雄史诗，以替代\Person{荷马}的史诗。他将这部作品分为二十四部分，每部分冠以希腊字母的名称，依其数量和顺序。他还模仿\Person{米南德}创作喜剧，模仿\Person{欧里庇得斯}创作悲剧，模仿\Person{品达}创作颂歌。简言之，他从圣经中选取整个知识领域的主题，在极短时间内创作出一套作品，这些作品在风格、表达、特色和布局上均与希腊文学相当，且数量和力度均等。若非人们对古代作品怀有极度的偏爱，我相信\Person{阿波利纳里}的著作会受到如同古人著作般的推崇。\par

他智识的广博尤其令人钦佩；因为他在各门文学上都出类拔萃，而古代作家仅精通一门。他写了一部非常卓越的著作，题为\Quote{\WorkTitle{真理}}，反对皇帝和异教哲学家，在其中他清晰证明，无需诉诸圣经权威，他们远未获得对天主的正确认识。皇帝为了嘲弄这类作品，写信给主教们说：\Quote{我读了，我明白了，我谴责了。}他们回复如下：\Quote{你读了，但你没有明白；因为如果你明白了，你就不会谴责。}\par

有些人将这封信归于卡帕多西亚教会的主教\Person{巴西略}，或许不无道理；但无论是由他口述还是由他人所作，这封信充分彰显了作者的宽宏大量和学问。\par

% structure: p0112 source=h2 target=subsection reason=relative-heading-rank; base=h2

\subsection{第十九章。\Person{尤利安}题为\WorkTitle{厌恶胡须}的作品。\Place{安提约基雅}的\Place{达芙妮}，其详尽描述。移葬圣殉道者\Person{巴比拉}的遗骸。}

\Person{尤利安}决定对波斯发动战争后，前往叙利亚的\Place{安提约基雅}。民众大声抱怨，尽管物资极其丰富，但定价过高。于是，皇帝出于对民众的宽厚，据我所见，将物价压得过低，以致商贩逃离了该城。\newline{}\par

因此发生了短缺，民众责怪皇帝；他们的愤恨通过嘲笑他的长胡须和钱币上铸造的公牛得以发泄；他们讽刺地说，他就像他的司祭献祭时摔翻牺牲一样，把世界搞得天翻地覆。\par

起初他被激怒，威胁要惩罚他们，并准备前往\Place{塔尔索}。但后来他抑制了愤怒之情，仅以言语回敬了他们的嘲笑；他写了一部非常优雅的作品，题为\Quote{\WorkTitle{厌恶胡须}}，送给了他们。他对待该城的基督徒与其他地方完全相同，并尽力促进异教的扩展。\par

在此我将叙述一些关于殉道者\Person{巴比拉}坟墓的细节，以及大约在这个时期发生在\Place{达芙妮}的\Person{阿波罗}神庙中的某些事件。\par

\Place{达芙妮}是\Place{安提约基雅}的一个郊区，种满了柏树和其他树木，树下各季花卉繁茂。这些树的枝条如此浓密交错，可以说形成了一个屋顶，而不仅仅是遮荫，阳光永远无法穿透它们照射到地面。由于水源丰富而美丽、空气温和、惠风和畅，这里变得宜人而极其可爱。希腊人编造了一个神话：\Place{拉冬河}的女儿\Person{达芙妮}在从\Place{阿卡狄亚}逃离，以躲避\Person{阿波罗}的爱恋时，在这里变成了一棵以她名字命名的树。他们说，\Person{阿波罗}的爱并未因这个变化而减少；他用他爱人的叶子做了一个花冠，并拥抱了那棵树。后来他常常定居于此，认为这里比其他任何地方都更令他喜爱。\par

然而，性情庄重的人都认为前往这座城郊是可耻的；因为该地的位置和性质似乎会激起淫欲之感，而故事本身充满情色，便给腐化的青年提供了可乘之机，使他们情欲倍增。那些以此为借口的人大为亢奋，毫无约束地纵情放荡，既无法自制，也无法容忍自制之人的存在。任何在达佛涅没有情妇的人都被视为麻木不仁、缺乏教养，人们像躲避可憎可恶之物一样回避他。异教徒也因这里矗立着一尊极其美丽的达佛涅阿波罗雕像，以及一座宏伟昂贵的庙宇——据说是由安提奥库斯之父、为安提奥基亚城命名的塞琉古所建——而对这地方表现出极大的崇敬。那些相信这类传说的人认为，一条溪流从卡斯塔利亚泉流出，能赋予人预言未来的能力，其名称和功效与德尔斐的泉相似。据说，哈德良在还是平民时，就在这里得到了关于他未来伟大成就的暗示；他将一片月桂叶浸入水中，发现叶上写有他的命运。据说他成为皇帝后，命令关闭该泉，以免有人能窥探未来的知识。但我将这个问题留给那些比我更熟悉神话的人。\par

当尤利安的兄弟加卢斯被君士坦提乌斯立为凯撒，并定居安提奥基亚后，他对基督宗教的热忱和对殉道者遗念的敬仰，促使他决心清除该地的异教迷信和放荡之徒的暴行。他认为达成此目的最便捷的方法是，在神庙中建造一座祈祷之所，并将殉道者巴彼拉斯的坟墓迁至那里。巴彼拉斯曾以卓越声誉主持安提奥基亚教会，并殉道身亡。据说自迁葬之日起，邪魔便不再发出神谕。起初，人们将这种沉默归因于对邪魔祭仪的忽视和以往崇拜的中断；但结果证明，这完全是因圣殉道者的临在所致。即使在尤利安成为罗马帝国的唯一统治者后，尽管向邪魔大量奉献了奠酒、香和牺牲，沉默仍未打破；因为最终神谕本身开口，说明了先前沉默的原因。皇帝亲自进入神庙以咨询神谕，并献上礼品和牺牲，恳求给予答复。邪魔并未公开承认障碍是由殉道者巴彼拉斯的坟墓所致，而是说该地充满了尸体，使神谕无法发声。\par

尽管达佛涅有许多埋葬，但皇帝察觉到，使神谕沉默的仅仅是殉道者巴彼拉斯的临在，于是他命令将坟墓迁走。因此，基督徒们聚集起来，将棺椁运往约四十斯塔迪昂之远的城中，安放在至今仍保存之处，该地也以殉道者的名字命名。据说，男女老少、青年少女、老人孩童一起拉动棺匣，沿途唱着圣咏互相鼓励，表面上是为减轻劳作，实际上是因为他们被热忱和为亲属的宗教信仰所激励，而这正是皇帝所反对的。最好的歌手先唱，群众齐声回应，他们唱的副歌是：\Quote{愿一切拜偶像者、自夸于神像的人蒙羞。}\par

% structure: p0121 source=h2 target=subsection reason=relative-heading-rank; base=h2

\subsection{第二十章  因迁葬之事，许多基督徒受虐待。宣信者德奥多若。达佛涅的阿波罗神庙被天火烧毁。}

上述之事激起了皇帝的愤怒，仿佛是对他的侮辱，他决定惩罚基督徒；但总督萨卢斯特虽为异教徒，却试图劝阻他。然而皇帝不肯平息，萨卢斯特被迫执行其命令，逮捕并监禁了许多基督徒。他最先逮捕的一个名叫德奥多若的年轻人，当即被拉上刑架；尽管他的肉被铁钉撕裂，但他并未向萨卢斯特哀求，也未祈求减轻痛苦；相反，他仿佛对自己的痛苦无动于衷，如同只旁观他人受苦一般，勇敢地承受着创伤；他唱起前一天曾参与咏唱的同一首圣咏，以表明他并不后悔因之受审的行为。总督对这年轻人的坚毅深感钦佩，去见皇帝并告诉他，除非他尽快停止所采取的措施，否则他和他的同党将受嘲笑，而基督徒将获得更大荣耀。这一陈述产生了效果，被捕的基督徒获得了释放。据说后来有人问德奥多若，在刑架上是否感到疼痛；他回答说，他并非完全没有痛苦，但有一位青年站在他身边，用最细的麻布擦去他的汗水，并为他提供最凉的水，以此减轻了他的炎症，使他得到舒缓。我相信，没有天主的特别助佑，任何人无论拥有何等豪迈气概，都无法对身体表现出如此完全的漠视。\par

殉道者\Person{巴比拉}的遗体，基于前述理由，被迁至\Place{达佛涅}，随后又被移往他处。遗体被移走不久，火突然降在\Place{达佛涅}的阿波罗神庙，屋顶和神像本身被焚毁，只有裸露的墙壁和支撑柱廊及建筑后部的柱子幸免于火灾。基督徒相信殉道者的祈祷从天上引火烧了那魔鬼；但异教徒却报告说基督徒放火烧了那地方。这一怀疑日益加深；阿波罗的祭司被带到法庭面前，要他交出胆敢纵火者的名字；但他虽被捆绑并遭受最残酷的折磨，却没有指认任何人。\par

因此基督徒比以前更确信，那火不是出于人的行为，而是天主的愤怒从天降下烧了那庙宇。这些就是当时发生的事。皇帝，据我猜想，听说\Place{达佛涅}的灾祸是由殉道者\Person{巴比拉}所致，又得知殉道者受尊崇的遗骸被保存在靠近\Place{米利都}城附近阿波罗·狄迪穆斯神庙的几座祈祷所中，便写信给\Place{卡里亚}总督，命令他用火摧毁所有有屋顶和祭台的此类建筑，并将那些尚未完备的祈祷所连根拆毁。\par

% structure: p0125 source=h2 target=subsection reason=relative-heading-rank; base=h2

\subsection{第二十一章 论\Place{公元尼斯}的基督雕像，\Person{尤利安}将其推倒并使之无价值；他竖立自己的雕像；此像被雷电击倒并毁灭。\Place{厄玛乌}泉，基督曾在此洗脚。关于\Place{埃及}的\Person{佩尔西斯}树，它敬拜基督，以及通过它施行的奇事。}

在\Person{尤利安}统治期间发生的如此多显著事件中，我不可忽略提及一件，它显明了基督的权能，并证明天主的愤怒临到皇帝身上。\par

他听说在\Place{腓尼基}的\Place{凯撒勒雅腓立比}（又称\Place{公元尼斯}）有一尊著名的基督雕像，是由一位主曾治好血漏病的妇人所立，便命令将其推倒，并在原地竖立他自己的雕像；但一阵猛烈的天火降下，打断了雕像胸前的部分；头和颈被抛在地上，脸朝下固定在胸像断裂处，从那天直到如今，它仍那样立着，布满闪电的铁锈。基督的雕像被异教徒拖过城中并加以损毁；但基督徒收回了碎片，将雕像保存在至今仍保存它的教堂中。\Person{欧瑟比}记述，在此雕像基部生长了一种草药，为医生和江湖医士所不知，却能有效治愈一切疾病。在我看来，既天主曾屈尊与人同住，他再降恩赐予人，便不足为奇。\par

显然，无数其他奇迹发生在不同城邑和村庄；这些事迹被当地居民准确保存，因为他们从先祖传说中得知；此事如何真实，我立刻说明。在\Place{巴勒斯坦}有一座如今称为\Place{尼科波利}的城，从前只是一个村庄，被称为\Place{厄玛乌}，神圣的福音书曾提及。\BibleRef{路 24:13} 罗马人在征服耶路撒冷、战胜犹太人后，将此城命名为尼科波利。就在城外三岔路口之处，是基督复活后向\Person{克罗帕}及其同伴告别的地方，仿佛他要往另一个村庄去；此处有一治愈泉，患各种疾病的人和其他活物在其中洗去痛苦；因为据说当基督与祂的门徒一同从旅途中来到此泉时，他们在其中洗了脚，从此那水便成了医治疾病的良药。\par

在\Place{底比斯}的\Place{赫尔莫波利斯}，有一棵称为\Person{佩尔西斯}的树，其枝、叶乃至最微小的树皮，据称当病人触摸时能治愈疾病；因为埃及人相传，当\Person{若瑟}带着基督与\Person{玛利亚}——天主之母——躲避\Person{黑落德}的愤怒时，他们来到\Place{赫尔莫波利斯}；当进入城门时，这棵最大的树，仿佛不能忍受基督的到来，便向地弯腰敬拜祂。我准确叙述我从许多来源听闻关于此树的事。我认为这一现象是天主临在此城的迹象；或者，更可能的是，这棵树曾被居民按异教习俗敬拜，它摇动是因为那曾被敬拜的魔鬼，一见到那位为净化此类势力而显现者，便惊跳起来。它自动移动了；因为在基督面前，埃及的偶像都摇动，正如\Person{依撒意亚}先知所预言。魔鬼被驱逐后，此树被留下作为所发生之事的纪念碑，并被赋予了医治信者的特性。\par

\Place{埃及}和\Place{巴勒斯坦}的居民证实这些发生在他们中间的事件之真实性。\par

% structure: p0131 source=h2 target=subsection reason=relative-heading-rank; base=h2

\subsection{第二十二章 因厌恶基督徒，\Person{尤利安}允许犹太人重建耶路撒冷圣殿；每当他们动手做工时，火焰向上喷出，烧死许多人。关于在此工中努力之人衣服上出现的十字架记号。}

皇帝虽憎恨并迫害基督徒，却对犹太人表现出仁慈与仁爱。他致函犹太人的族长、领袖及民众，请求他们为他及帝国的繁荣祈祷。我深信，他采取这一举措并非出于对其宗教的尊重；因为他知道，犹太教可谓是基督宗教之母，并且他知道两者都基于族长与先知的权威；但他以为，通过偏袒犹太人——基督徒最顽固的敌人——可以令基督徒伤心。然而，他或许也算计着说服犹太人改信异教与祭献；因为他们只懂得圣经的字面意思，不能像基督徒及希伯来人中最明智的少数人那样洞察隐藏的含义。\par

事实证明这才是他的真正动机；因为他召来一些犹太领袖，劝勉他们回归遵守梅瑟法律及祖先习俗。他们回答说，由于耶路撒冷的圣殿已被推翻，在他们被驱逐出京城的其他地方举行这种仪式既不合法，也不符合祖先传统。于是皇帝拨给公款，命令他们重建圣殿，并按照古代方式祭献，举行类似祖先的崇拜。犹太人开始这项工程，却没有想到，根据神圣先知的预言，这不可能完成。他们寻找最熟练的工匠，收集材料，清理地基，如此热心地投入工作，以致妇女们也搬运土堆，并献出项链及其他女性饰物以支付费用。皇帝、其他异教徒及所有犹太人都将此视为最重要的事业，轻视其他一切。虽然异教徒对犹太人并无好感，但他们仍协助这项工程，因为他们指望其最终成功，并希望借此推翻基督的预言。除此动机外，犹太人自己也因认为重建圣殿的时候已经来到而受到推动。当他们清除了旧建筑的废墟，挖开地基并清理其基础后，据说第二天正要将第一块基石安放时，发生了大地震。由于地面的剧烈震动，深处有石头被抛起，砸伤了正在工作的犹太人以及旁观者。靠近圣殿遗址的房屋与公共柱廊——他们曾在那里娱乐——突然倒塌；许多人被压住，有的当场死亡，有的被找到时半死不活，手脚残缺，其他部分受伤。当天主使地震停止后，幸存的工人再次返回工作，一方面是因为皇帝的谕令，另一方面也是因为他们自己对这项工程感兴趣。人们常常在试图满足自己的激情时，寻求有害之物，拒绝真正有利之物，并被一种错觉所迷惑，认为唯有令自己愉悦之物才是真正有用的。一旦被这种错误引入歧途，他们就不再能以有利于自身利益的方式行事，也无法从降临于他们的灾祸中汲取教训。\par

我相信，犹太人正处在这样的状态中；因为他们并未将这次突如其来的地震视为天主反对重建圣殿的明确征兆，反而继续开始工作。但各方都叙述说，他们刚一回去工作，圣殿地基中突然迸出火焰，烧死了几名工人。\par

这一事实被无所畏惧地陈述，并被所有人相信；叙述中唯一的差异是，有些人认为火焰是从圣殿内部迸出，当时工人们正努力强行进入；另一些人则说火直接来自地下。无论现象以何种方式发生，同样奇妙。随后发生了一个更显著、更非凡的奇事：从事这项工程的人的衣服上突然自发地出现了十字记号。这些十字架如星宿般排列，看似人工制作。许多人因此承认基督是天主，并且重建圣殿并不蒙祂喜悦；另一些人则来到教会，接受了入门仪式，并恳求基督以圣歌和祈祷赦免他们的过犯。若有人不愿相信我的叙述，让他去请教那些从目击者那里听到我所叙述的事实的人吧，他们仍然活着。也让他去询问那些放弃未完成之工，或更准确地说，能够开始它的人。\par

\SourceRecord{Translated by Chester D. Hartranft.；Nicene and Post-Nicene Fathers, Second Series；Vol. 2.；Edited by Philip Schaff and Henry Wace.；Buffalo, NY: Christian Literature Publishing Co.,；1890.；<http://www.newadvent.org/fathers/26025.htm>.}

% structure: p0001 source=h1 target=section reason=page-title

\section{教会史（第六卷）}

% structure: p0002 source=h2 target=subsection reason=relative-heading-rank; base=h2

\subsection{第一章 \Person{尤利安}进军波斯；他战败后悲惨地结束生命。\Person{利巴尼奥斯}写了一封信描述他的死亡。}

我在前书中叙述了\Person{尤利安}在位期间教会发生的事件。这位皇帝决定对波斯发动战争，于是在初春迅速渡过\Place{幼发拉底河}，由于对\Place{埃德萨}居民的仇恨（他们长期信奉基督教），他绕过了该城，前往\Place{卡雷}，那里有一座朱庇特神庙，他在此献祭并祈祷。然后他从军队中挑选了两万名武装士兵，派往\Place{底格里斯河}方向，以便守卫那些地区，并在需要时随时支援他。随后他写信给罗马的盟友之一——\Place{亚美尼亚}王\Person{阿尔萨西乌斯}，请求他在战争中提供援助。在这封信中，\Person{尤利安}表现出极度的傲慢；他吹嘘自己拥有高尚的品质，声称这些品质使他配得上皇位，并为他所关心的众神所悦纳；他辱骂他的前任\Person{君士坦提乌斯}是一位柔弱不敬神的皇帝，并以极其侮辱的方式威胁\Person{阿尔萨西乌斯}；由于他得知\Person{阿尔萨西乌斯}是基督徒，他便加重了侮辱，或是急切而大量地说出反对基督的亵渎之言，因为他惯于在任何情况下都这样大胆。他告诉\Person{阿尔萨西乌斯}，除非他按照他的指示行事，否则他所信靠的天主将无法保护他免受报复。当他考虑所有安排都已妥当后，他便率领军队穿过\Place{亚述}。\par

他攻占了大量城镇和堡垒，有些是通过背叛，有些是通过战斗，并且轻率地继续前进，没有考虑他必须沿原路返回。他洗劫了所到之处，拆除或烧毁了粮仓和仓库。当他沿\Place{幼发拉底河}上行时，抵达了\Place{泰西封}，这是一座非常大的城市，波斯君主如今已将王宫从\Place{巴比伦}迁至此地。\Place{底格里斯河}靠近这里。由于陆地的一部分将城市与河流隔开，他无法乘船到达城市，因此他判断必须要么走水路，要么弃船走陆路前往\Place{泰西封}；他就此事审问俘虏。从他们那里得知有一条运河因年久失修而被堵塞，他下令将其疏通，从而打通了\Place{幼发拉底河}与\Place{底格里斯河}之间的通道，随后他朝城市进发，他的舰队沿着军队侧旁行进。但波斯人在\Place{底格里斯河}岸上出现了，展示出令人畏惧的骑兵部队和许多武装士兵、大象和马匹；\Person{尤利安}意识到他的军队被两条大河包围，无论是留在当前位置还是撤退经过被他彻底摧毁以致无法获得补给的城镇和村庄，都有全军覆没的危险；因此他召集士兵观看赛马，并为最快的赛马提供奖励。与此同时，他命令舰船军官抛弃军队的补给和行李，以便士兵们看到自己因缺乏必需品而处于危险之中，从而能够大胆转向，更加拼命地对抗敌人。晚饭后，他召见将军和护民官，命令军队上船。他们沿着\Place{底格里斯河}航行了一整夜，立刻到达对岸并登陆；但他们的撤离被一些波斯人察觉，他们互相鼓励进行抵抗，但那些仍在睡觉的人被罗马人轻易制服。\par

天亮时，两军交战；双方伤亡惨重后，罗马人沿河返回，在\Place{泰西封}附近扎营。皇帝不再希望继续前进，烧毁了船只，因为他认为需要太多士兵来守卫它们；然后他开始沿着左面的\Place{底格里斯河}撤退。充当罗马人向导的俘虏把他们带到一片肥沃的土地，那里有丰富的补给。不久，一位决心为波斯自由而牺牲的老人故意被俘，并被带到皇帝面前。在受到关于路线的审问时，他似乎说了实话，说服他们跟随他，因为他能够非常迅速地将军队带到罗马边境。他指出，在三四天的路程中，这条路将很难走，在此期间必须携带补给，因为周围地区贫瘠。皇帝被这位智慧老人的话所欺骗，批准了这条路线。继续前进三天后，他们进入了一片未开垦的地区。那位老俘虏被施以酷刑。他承认自己为了国家自愿赴死，因此准备承受任何可以施加的苦难。\par

罗马军队现在因长途跋涉和补给匮乏而疲惫不堪，波斯人选择此时攻击他们。\par

在随后爆发的激战中，刮起了猛烈的风；天空和太阳完全被乌云遮蔽，空气中同时混杂着尘土。在这样的黑暗中，一骑兵全速奔驰，将长矛刺向皇帝，使他受了致命伤。将\Person{尤利安}掀下马后，这不知名的袭击者悄悄离去。有人猜测他是波斯人；也有人认为他是撒拉逊人。还有那些人坚称，动手的是罗马士兵，他对皇帝使军队暴露于如此危险的不慎与鲁莽感到愤慨。叙利亚本地人、诡辩家\Person{利巴尼乌斯}，\Person{尤利安}最亲密的朋友，关于行凶者表达了如下言辞：\Quote{你们想知道皇帝被谁所杀。我不知道他的名字。但我们有一个证据：凶手并非敌人；因为没有人前来领赏，尽管\Place{波斯}王曾使传令官宣告，要对完成此举者授予荣誉。我们确实应该感谢敌人没有将此事的光荣据为己有，而是让我们在自己人中寻找凶手。}\par

\Quote{那些寻求他死亡的人，是那些惯常违反法律、曾密谋反对他的人，因此他们一有机会就实施了此事。他们被渴望获得比在他统治下更多的自由所驱动；而且，也许他们主要是被对皇帝依附于他们所厌恶的众神侍奉的愤慨所刺激。}\par

% structure: p0009 source=h2 target=subsection reason=relative-heading-rank; base=h2

\subsection{第二章 \Person{尤利安}死于天主的义怒。不同的人所见皇帝之死的异象。木匠之子的答复；\Person{尤利安}将他的血抛向\Person{基督}。\Person{尤利安}给罗马人带来的灾祸。}

在上文引用的文件中，\Person{利巴尼乌斯}清楚表明皇帝死于一名基督徒之手；这或许是事实。并非不可能，当时在罗马军中服役的一些士兵可能就有了这个念头，因为希腊人以及所有人直到今日都赞美那些为了自由事业舍身赴死的弑君者，并精神抖擞地支持他们的国家、家庭和朋友。更不应受责备的是，谁为了天主和宗教而行此勇敢之举。除此之外，除了我所叙述的，我准确知道关于行凶者的事情就没了。然而，所有人都同意接受流传给我们的记载，该记载证明他的死亡是天主义怒的结果。一个证明是，他的一位朋友所得的天主异象，我现在就来描述。据传，他曾旅行到\Place{波斯}，打算与皇帝会合。在路上，他发现自己远离任何住处，以致有一夜，他被迫睡在一座教堂里。那一夜，他看见，或是在梦中或是在异象中，所有宗徒和先知聚集在一起，抱怨皇帝对教会施加的伤害，并商议应采取的最佳措施。经过许多商议和困窘，有两个人从集会中站起来，让其他人安心，然后匆匆离开，似乎要去剥夺\Person{尤利安}的皇权。这个奇观的旁观者没有继续前行，而是在可怕的悬念中等待这启示的结局。他再次在同一处躺下睡觉，又看见了同样的集会；那前一天晚上似乎离去以实施反对\Person{尤利安}目的的两个人，突然返回，向其他人宣布了他的死亡。\par

同一天，一个异象赐给了住在\Place{亚历山大里亚}的教会哲学家\Person{狄迪莫斯}；他为\Person{尤利安}的谬误和他对教会的迫害深感忧伤，因此禁食并不断为此向天主献上恳求。由于焦虑和前一晚缺乏食物，他坐在椅子上睡着了。然后，仿佛在神魂超拔中，他看见白马穿越天空，并听见一个声音对骑者说：\Quote{去告诉\Person{狄迪莫斯}，\Person{尤利安}刚在这时刻被杀了；让他把这消息传给主教\Person{阿塔纳修斯}，并让他起来吃饭。}我可靠地得知，\Person{尤利安}的朋友和那位哲学家都看到了那些事。结果证明，他们两人离见证真相都不远。但如果这些例子不足以证明\Person{尤利安}的死是因他迫害教会而致天主义怒的结果，那么请想起其中一位神职人员的预言。当\Person{尤利安}准备开始对\Place{波斯}人的战争时，他威胁说战争结束后他会严厉对待基督徒，并夸口说木匠之子将无力帮助他们；上述神职人员随即回答说，木匠之子当时正在为他准备一副木棺材，以备他死亡之用。\par

\Person{尤利安}自己很清楚致命打击从何而来，以及为何被施加；因为，据说，当他受伤时，他取了一些从伤口流出的血，抛向空中，好像他看见了\Person{耶稣基督}显现，打算将血抛向他，以责备基督使他被杀。另有人说，他是在生太阳的气，因为太阳偏袒了\Place{波斯}人，没有救他，尽管根据天文学家的学说，太阳在他出生时主掌他的命运；而他正是为了表达对这颗发光体的愤怒，才把血握在手中，向上抛向空中。\par

我不知道在死亡临近时，是否如通常那样，当灵魂正与身体分离并得以看见比给人类所定更为神圣的景象时，尤利安或许就看见了基督。对此论题提及甚少，然而我不敢完全否定这一假设；因为天主常常允许更不可能、更惊人的事件发生，以证明以基督命名的宗教并非由人的力量支撑。但很明显，在整个这位皇帝统治期间，天主明确显示出祂的不满迹象，并允许许多灾祸降临于罗马帝国的几个行省。祂以如此可怕的地震惩罚大地，以至于建筑摇动，屋内与露天一样不安全。据我所闻，我猜测正是在这位皇帝统治期间，或者至少在他位居政府第二把交椅时，在埃及亚历山大里亚附近发生了一场大灾难：海水后退，然后又因回潮波浪超越其边界，淹没大片土地，以致水退时，海中小艇发现搁在屋顶上。这次洪水的周年日，他们称之为地震诞辰，在亚历山大里亚仍以年度节日纪念；全城普遍点亮灯火；他们向天主献上感恩祈祷，并非常辉煌而虔诚地庆祝这一天。此次统治期间还发生了过度干旱；植物枯死，空气腐败；由于缺乏适当食物，人们不得不求助于通常其他动物食用的东西。\par

饥荒引发了特殊疾病，许多人因此丧生。这就是尤利安治下帝国的状况。\par

% structure: p0015 source=h2 target=subsection reason=relative-heading-rank; base=h2

\subsection{第三章 约维安统治；他颁布许多法律并在其统治中实施。}

尤利安去世后，军队一致同意将帝国统治权交给约维安。当军队要宣布他为皇帝时，他宣布自己是基督徒并拒绝君权，也不接受帝国象征；但当士兵们发现他拒绝的原因时，他们大声宣称他们自己也是基督徒。\par

尤利安策略所遗留的危险动荡局势以及军队在敌国因饥荒所受的痛苦，迫使约维安与波斯人缔结和约，并割让以前曾是罗马藩属的一些领土。他从经验中得知，其前任的不虔敬激起了天主的愤怒，并引发了公共灾祸，因此他立即致函各行省总督，指示人民应无所恐惧地聚集在教会中，恭敬地事奉天主，并接受基督信仰为唯一真正的宗教。他恢复了教会、圣职人员、寡妇及贞女在君士坦丁及其诸子所赐、后来被尤利安削夺的、为宗教的益处和荣耀而设的一切豁免权和所有先前拨款。他命令当时任禁军统领的塞孔都斯，将娶任何圣洁贞女为妻、甚或以不洁欲望看待她们并劫掠她们的行为定为死罪。\par

他因尤利安统治期间所盛行的邪恶而颁布此法律；因为许多人从圣洁贞女中娶妻，或凭武力或诡计完全败坏了她们；由此产生了那种可耻情欲得逞而免罚的放纵，当宗教被滥用时，这种情况总是发生。\par

% structure: p0019 source=h2 target=subsection reason=relative-heading-rank; base=h2

\subsection{第四章 教会中再起纷争；安提约基雅会议，尼西亚信德得到确认；这次重要会议关于此点写给约维安的内容。}

教会的主持者们现在恢复了教义问题和讨论的鼓动。在尤利安统治期间，当基督教本身面临危险时，他们保持安静，并一致祈求天主的怜悯。当人受到外敌攻击时，他们彼此和睦；但当外在麻烦消除时，内在纷争便潜入；然而，这不是引用历史上许多政府和国家中此类例证的适当场合。\par

此时，\Person{巴西略}（\Place{安基拉}主教）、\Person{西尔瓦诺}（\Place{塔尔索}主教）、\Person{索弗罗尼}（\Place{庞培城}主教）及其同党——他们极度厌恶所谓的异常派异端，并采用\Quote{相类本质}一词而不用\Quote{同质}一词——向皇帝写了一份文书；他们在表达对天主使其即位的感激之后，恳请他确认在\Place{阿里米农}和\Place{塞琉西亚}颁布的法令，并废除仅凭某些人的热忱和权力所确立的内容。\par

他们还恳求，如果因会议而存在的分裂仍在教会中盛行，各地主教应独自在皇帝指定的某处聚集，而不被允许在其他地方集会并发布相互矛盾的法令，如同\Person{君士坦提乌斯}统治时期所做的那样。他们补充说，他们没有去营地拜访他，是因为怕给他添麻烦；但如果他想见他们，他们愿意前往，并自行承担旅途的一切费用。这就是写给\Person{约维安}皇帝的文书。\par

此时，在\Place{叙利亚}的\Place{安提约基雅}召开了一次会议；\Place{尼西亚}会议所确立的信仰形式得到确认，并决定圣子与圣父是无可争议的同一实体。当时主理\Place{安提约基雅}教会的\Person{梅利提乌}、\Place{萨莫萨塔}主教\Person{欧瑟伯}、\Place{叙利亚劳狄刻雅}主教\Person{佩拉纠}、\Place{巴勒斯坦凯撒勒雅}主教\Person{阿卡西乌}、\Place{加沙}主教\Person{依勒内}以及\Place{安基拉}主教\Person{亚大纳削}参加了这次会议。\par

会议结束后，他们通过发送以下信函，将所发生之事禀报皇帝：\par

致最虔诚且天主所爱的奥古斯都、我们的君主\Person{约维安}凯旋者，由从各地聚集于\Place{安提约基雅}的主教们。\par

皇帝，天主所爱者，我们深知陛下的虔敬完全专注于维护教会的平安与和谐；我们也不忽视您已很好地接受了这一合一的关键，即真实而正统的信德。\par

\Quote{因此，免得我们被算作那些攻击这些真理教义的人，我们向陛下证实，我们接受并维护古代由圣\Place{尼西亚}会议所制定的信仰形式。尽管\InnerQuote{\OriginalTerm{同质}{consubstantial}}一词对某些人显得陌生，但教父们已安全地解释了它，意指圣子是由圣父的实体所生。这一词汇并不表示不间断的受生；也不与希腊人使用\InnerQuote{\OriginalTerm{实体}{substance}}一词的含义一致，而是旨在抵制\Person{亚略}不敬而鲁莽的主张，即圣子是从先前不存在之物产生的。刚刚兴起的\OriginalTerm{阿诺米派}{Anomians}无耻地坚持这一词汇，破坏了教会的和谐。我们随函附上在\Place{尼西亚}聚会的教父们所采纳的信经文本，我们也珍视它。}\par

这就是在\Place{安提约基雅}聚会的司铎们所作的决定；他们在信函后附上了\Place{尼西亚}信经的文本。\par

% structure: p0029 source=h2 target=subsection reason=relative-heading-rank; base=h2

\subsection{第五章　伟大的\Person{亚大纳削}深受皇帝敬重，并治理埃及的教会。\Person{安当}的异象。}

此时，治理亚历山大宗座的\Person{亚大纳削}和他的几位朋友认为，既然皇帝是基督徒，有必要前往他的宫廷。于是\Person{亚大纳削}去了\Place{安提约基雅}，并将他认为合适的事禀告皇帝。然而，另有人说，皇帝召见他，是为了咨询他关于宗教事宜和正确教义的问题。当教会事务尽可能处理完毕后，\Person{亚大纳削}开始考虑返回。\par

在\Place{安提约基雅}信奉亚略异端的主教\Person{欧佐约}，企图将怀有同样思想的阉人\Person{普罗巴提}安置在\Place{亚历山大}。\Person{欧佐约}的全党与他合谋实现这一计划；而\Place{亚历山大}的居民\Person{路济亚}（曾被\Person{乔治}祝圣为长老）试图在皇帝面前中伤\Person{亚大纳削}，声称他被指控犯有各项罪行，并被先前的皇帝判处永久流放，是教会关于天主的纷争与动荡的始作俑者。\Person{路济亚}同样恳求\Person{约维安}为\Place{亚历山大}教会另立一位主教。皇帝因为知道针对\Person{亚大纳削}的阴谋，所以对诽谤不予置信，并以威胁的口吻命令\Person{路济亚}安静退下；他也命令\Person{普罗巴提}和属于他宫廷的其他阉人（他视之为这些麻烦的始作俑者）行事更谨慎明理。从那时起，\Person{约维安}对\Person{亚大纳削}表现出了最大的友谊，并送他回埃及，指示他可以按他认为适当的方式管理该国的教会和人民。据说他还赞扬了主教的德性、他的生活、他的才智和口才。\par

这样，如同前几卷所叙述的，在长期遭受反对之后，\Place{尼西亚}信德在本届政府下完全得以重建；但不久之后，它面临着进一步的困扰。因为后来看来，隐修士\Person{安当}的整个预言并未在\Person{君士坦提乌斯}统治期间教会所遭遇的事中完全应验；部分预言直到\Person{瓦伦斯}统治时期才应验。据说，在亚略派于\Person{君士坦提乌斯}统治期间控制教会之前，\Person{安当}在梦中看到骡子用蹄子踢祭台并撞翻圣桌。醒来后，他立即预言教会将因引入虚假混杂的教义和异端者的反叛而受困扰。这一预言的真实性被当前所论时期前后发生的事件所证实。\par

% structure: p0033 source=h2 target=subsection reason=relative-heading-rank; base=h2

\subsection{第六章　\Person{约维安}之死；\Person{瓦伦提尼安}的生平及其对天主的信赖；他如何被推举登基并选立他的兄弟\Person{瓦伦斯}与他共治；两人的分歧。}

\Person{约维安}统治大约八个月后，在前往\Place{君士坦丁堡}的路上，突然死于\Place{比提尼亚}的\Place{达达斯塔纳}镇。有人说他的死因是晚餐吃得过饱；另有人归因于他睡觉的房间潮湿；因为房间最近用生石灰粉刷过，而且冬季在里面烧了大量煤炭以作防护；墙壁变得潮湿，非常湿润。\par

军队到达比提尼亚的\Place{尼西亚}时，拥立\Person{瓦伦提尼安}为皇帝。他是个好人，能掌控帝国。他刚从流放中回来不久；据说\Person{尤利安}一登基，就把\Person{瓦伦提尼安}从所谓的约维安军团除名，并以他未尽职责率领部下士兵攻打敌人为借口，判他永久流放。但他被判刑的真正原因如下：\Person{尤利安}在高卢时，有一天去神庙献香。\Person{瓦伦提尼安}陪着他，根据一条古罗马法律（至今仍有效），该法律规定约维安军团和赫丘利军团（即以朱庇特和赫丘利命名的军团）的首领应始终作为皇帝的贴身侍卫。当他们正要进入神庙时，祭司按照异教习俗，用树枝向他们洒水。一滴水落在\Person{瓦伦提尼安}的袍子上；他几乎不能自制，因为他是基督徒，他斥责了洒水者；甚至据说他在皇帝面前，剪下了沾水的那块衣料，扔掉了。从那一刻起，\Person{尤利安}就对他怀有敌意，不久就借口他在军事上的不当行为，将他流放到亚美尼亚的\Place{梅利蒂尼}；因为他不愿宗教被视为判决的原因，免得\Person{瓦伦提尼安}被视为殉道者或宣信者。\Person{尤利安}对其他基督徒也是同样对待，正如我们已说过的；因为如前所述，他意识到让他们冒险只会增加他们的声誉，并有助于巩固他们的宗教。\Person{尤利安}一登基，\Person{瓦伦提尼安}就被从流放中召回\Place{尼西亚}；但皇帝在此期间去世，\Person{瓦伦提尼安}经军队和政府首要人物一致同意，被指定为继承人。当他被授予帝国权力的象征时，士兵们喊道必须选一个人分担政府的负担。对于这个提议，\Person{瓦伦提尼安}做了如下回答：\Quote{士兵们，选我为皇帝只取决于你们；但现在你们选了我，执行你们的要求就不取决于你们，而取决于我了。像臣民应该做的那样保持安静，让我作为皇帝来处理公共事务。}\par

拒绝士兵的要求后不久，他前往\Place{君士坦丁堡}，立他的弟弟为皇帝。他将东方作为帝国的一部分分给他，自己保留了从\Place{伊利里亚}到利比亚最远海岸的西大洋地区。两兄弟都是基督徒，但意见和性情不同。因为\Person{瓦伦斯}受洗时，由\Person{欧多克西乌斯}作为他的施洗者，并热烈拥护\Person{阿里乌}的教义，并会以武力迫使所有人屈服。但\Person{瓦伦提尼安}坚持尼西亚大公会议的信德，支持持有同样观点的人，而不骚扰持有其他意见的人。\par

% structure: p0037 source=h2 target=subsection reason=relative-heading-rank; base=h2

\subsection{第7章 教会再次出现纷争，\Place{兰普萨库斯}会议召开。支持\Person{欧多克西乌斯}的阿里乌派得势，将正统派逐出教会。被逐者中有\Place{安提约基雅}的\Person{梅利提乌斯}。}

当\Person{瓦伦提尼安}从\Place{君士坦丁堡}前往\Place{罗马}时，他必须经过\Place{色雷斯}；\Place{赫勒斯滂}和\Place{比提尼亚}的主教，以及其他坚持圣子与圣父同质的主教，派遣\Place{佩林托斯}的\Place{赫拉克利亚}主教\Person{希帕提安}去见他，请求允许他们聚集在一起审议教义问题。\par

\Person{希帕提安}转达了委托给他的消息后，\Person{瓦伦提尼安}做了如下回答：\Quote{我只是平信徒之一，因此无权干预这些事务；让司铎们，这些事属于他们，在他们愿意的地方聚集。}通过他们的代表\Person{希帕提安}得到这个答复后，主教们在\Place{兰普萨库斯}集会。\par

他们商议了两个月后，废除了在\Place{君士坦丁堡}通过\Person{欧多克西乌斯}和\Person{阿卡西乌斯}同党的阴谋所颁布的一切。他们还宣布，那份以虚假声明流传的、据称是西方主教编撰的信德公式无效，该公式曾以承诺谴责本质不相似的教义而获得许多主教的签名——但这一承诺从未兑现。\par

他们规定圣子在本质上与圣父相似的教义应占主导地位；因为他们说必须使用\Quote{相似}一词来表明天主位格中的特性。他们同意所有教会应维持曾在\Place{塞琉西亚}宣认并在\Place{安提约基雅}圣堂奉献时所颁布的信德形式。\par

他们指示，所有被那些主张圣子与圣父不相似的人所罢黜的主教，应立即复职，因为他们被不公正地逐出教会。他们宣布，如果有人想控告他们，可以这样做，但若控告证明为虚假，则控诉者将承担与所控罪行相同的惩罚。本省和邻省的正统主教应担任法官，并在教堂内与提供证词的证人一起集会。\par

做出这些决定后，主教们传唤\Person{欧多克西乌斯}的同党，劝他们悔改；但他们不听劝诫，于是会议颁布的法令被送往所有教会。主教们判断\Person{欧多克西乌斯}可能会试图说服皇帝支持他，并诽谤他们，于是决定抢先一步，将他们与\Place{兰普萨库斯}的会议纪要送往朝廷。\par

他们的代表在皇帝\Person{瓦伦斯}从\Place{赫拉克利亚}返回\Place{色雷斯}的途中遇见了他，当时他正与他的兄弟同行，他的兄弟已经前往\Place{旧罗马}。\par

然而，\Person{欧多克修斯}事先已经争取了皇帝及其朝臣支持他的观点；所以，当\Place{兰普萨科斯}会议的代表谒见\Person{瓦伦斯}时，他只是劝告他们不要与\Person{欧多克修斯}不和。代表们回答时提醒他\Person{欧多克修斯}在\Place{君士坦丁堡}所使用的诡计，以及他力图废除\Place{塞琉西亚}会议决议的阴谋；这些陈述激起了\Person{瓦伦斯}的极大愤怒，以至于他将代表们判为流放，并将教会交与\Person{欧多克修斯}的党羽。随后他越过\Place{叙利亚}，因为他担心\Place{波斯人}会破坏他们与\Person{约维安}缔结的三十年休战。然而，当他发现\Place{波斯人}并无反叛之意时，便在\Place{安提约基雅}定居下来。他将主教\Person{梅利提乌斯}流放，但却饶过了\Person{保禄}，因为他钦佩他生活的圣洁。那些不与\Person{犹佐修斯}共融的人，不是被逐出教会，就是受到虐待或其他形式的骚扰。\par

% structure: p0046 source=h2 target=subsection reason=relative-heading-rank; base=h2

\subsection{第八章 普罗科比乌斯的叛乱与异常死亡。\Place{基齐库斯}主教\Person{埃莱乌修斯}与异端者\Person{欧诺米乌斯}。\Person{欧诺米乌斯}继任\Person{埃莱乌修斯}。}

在这个关键时刻，很可能发生一场严厉的迫害，如果不是\Person{普罗科比乌斯}发动了一场内战。当他在\Place{君士坦丁堡}开始以暴君自居时，他迅速召集了一支庞大的军队，并向\Person{瓦伦斯}进军。\par

后者离开\Place{叙利亚}，在\Place{弗里吉亚}的城市\Place{纳科利亚}附近与\Person{普罗科比乌斯}相遇，并因他的两位将军\Person{阿格隆}和\Person{戈马留斯}的背叛而活捉了他。\par

\Person{瓦伦斯}将他和他的背叛者处以残酷的死刑；虽然据说他曾发誓要优待这两位将军，但他却下令将他们锯成两半。\par

他命令将\Person{普罗科比乌斯}双腿绑在两根被弯到地上的树上，然后让这些树弹起；当树木恢复其自然位置时，受害者被撕成两半。\par

这场战争结束后，\Person{瓦伦斯}退隐到\Place{尼西亚}，发现自己处于极度平静之中，他便又开始骚扰那些在关于天主性体的观点上与他不同的人。\par

他对\Place{兰普萨科斯}会议的主教们愤怒无比，因为他们谴责了亚略派主教们以及在\Place{阿里米努姆}制定的信仰宣言。\par

在受这些怨恨情绪影响下，他从\Place{叙利亚}召来\Person{埃莱乌修斯}，并召集了一个与他观点一致的主教会议，企图强迫他同意他们的教义。\Person{埃莱乌修斯}起初勇敢地拒绝服从。但后来，由于惧怕皇帝威胁的流放和剥夺财产，他屈服于命令。他很快为自己的软弱感到后悔，回到\Place{基齐库斯}后，他在教会中公开承认了自己的过错，并敦促民众另选一位主教，因为他说自己在背叛了自己的教义之后，无法履行司铎的职责。市民们尊敬他的行为，特别善待他，所以他们不愿另选主教。然而，在\Place{君士坦丁堡}的亚略派领袖\Person{欧多克修斯}却祝圣\Person{欧诺米乌斯}为\Place{基齐库斯}主教；因为他期望凭借\Person{欧诺米乌斯}卓越的雄辩口才，能轻易地将\Place{基齐库斯}的民众吸引到自己的观点上来。\Person{欧诺米乌斯}抵达该城后，逐出了\Person{埃莱乌修斯}，因为他有一道帝国敕令支持此举，并自己占据了教会。\par

\Person{埃莱乌修斯}的追随者在城墙外建造了一座祈祷堂，他们在那里举行集会。我很快将有机会再次谈及\Person{欧诺米乌斯}以及以他命名的异端。\par

% structure: p0055 source=h2 target=subsection reason=relative-heading-rank; base=h2

\subsection{第九章 维护尼西亚信德者的苦难。\Person{阿格利乌}，诺瓦提安派的领袖。}

在\Place{君士坦丁堡}城，代表尼西亚教义的基督徒和诺瓦提安观点的追随者受到同样严厉的对待。\par

他们最终全都被逐出城；诺瓦提安派的教会根据皇帝的命令被关闭。另一方则没有可关闭的教会，因为在\Person{君士坦提乌斯}统治期间他们已经失去了所有教会。\par

在这个时期，从\Person{君士坦提乌斯}时代起就管理\Place{君士坦丁堡}诺瓦提安派教会的\Person{阿格利乌}被判处流放。据说他尤其以按照教会法规生活而著称。就他的生活方式而言，他已经达到了哲学的最高境界，即摆脱世俗的财产；这从他的日常行为可以证明：他只有一件外衣，并且总是赤脚行走。他被流放后不久就被召回，接收了他属下的教会，并通过\Person{马尔西安}的影响大胆地召集教会；\Person{马尔西安}是一个德行和口才非凡的人，他曾经是宫廷卫队的一员，但此时是诺瓦提安异端的司铎，也是皇帝女儿\Person{阿纳斯塔西娅}和\Person{卡罗莎}的语法教师。\Place{君士坦丁堡}至今仍有以这两位公主命名的浴场。仅仅是为了\Person{马尔西安}的缘故，上述特权才被赐予诺瓦提安派。\par

% structure: p0059 source=h2 target=subsection reason=relative-heading-rank; base=h2

\subsection{第十章 关于\Person{小瓦伦提尼安}和\Person{格拉提安}。\Person{瓦伦斯}统治下的迫害。同质派受亚略派和马其顿派压迫，派遣使团前往\Place{罗马}。}

大约在这个时期，\Person{瓦伦提尼安}在西方得了一个儿子，皇帝给他取了自己的名字。不久之后，他宣布他的儿子\Person{格拉提安}为皇帝；这位王子是在他父亲掌权之前出生的。\par

同时，尽管异常巨大的冰雹降落在各处，许多城市，特别是\Place{俾提尼亚}的\Place{尼西亚}，被地震摇撼，然而皇帝\Person{瓦伦斯}和主教\Person{欧多克修斯}并未停下脚步，继续迫害所有与他们意见不同的基督徒。他们在对付那些坚持尼西亚教义的基督徒的阴谋中，达到了预期的最大成功；因为在瓦伦斯统治的大部分时间里，特别是在\Place{色雷斯}、\Place{俾提尼亚}和\Place{赫勒斯滂}，甚至更远的地方，这些基督徒既没有教堂也没有司铎。然后\Person{瓦伦斯}和\Person{欧多克修斯}将他们的怨恨转向马其顿派，那些人在该地区比上述基督徒人数更多，并毫无节制地迫害他们。\par

马其顿派预感将遭受更多苦难，便派遣代表到各城市，最后同意求助于\Person{瓦伦提尼安}和罗马主教，而不愿分享\Person{欧多克修斯}和\Person{瓦伦斯}及其追随者的信仰；当此事似乎可行时，他们从自己中间选出三人——\Place{塞巴斯特}的主教\Person{尤斯塔修斯}、\Place{塔尔苏斯}的主教\Person{西尔瓦努斯}和\Place{卡斯塔巴利斯}的主教\Person{狄奥菲鲁斯}——派遣他们到皇帝\Person{瓦伦提尼安}那里；他们也托付他们一封写给罗马主教\Person{利贝里乌斯}及其他西方司铎的信，信中恳求他们，作为坚守宗徒们认可并确认的信仰的主教，应当比别人更关心宗教，以完全的认可接待他们的代表，并与他们商议在教会事务得以妥善安排之前应做之事。\par

代表们到达意大利后，发现皇帝在\Place{高卢}与蛮族作战。他们认为前往高卢的战地会有危险，便将信件交给了\Person{利贝里乌斯}。在与他就其使命的目的进行商议后，他们谴责了\Person{亚略}及其持有并教导其教义的人；他们弃绝了所有与\Place{尼西亚}所确立的信仰相反的异端；并接受了\Quote{同质}这一术语，认为这个词与\Quote{本质相似}表达相同的意思。当他们向\Person{利贝里乌斯}呈交了与上述类似的信仰宣认后，他接纳他们与自己共融，并写信给东方的主教们，赞扬他们信仰的正统，并详细讲述了他与他们会议中所发生的事。\Person{尤斯塔修斯}及其同伴所作的信仰宣认如下：\par

% structure: p0064 source=h2 target=subsection reason=relative-heading-rank; base=h2

\subsection{第11章 马其顿派代表\Person{尤斯塔修斯}、\Person{西尔瓦努斯}和\Person{狄奥菲鲁斯}致罗马主教\Person{利贝里乌斯}的信仰宣认}

致我们的主、兄弟和同工\Person{利贝里乌斯}——\Person{尤斯塔修斯}、\Person{西尔瓦努斯}和\Person{狄奥菲鲁斯}在主内问安。\par

由于异端者的疯狂意见，他们不停地为天主教会散布丑闻，我们这些废除他们每次攻击的人，宣认由正统主教们在\Place{兰普萨库斯}举行的主教会议、在\Place{士每拿}举行的会议以及在别处举行的其他会议。我们已经备好信件，并派遣使团前往尊下，同样也前往意大利和西方所有其他主教，以确认并保存天主教信仰，这信仰是由真福\Person{君士坦丁}和三百一十八位敬畏天主的教父在神圣的\Place{尼西亚}大公会议上所确立的。\par

\Quote{这信仰藉着一纯粹且不可动摇的规定，至今保持不变，并将永远保留；其中\InnerQuote{同质}一词以全部的神圣和虔诚被确立，以见证反对\Person{亚略}的乖谬。我们各人亲手宣认，我们与上述人士始终持守这同一信仰，现今仍持守，并将坚持到底。我们谴责\Person{亚略}、他的不敬教义及其门徒。我们也谴责\Person{帕特罗帕西安努斯}、\Person{佛提努}、\Person{玛尔塞鲁}、\Person{萨莫萨塔的保禄}的异端，以及所有自行持守此类教义者。我们绝罚所有反对神圣教父们在\Place{尼西亚}所确立的上述信仰的异端。我们特别绝罚\Person{亚略}，并谴责在\Place{阿里米努姆}所颁布的一切与神圣\Place{尼西亚}大公会议所确立的上述信仰相悖的法令。我们从前曾被某些人的诡诈和伪证所欺骗，当这些法令从\Place{色雷斯}的\Place{尼西亚}传至\Place{君士坦丁堡}时，我们签署了它们。}\par

在这次宣认之后，他们将尼西亚信经全文的副本附在自己的信经之后，并从\Person{利贝里乌斯}那里收到了他们所办理的一切的书面记录，然后乘船前往\Place{西西里岛}。\par

% structure: p0069 source=h2 target=subsection reason=relative-heading-rank; base=h2

\subsection{第12章 \Place{西西里}和\Place{提亚纳}的会议。原定在\Place{基利家}举行的主教会议被\Person{瓦伦斯}解散。当时的迫害。伟大的\Person{阿塔纳修斯}再次逃亡并隐藏；因\Person{瓦伦斯}的信他重新出现，并治理\Place{埃及}的教会。}

在\Place{西西里}召开了一次会议；在确认了与代表们的宣认中所阐述的相同教义之后，会议解散。\par

同一时期，在\Place{提雅纳}召开了一次大公会议。出席会议的有\Place{凯撒勒雅（卡帕多细雅）}的主教\Person{尤西比乌斯}、\Place{安基拉}的主教\Person{阿塔纳修斯}、\Place{劳狄刻雅}的主教\Person{佩拉吉乌斯}、\Place{提洛}的主教\Person{则农}、\Place{厄梅撒}的主教\Person{保禄}、\Place{默利特内}的主教\Person{奥特雷乌斯}以及\Place{纳齐盎}的主教\Person{额我略}，还有许多其他主教。他们在\Person{约维安}皇帝在位期间曾聚集在\Place{安提约基雅}，并决定坚持圣子与圣父同体的教义。此次会议上宣读了\Person{利贝里乌斯}及西方主教们的信函。这些信函令与会者极为满意；他们致信所有教会，请他们阅读亚洲主教们所颁布的法令，以及\Person{利贝里乌斯}与意大利、非洲、高卢和西西里主教们所撰写的文件——这些文件已委托给\Place{兰普萨库斯}会议的使者。他们敦促众人思考，这些文件的制定者人数众多，远超\Place{阿里米努姆}会议的成员，并劝勉他们同心合意，与他们共融，以书面形式表明这一点，最后在春季结束前齐集于\Place{西利济雅}的\Place{塔尔索斯}。他们在指定的固定日期彼此催促召开会议。待到约定的日子临近，当会议即将在\Place{塔尔索斯}召开时，约三十四位亚洲主教聚集在\Place{亚洲}行省的\Place{卡里雅}，他们赞同在教会内统一信仰的计划，但反对使用\Quote{同体}\OriginalTerm{同体}{consubstantial}一词，坚持认为由\Place{安提约基雅}会议和\Place{塞琉西亚}会议制定、并由殉道者\Person{路济安}及其许多前任在危险与紧张中维护的信条，应超越其他一切信条而居主导地位。\par

皇帝在\Person{欧多克修斯}的唆使下，写信阻止在\Place{西利济雅}召开会议，甚至以严厉惩罚加以禁止。他还致函各省总督，命令他们将所有曾被\Person{君士坦丁}流放、后又在于\Person{利安}皇帝治下重新担任司铎职务的主教逐出教会。根据这一命令，埃及政府首脑急于剥夺\Person{阿塔纳修斯}的主教职务并将其逐出城市；因为皇帝的谕令中规定了不轻的惩罚；若不执行命令，所有官员、其下属士兵及参事均须缴纳巨额罚款，并面临身体刑罚的威胁。\par

然而，城内大多数基督徒聚集起来，恳求总督不要不加考虑帝国皇帝信中的条款就放逐\Person{阿塔纳修斯}；该信仅指明所有曾被\Person{君士坦提乌斯}流放、而后被\Person{于利安}召回的主教；显然\Person{阿塔纳修斯}不在此列，因为他已被\Person{君士坦提乌斯}召回并恢复了主教职务；但\Person{于利安}却在其他所有主教均被召回之时迫害他，最后\Person{约维安}又将他召回。总督对这些论点毫不信服；不过他克制自己，没有诉诸武力。民众从四面八方涌来；全城陷入巨大的骚动与不安；一场暴动似乎即将爆发；因此他将实情禀报皇帝，并允许主教留在城内。数日后，当民众的激动似乎有所平息时，\Person{阿塔纳修斯}在黄昏时分秘密离开城市，藏匿于某处。就在当天夜里，埃及总督与军事首领占领了\Person{阿塔纳修斯}通常居住的教堂，搜遍了整个建筑，甚至屋顶，却一无所获；因为他们本想趁民众骚动部分平息、全城沉睡之际，执行皇帝的命令，悄悄将\Person{阿塔纳修斯}押送出城。\par

未能找到\Person{阿塔纳修斯}自然引起普遍惊讶。有人将他的逃脱归因于上天的特殊启示；有人则归因于其追随者的建议；两者结果相同；但预知并避开如此阴谋似乎需要超乎人的明智。有人说，当民众一有骚动迹象时，他因担心被视为任何可能发生的骚乱的起因，便藏匿于其祖先的坟墓中；随后又退至另一藏身之处。\par

不久之后，\Person{瓦伦斯}皇帝致信准许他返回并主持其教会。\Person{瓦伦斯}做出这一让步是否出于本意，十分令人怀疑。我更倾向于认为，他考虑到\Person{阿塔纳修斯}普遍受人敬重，担心引起\Person{瓦伦提尼安}皇帝的不满——\Person{瓦伦提尼安}以拥护尼西亚教义而闻名；或者他担心这位主教众多敬仰者会发动骚乱，以免某种变革损害公共事务。\par

我也认为，亚流派的主教们此次并未极力反对\Person{阿塔纳修斯}；因为他们认为，若他被逐出城市，他很可能会向皇帝们诽谤他们，从而获得与他们辩论的机会，或许还能说服\Person{瓦伦斯}采纳他自己的观点，并激起与之志同道合的\Person{瓦伦提尼安}对他们的愤怒。\par

\Person{君士坦提乌斯}在位期间所发生的事件证明了\Person{阿塔纳修斯}的德行与勇气，这使他们极为困扰。事实上，他如此巧妙地避开了敌人的阴谋，以致他们不得不同意他重新掌管埃及教会；然而，尽管\Person{君士坦提乌斯}已为此发出信函，他却几乎不愿从意大利返回。\par

我相信，正是由于这些原因，\Person{阿塔纳修斯}才没有像其他主教那样被逐出教会，而那些主教遭受了比异教徒施加的任何迫害更残酷的迫害。\par

那些不肯改变自己教义主张的人被放逐；他们的祈祷所被夺走，转归持相反意见者所有。唯独\Place{埃及}，在\Person{亚大纳削}在世期间，免于这场迫害。\par

% structure: p0080 source=h2 target=subsection reason=relative-heading-rank; base=h2

\subsection{第十三章 \Person{德摩斐洛}，一位亚略派人士，在\Person{欧多希乌斯}后成为\Place{君士坦丁堡}主教。虔诚者选举\Person{埃瓦格里乌斯}。记述随之而来的迫害。}

大约此时，皇帝\Person{瓦伦斯}前往\Place{奥龙特斯河}畔的\Place{安提约基雅}；在他旅途中，\Person{欧多希乌斯}去世了，他在治理\Place{君士坦丁堡}教会十一年之后。\Person{德摩斐洛}立即被亚略派主教们祝圣为他的继任者。尼西亚信经的追随者，认为事态的发展尽在他们掌握之中，便选举\Person{埃瓦格里乌斯}作为他们的主教。他是由\Person{欧斯塔修}祝圣的，\Person{欧斯塔修}曾治理\Place{叙利亚}的\Place{安提约基雅}教会，并由\Person{约维安}从流放地召回后，以私人身份居住在\Place{君士坦丁堡}，专心教导那些持守他见解的人，劝勉他们坚持对天主性的观点。亚略派异端者被鼓动造反，开始对参与\Person{埃瓦格里乌斯}祝圣的人发动猛烈迫害。皇帝\Person{瓦伦斯}当时在\Place{尼科美底亚}，得知\Person{欧多希乌斯}死后\Place{君士坦丁堡}发生的事件后，担心任何城市利益因骚乱受损，便派去了他认为足以维持安宁的军队。\par

\Person{欧斯塔修}奉他的命令被逮捕，放逐到\Place{比西亚}，一座\Place{色雷斯}的城市；\Person{埃瓦格里乌斯}被放逐到其他地区。此事便是这样发生的。\par

% structure: p0083 source=h2 target=subsection reason=relative-heading-rank; base=h2

\subsection{第十四章 记述\Place{尼科美底亚}的八十位虔诚代表，\Person{瓦伦斯}将他们连同船只在海中烧死。}

亚略派人士，如同得势者惯常所为，因更加傲慢，毫不怜悯地迫害所有宗教信仰与他们对立的基督徒。\par

这些基督徒遭受身体伤害，被出卖给长官和监狱，又因频繁罚款而逐渐贫困，终于被迫向皇帝上诉求援。皇帝虽极其愤怒，却没有公开表现任何怒气，而是秘密命令总督逮捕并处死整个使团。但总督担心，若不经任何司法程序处死这么多善良虔诚的人，会激起全民暴动，便假称流放他们，以此借口强迫他们登上一艘船，他们以最彻底的顺从同意了。当他们航行到海湾中央（称为\Place{阿斯塔基乌斯}）时，水手们按照接到的命令点燃了船，跳上小艇。一阵风起，船被吹到\Place{达西比扎}，\Place{比提尼亚}海岸的一个地方；但刚一靠近岸边，船就连同船上所有人完全烧毁了。\par

% structure: p0086 source=h2 target=subsection reason=relative-heading-rank; base=h2

\subsection{第十五章 \Place{凯撒勒雅}主教\Person{欧瑟伯}与\Person{巴西略}大帝之间的纷争。因此亚略派人士壮胆来到\Place{凯撒勒雅}，并被击退。}

\Person{瓦伦斯}离开\Place{尼科美底亚}后，前往\Place{安提约基雅}；在经过\Place{卡帕多细亚}时，他照常竭尽全力损害正统派，并将教会交给亚略派。他以为因\Person{巴西略}与当时治理\Place{凯撒勒雅}教会的\Person{欧瑟伯}之间正发生的争执，他的计划会更容易实现。这场不和曾导致\Person{巴西略}离开\Place{本都}，他在那里与一些修道的隐修士共同生活。城里的人民和一些最有势力、最有智慧的人开始怀疑\Person{欧瑟伯}，尤其认为他是使一位同样以虔诚和口才著称的人离开的原因；于是他们开始策划脱离教会并另立分开的教会。与此同时，\Person{巴西略}害怕成为教会更大麻烦的根源，教会已因异端者的纷争而分裂，便留在\Place{本都}的隐修院中隐居。皇帝和亚略派异端的主教们（他们常跟随他）因\Person{巴西略}的缺席和人民对\Person{欧瑟伯}的仇恨而更加得意洋洋。但结果却与他们的判断相反。一听说皇帝打算经过\Place{卡帕多细亚}，\Person{巴西略}便离开\Place{本都}回到\Place{凯撒勒雅}，在那里他与\Person{欧瑟伯}和解，并以他的口才适时帮助了教会。\Person{瓦伦斯}的计划因此受挫，他带着他的主教们无功而返。\par

% structure: p0088 source=h2 target=subsection reason=relative-heading-rank; base=h2

\subsection{第十六章 \Person{巴西略}在\Person{欧瑟伯}后成为\Place{凯撒勒雅}主教；他对皇帝和总督的胆色。}

此后不久，皇帝再次访问\Place{卡帕多细亚}，发现\Person{巴西略}在\Person{欧瑟伯}死后治理那里的教会。他本想驱逐\Person{巴西略}，却不情愿地克制了意图。据说，他制定计划的当晚，他的妻子被一个可怕的梦惊扰，他唯一的儿子\Person{加拉特}被一种急病夺去性命。这个儿子的死普遍被认为是天主的报复，惩罚他父母对\Person{巴西略}的阴谋。\Person{瓦伦斯}自己也这样认为，并在儿子死后不再骚扰主教。\par

当皇子因病日渐衰弱，濒临死亡时，皇帝派人请来巴西略，请求他为儿子的康复向天主祈祷。因为瓦伦斯一到凯撒勒雅，总督就派人请来巴西略，命令他接受皇帝的宗教观点，并以不服从就处死相威胁。巴西略回答说，能尽快从身体的束缚中解脱出来，对他来说是巨大的收获和至高恩惠的赐予。总督给了他当天剩余的时间和即将来临的夜晚来考虑，并劝告他不要鲁莽地冲入明显的危险，而应当在第二天来发表他的意见。\Quote{我不需要思考，}巴西略回答说。\Quote{我的决定明天将与今天相同；因为作为受造物，我永远不能被诱导去崇拜类似我自身的东西并将其当作天主来崇拜；我也不会顺从你们的宗教，也不会顺从皇帝的宗教。虽然你的显赫地位可能很高，虽然你有幸统治帝国不小的一部分，但我却不应该因此追求取悦于人，同时贬低那神圣的信仰，这信仰无论是丧失财产、流放还是被判死刑，都永远不会促使我背叛。此类折磨从未在我心中引起一丝悲伤。我除了一件斗篷和几本书之外一无所有。我在地上如同过客。身体因其软弱，在第一击之后将超脱一切感觉和折磨。}\par

总督钦佩这大胆回答中所显示出的勇气，并将情况报告了皇帝。在主显节的庆日，皇帝与众官员和卫队一同前往教堂，在圣桌上献上礼物，并与巴西略举行了一次会谈。巴西略在司铎职和教会管理方面的智慧、秩序和安排赢得了皇帝的赞扬。\par

然而不久之后，他敌人的诽谤得逞，巴西略被判处流放。执行法令的夜晚即将来临；皇帝的儿子突然患了紧急而危险的高烧。父亲俯伏在地，为仍然活着的儿子哭泣，不知还有什么其他方法能治愈儿子，于是派遣几名侍从去请巴西略来探望这倒卧的孩子；因为他自己因刚对巴西略施加的伤害而不敢召请主教。巴西略一到，男孩就开始好转；以致许多人认为，若不是有一些异端者被召来与巴西略一起为男孩的康复祈祷，他本可完全康复。据说总督也同样患病；但在他悔改并向天主祈祷后，恢复了健康。以上事例远不足以表达巴西略奇妙的禀赋；他对哲学生活的极度投入和惊人的口才吸引了极大的声誉。\par

% structure: p0093 source=h2 target=subsection reason=relative-heading-rank; base=h2

\subsection{第十七章 \Person{巴西略}与\Person{额我略}（神学家）的友谊；智慧同等，他们捍卫尼西亚教义。}

\Person{巴西略}和\Person{额我略}是同时代的人，可以说，他们都被公认为同样专注于修德。他们年轻时都在\Place{雅典}，师从\Person{希麦里乌斯}和\Person{普洛埃雷西乌斯}，这两位是当时最受推崇的诡辩家；后来在\Place{安提约基雅}，师从叙利亚人\Person{利巴尼乌斯}。但后来他们鄙视诡辩术和法律研究，决定按照教会的法律研究哲学。在花了一些时间追求异教哲学家所传授的科学之后，他们开始研究\Person{奥力振}以及在他前后时期最受认可的作者为解释\WorkTitle{圣经}而写的注释。\par

他们为那些像自己一样坚持尼西亚教义的人提供了巨大帮助，因为他们勇敢地反对亚略派异端者的教条，证明这些异端者既没有正确理解他们所依据的资料，也没有正确理解他们主要依赖的\Person{奥力振}的意见。这两位圣人通过彼此协议，或者如我所获悉的，通过抽签，分担了他们事业中的危险。\Place{彭图斯}邻近的城市分给了\Person{巴西略}；他在那里建立了许多隐修院，并通过教导百姓，说服他们接受与他相同的观点。\Person{额我略}在他父亲去世后，担任了\Place{纳齐盎}小城的主教，但因此住在许多地方，尤其是\Place{君士坦丁堡}。不久之后，他被许多司铎投票任命为那里人民的领袖；因为当时\Place{君士坦丁堡}既没有主教也没有教堂，\Place{尼西亚}大公会议的信理几乎绝迹。\par

% structure: p0096 source=h2 target=subsection reason=relative-heading-rank; base=h2

\subsection{第十八章 发生在\Place{安提约基雅}（位于奥龙特斯河畔）的迫害。\Place{厄德萨}的祈祷所，以宗徒\Person{多默}命名；那里的集会，以及\Place{厄德萨}居民的宣认。}

皇帝前往\Place{安提约基雅}，将城中及邻近各城所有拥护尼西亚信经的人悉数逐出教会；此外，他还用各种刑罚迫害他们；据一些人断言，他下令以各种方式处死许多人，并使其他人被投进\Place{奥龙特斯河}。他听说\Place{厄德萨}有一座以\Person{宗徒多默}命名的宏伟祈祷所，便去参观。他看见天主教会的成员在城外的平原上聚集敬拜；因为在那里，他们也被剥夺了祈祷之所。据说皇帝严厉斥责了长官，并挥拳打他的下巴，因为他容许这些集会，违背了他的谕令。\Person{莫德斯特}（长官的名字）虽然自己是个异端者，却暗中警告\Place{厄德萨}的百姓次日不要到惯常的地点聚集祈祷；因为他已接到皇帝的命令，要惩罚所有被抓获的人。他说出这样的威胁，是出于先见之明，好使没有人或至少只有少数人遭遇危险，并想平息君王的愤怒。但\Place{厄德萨}的百姓全然不顾威胁，比平常更加热切地聚集在一起，挤满了惯常的聚会地点。\par

\Person{莫德斯特}得知他们的行动后，不知该采取何种措施，便尴尬地随着人群来到平原。一个妇人牵着一个孩子的手，披风拖曳，不合妇女的体统，仿佛有要事似的，从长官率领的士兵队列中挤了过去。\Person{莫德斯特}注意到她的行为，下令逮捕她，并召她到面前，询问她奔跑的原因。她回答说，她正赶往平原上天主教会成员聚集的地方。\Quote{你不知道，}\Person{莫德斯特}回答说，\Quote{长官正在前往那里，要将所有在场的人处死吗？}\Quote{我听说了，}她回答说，\Quote{这正是我匆忙的原因；因为我怕去得太晚，从而失去为天主殉道的荣耀。}长官问她为何带着孩子，她回答说：\Quote{为了让他也分担共同的苦难，同享同样的赏报。}\Person{莫德斯特}对这妇人的勇气大为惊讶，便去见皇帝，将所发生的事告诉他，并劝他不要实行他显示为可耻且灾难性的计划。就这样，整个\Place{厄德萨}城宣认了基督信仰。\par

% structure: p0099 source=h2 target=subsection reason=relative-heading-rank; base=h2

\subsection{第十九章 伟大的\Person{亚大纳削}之死；拥戴亚流思想的\Person{路济乌斯}晋升主教座；他为埃及教会带来的众多灾祸；继\Person{亚大纳削}之后的\Person{伯多禄}前往罗马。}

\Person{亚大纳削}，\Place{亚历山大里亚}教会的主教，大约在这一时期去世，完成了约四十六年的大司祭职务。亚流派人士及早得知他的死讯，\Place{安提约基雅}的亚流派首领\Person{欧佐依}和财务总管\Person{玛格努斯}被皇帝派遣，毫不迟延地抓获并监禁了\Person{亚大纳削}指定继承主教职位的\Person{伯多禄}；他们立即将教会的管理权移交给\Person{路济乌斯}。\par

因此，在埃及的人比在其他地方的人遭受更重的苦难，灾祸连连，压迫着天主教会的成员；因为\Person{路济乌斯}一在\Place{亚历山大里亚}定居，就试图夺取教堂；他遭到人民的反对，而神职人员和圣童贞女被指控为骚乱的煽动者。有些人如同城市落入敌手一般逃脱了；另一些人被抓获并投入监狱。后来，一些囚犯从地牢中被拖出，用钩子和皮鞭撕裂；另一些人则被燃烧的火把焚烧。他们所受的折磨令人难以置信地存活下来。放逐甚至死亡本身都比这种痛苦更可忍受。主教\Person{伯多禄}从监狱逃脱，登船前往罗马，该教会的主教持有与他相同的见解。就这样，亚流派人士虽然人数不多，却占据了教堂。同时，皇帝颁布谕令，规定按照\Person{路济乌斯}的指示，将许多信奉尼西亚信经的人逐出\Place{亚历山大里亚}和埃及其余地方。\Person{欧佐依}于是完成了所有计划，返回\Place{安提约基雅}。\par

% structure: p0102 source=h2 target=subsection reason=relative-heading-rank; base=h2

\subsection{第二十章 对埃及隐修士及\Person{圣安当}门徒的迫害；他们因正统信仰而被围困在某岛上；他们所行的奇迹。}

\Person{路济乌斯}同埃及军队的将军前往旷野攻打隐修士；因为他认为，如果他能打破他们的抵抗，中断他们所爱的宁静，那么他在争取居住在城市中的基督徒加入他的阵营时遇到的障碍就会更少。这个国家的隐修院由几位圣德卓越的人管理，他们坚决反对亚流的异端。人们既不愿意也无能力探讨教义问题，他们从这些人那里接受观点，并与他们意见一致；因为他们相信，那些以行为彰显美德的人掌握着真理。我们听说，这些埃及苦修者的领袖是两位名叫\Person{玛加略}的人（前面已经提到过）、\Person{庞博}、\Person{赫拉克利德}和\Person{安当}的其他门徒。\par

思及亚略派永远无法在天主教会中成功建立优势，除非将隐修士拉拢到他们一边，\Person{路济乌斯}决定诉诸武力强迫隐修士支持他，因为他无法说服他们。然而，他的计划再次失败；因为隐修士已准备引颈受刃，也不愿偏离尼西亚教义。据传，正当士兵要攻击他们时，一个四肢枯萎、无法站立的人被抬到他们面前；他们用油傅抹了他，并奉基督（\Person{路济乌斯}所迫害的那位）的名命令他起来回家，他顿时痊愈。这一奇迹般的治愈公开显明了采纳那些天主亲自证明拥有真理之人的见解的必要性，而\Person{路济乌斯}被定罪，因为天主垂听了他们的祈祷并治愈了病人。\par

但密谋反对隐修士的人并未因这奇迹而悔改；相反，他们在夜间逮捕了这些圣人，并将他们送到埃及一个隐藏在沼泽中的岛上。这岛上的居民从未听说过基督信仰，而是侍奉魔鬼：岛上有一座极其古老且备受尊崇的神庙。据说，当隐修士登上该岛时，司祭的女儿（被魔鬼附身）走向他们。那女孩尖叫着跑向他们；岛上的人因她突然而奇怪的举动感到惊讶，便跟随着。当她走近载有圣使者的船时，她恳求地扑倒在地，大声哀求道：\Quote{你们为何到我们这里来，啊，伟大天主的仆人？因为我们长期居住在这岛上，未曾打扰任何人。不为人知，我们隐藏在此，四周都是沼泽。但如果你们愿意，请接受我们的财产，在此定居；我们将离开此岛。}\par

这就是她所说的话。\Person{玛加略}和他的同伴斥责了邪魔，女孩便恢复了神智。她的父亲和全家，以及岛上的居民，立即接受了基督信仰，拆毁了他们的神庙，将其改为一座教堂。这些事报告到\Place{亚历山大城}后，\Person{路济乌斯}被过度的忧伤所压倒；并且，担心招致同党仇恨，被指控与天主作战而非与人作战，他秘密下令将\Person{玛加略}和他的同伴送回他们旷野中的居所。这样，\Person{路济乌斯}就在埃及引起了麻烦和骚动。\par

大约同一时期，\Person{狄迪慕斯}哲学家和其他几位杰出人士获得了巨大声誉。民众被他们的德行以及隐修士的德行所打动，于是追随他们的教义，反对\Person{路济乌斯}的党羽的教义。\par

亚略派虽然在人数上不如另一方，却严重迫害了埃及的教会。\par

% structure: p0109 source=h2 target=subsection reason=relative-heading-rank; base=h2

\subsection{第二十一章 记载尼西亚教义被代表的各地清单；斯基泰人所显明的信德；维特拉尼奥，该民族的首领。}

\Emph{亚略主义}在同一时期在\Place{奥斯若恩}遇到了类似的反对；但在\Place{卡帕多西亚}地区，上智安排了一对神圣且最有学识的人——\Place{该地凯撒勒雅}的\Person{巴西略}主教和\Place{纳齐盎}的\Person{额我略}主教。\Place{叙利亚}及邻近省份，尤其是\Place{安提约基雅}城，陷入了混乱和无序；因为亚略派在这些地区人数众多，并占据了教堂。但天主教会成员人数也不少。他们被称为欧斯塔修斯派和保利诺派，由\Person{保利诺}和\Person{梅利提乌斯}领导，如前所述。正是通过他们的帮助，安提约基雅教会才得以免受亚略派的侵扰，并能够抵抗皇帝及其当权亲信的热忱。事实上，在所有由勇敢之人治理的教会中，民众都没有偏离他们先前的见解。\par

据说，这就是斯基泰人坚守其信德的原因。该国有很多城市、村庄和要塞。城都称为\Place{托米}，是一个大型且人口稠密的城市，位于海边的左侧，对航行到称为\Place{攸克辛海}的人来说。\par

按照至今仍盛行的古老习俗，全国所有教会都受一位主教管辖。\par

\Person{维特拉尼奥}在皇帝访问\Place{托米}时统治这些教会。\Person{瓦伦斯}前往教堂，并试图按照他的惯常做法将主教争取到亚略的异端；但后者勇敢地反驳了他的论点，在勇敢地为尼西亚教义辩护之后，离开了皇帝，前往另一座教堂，民众跟随了他。几乎全城的人都挤来看皇帝，因为他们期待这次与主教的面谈会产生一些不寻常的结果。\par

\Person{瓦伦斯}因与随从独自留在教堂而极为恼怒，并在怨恨中判决\Person{维特拉尼奥}流放。然而不久后，他召回了维特拉尼奥，我相信是因为他担心暴动；因为斯基泰人对主教的缺席感到不满。\par

他深知斯基泰人是一个勇敢的民族，他们的国家由于地理位置而拥有许多天然优势，对罗马帝国来说是必要的，因为它充当了抵御野蛮人的屏障。\par

这样，统治者的意图就被\Person{维特拉尼奥}公开挫败了。斯基泰人自己见证，他在其他方面都是良善的，并以生活的德行而卓越。\par

皇帝的愤怒波及所有圣职人员，但西方教会除外；因为统治西方地区的\Person{瓦伦提尼安}是尼西亚教义的仰慕者，他对宗教充满敬意，从未对司铎施加任何命令，也从未试图在教会规章中引入任何或好或坏的改变。虽然他已成为最好的皇帝之一，并展示了他治理事务的能力，但他认为教会事务超出他的管辖范围。\par

% structure: p0118 source=h2 target=subsection reason=relative-heading-rank; base=h2

\subsection{第二十二章 那时，关于圣神的教义引起争论，并决定祂应被视为与圣父和圣子\OriginalTerm{同性同体}{consubstantial}。}

当时一个问题重新被提出，它先前已引起许多探讨，现在更多；即圣神是否应被视为与圣父和圣子\OriginalTerm{同性同体}{consubstantial}。\par

关于这个主题发生了许多争执和辩论，类似于先前关于天主圣言本性的争论。那些断言圣子\OriginalTerm{与圣父不相似}{dissimilar from the Father}的人，以及那些坚持圣子\OriginalTerm{在本质上与圣父相似}{similar in substance}的人，在关于圣神的问题上达成了一致意见；因为双方都主张圣神在本质上不同，祂只是仆役，在秩序、尊荣和本质上位居第三。相反，那些相信圣子与圣父同性同体的人，对圣神也持相同的观点。在叙利亚，劳底迦主教\Person{阿波利纳里乌斯}尊贵地维护了这一教义；在埃及，是主教\Person{亚大纳削}；在卡帕多细亚和彭都斯的教会中，是\Person{巴西略}和\Person{额我略}。罗马主教得知这个问题被激烈地争论，并且每天都因争论而加剧，便写信给东方教会，敦促他们接受西方圣职人员所持的教义；即天主圣三的三位是同一实体、同等尊荣。这个问题既已由罗马教会如此决定，和平得以恢复，探讨似乎也告结束。\par

% structure: p0121 source=h2 target=subsection reason=relative-heading-rank; base=h2

\subsection{第二十三章 罗马主教\Person{里贝里乌斯}之死。\Person{达玛苏斯}和\Person{西里修斯}继任。正统教义在西方各处盛行，除米兰外，那里是\Person{奥克森提乌斯}担任大祭司。在罗马召开会议，\Person{奥克森提乌斯}被废黜；会议通过信函发送的定义。}

大约此时，\Person{里贝里乌斯}去世，\Person{达玛苏斯}继任罗马教区。一位名叫\Person{乌尔西奇乌斯}的执事，曾获得一些支持他的选票，但不甘失败，于是让一些名不见经传的主教秘密地祝圣了他，并试图在人民中制造分裂，成立一个独立的教会。他成功造成了分裂，一些人尊崇他为主教，而其余的人则拥护\Person{达玛苏斯}。这引发了人民中许多争斗和叛乱，最终发展到伤害和谋杀的地步。罗马长官不得不干预，惩罚了许多人民和圣职人员，并终结了\Person{乌尔西奇乌斯}的企图。\par

然而在教义方面，无论是在罗马还是在其他西方教会中，都没有发生分歧。人民一致坚持在尼西亚所确立的信仰形式，视圣三的三位在尊严和权能上平等。\par

\Person{奥克森提乌斯}及其追随者在意见上与其他人不同；他当时是米兰教会的主席，与少数同谋者一起，意图引入革新，并维护\OriginalTerm{阿里乌斯派的教义}{Arian dogma}，即\OriginalTerm{圣子和圣神的不相似}{dissimilarity of the Son and of the Holy Ghost}，根据最新兴起的探讨，这与西方司铎的一致意见相悖。高卢和威尼蒂亚的主教们报告说，他们之中也有其他人试图破坏教会的和平，于是几个省的主教们不久后在罗马集会，并下令将\Person{奥克森提乌斯}及持有其观点的人逐出共融。他们确认了尼西亚会议所确立的传统信仰，并废除了在里米尼会议上发布的与那信仰相悖的所有法令，理由是这些法令没有得到罗马主教以及同意他们的其他主教的认可，而且许多出席该会议的人不赞成他们在那里制定的条例。罗马主教\Person{达玛苏斯}和其余与会者致伊利里亚主教的信函证明了会议确实作出了这样的决定。信函如下：\par

\Quote{\Person{达玛苏斯}、\Person{瓦莱里乌斯}及在罗马举行的神圣会议的其他主教，致定居在伊利里亚的亲爱弟兄们，在主内问候。}\par

\Quote{我们相信你们坚持并向人民教导我们神圣的信仰，这信仰建立在宗徒的教义上。这信仰与教父们所定义的毫无差异；天主的司铎，其职责是教导智者，也不允许有其他想法。然而，我们从高卢和威尼斯的几位弟兄那里得知，某些人执意要引入异端。}\par

\Quote{所有主教都应勤勉防备这邪恶，免得他们中的一些羊群因缺乏经验，或另一些因单纯，而反对正确的解释。}\par

\Quote{不应追随那些发明新奇教义的人；而应保留我们父辈的意见，无论在我们周围有多少不同的判断。}\par

\Quote{因此，米兰主教\Person{奥克森提乌斯}已被公开宣布在这件事上首先被定罪。因此，罗马世界的所有教师应当同心一意，不因各种冲突的教义而玷污信仰。}\par

\Quote{因为当异端的恶意开始成熟时，正如阿里乌斯派的亵渎至今所做的——愿它远离我们——我们选出的三百一十八位蒙选者，在尼西亚进行研究后，筑起围墙对抗魔鬼的武器，并以这解毒剂击退致命的毒药。}\par

这份解毒剂在于相信：圣父与圣子具有同一\OriginalTerm{天主性}{Godhead}、同一德能、同一\OriginalTerm{本体}{\GreekText{χρῆμα}}。也必须相信圣神属同一\OriginalTerm{位格}{hypostasis}。我们规定，凡持其他道理者，应与我们的共融隔绝。\par

有人曾议定要玷污这救恩的定义和可钦崇的观点；但起初，在\Place{阿里米努姆}会议上推行革新的人，或被迫投票支持改变的人，后来多少弥补了过失，他们承认自己被某些似是而非的论点所欺骗，而这些论点在他们看来，与我们祖先在\Place{尼西亚}所确立的原则并不相悖。在\Place{阿里米努姆}聚会的个人数目，并不能作为反对正统教义的论据；因为那次会议既未经\Place{罗马}主教们的认可（\Place{罗马}主教们的意见本应最优先接受），也未获得\Person{味增爵}（他长期纯洁地守护主教职位）以及许多与上述主教意见相同的其他主教的赞同。\par

\Quote{此外，如前所述，那些似乎倾向某种虚幻的人，一旦运用更佳判断，就立即表明不赞同自己的作为。因此，你的纯洁应当看出，唯独这是在\Place{尼西亚}基于宗徒权威所确立的信仰，必须永远保持不受侵犯，并且所有东西方的、宣认大公宗教的主教，都应以与我们共融为荣。我们相信，不久之后，那些持有其他主张的人将被排除出共融，并被剥夺主教的名称和尊位；这样，现今被这些有害而欺骗性原则的轭所压迫的人民，可以得享自由。因为这些主教没有能力纠正人民的错误，因为他们自己也被错误所束缚。因此，让你的意见也与所有天主的司铎一致，我们相信在此事上你是圣善而坚定的。我们应当与你们一同如此相信，这可以通过与你的爱德交换信件来证明。}\par

% structure: p0134 source=h2 target=subsection reason=relative-heading-rank; base=h2

\subsection{第二十四章。论\Person{圣盎博罗削}及其被提升至大司祭职；他如何劝勉人民修持虔诚。\Place{夫里家}的诺瓦提安派与逾越节。}

西方的圣职人员既已预先阻止了那些试图在他们中推行革新之人的计划，便谨慎地继续保全那从起初就传授给他们的信仰的不可侵犯性。除了\Person{奥克森提乌斯}及其同党外，他们中间没有一个人持有异端意见。然而，\Person{奥克森提乌斯}此后不久便去世了。他死后，人民因\Place{米兰}教会主教的选择而发生了骚动，城市陷入危险。那些曾觊觎主教职位却未能如愿的人，像这类骚动中常见的那样，大声威胁。\par

当时担任该省总督的\Person{盎博罗削}，因担忧民众的骚动，便来到教堂，劝勉人民停止争论、遵守法律、重建和谐与和平所带来的繁荣。他讲话尚未结束，所有听众立刻抑制了彼此激起的愤怒情绪，并将主教职位的选票投向他，视之为他劝人和睦的成果。他们劝他领受洗礼，因为他尚未领洗，并恳求他接受司铎之职。在他拒绝、推辞、并真心逃避此事之后，人民仍坚持，声称除非他同意他们的愿望，否则纷争永不会平息；最后，这些消息传到了朝廷。据说，\Person{瓦伦提尼安}皇帝祈祷感谢天主，因为他所任命的总督正被选来担任司铎之职。当他得知人民的迫切愿望和\Person{盎博罗削}的拒绝时，他推断这些事件是天主为恢复\Place{米兰}教会的和平而安排的，于是下令尽快祝圣\Person{盎博罗削}。他在同一时刻领洗并被祝圣，随后立即着手使他所管辖的教会对神性本质达成一致意见；因为该教会在\Person{奥克森提乌斯}领导下，长期以来因这一议题而分裂。关于\Person{盎博罗削}在祝圣之后的行为，以及他如何勇敢而神圣地履行司铎职责，我们将在以后有机会提及。\par

大约此时，\Place{夫里家}的诺瓦提安派一反其古老习俗，开始与犹太人在同一日庆祝逾越节。其异端的创始人\Person{诺瓦提安}拒绝接受那些悔改自己罪恶的人进入共融，唯有在这一方面，他创新了既定的教义。但他及其继承者按照罗马教会的习俗，在春分后庆祝逾越节。然而，大约此时，一些诺瓦提安派主教在\Place{夫里家}靠近\Place{桑加鲁斯}河源的\Place{帕齐}镇集会，决定在纪律上不再遵循那些在教义上与他们不同者的做法，设立了一条新律法；他们决定与犹太人在同一日守无酵节并庆祝逾越节。\Place{君士坦丁堡}的诺瓦提安派主教\Person{阿格利乌斯}，以及\Place{尼西亚}、\Place{尼科美底亚}和\Place{夫里家}著名城市\Place{科提埃乌姆}的诺瓦提安派主教，都没有参加这次主教会议，尽管诺瓦提安派视他们为有关其异端及其教会事务的主宰与标志，可以这么说。由于这个原因，这些革新者深入分歧，并切断联系，形成了一个独立的教会，我将在适当的时机讲述此事。\par

% structure: p0138 source=h2 target=subsection reason=relative-heading-rank; base=h2

\subsection{第二十五章。论\Person{阿波利纳里}：父子同名。司铎\Person{维塔利安}。在被逐出一种异端后，他们转向其他异端。}

大约在此时，\Person{Apolinarius}公开炮制了一种异端，此后便以其名字命名。他引诱许多人脱离教会，并组成独立的聚会。\Place{Antioch}的一位司铎\Person{Vitalius}，也是\Person{Meletius}的神职人员之一，赞同并确认了他的独特见解。在其他方面，\Person{Vitalius}的生活与品行出众，热忱地看顾那些托付给他牧灵管理的人；因此他深受民众敬重。他脱离了与\Person{Meletius}的共融，加入\Person{Apolinarius}，并在\Place{Antioch}主持那些接纳相同意见的人；因他圣洁的生活，他吸引了大批追随者，这些人至今仍被\Place{Antioch}的市民称为\OriginalTerm{维塔利安派}{Vitalians}。据说他脱离教会是由于他对\Person{Flavian}（当时是他的同僚司铎，后来被擢升为\Place{Antioch}的主教）所表现出来的轻蔑感到愤慨。\Person{Flavian}阻止了他与主教进行惯常的会面，他便觉得自己被轻视，于是加入了与\Person{Apolinarius}的共融，并视其为朋友。从那时起，这一教派的成员在各城市中建立了独立的教会，有自己的主教，并制定了与天主教会不同的法规。除了常规的圣秩外，他们还咏唱一些由\Person{Apolinarius}创作的韵律诗歌；因为除了其他学识外，他还是一位诗人，精通多种格律，并以其甜美吸引许多人依附于他。男人在宴饮和日常劳作时唱他的诗，女人在织布时也唱。然而，无论是他的柔情诗歌适用于节日、庆节还是其他场合，它们都同样赞美和光荣天主。\Place{Rome}的主教\Person{Damasus}和\Place{Alexandria}的主教\Person{Peter}最先得知这一异端在民众中蔓延，并在\Place{Rome}举行的会议上投票将其视为与天主教会无关。据说，\Person{Apolinarius}在教义上创新，既有心胸狭隘的原因，也有其他原因。因为当治理\Place{Alexandria}教会的\Person{Athanasius}在从被\Person{Constantinus}放逐之地返回\Place{Egypt}途中，必须经过\Place{Laodicea}，他在该城与\Person{Apolinarius}建立了亲密关系，最终结成最牢固的友谊。然而，由于异端派认为与\Person{Athanasius}共融是可耻的，该城亚流派的主教\Person{George}以一种极具侮辱性的方式将\Person{Apolinarius}逐出教会，借口是他违反教规和神圣法律接待了\Person{Athanasius}。这位主教并未就此止步，还以他许久以前所犯并已悔改的罪行来指责他。因为当\Person{George}的前任\Person{Theodotus}管理\Place{Laodicea}教会时，诡辩家\Person{Epiphanius}朗诵了一首他为赞美\OriginalTerm{狄俄尼索斯}{Dionysus}而创作的颂歌。那时尚是青年的\Person{Apolinarius}，作为\Person{Epiphanius}的学生，由他的父亲（其名亦为\Person{Apolinarius}，一位著名的文法家）陪同去听朗诵。开场白之后，\Person{Epiphanius}按照公开朗诵颂歌时一贯的惯例，吩咐未受启蒙者和亵渎者离开。但年轻的\Person{Apolinarius}和年长的\Person{Apolinarius}，以及在场的任何基督徒，都没有离席。当主教\Person{Theodotus}听说他们曾在朗诵时在场，他极其不悦；然而，他赦免了犯此过错的平信徒，仅给予温和的责备。至于\Person{Apolinarius}父子，他公开定了他们的罪，并将他们逐出教会；因为他们二人原属神职人员，父亲是司铎，儿子是圣经的读经员。过了一段时间，当父子二人以眼泪和禁食表现出与罪行相称的悔改后，\Person{Theodotus}恢复了他们在教会中的职务。当\Person{George}接任同一主教职位时，他绝罚了\Person{Apolinarius}，视其为教会之外的人，原因如前所述，他接待了\Person{Athanasius}共融。据说\Person{Apolinarius}多次恳求他恢复共融，但他毫不宽容。\Person{Apolinarius}悲伤过度，扰乱了教会，并通过教义创新引入了前述异端；他企图凭自己的口才报复敌人，证明\Person{George}罢免了一个比自己更精通圣经的人。就这样，神职人员之间的私人仇怨时常严重损害教会，并将宗教分裂为许多异端。这便是证明；因为倘若\Person{George}像\Person{Theodotus}一样，在\Person{Apolinarius}悔改后接纳他共融，我相信我们绝不会听到以他命名的异端。人一旦遭受侮辱和轻视，便倾向于竞争和创新；而若受到公正对待，就会变得温和，安于原位。\par

% structure: p0140 source=h2 target=subsection reason=relative-heading-rank; base=h2

\subsection{第二十六章。\Person{Eunomius}及其师\Person{Aëtius}，他们的事务与教义。他们是首先提出圣洗圣事只进行一次浸洗的人。}

\Person{欧诺米}接替\Person{厄肋乌西}掌管\Place{基齐库}教会并领导亚略异端，又在此之外另创了一个异端，有些人以他的名称之为欧诺米异端，但有时也被称为\OriginalTerm{非类者}{Anomian}异端。有人断言，\Person{欧诺米}是第一个敢于主张神圣的圣洗应只进行一次浸洗，并以此方式败坏那慎重传承至今的传统的人。据说他发明了一种与教会相反的训练方式，并以庄重和更严格的表象掩饰这一创新。他是言辞与争辩的巧匠，喜爱辩论。持他见解的人大多也有同样的偏好。他们不赞赏善行或善德，或对穷人的怜悯，除非由他们本派的人表现出来，他们更赞赏辩论的技巧和在争论中获胜的能力。拥有这些才能的人在他们中间被视为最虔诚者。另一些人，我相信更真实地断言，\Place{卡帕多细亚}人\Person{特奥弗罗尼}和\Person{欧提基}，都是这一异端的热心传播者，在随后的朝代中脱离了与\Person{欧诺米}的共融，并在\Person{欧诺米}的其他教义和神圣的圣洗上创新。他们主张圣洗不应因天主圣三之名施行，而应因基督的死亡之名施行。看来\Person{欧诺米}在这个问题上并未提出新见解，而是从一开始就坚定地依附于亚略的观点，并一直如此。他晋升为\Place{基齐库}主教后，被自己的神职人员指控在教义上引入创新。\Place{君士坦丁堡}亚略异端的统治者\Person{欧多西乌}，召见他并迫使他向民众解释其教义；但发现他并无过错后，\Person{欧多西乌}劝他返回\Place{基齐库}。然而\Person{欧诺米}回答说，他不能与怀疑他的人待在一起；据说他抓住了分离的机会，尽管看起来他采取这一步骤实际上是出于对拒绝接纳他的老师\Person{阿埃齐}进入共融的愤恨。还有人补充说，\Person{欧诺米}与\Person{阿埃齐}住在一起，从未偏离他最初的见解。以上就是各人相互矛盾的陈述；一些人以这种方式叙述情况，另一些人则以另一种方式叙述。但无论是\Person{欧诺米}还是任何其他人首先对这些关于圣洗的传统做了这些创新，在我看来，这样的创新者，无论他们是谁，按照他们自己的说法，只有他们自己面临在未领受神圣的圣洗就离开此生的危险；因为如果他们在按最初推荐的方式受洗之后，发现自己无法重洗自己，那么就必须承认他们引入了自己并未遵守的做法，从而向他人施行了从未由他们自己或他人向他们施行过的圣洗。这样，在通过某种不存在的原则和私自的假定确立了教条之后，他们就把自己未曾领受的施行给了别人。这种假定的荒谬性从他们自己的承认中显而易见；因为他们承认未领洗者没有能力为他人施洗。那么，按照他们的意见，没有按他们的传统受洗的人就是未受洗者，因为未被适当地入门，而他们通过自己的实践证实了这一意见，因为他们重洗所有加入他们派别的人，尽管这些人先前已按照天主教会传统入门。这些不同的教条给信仰带来了无数的麻烦；许多人因在教义问题上普遍存在的意见分歧而不敢信奉基督教。\par

争论日益激烈，并且如同异端初期那样，它们日益增长；因为它们的领袖并不缺少热忱或言辞的能力；事实上，似乎天主教会的大部分会被这一异端颠覆，若不是遇到了\Place{卡帕多细亚}的\Person{巴西略}和\Person{额我略}作为对手的话。\Person{德奥多西}的统治稍后不久开始；他将异端派别的创始人从人口稠密的帝国区域放逐到更偏远的地区。\par

但是，为了使我历史的读者不至于不知道我所特别指出的两个异端的确切性质，我认为有必要说明，\Place{叙利亚}人\Person{阿埃齐}是通常归诸于\Person{欧诺米}的异端的创始人；他像\Person{亚略}一样，主张圣子与圣父不相似，圣子是被造物，从无中受造。持这些观点的人过去被称为\OriginalTerm{阿埃齐派}{Aëtians}；但后来，在\Person{君士坦提}统治期间，正如我们所说，有些派别主张圣子与圣父\OriginalTerm{同质}{homoousios}，另一些派别主张圣子与圣父\OriginalTerm{相似本质}{homoiousios}，而\Place{阿里米农}会议曾下令只认为圣子与圣父相似，\Person{阿埃齐}因对天主的亵渎和不敬被判处放逐。此后一段时间，他的异端似乎被压制了；因为既没有任何其他有名望的人，甚至连\Person{欧诺米}也不敢公开为其辩护。但当\Person{欧诺米}被提升到\Place{基齐库}教会接替\Person{厄肋乌西}时，他再也不能安静地约束自己，在公开辩论中又重新提出了\Person{阿埃齐}的教义。因此，正如经常发生的那样，异端派别最初创始人的名字被遗忘，\Person{欧诺米}的追随者以其自己的名字被称呼，尽管他仅仅复兴了\Person{阿埃齐}的异端，并且比最初传授它的人更大胆地传播了它。\par

% structure: p0144 source=h2 target=subsection reason=relative-heading-rank; base=h2

\subsection{第27章 \Person{神学家额我略}在致\Person{内克塔里}书信中论及\Person{阿波利纳里}和\Person{欧诺米}。当时的隐修士哲学使他们的异端显得突出，因这两人的异端几乎控制了整个东方。}

显然，\Person{欧诺米乌斯}和\Person{阿埃提乌斯}持相同意见。在著作的几处段落中，\Person{欧诺米乌斯}夸耀并屡次作证，\Person{阿埃提乌斯}是他的导师。\Person{额我略}，\Place{纳西盎}的主教，在写给君士坦丁堡教会的首领\Person{奈克塔里乌斯}的信中，以下列言辞谈及\Person{阿波利纳里}：\Quote{\Person{欧诺米乌斯}是我们中不断的麻烦之源，他自己成为我们的负担尚不满足，还要认为若不竭力拖每个人同他一起奔向他所急趋的毁灭，就是自己的过错。不过，这种行为在某种程度上或许可以容忍。教会现在必须对抗的最严重灾难源于阿波利纳里派的狂妄。我不知道您的圣座如何能同意他们像我们一样自由举行集会。您已因天主的恩宠在天主的奥秘中受到充分的教导，不仅明白天主的圣言的辩护，也明白异端者所做出的反对纯正信仰的一切创新；然而，您的可敬阁下从我们这些狭隘之人听到以下情况或许并无不妥：我得到了一本\Person{阿波利纳里}所写的书，其主张超出了所有异端邪恶的形式。他断言，为转化我们的本性而取的血肉，在天主独生子的救世计划中，并非为此目的而获得；而是这血肉之性从起初就存在于圣子内。他通过滥用以下圣经言辞来证实这一邪恶假设：\BibleQuote{没有人上过天}\BibleRef{若 3:13}。他根据这段经文断言，基督在未从天降下之前就是人子，并且当他降下时，他带来自己的血肉，这血肉是他已在万世之前就拥有且本质合一的。他还提出了另一句宗徒的话：\BibleQuote{第二个人是出于天}\BibleRef{格前 15:47}。此外，他主张那天降之人缺乏\OriginalTerm{理智}{\GreekText{νοῦς}}，但独生子的天主性实现了理智的本质，并构成了人类复合体的第三部分。身体和\OriginalTerm{灵魂}{\GreekText{ψυχὴ}}构成了两部分，如同在其他人类中一样，但没有理智，而是天主的圣言填补了理智的位置。这还不是可怕景象的结束；因为该异端最严重的一点是，他断言那独生天主，众人的审判者，生命的赐予者，死亡的毁灭者，自己竟屈服于死亡；他在自己的天主性内受苦，并且，在第三天身体的复活中，天主性也与身体一起从死者中复活；而它是由圣父从死者中复活的。要详述这些异端者所提出的其他一切荒谬教义，时间太长。}我想，以上所述足以表明\Person{阿波利纳里}和\Person{欧诺米乌斯}所持见解的性质。若有人渴望更详细的信息，我只能请他参考由他们或由别人针对这些人所写的相关著作。我不自称能轻易理解或阐述这些事，因为在我看来，这些教义未能盛行并进一步推进，除了提到的原因外，尤其应归因于那一时期的隐修士们；因为叙利亚、卡帕多细亚及邻近省份的所有那些哲学家都真诚地依附尼西阿信经。然而，从基里基雅到腓尼基的东部地区却处于\Person{阿波利纳里}异端的危险中。\Person{欧诺米乌斯}的异端从基里基雅和托鲁斯山脉一直蔓延到赫勒斯滂和君士坦丁堡。这两个异端者发现很容易吸引他们所居住和邻近地区的人加入各自的团体。但他们也遭遇了与亚略派相同的命运；因为人民钦佩那些以工作彰显德行并被认为持正确见解的隐修士，而回避那些持其他见解的人，将他们视为不虔敬并持有谬误教义的人。同样，埃及人被隐修士引导反对亚略派。\par

% structure: p0146 source=h2 target=subsection reason=relative-heading-rank; base=h2

\subsection{第二十八章 论这一时期在埃及兴盛的神圣人物。\Person{若望}，或\Person{阿孟}、\Person{贝努斯}、\Person{特奥纳斯}、\Person{科普勒斯}、\Person{赫勒斯}、\Person{埃利亚}、\Person{阿佩勒斯}、\Person{伊西多尔}、\Person{塞拉皮翁}、\Person{迪奥斯库罗斯}和\Person{欧洛吉乌斯}。}

由于这一时期以许多献身于哲学生活的圣洁人物而闻名，似乎有必要对他们作一些介绍，因为在那个时期涌现了大量天主所爱的人。看来，在埃及没有比\Person{若望}更著名的人物了。他从天主那里获得了如同古代先知一样清楚辨别未来和最隐秘之事的能力，此外，他还拥有治愈那些患有不治之症和疾病的人的恩赐。\Person{奥尔}是这一时期的另一位杰出人物；他自幼生活在孤独中，不断从事于咏唱赞美天主。他靠草药和根茎为生，有水时便以水为饮。晚年时，他奉天主的命令前往\Place{特拜}，管理几座隐修院，并且他也参与了神圣的工作。仅凭祈祷他就驱除了疾病和魔鬼。他不识字，也不需要书籍来帮助记忆；因为他心中所接受的任何事物都不会随后被遗忘。\par

\Person{阿孟}，被称为\OriginalTerm{塔本尼西奥特派}{Tabennesiotians}的隐修士们的领袖，住在同一地区，有大约三千门徒跟随。\Person{贝努斯}和\Person{特奥纳斯}同样管理隐修团体，并拥有预知和预言的恩赐。据说，尽管\Person{特奥纳斯}通晓埃及人、希腊人和罗马人的一切学问，他却实行了三十年的沉默。\Person{贝努斯}从未被看见表现出任何愤怒的迹象，也从未被听见发誓，或说出虚假、无益、轻率或无用的言语。\par

科普雷斯、赫勒斯和埃利亚斯也在这一时期兴盛。据说科普雷斯从天主获得了治愈疾病和各种病症以及战胜魔鬼的能力。赫勒斯自幼在隐修生活中受训，行了许多奇事。他能怀中抱火而不烧衣服。他激励同会隐修士修德，声称品行端正便会伴随奇迹。埃利亚斯在安提努城附近实践哲学，此时约一百一十岁；在此之前他说自己在荒漠独居了七十年。尽管年事已高，他仍不懈地守斋和勇敢的苦修。\par

阿佩莱斯在同一时期兴盛，在埃及的隐修院、靠近阿科里斯城的地方行了无数奇迹。他曾一度作铁匠，这是他本行；一天夜里魔鬼企图引诱他不洁，以美女形象出现；但阿佩莱斯抓住炉中烧红的铁，烫了魔鬼的脸，魔鬼像野鸟般尖叫逃走。\par

伊西多尔、塞拉皮翁和狄奥斯科普斯在这一时期是隐修士中最著名的教父。伊西多尔关闭他的隐修院，无人能出入，并供给院内所需。塞拉皮翁住在阿尔塞诺伊特附近，带领约一千名隐修士。他教导所有人通过劳动赚取生活所需，并供应其他穷人。在收获季节，他们忙于收割赚取酬劳；留足自己的粮食，与其余隐修士分享。狄奥斯科普斯只有不到一百位门徒；他是司铎，极其认真地履行司铎职责；他仔细查问那些前来申请参与神圣奥迹的人，使他们净化心灵，不忽视任何所犯的罪。司铎欧洛吉乌斯在分施神圣奥迹时更加谨慎。据说，当他执行司祭职务时，能看出前来者心中所想，清楚察觉每个人的罪和隐秘念头。他将所有犯罪或怀有恶意的人逐出祭台，公开揭露他们的罪；但当他们通过悔改净化自己后，他又重新接纳他们共融。\par

% structure: p0152 source=h2 target=subsection reason=relative-heading-rank; base=h2

\subsection{第29章 关于底比斯地区的隐修士：阿波罗、多罗特奥斯；关于皮亚蒙、若望、马尔谷、玛加略、阿波罗多洛斯、梅瑟、保禄（在菲尔马）、帕科、斯德望、皮奥尔。}

阿波罗在差不多同一时期在底比斯兴盛。他早年献身于哲学生涯；在荒漠度过四十年后，遵照天主的命令，将自己关闭在山脚下一个洞穴里，靠近一个人口众多的地区。由于许多奇迹，他很快声名显赫，成为许多隐修士的领袖；因为他用教诲有益地指导他们。领导亚历山大里亚教会的弟茂德记述了他的纪律方法以及他所行的神圣奇妙事迹；他也叙述了其他受称赞的隐修士的生平，其中许多我已提到。\par

那时，约有两千名善良的隐修士在亚历山大里亚附近宣讲哲学；有些在被称为隐修院的地区，另一些更靠近玛勒奥提斯和利比亚。多罗特奥斯是底比斯人，是这些隐修士中最著名的一位。他白天在海岸收集石头，用来建造小屋送给无力建造的人。夜间，他用棕榈叶编篮子出售以维持生计。他每天只吃六盎司面包配少量蔬菜，只喝水。从年轻时起就习惯这种极端的节食，并一直坚持到老年。从未见他躺卧在席子或床上，也不以舒适姿势摆放四肢，或愿意入睡。有时，由于自然倦怠，在日间劳作或吃饭时眼睛会不由自主闭上；吃饭打盹时食物会从口中掉落。一天，他完全被睡意压倒，跌倒在席上；他对自己处在此姿势感到不悦，低声说：\Quote{如果天使被说服去睡觉，你也会说服热心者。}或许他是对自己说的，或许是对那个阻碍他热心操练的魔鬼说的。有一次，一个人来到他那里，见他消耗自己，问他为何毁坏身体。他回答：\Quote{因为它毁坏我。}\par

皮亚蒙和若望主持了两座著名的埃及隐修院，靠近迪奥尔库斯。他们是司铎，非常谨慎恭敬地履行司祭职务。据说有一天，当皮亚蒙执行司祭职务时，他看见一位天使站在祭台旁，在一本书上记下在场的隐修士的名字，同时抹去缺席者的名字。若望从天主那里得到了战胜痛苦和疾病的权能，以致他能治愈痛风者，使瘫痪者康复。\par

其时，在\Place{斯刻提斯}附近旷野中，有一位年事已高、名叫\Person{本雅明}的长者，正以卓越的方式修行。天主赐予他能力，不用药物，仅凭手触摸或藉由经祈祷祝圣的一点油，就能治愈各种疾病。据说他患了水臌，身体肿胀到必须扩大门才能将他从斗室中抬出的地步。由于病情不允许他躺卧，他便在一张极大的皮褥上坐了八个月，仍继续治愈病人，并不因自己未得痊愈而懊恼。他安慰来访者，并请他们为自己的灵魂祈祷；又说他并不在意自己的身体，因为健康时它未曾有益于他，如今患病，也不能加害于他。\par

大约同时，著名的\Person{玛尔谷}、\Person{小玛加略}、\Person{阿颇罗尼}，以及埃及人\Person{梅瑟}，都住在\Place{斯刻提斯}。据说\Person{玛尔谷}自幼就以极端温和及谨慎著称；他熟记圣经，并显出了如此卓越的虔诚，以至\Place{塞利亚}的司铎\Person{玛加略}本人宣称，他从未将司铎在圣桌上给予领洗者之物交给\Person{玛尔谷}，而是一位天使将之递给\Person{玛尔谷}，玛加略说他看见了那只手，直到前臂。\par

\Person{玛加略}从天主领受了驱魔的能力。他无意中犯下的一桩杀人案，原是他投入修行生活的起因。他曾是牧羊人，带领羊群到\Place{马勒奥提斯湖}岸边吃草时，在戏耍中杀死了一个伙伴。因害怕被送交官府，他逃往旷野。他在那里隐藏了三年，后来在原地筑了一所小屋，住了二十五年。他惯常说，自己深受早年遭遇的灾祸之恩，甚至称那无意杀人为有裨益之举，因为它是他投身修行与蒙福生活的原因。\par

\Person{阿颇罗尼}在商业生涯中度过后，年老时退隐\Place{斯刻提斯}。他反思自己已太老，无法学习文字或任何技艺，便自费购买了各类药品及适合病人的食物，每天到第九时辰为止，将部分药品和食物送到每座隐修院门口，以救济患病者。他发觉这种做法对自己有益，便采用这种生活方式；当感到死亡临近时，他把药品交给一个人，并劝勉那人也去照他所做的去做。\par

\Person{梅瑟}原为奴隶，因道德败坏被主人赶出家门。他加入盗匪，成为匪帮头目。在犯下许多恶行并胆敢杀人之后，藉着某种突如其来的皈依，他度起了隐修生活，达到了修行的最高境界。他先前的生活方式所带来的健康而强健的体魄，刺激了他的想象力，激起对享乐的欲望，于是他采取一切可能的办法来折磨身体：他仅靠一点无烹饪的面包为生，承受艰苦劳作，每日祈祷五十次；他站着祈祷，不弯膝，也不闭眼入睡。有时，他夜间去隐修士的斗室，悄悄把他们的水罐装满，这十分费力，因为他有时要走十、二十、三十甚至更多的斯塔狄去寻水。尽管他竭尽全力折磨身体，但很久之后才能制服天生的体质健壮。据说有一次，盗匪闯入他修行的住所；他抓住并捆绑了他们，将这四个男子扛在肩上，带到教堂，让聚集在那里的隐修士照他们认为合适的方式处置他们，因为他认为自己无权惩罚任何人。他们说，从未见过如此突然的从恶到善的转变，也从未见过在隐修修行方面如此迅速的进步。因此，天主使他在魔鬼面前成为可畏的对象，他被祝圣为\Place{斯刻提斯}隐修士的司铎。这样度过一生后，他在七十五岁时去世，留下许多杰出的门徒。\par

\Person{保禄}、\Person{帕科}、\Person{斯德望}、\Person{梅瑟}（后两者是利比亚人）以及埃及人\Person{皮奥尔}，都在这个朝代兴盛。\Person{保禄}住在\Place{斯刻提斯}的\Place{费尔默山}，管理五百位苦修者。他不动手劳作，也不接受任何人施舍的衣食，只接受维持生存所需的食物。他除了祈祷什么也不做，每天向天主献上三百次祈祷。他把三百块小石子放在怀里，以免忘记其中任何一次祈祷；每次祈祷完毕，他就取出一块石子。当石子不再剩余时，他就知道已经完成了全部规定的祈祷。\par

\Person{帕科}也在这一时期生活在\Place{斯刻提斯}。他从青年直到耄耋之年，一直遵循这条道路，从未发现他在任何方面缺乏自制——无论是身体的食欲、灵魂的激情，还是魔的诱惑——简言之，在修行者应当战胜的一切事上，他都未曾失却刚毅。\par

斯德望住在于靠近马尔马里卡的马里奥提斯。六十年来，他藉着严谨达到了苦修成全，成为一位非常著名的隐修士，并与安当大帝亲密。他极其温和而审慎，其惯常交谈的方式甜美且有益，极能使忧苦者的灵魂得安慰，将他们转变为良好的心情，即使他们先前被似乎不可避免的悲伤所压抑。他对自己的痛苦也同样如此。他受着严重而无法治愈的溃疡之苦，外科医生被请来对患处进行手术。手术期间，斯德望忙于编织棕榈叶，并劝慰周围的人都不要为他的痛苦担忧。他告诉他们，除了思想天主所做的一切都是为了我们的益处，以及他的痛苦将有助于他真正的福祉之外，不要再有其他想法，因为这痛苦或许会补赎他的罪过，在此生受审判胜过来生。\par

摩西以其温和、仁爱以及藉祈祷治愈痛苦的能力而闻名。皮奥尔自幼便决心献身于哲学的生活；为此，他离开父家，并发誓永不再见任何亲戚。五十年后，他的一位姐妹听说他还活着，对这意外的消息喜不自胜，非要见到他不可。她所住地方的主教被这老妇人的叹息和眼泪所感动，便写信给斯刻提斯旷野中隐修院的院长们，请求他们派皮奥尔到他那里去。院长们于是命他前往他出生的城市，他不能拒绝，因为在埃及的隐修士看来，不顺服是不合法的，我想其他隐修士也是如此。他和另一位隐修士来到父亲的门口，让人通报。当他听见开门的声音时，便闭上眼睛，呼唤他姐妹的名字，对她说：\Quote{我是皮奥尔，你的兄弟；你尽可能看我看个够吧。}他的姐妹再次见到他，喜出望外，感谢天主。他在所站的门前祈祷，然后返回他所住的地方；他在那里挖了一口井，发现水是苦的，但他坚持饮用直到去世。这时人们才知道他克己的程度；因为他去世后，有好几个人试图在他住过的地方实践哲学，却发现无法留在那里。我深信，若不是因为他所奉行的哲学原则，他本可轻易藉祈祷将水变为甘甜；因为他在一个原本无水的地方引出了水。据说，在摩西的带领下，一些隐修士曾挖一口井，但预期的水脉并未出现，无论挖多深都不见水，他们正要放弃时，大约中午时分，皮奥尔加入了他们；他先拥抱了他们，然后责备他们缺乏信心和心量狭小；接着他下到他们所挖的坑中，祈祷后，用棍子击地三次。不久，一股泉水涌出地面，充满了整个坑。祈祷后，皮奥尔离开了；尽管隐修士们恳求他同他们一起开斋，但他拒绝了，声称他来到他们这里并非为此，而只是为了完成他所作的事。\par

% structure: p0165 source=h2 target=subsection reason=relative-heading-rank; base=h2

\subsection{第三十章  斯刻提斯的隐修士：奥利振、狄迪莫斯、克罗尼翁、奥西修斯、普图巴图斯、阿尔西翁、塞拉皮翁、阿蒙、欧瑟比、狄奥斯克斯，以及被称为\Quote{长人}的弟兄们，与哲学家埃瓦格里乌斯。}

在此时期，圣安当的门徒之一奥力振仍年事甚高，居于塞提斯隐修院中。另有狄迪莫斯、克罗尼翁（约一百一十岁）、伟大的阿尔西修斯、普图巴图斯、阿尔西翁及塞拉皮翁，他们皆曾与圣安当同时。在修习哲学中已臻年老，且于此时期主持各隐修院。其中有些青年和中年圣德之人，因卓越的善行而闻名。其中有阿蒙尼乌斯、欧瑟比乌和狄奥斯库鲁斯。他们是兄弟，但因身材高大，被称为\Quote{长兄弟}。据说阿蒙达到了哲学的顶峰，因而克服了对安逸与享乐的爱慕。他极好学，读过奥力振、狄迪莫斯及其他教会作家的著作。自幼年至逝世，除面包外未尝吃过任何经火烹制之物。他曾被选为主教，在竭力推辞无效后，割去自己一只耳朵，并对前来迎请的人说：\Quote{去吧。从此司铎律法禁止我受祝圣，因为司铎的形体应当完美无缺。}前来迎请的人便离去；但得知教会并不恪守犹太律法要求司铎肢体完全，而只要求其品性无瑕疵后，他们返回阿蒙处，试图强逼他。他向他们声明，若施以任何强暴，他将割掉自己的舌头；他们被此威胁吓住，立即离去。从此阿蒙被称为帕罗特斯。此后不久，在随后统治时期，智者埃瓦格里乌斯与他结为密友。埃瓦格里乌斯是一位智者，在思想和言语上都有力量，善于辨别导向美德和恶习的论据，并能够劝勉他人效法前者、避免后者。他的辩才由留下的著作充分证明。至于他的品德，据说他完全无骄傲或傲慢，故受公正赞誉时不骄，遭不公指责时不愠。他是靠近攸克辛海的伊比利亚的公民。他曾在纳齐安的额我略主教门下修习哲学并研读圣经，并在额我略管理君士坦丁堡教会时担任总执事。他容貌俊美，衣着考究；因此与某位妇人的相识引起了其丈夫的嫉妒，后者图谋杀害他。就在阴谋即将实施时，天主在他睡眠中赐予他一个可怕而拯救的异梦。他梦见自己因犯某罪而被捕，手足被铁链捆绑。当他被带到法官面前接受判决时，一位手持福音书的人向他发言，许诺将他从捆绑中解放，并起誓确认，条件是他离开该城。埃瓦格里乌斯触摸了经书，起誓他会这样做。立即他的锁链似乎脱落，他醒了。他被此神圣之梦说服，逃离了危险。他决心献身于苦修生活，从君士坦丁堡前往耶路撒冷。过了一段时期，他去探望塞提斯的哲学家们，并欣然决定住在那里。\par

% structure: p0167 source=h2 target=subsection reason=relative-heading-rank; base=h2

\subsection{第三十一章 论尼特里亚的隐修士及称为\Quote{塞利亚}的隐修院；关于林科科鲁拉的一处（隐修院）；关于梅拉斯、狄奥尼修斯和索隆。}

他们称此地为尼特里亚。此处住有大量献身于修习哲学者，其名源自附近一个采集天然碱的村庄。该地约有五十座隐修院，彼此相距不远，有些由共居的隐修士居住，另一些则由独居的隐修士居住。在沙漠更深处，距此地约七十斯塔迪翁处，另有一地称为塞利亚，其中散布着许多小屋，故名；但彼此相隔甚远，以至于其中的人既不能看见也不能听见对方。他们在每周的第一日和最后一日集会；若有隐修士缺席，显然是非自愿地被落下，因患病而受阻；他们并不全体即刻前去看望和照料他，而是轮流在不同时间前往，各人携带适合疾病之物。除此类原因外，他们很少交谈，除非其中有人能够传授关于天主及灵魂得救的更深入知识。居住在塞利亚中的是那些已达哲学顶峰、因而能自主行为、独居生活的隐修士，他们为求宁静而与其他人分离。关于塞提斯及其哲学家，我简要陈述如此。若我再详述其生活方式，恐有人会指责我写作冗长；因为他们建立了各自的生活、劳作、习俗、锻炼、节制和时间，按照个人年龄自然划分。\par

此时，\newline{}\Place{Rhinocorura}也因圣人们而闻名，他们并非来自外地，而是本地人。\newline{}我听说他们中最杰出的修行者包括\Person{梅拉斯}，他当时治理该地的教会；\Person{狄约尼削}，他主持了城北的一座隐修院；以及\Person{索隆}，他是\Person{梅拉斯}的兄弟，也是主教位的继任者。据说，当驱逐所有反对亚略主义之司铎的法令颁布后，奉命逮捕\Person{梅拉斯}的官员发现他正以最卑微仆人的身份，在教堂修剪灯芯，斗篷上系着沾满油的腰带，并拿着灯芯。当他们向他询问主教时，他回答说主教在里面，并愿意带他们去见他。由于旅途疲惫，他带他们到主教住所，请他们坐在桌前，用他所有之物给他们吃。用餐后，他给他们水洗手，因为他服侍了客人，然后告诉他们他是谁。他们对他的行为感到惊奇，承认了他们的使命，但出于对他的尊敬，给了他完全的自由，可去任何地方。然而，他回答说，他不会逃避与他持相同信念的其他主教所面临的苦难，他愿意去流亡。他自幼便修行，已锻炼自己具备一切隐修美德。\par

\Person{索隆}放弃了商业追求，选择了隐修生活，这对他的福祉大有裨益。因为在他兄弟和其他苦修者的指导下，他在对天主的虔诚和对邻人的仁爱方面进步迅速。\Place{Rhinocorura}的教会从一开始就在如此模范的主教领导下，从未偏离他们的训诲，并培养出许多善人。该教会的圣职人员同住一屋，同席而坐，一切共享。\par

% structure: p0171 source=h2 target=subsection reason=relative-heading-rank; base=h2

\subsection{第32章　巴勒斯坦的隐修士：\Person{赫西加斯}、\Person{厄丕法尼}（后来在塞浦路斯）、\Person{阿孟}与\Person{西尔瓦诺}}

巴勒斯坦有许多隐修团体兴盛。我在\Person{君士坦提乌斯}统治下所列举的许多人，仍在从事修行科学。他们及其同伴达到了修行生活的顶峰，为自己的隐修院增添了更大的声誉；其中\Person{赫西加斯}——\Person{希拉里翁}的同伴，以及\Person{厄丕法尼}——后来成为塞浦路斯\Place{Salamis}的主教，尤其值得注意。\Person{赫西加斯}在他老师先前居住的同一地点专务修行生活；而\Person{厄丕法尼}则在\Place{Eleutheropolis}辖区的村庄\Place{Besauduc}附近定居，那是他的出生地。他自幼受最著名苦修者的教导，为此大部分时间在埃及度过，\Person{厄丕法尼}因其在隐修修行方面的造诣，在埃及和巴勒斯坦最为著名，并被塞浦路斯居民选为岛上首府的主教。因此，我认为他是普天之下最受尊敬的人，可以说；因为他是在一个大城市和港口履行司祭职；当他投身于世俗事务时，他以如此的美德处理它们，以至于不久便为所有市民和各种各样外来人所知；一些人是因为见过他本人，体验过他的生活方式；另一些人则是从这些见证者那里听说的。在去塞浦路斯之前，他曾在当今皇帝统治期间，在巴勒斯坦居住了一段时间。\par

同时期，在隐修院中，四兄弟——\Person{撒拉米内斯}、\Person{普斯科}、\Person{玛拉基翁}和\Person{克里斯丕翁}——非常杰出；他们在\Place{Gaza}的一个村庄\Place{Bethelia}附近修行；他们出身于一个居住在此的贵族家庭，并由\Person{希拉里翁}教导修行。据说有一次兄弟们正在回家途中，\Person{玛拉基翁}突然被带走，变得看不见；不久之后，他又重新出现，继续与兄弟们同行。此事后不久，他并未存活多久，而是在青春年华去世。他在美德生活与虔诚的修行方面，不亚于年长之人。\par

\Person{阿孟}住在距离上述之地十斯塔狄的地方；他住在\Place{Gaza}的一个城镇\Place{Capharcobra}附近，那是他的出生地。他在推行苦修主义方面非常严格且勇敢。我认为\Person{西尔瓦诺}——巴勒斯坦本地人，因他的崇高德行，曾有一次见到天使服侍他——大约同时期在埃及修行。之后他住在西乃山，后来又曾在\Place{Gerari}的洼地建立了一个非常广阔且著名的隐修院，容纳许多善人，杰出的\Person{匝加利亚}后来管理该院。\par

% structure: p0175 source=h2 target=subsection reason=relative-heading-rank; base=h2

\subsection{第33章　叙利亚和波斯的隐修士：\Person{巴特乌斯}、\Person{欧瑟比}、\Person{巴尔革斯}、\Person{哈拉斯}、\Person{阿博}、\Person{拉匝禄}、\Person{阿巴德乌斯}、\Person{则诺}、\Person{赫略多洛}、\Place{Carræ}的\Person{欧瑟比}、\Person{普罗托革内斯}和\Person{阿奥内斯}}

讓我們由此轉向敘利亞和波斯，即敘利亞鄰近的地區。我們將發現這些國家的隱修士在實踐修道哲學方面效法了埃及的隱修士。\Person{巴特烏斯}、\Person{歐瑟伯}、\Person{巴爾格斯}、\Person{哈拉斯}、\Person{阿波斯}、\Person{拉匝祿}（後來獲得主教尊位）、\Person{阿巴德烏斯}、\Person{則諾}和\Person{赫略多洛}，均在尼西比斯附近稱為西戈隆的山上興盛一時。當他們最初開始修行生涯時，被稱為牧羊人，因為他們沒有房屋，不食麵包、肉，也不飲酒，而是常居山上，按照教會的規矩，以祈禱和讚美詩讚美天主來度過時光。在通常的用餐時間，他們每人拿起一把鐮刀，到山上割些草，如同在牧場上吃草的羊群，這便成了他們的餐食。這就是他們修道哲學的歷程。\Person{歐瑟伯}自願將自己關在卡雷附近的一間小室內修行。\Person{普羅托革內}居住在同一地點，並在\Person{維托}（當時的主教）之後管理那裡的教會。這就是著名的維托，據說君士坦丁皇帝初次見他時，承認天主曾多次在顯現中向他顯示這人，並命令他絕對服從這人所說的話。\Person{阿奧內斯}在法達納有一座隱修院；那裡是亞巴郎的孫子雅各伯從帕勒斯坦旅行途中遇見日後所娶少女的地方，也是他滾開石頭使她的羊群能喝井水的地方。據說阿奧內斯是第一個將離群索居的生活和嚴格的修道哲學引入敘利亞的人，正如安當首先將其引入埃及一樣。\par

% structure: p0177 source=h2 target=subsection reason=relative-heading-rank; base=h2

\subsection{第三十四章 埃德薩的隱修士：\Person{猶利安}、\Person{敘利亞的厄弗冷}、\Person{巴魯斯}和\Person{厄烏洛吉}；以及科勒敘利亞的隱修士：\Person{瓦倫提諾}、\Person{德奧多羅}、\Person{梅羅薩斯}、\Person{巴蘇斯}和\Person{巴索尼烏}；又加拉太、卡帕多西亞及其他地方的聖人；為何這些聖人直到最近都長壽}

加達納斯和阿濟祖斯與\Person{阿奧內斯}同住，並效法他的德行。\Person{敘利亞的厄弗冷}是位史學家，在我們敘述君士坦提烏斯統治時期的事件時已經提及，他與\Person{猶利安}一同成為當時最著名的修行者，活躍在埃德薩及其鄰近地區。\Person{巴爾塞斯}和\Person{厄烏洛吉}，在比我們此刻所論較晚的時期，均被祝聖為主教，但並非任何城市的主教；因為這只是榮譽頭銜，作為對其卓越行為的補償而授予，並在他們自己的隱修院內祝聖。\Person{拉匝祿}，我們已提及，也以同樣方式被祝聖為主教。以上便是敘利亞、波斯及鄰近地區最著名的苦修主義修行者，就我所能查實而言。他們共同的途徑，可以這樣說，在於勤勉關注靈魂的狀態，通過禁食、祈禱和讚美天主的聖詩，使之恆常準備脫離此世的事物。他們將大部分時間用於這些神聖的操練，並且完全鄙視世俗的財產、暫時的事務以及身體的安逸和修飾。有些隱修士將克己推至非凡的高度。例如，\Person{巴特烏斯}因過度的節制和禁食，竟有蠕蟲從他的牙齒中爬出；\Person{哈拉斯}又有八十年未嘗麵包；\Person{赫略多洛}則有許多夜晚不眠，並加上七天的禁食。\par

雖然科勒敘利亞和上敘利亞，除安提約基雅城外，緩慢地皈依了基督教，但並不缺乏教會的修行者，他們的行為因必須面對當地居民的敵意和仇恨而顯得更為英勇。他們高尚地克制自己，不抵抗，也不訴諸法律，而是勇敢地忍受外教人的侮辱和毆打。我發現，\Person{瓦倫提安}採取了這樣的途徑，根據某些記載，他生於厄梅撒，但根據其他記載，生於阿雷圖薩。另一個同名的個人以類似行為著稱，\Person{德奧多羅}也一樣。兩人都來自提提，位於阿帕梅亞地區。\Person{梅羅薩斯}也毫不遜色，他是內基利斯本地人，\Person{巴蘇斯}、\Person{巴索內斯}和\Person{保祿}亦如此。後者來自特爾米松村。他在許多地方建立了許多團體，並引入了修道哲學知識所必需的方法，最終在一個稱為尤加圖姆的地方建立了最大、最著名的隱修士團體。在那裡，經過長期而光榮的生活後，他去世並被安葬。有些以卓越和神聖方式實踐修道哲學的隱修士存活到了我們自己的時代；確實，提到的多數人都享有很長的壽命；我深信天主延長他們的日子，正是為了促進信仰的利益。他們引導幾乎整個敘利亞民族、以及大部分波斯人和撒拉森人皈依了正確的信仰，並使他們停止了崇拜偶像。在那裡開始修道哲學之後，他們培養了許多像他們自己的人。\par

我想，加拉太、卡帕多西亞及鄰近省份在那時也有許多其他教會的修行者，因為這些地區先前已熱切地接受了我們的教義。這些隱修士大多居住在城鎮和村莊的團體中，因為他們不習慣前輩的傳統。該國冬季的嚴寒（通常是一自然特徵）可能使獨修生活變得不切實際。據我所知，\Person{萊昂提烏}和\Person{普拉皮狄烏}是這些隱修士中最著名的。前者後來管理安基拉教會，後者年事極高，在幾個村莊執行主教職務。他還主管了巴西利亞斯，即最著名的窮人收容所。它是由凱撒勒雅主教\Person{巴西略}建立的，因而最初得名，並保留至今。\par

% structure: p0181 source=h2 target=subsection reason=relative-heading-rank; base=h2

\subsection{第三十五章 木製三腳架與通過其字母知識得知的皇位繼承；哲學家的毀滅；天文學}

关于那个时代的教会哲学家，我已能收集到的信息就是如此。至于外教人，在我们所指的那个时期，他们几乎全被消灭。他们中有些人，以精通哲学著称，对基督宗教的进展极为不满，正在谋划谁将成为\Person{瓦伦斯}在罗马帝国皇位上的继承人，并为了获取这种对未来的洞察，诉诸各种占卜术。经过各种咒语后，他们用月桂木建造了一个三脚架，并以他们惯用的祈求与话语结束；这样，皇帝的名字就可以通过三脚架和预言机制逐一指示的字母组合显示出来。他们张口结舌地等待着\Person{狄奥多若}，此人在宫中担任显赫的军职。他是外教人，也是一位有学问的人。字母的排列，一直排到他名字的\GreekText{Δ}，欺骗了这些哲学家。他们因此预期狄奥多若很快会成为皇帝。当他们的计划被告发时，\Person{瓦伦斯}怒不可遏，仿佛有人阴谋反对他的安全。因此，所有人被捕；\Person{狄奥多若}和三脚架的建造者被下令处死，有的用火，有的用剑。同样，由于同样的原因，帝国中最杰出的哲学家被杀；由于皇帝的愤怒不受约束，死刑甚至波及那些不是哲学家、但穿着类似服装的人；因此，从事其他职业的人不再穿\OriginalTerm{克罗科提翁袍}{crocotium}或\OriginalTerm{特里波尼翁袍}{tribonium}，出于怀疑和恐惧危险，以免显得是在追求巫术和邪术。我丝毫不认为，明智之士会因皇帝的愤怒和残忍而责备他更甚于责备那些哲学家，因他们的鲁莽和非哲学性的行为。皇帝荒谬地以为他可以处死自己的继任者，因此既不宽恕那些预言者，也不宽恕他们所预言的对象，据说他也不宽恕那些与\Person{狄奥多若}同名的人——其中一些是显贵——无论他们的名字是否完全一致，还是以\GreekText{θ}开头、以\GreekText{δ}结尾的相似名字。另一方面，这些哲学家的行为仿佛皇帝的被废黜和复位完全取决于他们；因为如果皇位继承应被视为取决于星辰的排列，那么除了等待未来皇帝（无论他是谁）登基之外，还需要什么呢？或者如果继承被视为取决于天主的旨意，人有什么权利干预？因为理解天主的想法不是人的预知或热忱所能及的；即使这是正确的，人，即使是最聪明的人，以为他们可以比天主计划得更好，也是不妥的。如果仅仅出于鲁莽的好奇心去探究未来的事，以致他们表现出如此缺乏判断力，以至于准备陷入危险，并藐视罗马人自古设立的法律，而且是在进行外教崇拜和祭献并不危险的时期；在这方面，他们的想法与\Person{苏格拉底}不同；因为当\Person{苏格拉底}被不公正地判处饮毒酒时，他拒绝通过违反他出生和受教育所在的法律来拯救自己，也不愿越狱，尽管他有权这样做。\par

% structure: p0183 source=h2 target=subsection reason=relative-heading-rank; base=h2

\subsection{第三十六章 远征\Place{萨马提亚人}；\Person{瓦伦提尼安}在\Place{罗马}驾崩；\Person{小瓦伦提尼安}被立为帝；对司铎的迫害；哲学家\Person{忒米斯提乌斯}的演说，因之\Person{瓦伦斯}倾向于更人道地对待与他意见不同的人。}

然而，上述这类主题最好留给个人判断去审查和决定。\par

\Place{萨马提亚人}入侵帝国西部，\Person{瓦伦提尼安}征召军队对抗他们。然而，他们一听说征召来对抗他们的军队的数量和实力，就派遣使团求和。当使节被引入\Person{瓦伦提尼安}面前时，他问他们是否所有\Place{萨马提亚人}都像他们一样。他们回答说，该国的主要人物被选为使团成员，皇帝大怒，喊道：\Quote{我们的臣民遭受着可怕的事，一场灾难正笼罩着罗马政府，如果\Place{萨马提亚人}——这个野蛮种族，这些人是你们的最优秀者——不喜欢安分守己，却胆敢入侵我的政府，并且竟敢对罗马人发动战争。}他用非常高的音调这样说了好一会儿，他的愤怒如此猛烈和无度，以至于最终他同时爆裂了一根血管和一根动脉。结果他失血过多，不久后在\Place{高卢}的一座堡垒中去世。他大约五十四岁，在十三年间，以良好的结果和卓越的成就驾驭着政府的缰绳。他死后六天，他的幼子与他同名的，被士兵们立为皇帝；不久之后，\Person{瓦伦斯}和他的兄弟\Person{格拉提安}正式同意了这个选举，尽管他们起初对士兵在未事先征得他们同意的情况下就将政府的象征移交给他而感到恼怒。\par

在此期间，\Person{瓦伦斯}已将居所定于叙利亚的\Place{安提约基雅}，对在有关天主性问题上与他意见相左之人更为敌视，更严厉地折磨并迫害他们。哲学家\Person{德米斯提乌斯}在他面前发表了一篇演说，劝诫他不应对教会教义上的分歧感到惊异，因为与外教人相比，这已算是较为温和且较轻微的了；外教人中间的意见是多形的。在众多导致无尽争论的信理中，其差异必然引发更多的争辩与讨论。因此，或许这正是天主所喜悦的：祂不愿被轻易认识，甚至允许意见分歧，这样每个人反而更敬畏祂，因为对祂的精确认识是如此不可企及。而试图归纳这浩瀚时，人便会倾向于得出祂是何等伟大、何等美善的结论。\par

% structure: p0187 source=h2 target=subsection reason=relative-heading-rank; base=h2

\subsection{第37章 论多瑙河以外的蛮族被匈奴驱逐、进逼罗马人及其皈依基督信仰；\Person{乌尔菲拉}与\Person{阿塔纳里克}；他们之间发生的事件；哥特人由此接受亚略主义。}

特米斯提乌斯这篇杰出的演说使皇帝变得较为人道，因此惩罚比以前减轻了。若不是国家大事的忧虑接踵而至，使他无暇继续追究，他决不会完全收回对司铎们的愤怒。原来，居住在依斯特河对岸、曾征服其他蛮族的哥特人，被匈奴人击败并逐出家园后，越过边界进入罗马领土。据说，在此之前，依斯特河的色雷斯人和哥特人都不知道匈奴人；因为虽然他们彼此隐秘地相邻而居，但中间隔着一个巨大的湖泊，湖两岸的居民各自以为自己的国家位于大地尽端，以外别无他物，只有大海和水。然而，恰巧有一头受昆虫折磨的牛跳入湖中，牧人追赶它；他第一次发现对岸有人居住，便将此事告知本族同胞。然而，有些人叙述说，一只鹿在逃跑，向几个属于匈奴族的猎人指出了那条被水面浅浅掩盖的道路。到达对岸后，猎人们对那片土地的美丽、空气的宁静以及适宜耕种感到惊叹，并将所见报告给他们的国王。于是匈奴人尝试以少数士兵攻击哥特人；但随后他们组建了一支强大的军队，在战斗中征服了哥特人，并夺取了他们的全部领土。被征服的民族，在敌人的追击下，越过边界进入罗马领土。他们渡过大河，派遣使团到皇帝那里，向他保证愿意参与他可能发动的任何战争，只要他划出一块土地供他们居住。该民族的\Person{乌尔菲拉}主教是使团的首领。他的使命完全达成了，哥特人被允许在\Place{色雷斯}定居。不久，他们内部爆发了争执，导致他们分裂为两部分，一个由\Person{阿塔纳里克}领导，另一个由\Person{弗里提格恩}领导。他们互相拿起武器对抗，弗里提格恩战败，恳求罗马人援助。皇帝命令色雷斯的军队援助并与弗里提格恩结盟，于是发生了第二次战斗，阿塔纳里克及其党羽被击溃。为感谢\Person{瓦伦斯}及时提供的援助，并证明自己对罗马人的忠诚，弗里提格恩接受了皇帝的宗教，并说服他所统治的蛮族效仿他。然而，我认为这并不是唯一可以用来解释哥特人为何直到今日仍保留\OriginalTerm{阿里乌斯主义}{Arianism}教义的理由。因为他们的主教乌尔菲拉最初所持的见解与天主教会并无分歧；因为在君士坦提乌斯统治时期，尽管他（我相信是出于轻率）与\Person{欧多克西乌斯}和\Person{阿卡基乌斯}一起参加了君士坦丁堡会议，但他并未偏离\OriginalTerm{尼西亞大公会议}{Nicæan council}的信理。后来他似乎回到了\Place{君士坦丁堡}，据说与阿里乌斯派的首领们就教义问题进行了辩论；他们承诺，如果他同意他们的意见，就会把他的请求上呈皇帝，并促成他使命的目标。迫于形势的紧迫，或者可能认为持有关于天主性体这样的见解更好，乌尔菲拉便与阿里乌斯派共融，并使自己及其整个民族与天主教会断绝了一切关系。因为他曾教导哥特人宗教的基本道理，并通过使他们分享了一种更温和的生活方式，所以他们对他的指导深信不疑，坚信他不能也不会做或说任何邪恶之事。事实上，他曾多次出示其伟大德性的显著证明。在上述蛮族信奉异教崇拜的时期，他为了保卫信仰曾使自己遭受无数危险。他教他们使用文字，并将\WorkTitle{圣经}翻译成他们自己的语言。正因如此，依斯特河岸的蛮族追随了\Person{阿利乌}的教义。同一时期，弗里提格恩的许多臣民为基督作证，殉道而死。阿塔纳里克怨恨他的臣民在乌尔菲拉的劝说下成为基督徒；因为他们放弃了祖传的敬拜，他对许多人施以种种刑罚；有些人被拖到法庭前，高贵地宣认了信仰后被他处死，还有一些人被杀害，却不容他们说一句话为自己辩护。据说，阿塔纳里克任命执行其残酷命令的官员们让人建造了一尊雕像，放在一辆战车上，运到那些被怀疑已经接受基督教的帐篷前，命令他们敬拜雕像并献祭；如果他们拒绝，就连人带帐篷一起烧掉。但我还听说，这一时期发生了一起更加残暴的暴行。许多人拒绝服从那些强迫他们献祭的人。其中有男有女；妇女中有的牵着小孩子，有的正在哺乳新生婴儿；他们逃到他们的教堂，那是一个帐篷。异教徒放火烧了帐篷，所有人都被烧死了。\par

哥特人不久便彼此讲和；接着在失去理性的激动中，他们开始蹂躏\Place{色雷斯}，洗劫城镇和村庄。\Person{瓦伦斯}经查问后，从经验中认识到自己犯了多大的错误；因为他原以为哥特人将始终有利于帝国并令敌人畏惧，因此疏忽了对罗马军队的增援。他收取了罗马治下城镇和村庄的金子，以代替通常的兵员补充。期望落空后，他离开\Place{安提约基雅}，匆忙赶往\Place{君士坦丁堡}。因此，他先前对信仰不同意见的基督徒所进行的迫害暂告停止。阿里乌派的首领\Person{欧佐约}去世，\Person{多罗泰}被提议接替其治理。\par

% structure: p0190 source=h2 target=subsection reason=relative-heading-rank; base=h2

\subsection{第三十八章. 关于萨拉森人的酋长\Person{玛尼亚}。当与罗马人的和约解除时，由基督徒祝圣的他们的主教\Person{摩西}重新订立了和约。关于\Person{依市玛耳人}和萨拉森人及其货物的叙述；以及他们如何通过他们的酋长\Person{佐科莫斯}开始归化于基督教。}

大约此时，萨拉森人的国王去世了，先前存在于该民族与罗马人之间的和平被解除。已故君主的遗孀\Person{玛尼亚}，在执掌了本族统治权后，率领她的军队进入\Place{腓尼基}和\Place{巴勒斯坦}，远至\Place{埃及}那些位于驶向\Place{尼罗河}源头者左侧的地区，这些地区通常被称为\Place{阿拉伯}。这场战争绝不可轻视，尽管是由一位妇女指挥。据说罗马人认为它如此艰巨和危险，以至于腓尼基军队的将军向整个东部地区骑兵和步兵的将军请求援助。后者嘲笑这个请求，并独自承担战斗。于是他进攻了亲自指挥自己军队的\Person{玛尼亚}；他好不容易才被巴勒斯坦和腓尼基军队的将军救出。这位将军察觉到极端的危险，认为没有必要遵守他所收到的远离战斗的命令；因此他冲向蛮族，为他的上级提供了安全撤退的机会，而他本人则边撤退边向逃跑者射箭，并用箭击退了紧逼的敌人。这件事仍在当地人民中流传，并被萨拉森人编成歌曲传唱。\par

由于战争仍在激烈进行，罗马人认为有必要派使团向\Person{玛尼亚}求和。据说她拒绝同意使团的请求，除非同意祝圣一个名叫\Person{摩西}的人（他在邻近的沙漠中修行）作她臣民的主教。这\Person{摩西}是一个德行高尚的人，以施行神圣奇迹而闻名。当这些条件禀告给皇帝时，军队的首领奉命捉拿\Person{摩西}，把他带到\Person{路济伍斯}面前。这位隐修士在统治者与会众面前大声说：\Quote{我不配承受主教名号和尊位的荣誉；但尽管我不配，若天主命定我担任此职，我以创造天地的上主为证，我绝不受\Person{路济伍斯}的按手礼，那双手被圣人们的血玷污了。}\Person{路济伍斯}立即回答说：\Quote{你若不认识我信经的性质，你在了解案件全部真相之前就判断我，这是错误的。如果你因流传的诽谤对我抱有偏见，至少容我向你表明我的观点；请你评判它们。}\Person{摩西}回答说：\Quote{你的信经我早已熟知；其性质由在流放中和矿场中受苦的诸多主教、司铎和执事所见证。显然，你的观点与基督的信仰以及有关天主性的所有正统教义相悖。}他再次发誓抗议，绝不接受他们的祝圣，便去了萨拉森人那里。他使他们与罗马人言和，并使许多人皈依了基督教，在他们中间作为司铎度过一生，尽管他找到的与他共信的人很少。\par

这就是那支从亚巴郎之子依市玛耳起源并得名的部落；古人根据其始祖称他们为依市玛耳人。由于他们的母亲哈加尔是女奴，后来为了掩盖其起源的耻辱，他们自称为撒拉逊人，仿佛他们是亚巴郎的妻子撒辣的后裔。如此起源的他们，如同犹太人一样行割损礼，戒食猪肉，并遵守许多其他犹太人的仪式和习俗。倘若他们在某些方面偏离了该民族的惯例，那必是因时间的流逝及与邻近民族的交往所致。梅瑟生活在亚巴郎之后许多世纪，只为他从埃及领出的那些人立法。邻近的居民热衷于迷信，很可能很快就败坏了他们的祖先依市玛耳所传下的律法。古代的希伯来人在梅瑟立法之前仅凭此法度集体生活，因此使用不成文的习俗。这些人无疑与邻近民族敬拜相同的神，以相似的方式尊崇并称呼它们，以致在宗教上与祖先的这种相似性，显明了他们背离了祖先的律法。一如常情，随着时间流逝，他们古老的习俗被遗忘，其他做法逐渐在他们中占据上风。后来，该部落的一些人偶然与犹太人接触，从他们那里得知自己真正的起源，便回到同族人中，倾向于希伯来的习俗和律法。从那时起直到如今，他们中有许多人按犹太人的规条生活。就在本朝不久之前，一些撒拉逊人皈依了基督信仰。他们通过与居住在附近、在邻近旷野中实践哲学、并以卓越的生活和奇迹般的工作著称的司铎和隐修士的交往，分享了基督的信德。据说大约在此时，整个部落及其首领佐科穆斯在以下情形下皈依了基督信仰并受了洗：佐科穆斯没有子女，便去见一位极负盛名的隐修士，向他倾诉这一不幸；因为在撒拉逊人中间，而且我相信在其他蛮族中也一样，有子女被视为极其重要。那隐修士让佐科穆斯放心，为他祈祷，并给他一个应许后送他离去：如果他相信基督，就会得到一个儿子。当天主证实了这一应许，且一个儿子为他降生时，佐科穆斯和他所有的臣民都受了洗。从那时起，这一部落格外幸运，人数壮大，对波斯人和其他撒拉逊人都变得可怕。关于撒拉逊人的皈依及其首位主教，这就是我所收集到的详情。\par

% structure: p0194 source=h2 target=subsection reason=relative-heading-rank; base=h2

\subsection{第三十九章 \Person{伯多禄}从\Place{罗马}返回，在\Person{路基乌}让步后重获\Place{埃及}教会；\Person{瓦伦斯}向西远征对抗斯基泰人。}

各城持守尼西亚教义的人如今开始鼓起勇气，尤其是在\Place{埃及}的\Place{亚历山大}城的居民。\Person{伯多禄}从\Place{罗马}带着\Person{达玛苏}的信函返回，信函确认了尼西亚的信条和他自己的祝圣；他取代了\Person{路基乌}掌管教会，\Person{路基乌}在被逐后乘船去了\Place{君士坦丁堡}。\Person{瓦伦斯}皇帝自然被其他事务缠身，无暇顾及这些事。他刚到\Place{君士坦丁堡}，就引起了百姓的猜疑和憎恨。蛮族正在劫掠色雷斯，甚至已逼近城郊，企图在无人阻挡之下攻打城墙。全城对这种冷漠感到愤慨；百姓甚至指控皇帝是他们攻击的同谋，因为他没有出击，而是拖延应战。最后，当他在赛马场观看竞赛时，百姓公开大声指责他忽视国事，并要求提供武器以便自卫。\Person{瓦伦斯}被这些指责所触怒，立即起兵征讨蛮族；但他威胁说，回来后将惩罚百姓的狂妄，并要向他们报复，因为他们先前曾支持篡位者\Person{普罗科皮乌斯}。\par

% structure: p0196 source=h2 target=subsection reason=relative-heading-rank; base=h2

\subsection{第四十章 \Person{圣依撒格}隐修士预言\Person{瓦伦斯}之死。\Person{瓦伦斯}逃入草棚被烧，因而丧命。}

当瓦伦斯即将离开君士坦丁堡时，一位品德高尚、为天主的事业不惧危险的隐修士依撒格出现在他面前，对他说了以下的话：\Quote{皇帝啊，请将你剥夺的正统派和拥护尼西亚信理者的教会归还他们，胜利就必属于你。}皇帝被这种大胆行为触怒，下令逮捕依撒格，将他囚禁起来，直到他回来，那时他打算因他的鲁莽而审判他。依撒格，然而，回答说：\Quote{你若不放还教会，你必不能返回。}事实果然如此。因为当瓦伦斯率军出征时，哥特人边战边退。他行军经过色雷斯，来到哈德良堡。当离蛮族不远时，发现他们驻扎在安全的位置；然而他竟鲁莽地在未将自己的军团布置妥当之前就发动进攻。他的骑兵被击溃，步兵被迫撤退；被敌人追击时，他跳下马，带着几个随从进入一座小屋或塔楼，藏身其中。蛮族全力追击，越过了塔楼，未怀疑他选择那里作为藏身之地。当蛮族的最后一支队伍经过塔楼时，皇帝的随从从隐蔽处射出一阵箭雨，这立刻导致有人呼喊说瓦伦斯藏在塔楼内。稍前一点的人听到这呼喊，便用喊声将消息告诉前面的人；就这样，消息一直传到最前面的追兵。他们返回并包围了塔楼。他们从周围乡下收集大量木柴，堆在塔楼旁，最后点燃了木柴堆。恰巧刮起的一阵风助长了火势；不久，塔楼及其所有一切，包括皇帝和他的随从，完全被烧毁。瓦伦斯享年五十岁。他与兄弟共同统治十三年，又独自统治三年。\par

\SourceRecord{Translated by Chester D. Hartranft.；Nicene and Post-Nicene Fathers, Second Series；Vol. 2.；Edited by Philip Schaff and Henry Wace.；Buffalo, NY: Christian Literature Publishing Co.,；1890.；<http://www.newadvent.org/fathers/26026.htm>.}

% structure: p0001 source=h1 target=section reason=page-title

\section{\WorkTitle{教会史（第七卷）}}

% structure: p0002 source=h2 target=subsection reason=relative-heading-rank; base=h2

\subsection{第一章 当罗马人受蛮族压迫时，玛维亚派遣援军，部分民众取得胜利。格拉提安命令各人按自己意愿信仰。}

瓦楞斯的下场如此。蛮族因胜利而气焰高涨，蹂躏了色雷斯，并进逼君士坦丁堡城门。在此危急关头，玛维亚派来的少数撒拉逊盟军，连同许多民众，发挥了巨大作用。据传，瓦楞斯的妻子多米尼加从国库中拨出款项，部分民众仓促武装后，攻击蛮族，将其逐出城外。\par

此时与弟弟共治整个罗马帝国的格拉提安，不赞成先前为遏制宗教信条分歧而进行的迫害，召回了所有因信仰被放逐者。他还颁布了一项法律，规定每个人可自由奉行自己的宗教，并允许举行集会，但摩尼教徒以及福提努斯和欧诺米乌斯的追随者除外。\par

% structure: p0005 source=h2 target=subsection reason=relative-heading-rank; base=h2

\subsection{第二章 格拉提安选立西班牙人狄奥多西与他共治，亚略主义盛行于东方各教会，唯耶路撒冷教会除外。安提约基亚公会议。各教会首席职位之安排。}

格拉提安考虑到，既要遏制伊斯特河一带伊利里亚和色雷斯的蛮族入侵，又必须亲自前往高卢抵御阿勒曼尼人的进犯，于是在西尔米乌姆将狄奥多西立为共治皇帝。狄奥多西出身于伊比利亚比利牛斯山区的显赫家族，在战事中声誉卓著，因此在他登基之前，人们已普遍认为他足以执掌帝国大权。\par

此时东方各教会，除耶路撒冷教会外，尽在亚略派手中。马其顿派的见解与持守尼西亚教义者相差无几，在各城与后者交往共融，尤其在君士坦丁堡的马其顿派自与利贝里乌斯和好以来更是如此。但格拉提安法令颁布后，马其顿异端的一些主教壮起胆子，重新占据了被瓦楞斯驱逐时所失去的教堂。他们在卡里亚的安提约基亚集会，声称圣子不应被宣称为与圣父\Quote{同质}，而只是本质相似。自那时起，许多马其顿派脱离其他派别，自设教堂；另一些人则谴责作出这些决议者的对立和争执，更坚定地与持守尼西亚教义者联合。\par

许多曾被瓦楞斯放逐、此时因格拉提安法令而被召回的主教，并无意恢复教会最高职位；他们更看重民众的合一，因此恳请亚略派主教保留现有职位，不要因纷争撕裂那由天主和宗徒传承为合一的教会，而纷争与争权已使之四分五裂。本都阿马西亚主教欧拉利乌斯便是持此态度者之一。据说他自放逐归来时，发现自己的教会由一位亚略派主教掌管，城中几乎只有五十人服从这位新主教。欧拉利乌斯将合一置于一切之上，主动提出与那位亚略主教共同管理教会，并明确同意让其居先。但因亚略主教不允此议，不久后，原本跟随他的少数人便离他而去，投靠了另一方。\par

% structure: p0009 source=h2 target=subsection reason=relative-heading-rank; base=h2

\subsection{第三章 论圣默肋提乌斯与安提约基亚主教保利努斯。二人关于主教之位的誓言。}

因着这项法律，默肋提乌斯大约此时回到叙利亚的安提约基亚，他的出现引起了民众的巨大纷争。保利努斯仍健在，瓦楞斯因敬重其虔诚，未敢放逐他。默肋提乌斯的支持者提议让默肋提乌斯与保利努斯共治，但保利努斯谴责默肋提乌斯的祝圣，因其由亚略派主教施行；然而默肋提乌斯支持者人数不少，强行推进其计划，将默肋提乌斯安置在城外某教堂的主教座上。双方仇恨加剧，几乎酿成骚乱，幸而一项促进和睦的非凡计划占了上风。他们决定从被认为有资格、或有望继承该地主教职位的人中取誓。除弗拉维安外，另有五人立誓：在保利努斯和默肋提乌斯在世期间，若他们中间有人被祝圣，他们既不争取也不接受主教职；若两位伟人中有一位去世，另一人独自主教区。他们以誓言批准此承诺后，几乎全体民众恢复了和睦；只有少数路济弗尔派仍持异议，因默肋提乌斯由异端者祝圣。此争端结束后，默肋提乌斯前往君士坦丁堡，那里已有许多主教聚集，商议是否应将国瑞从纳西盎主教之职调任至该城主教位。\par

% structure: p0011 source=h2 target=subsection reason=relative-heading-rank; base=h2

\subsection{第四章 伟大狄奥多西的统治；他由得撒洛尼主教阿斯科利乌斯授以神圣洗礼。他致函不持守尼西亚公会议定义者。}

当此时高卢正受阿拉曼人侵扰，\Person{格拉提安}返回其父辈的领地——他为自己和弟弟保留了该领地——而将伊里利亚和东方各省的治理交给\Person{德奥多西}。他对蛮族达成了目的；\Person{德奥多西}对伊斯特河两岸的部落也同样成功：他击败了他们，迫使他们求和，并在接受人质后前往\Place{得撒洛尼}。他在该城患病，在接受主教\Person{阿斯科利乌}的训诲后领受入门圣事，不久便康复。\Person{德奥多西}的父母是基督徒，信奉尼西亚信理；他对同样持守此信理、且司铎一切德行兼备的\Person{阿斯科利乌}感到满意。他也欣喜发现伊里利亚未曾参与亚略异端。他询问其他省份盛行的宗教信仰，得知直至\Place{马其顿}，所有教会心思一致，所有教会都认为应当对天主圣言和圣神与天主圣父同等敬拜；但东方，尤其是\Place{君士坦丁堡}，人民分裂为许多不同的异端。\Person{德奥多西}思忖：最好向臣民陈明自己的宗教观点，以免看来是以强迫命令不愿服从者违背自己的判断去崇拜。于是他在\Place{得撒洛尼}颁布一道法律，并使之在\Place{君士坦丁堡}公布；他深知，敕令一旦从该城——如同整个帝国的堡垒——发出，便会迅速传至所有其他城市。通过此法律，他表明了自己的意图：引导所有臣民接受那从起初由宗徒之长\Person{伯多禄}向罗马人所传、并由\Place{罗马}主教\Person{达玛苏斯}和\Place{亚历山大里亚}主教\Person{伯多禄}所宣认的信仰。他颁布法令：只有那些对天主圣三的三位同等敬拜的人，才能被单独称为\Quote{天主教会}，而那些持相反见解的人应被视作异端，遭轻蔑并被交付处罚。\par

% structure: p0013 source=h2 target=subsection reason=relative-heading-rank; base=h2

\subsection{第五章。神学家\Person{格列高利}从\Person{德奥多西}接受教会治理。\Person{德莫斐鲁斯}及所有否认圣子与圣父\OriginalTerm{同性同体}{Consubstantial}者被驱逐。}

此法令颁布后不久，\Person{德奥多西}前往\Place{君士坦丁堡}。亚略派在\Person{德莫斐鲁斯}领导下仍占据教堂。纳齐盎的\Person{格列高利}则主持那些持守天主圣三\OriginalTerm{同性同体}{consubstantiality}的人，并将他们聚集在一间改为祈祷所的小屋内——这屋是由持同样见解和同样礼仪的人改建的。该地后来成为城中显眼的教堂之一，直至今日，不仅因其建筑之美和众多，更因天主可见显现带来的恩惠。因为天主的权能在此显现，无论在醒着的异象或梦境中都大有帮助，常为许多疾病和遭遇突然变故者带来缓解。这权能归于天主之母、童贞\Person{玛利亚}，因为她确实以这种方式显现自己。这座教堂被称为阿纳斯塔西亚，我相信，这是因为在君士坦丁堡已废置、可谓被异端势力埋葬的尼西亚信理，通过格列高利的讲论从死中复活并被重新唤醒；或者，如我所听说的，有人确证有一天，当人们在这座建筑里聚集敬拜时，一位怀孕的妇女从最高的楼座跌落，当场被发现死亡；但全团体祈祷后，她复活了，她和婴儿都得了救。由于这一神圣奇迹的发生，据一些人断言，该地得到了这个名字。\par

皇帝派人命令\Person{德莫斐鲁斯}服从\Place{尼西亚}的信理，引导人民接受同样观点，否则就离开教堂。\Person{德莫斐鲁斯}召集众人，告知他们皇帝的谕令，并通知他们他打算第二天在城外举行聚会，他说，这是根据神律，那律命令我们当在一城被迫害时，\BibleQuote{逃往另一座}。\BibleRef{玛 10:23}从那天起，他总与\Person{路济乌斯}在城外举行聚会——\Person{路济乌斯}原为亚历山大里亚的亚略派主教，如前所述，被逐出该城后逃往君士坦丁堡并定居于此。当\Person{德莫斐鲁斯}及其追随者离开教堂后，皇帝进入其中并祈祷；从那时起，那些持守天主圣三同性同体的人占据了祈祷所。这些事发生在\Person{格拉提安}执政第五年、\Person{德奥多西}执政第一年，并在教会落入亚略派手中四十年之后。\par

% structure: p0016 source=h2 target=subsection reason=relative-heading-rank; base=h2

\subsection{第六章。论亚略派；以及\Person{欧诺米乌斯}的成功。\Person{圣盎斐洛基乌斯}对皇帝的胆量。}

亚略派在人数上仍然十分强大，并且由于曾得到\Person{君士坦提乌斯}和\Person{瓦伦斯}的保护，仍然无所畏惧地集会，公开谈论天主和天主性。现在他们决定通过其教派中在朝廷任职的人，试图争取皇帝加入他们的一方；他们期望在这个计划上能够成功，就像在\Person{君士坦提乌斯}身上成功那样。这些阴谋在天主教会成员中引起了极大的焦虑和恐惧；但他们最害怕的是\Person{欧诺弥}的辩才。看来，在\Person{瓦伦斯}统治时期，\Person{欧诺弥}曾在\Place{齐济库斯}与自己的神职人员发生争执，因而脱离了亚略派，退居到\Place{君士坦丁堡}附近的\Place{比提尼亚}。许多人聚集到他那里；也有一些人从各地而来，少数人是为了测试他的原则，另一些人仅仅是为了听他讲论。他的名声传到了皇帝耳中，皇帝很愿意与他交谈。但皇后\Person{弗拉基拉}有意阻止了他们之间的会面；因为她是尼西亚信理最忠实的捍卫者，她担心\Person{欧诺弥}会凭其辩才改变皇帝的观点。\par

与此同时，当双方各自进行这些阴谋时，据说当时住在\Place{君士坦丁堡}的主教们前往皇帝处，向他作惯常的问候。有一位来自一个小城的老司铎，他朴实而淡泊世俗，但在神圣事物上却很有学识，他也在这一行人中。其他人都是脱帽恭恭敬敬地向皇帝致敬。这位老司铎也以同样的形式问候了皇帝；但是，他没有向坐在父亲身旁的王子给予同等的尊敬，而是走近他，亲切地拍了拍他，称他为亲爱的孩子。皇帝因其子所受的侮辱而愤怒，因为老司铎没有同样尊敬王子，便命令粗暴地将老人推出去。然而，在被推来推去时，老司铎转过身来大声说道：\Quote{请思考，皇帝啊，天上的父对那些不尊敬祂的圣子如同尊敬祂自己、并且胆敢断言圣子低于圣父的人所发的义怒。}皇帝感到这话的分量，召回司铎，为所发生的事向他道歉，并承认他说的是真理。从此，皇帝更不愿意与异端者交往，并禁止在市场中进行辩论和集会。他颁布法律，明确对此类关于天主本质和性体的讨论的惩罚，使这类讨论变得危险。\par

% structure: p0019 source=h2 target=subsection reason=relative-heading-rank; base=h2

\subsection{第七章 关于第二次神圣大公会议及其召开的地点和原因。神学家格列高利的退位。}

皇帝不久后召开了一次正统主教会议，旨在确认尼西亚法令，并选举一位主教担任\Place{君士坦丁堡}的空缺主教座。他也邀请了马其顿派参加这次会议；因为他们的教义与天主教会相去不远，他判断容易与他们重新合一。大约一百五十位坚持天主圣三同体的主教出席了会议，另外有三十六位马其顿派主教，主要来自\Place{赫勒斯滂}的城市；其中主要的有\Place{齐济库斯}主教\Person{厄肋乌西乌}、\Place{拉姆普萨库斯}主教\Person{玛尔西安}。另一方由\Person{提摩太}带领，他接替了其兄\Person{伯多禄}担任\Place{亚历山大里亚}主教座；还有\Place{安提约基亚}主教\Person{默利提乌}，他因\Person{格列高利}的选举而于不久前来到了\Place{君士坦丁堡}；以及\Place{耶路撒冷}主教\Person{济利禄}，他此时已放弃了先前所持的马其顿派主张。\Place{得撒洛尼}主教\Person{阿斯科利乌}、\Place{塔尔索}主教\Person{狄奥多若}、\Place{贝雷亚}主教\Person{阿卡基乌}也出席了会议。后者一致坚持尼西亚法令，并敦促\Person{厄肋乌西乌}及其同党遵从这些意见，同时提醒他们以前派往\Person{利贝里乌}的使团，以及他们通过\Person{欧斯塔修斯}、\Person{西尔瓦诺斯}和\Person{德奥斐洛}向\Person{利贝里乌}所作的宣认，如前面所述。然而，马其顿派公开声明，无论他们以前向\Person{利贝里乌}作过什么宣认，他们绝不会承认圣子与圣父同体，于是立即退席。然后他们写信给各地依附他们的人，劝他们不要顺从尼西亚教义。\par

当时留守\Place{君士坦丁堡}的主教们，转而关注为该城选举一位主教。据说，皇帝因深敬佩\Person{额我略}的圣德与口才，认为他堪当此主教之职，且因敬畏其德行，多数主教会议的成员也持相同意见。\Person{额我略}起初同意接受\Place{君士坦丁堡}教会的主席职务；但后来得知有些主教，特别是\Place{埃及}的主教们反对这一选举，便撤回了同意。在我看来，这位最智慧之人不仅因各方面的才华值得钦佩，尤其因他在当前情况下的行为而值得赞赏。他的口才没有使他骄傲，虚荣也没有驱使他渴望执掌一个教会——他曾在教会脱险之后才接受它。当主教们要求他时，他交出了任命，从未抱怨自己付出的许多辛劳，或是在镇压异端时遭遇的危险。倘若他保留\Place{君士坦丁堡}的主教职，也不会对任何人造成损害，因为\Place{纳西盎}已另立了一位主教代理其位。然而，主教会议严格遵从教父的法律与教会秩序，在他本人的默许下，收回了原先托付给他的这份职守，甚至没有为这样一位杰出的人物开例。因此，皇帝与司铎们转而选举另一位主教，认为这是当时最需要关注的事；皇帝迫切要求进行审慎调查，以便将最优秀、最卓越之人托付给这座伟大皇家之城的大司祭之职。然而，主教会议意见分歧，因为每位成员都希望看到自己的一位朋友被祝圣为教会首领。\par

% structure: p0022 source=h2 target=subsection reason=relative-heading-rank; base=h2

\subsection{第八章 \Person{纳克塔利}当选\Place{君士坦丁堡}主教；其出生地与教育。}

\Emph{有一位}\Place{塔尔索}人，属\Place{基里基雅}，出身显赫的元老阶层，当时正居住在\Place{君士坦丁堡}。他即将返回本国，便拜访\Place{塔尔索}主教\Person{狄奥多若}，询问是否有信件需他转交。\Person{狄奥多若}正专心于当时众人关注的祝圣事宜。他一见到\Person{纳克塔利}，便认为此人堪当主教之职，并立即在心中定意，思及此人的可敬年岁、适合于司铎的仪态以及温和的举止。他仿佛为其他事一般，将\Person{纳克塔利}带到\Place{安提约基雅}主教那里，请求其运用影响力促成这一选举。\Place{安提约基雅}主教对此请求不以为然，因为已有最杰出之人的名字被提出来考虑。但他仍召来\Person{纳克塔利}，请他暂时留下。过了一段时日，皇帝命司铎们列出他们认为堪当祝圣之人的名单，保留自己从中挑选一人的权利。所有主教都遵命照办；\Place{安提约基雅}主教也写下他举荐候选主教的名字，并在名单末尾，出于对\Person{狄奥多若}的考量，加上了\Person{纳克塔利}之名。皇帝阅读所录名单，目光停在文书末尾的\Person{纳克塔利}之名上，用手指按住，沉思片刻；又从头看起，再次检视整份名单，最终选择了\Person{纳克塔利}。这一任命引起极大惊讶，所有人都急于查知\Person{纳克塔利}是谁、他的生活方式及出生地。当听说他尚未领受入门圣事时，人们对皇帝的决定更加惊异。我相信，\Person{狄奥多若}自己也不知道\Person{纳克塔利}未受洗；因为若他知晓此事，必不敢投票支持一位未入门者领受司祭职。合理推测是，当他见\Person{纳克塔利}已年长，便想当然认为他早已领洗。但这些事并非没有天主的干预。因为当皇帝得知\Person{纳克塔利}未受洗时，尽管许多司铎反对，他仍坚持己见。最后，当皇帝的命令得到同意后，\Person{纳克塔利}领受了入门圣事，并且仍穿着入门白衣时，便由主教会议一致宣告为\Place{君士坦丁堡}主教。许多人猜测，皇帝作此选择是出于天主的启示。我无法判断此猜测是真是假；但我深信，反思这场祝圣的非凡经历时，这些事绝非没有天主的德能；是天主指引这位温良、有德、卓越之人进入了司祭职。以上是我所能查知的关于\Person{纳克塔利}祝圣的细节。\par

% structure: p0024 source=h2 target=subsection reason=relative-heading-rank; base=h2

\subsection{第九章 第二次大公会议的法令。犬儒哲学家\Person{玛克西穆}。}

在这些事之后，内克塔里乌斯（Nectarius）与其他司铎聚集一起，下令规定：尼西亚大公会议所确立的信仰应居主导地位，所有异端都应受谴责；各地教会应按照古代规条管理；每位主教应留守自己的教会，不得以任何轻率借口前往他处，或未经邀请而施行无权干涉的圣秩授予，正如在迫害时期天主教会中常见的情况那样。他们同样规定，各教会的事务应受省议会的审查与管控；且君士坦丁堡主教在地位顺序上应仅次于罗马主教，因其占据新罗马的宗座；因为君士坦丁堡不仅已享有此称谓，还拥有许多特权——例如拥有自己的元老院，公民分为不同等级与阶层；它也由自己的行政官治理，并在契约、法律及豁免权方面与意大利的罗马享有同等地位。\par

该大公会议同样规定：马克西穆斯（Maximus）过去不是、现在也不是主教；他所祝圣的那些人不属于圣职人员；凡由他或以其名义所做的一切均属无效。马克西穆斯原籍亚历山大里亚，职业是犬儒派哲学家。他热诚拥护尼西亚教义，并由从埃及聚集到该城的主教们秘密祝圣为君士坦丁堡主教。\par

这些便是大公会议的法令。它们得到皇帝的确认，皇帝颁布法令规定：尼西亚所确立的信仰应居主导地位，各地教会应置于那些承认在三个同等尊荣、同等权能的\OriginalTerm{位格}{hypostasis}——即圣父、圣子及圣神——中具有同一神性的人手中。为更精确地指明这些人，皇帝宣布他所指的是那些与君士坦丁堡的内克塔里乌斯、埃及亚历山大里亚主教提摩泰（Timothy）共融的人；在东方的教会中，与塔尔苏斯主教狄奥多鲁斯（Diodorus）共融；在叙利亚，与劳迪塞亚主教佩拉吉乌斯（Pelagius）共融；在亚细亚，与以哥念教会首领安菲洛基乌斯（Amphilochius）共融；在本都（从比提尼亚到亚美尼亚）的城市中，与卡帕多西亚凯撒利亚教会主教赫拉迪乌斯（Helladius）、尼撒主教额我略（Gregory）、梅利泰内主教奥特雷努斯（Otreinus）共融；在色雷斯和西徐亚的城市中，与托弥主教特伦提乌斯（Terentius）、马尔恰诺波利斯主教玛尔提里乌斯（Martyrius）共融。皇帝本人认识所有这些主教，并确认他们贤明而虔诚地管理各自的教会。这些事之后，大公会议解散，每位主教返回家乡。\par

% structure: p0028 source=h2 target=subsection reason=relative-heading-rank; base=h2

\subsection{第十章 论基利家的玛尔提里乌斯。圣保禄（Confessor）及安提约基雅主教梅肋提乌斯遗骸的迁葬。}

内克塔里乌斯在阿达纳主教基里亚库斯（Cyriacus）的指导下熟悉了司祭礼仪的程序；他曾请求塔尔苏斯主教狄奥多鲁斯让基里亚库斯与他共处一段时间。内克塔里乌斯还留了几位基利家人，其中包括他的医生玛尔提里乌斯，此人曾目睹他少年时期的放荡行为。内克塔里乌斯想祝圣他为执事，但玛尔提里乌斯以自己不配从事如此神圣的服事为由拒绝了这项荣誉，并请内克塔里乌斯本人作证他过去的生活。内克塔里乌斯回复如下：\Quote{虽然我现在是司铎，难道你不知道我过去的行径比你更有罪吗？因为你在我无数放荡行为中不过是一个工具。} \Quote{但是，有福的人，}玛尔提里乌斯回答说，\Quote{你已藉着圣洗被洁净，然后被堪当授予司铎职。这两件圣事都是由神的律法规定为除罪之途；在我看来，你现在与初生婴儿毫无区别；而我很久以前就领受了神圣的圣洗，此后却一直继续同样的放荡恶行。} 他以此为由推辞了接受圣秩；我赞赏这个人的拒绝，因此愿在我的历史中给他一个位置。\par

狄奥多西皇帝得知与君士坦丁堡前任主教保禄有关的各项事件后，命人将其遗体移至由其敌人马其顿尼（Macedonius）所建的教堂中安葬。这座圣殿宽敞而极其显赫，至今仍以保禄命名。因此许多不明事实真相的人，尤其是妇女和民众，以为宗徒保禄安葬于此。与此同时，梅肋提乌斯的遗骸被运往安提约基雅，安放在殉道者巴比拉（Babylas）墓旁。据说，皇帝命令沿途每座城市均在城墙内迎接这些遗骸，这违反罗马习俗，并在这些地方以对唱圣咏的方式予以尊崇，直到它们被运至安提约基雅。\par

% structure: p0031 source=h2 target=subsection reason=relative-heading-rank; base=h2

\subsection{第十一章 弗拉维亚努斯被祝圣为安提约基雅主教，以及因誓言而发生的后续事件。}

在隆重安葬\Person{默莱提乌斯}遗体之后，\Person{弗拉维安}被祝圣代替他，而这公然违背了他所发的誓言，因为\Person{保利努斯}仍然在世。这给\Place{安提约基雅}教会带来了新的麻烦。许多人拒绝与\Person{弗拉维安}保持共融，而与\Person{保利努斯}一起保持分离的教会团体。甚至连司铎们之间在这件事上也意见相左。埃及人、阿拉伯人和塞浦路斯人对于向\Person{保利努斯}表现出的不公感到愤慨。另一方面，叙利亚人、巴勒斯坦人、腓尼基人，以及\Place{亚美尼亚}、\Place{卡帕多细亚}、\Place{加拉提亚}和\Place{本都}的大部分，都支持\Person{弗拉维安}。罗马的主教和所有西方的司铎对\Person{弗拉维安}的行为极为不满。他们按照惯例，将称为主教会议信函的书信寄给\Person{保利努斯}作为\Place{安提约基雅}的主教，而对\Person{弗拉维安}不予理会。他们还断绝了与\Place{塔尔索}主教\Person{迪奥多鲁斯}和\Place{贝雷亚}主教\Person{阿卡西乌斯}的共融，因为他们祝圣了\Person{弗拉维安}。为了进一步处理此事，西方的主教们和\Person{格拉提安}皇帝写信给东方的主教们，并召集他们到西方参加主教会议。\par

% structure: p0033 source=h2 target=subsection reason=relative-heading-rank; base=h2

\subsection{第十二章 \Person{狄奥多西}统一所有异端的计划。\Person{阿格利乌斯}和\Person{西西尼乌斯}（诺瓦替安派）提出的提议。在另一次主教会议上，皇帝只接待了那些代表同质论的人；持不同观点的人，他将他们逐出教会。}

虽然在此时期，所有祈祷之所都归天主教会所有，但在帝国各地发生了许多由亚略异端者煽动的动乱。因此，戴奥陶皇帝在上述大公会议后不久，再次召集当时兴盛的各个派别的领袖，希望他们要么使他人对争议话题达到自己的确信状态，要么自己被说服；因为他设想，如果就教义中含糊之点展开自由讨论，所有人都能达成一致意见。于是，会议被召集。此事发生在梅罗保德第二次担任执政官、萨图尼努斯第一次担任执政官的那一年，同时阿尔卡狄乌与其父共同治理帝国。戴奥陶派人召见内克塔里，与他商议即将举行的主教会议，并命他引发所有导致异端的问题的讨论，以便基督信徒的教会合而为一，并在应当遵守虔敬的教义上达成一致。内克塔里回到家中，对所托之事感到忧虑，便将皇帝的命令告知诺瓦派的教会领袖阿格利乌，此人持有与他相同的宗教信仰。阿格利乌以行为证明了他德行的卓越，但不习惯言语的诡计和欺骗；因此他提议由他的一位读经者——名叫西西尼乌——代替他，此人后来继任为主教，既能看清实际，又能必要时辩论。西西尼乌拥有智力和表达能力；他精通圣经的解释，并熟悉世俗与教会文学。他建议避免与异端者的一切争论，因其是纷争与战争的丰富源泉；而应当调查异端者是否承认在教会分裂之前生活过的圣经的诠释者与教师们的见证。\Quote{如果他们拒绝这些伟人的见证，}他说，\Quote{他们将被自己的追随者所谴责；但如果他们承认这些人的权威足以解决教义中含糊之点，我们就拿出他们的著作。}因为西西尼乌深知，正如古人承认圣子与圣父一样永恒，他们从未妄称圣子有某个起始的开始。这一建议得到了内克塔里的赞同，后来也得到了皇帝的赞同；于是开始调查异端者关于古代圣经诠释者的意见。当发现异端者声称对这些早期作者极为钦佩时，皇帝公开询问他们是否愿意在争议话题上遵从上述作者的权威，并以那些著作中所提出的见解来检验自己的教义。这一提议在各异端派别的领袖中引起了极大的争论，因为他们对古代著作的看法并不一致；皇帝知道他们仅凭自己的言语辩论就被定罪，于是撤回了这一提议。他责备他们的意见，并命令每一派别写出一份自己信条的书面陈述。在为此类文件的提交指定的那天，内克塔里和阿格利乌作为维护天主圣三同一实体的人的代表出现在皇宫；亚略派领袖德莫斐罗作为亚略派的代表出现；欧诺弥乌代表欧诺弥派；而基齐库的主教厄肋乌西乌代表称为马其顿派的宗派分子出现。皇帝在收到他们的信纲后，只赞成其中承认天主圣三同一实体的那一份，并销毁了其他信纲。诺瓦派的利益未受此事件的影响，因为他们关于天主本性持有与天主教会相同的教义。其他派别的成员对司铎们在皇帝面前进行不明智的辩论感到愤慨。许多人放弃了他们先前的意见，拥抱了公认的宗教形式。皇帝颁布了一项法律，禁止异端者拥有教堂、公开讲授信仰、以及祝圣主教或其他人。一些异端者被逐出城市和村庄，而另一些人则被剥夺名誉和帝国其他臣民所享有的特权。尽管法律所定的对异端者的惩罚十分严厉，但并不总是被执行，因为皇帝并不想迫害他的臣民；他只想通过威慑来强制达成关于天主的统一观点。那些自愿放弃异端意见的人得到了他的赞扬。\par

% structure: p0035 source=h2 target=subsection reason=relative-heading-rank; base=h2

\subsection{第十三章 暴君马西莫斯。关于尤斯蒂娜皇后与圣盎博罗削之间的事件。格拉先皇帝被诡计杀害。瓦伦提尼安及其母亲逃往得撒洛尼投靠戴奥陶。}

当格拉先皇帝此时正忙于与阿勒曼尼人作战时，马西莫斯离开不列颠，企图篡夺皇权。瓦伦提尼安当时居住在意大利，但由于他尚未成年，国事由近卫长官普罗布斯处理，他此前曾任执政官。\par

皇帝之母尤斯蒂娜信奉亚略异端，迫害米兰主教盎博罗削，并竭力引入对《尼西亚信理》的篡改，试图使在阿里米努姆颁布的信条形式占据主导地位，从而扰乱各教会。她对盎博罗削怀恨在心，因他坚决反对她革新之举，她便向儿子进谗，说这位主教侮辱了她。瓦伦提尼安听信了这诬告，决意为母亲报仇，便派一队士兵前往教堂。他们抵达圣殿后，强行闯入，逮捕了盎博罗削，正要立即将他流放之际，民众聚集在教堂内，决心宁死也不从司铎被逐。此事更激怒了尤斯蒂娜；为依法推行其计划，她召来一名叫梅尼沃卢斯的法律秘书，命他尽快起草一道确认阿里米努姆法令的敕令。梅尼沃卢斯因坚定依附天主教会，拒绝撰写该文件；女皇又许以更高官职来贿赂他。但他仍拒不从命，解下腰带扔在尤斯蒂娜脚前，声明自己宁可不保留现有官职，也不接受晋升以换取不虔敬之报酬。因他坚持拒绝，其他人便受委托编纂此法。此法鼓励所有遵行阿里米努姆颁布并在君士坦丁堡批准的信条之人，大胆集会；并规定凡阻碍或违抗皇帝此法者，处以死刑。\par

当皇帝之母正策划施行这严酷法律时，传来了格拉提安被杀的消息，凶手是马克西穆斯的将军安德拉加提乌斯。安德拉加提乌斯夺取了皇帝的仪仗马车，派人通知皇帝说皇后正前往他的军营。格拉提安新婚不久，年轻且深爱妻子，便轻率地渡河，急于迎见她，却因轻率而落入安德拉加提乌斯之手；他被捉住，不久即被处死。当时他年仅二十四岁，在位十五年。这场灾难平息了尤斯蒂娜对盎博罗削的愤怒。\par

与此同时，马克西穆斯召集了一支由不列颠人、邻近的高卢人、克尔特人及其他民族组成的大军，向意大利进军。他表面上宣称此举是为了阻止古代宗教和教会秩序的变革，实际上却是出于消除对自己觊觎暴政的怀疑。他暗中窥伺并谋划夺取帝国统治权，以便显得他是依法而非凭武力取得罗马政权。瓦伦提尼安迫于时势，承认了他的统治标志；但不久之后，因害怕受害，他与母亲尤斯蒂娜以及意大利禁卫军长官普罗布斯一同逃往得撒洛尼。\par

% structure: p0040 source=h2 target=subsection reason=relative-heading-rank; base=h2

\subsection{第14章 霍诺留之诞生。狄奥多西将阿尔卡迪乌斯留在君士坦丁堡，前往意大利。诺瓦替安派及其他宗主教之继任。亚略人的胆大妄为。狄奥多西消灭僭主后，在罗马举行盛大凯旋。}

当狄奥多西筹备对马克西穆斯的战争时，他的儿子霍诺留出生了。战争准备就绪后，他将儿子阿尔卡迪乌斯留在君士坦丁堡治理，自己前往得撒洛尼，在那里接待了瓦伦提尼安。他既未公开遣返马克西穆斯派来的使团，也未予以接见，而是率领军队继续向意大利进发。\par

大约此时，君士坦丁堡诺瓦替安派的主教阿格利乌斯感到时日无多，任命了其教会的一位司铎西西尼乌斯为继任者。然而，民众因未能优先选择以虔诚著称的马尔西安而感到不满，阿格利乌斯因此祝圣了马尔西安，并对聚集在教堂的群众说了以下的话：\Quote{在我之后，你们将有马尔西安作为主教；在他之后，将是西西尼乌斯。}阿格利乌斯说完这些话后不久便去世了。他治理自己的教会四十年，深得本派异端团体的赞许；有人断言，在异教迫害时期，他曾公开宣认基督之名。\par

不久之后，弟茂德和济利禄去世；德奥斐洛继任亚历山大主教，若望继任耶路撒冷主教。君士坦丁堡亚略派的领袖德莫菲卢斯也去世了，由色雷斯的马里努斯继任；但他很快被多罗特乌斯取代，后者不久从叙利亚的安提约基雅来到，其派别认为他比马里努斯更胜任此职。\par

与此同时，狄奥多西进入意大利，关于他军事成败的消息众说纷纭。亚略人中谣传其军队大部分在战斗中覆灭，他本人也被僭主俘获；这些教派人士信以为真，变得大胆起来，跑到内克塔里乌斯的住所，放火焚烧，因恼恨这位主教在教会中所获得的权力。然而，皇帝在战争中实现了他的目标，因为马克西穆斯的士兵或因惧怕针对他们的备战，或因背叛，便捉住并杀死了僭主。杀害格拉提安的凶手安德拉加提乌斯一听到马克西穆斯死亡的消息，便披甲跳入河中淹死。战争就此结束，格拉提安之仇已报，狄奥多西偕同瓦伦提尼安在罗马举行凯旋，并整顿了意大利各教会，因为女皇尤斯蒂娜已去世。\par

% structure: p0045 source=h2 target=subsection reason=relative-heading-rank; base=h2

\subsection{第15章 安提约基雅主教弗拉维安与埃瓦格里乌斯。亚历山大城在摧毁狄俄尼索斯神庙时发生的事件。塞拉皮雍神庙及其他被摧毁的偶像神庙。}

\Place{安提约基雅}的主教\Person{保利努斯}大约在这个时期去世，那些与他聚集在教会中的人仍然坚持对\Person{弗拉维安}的厌恶，尽管他的宗教信仰与他们完全一致，因为他违背了从前向\Person{梅利提乌斯}所作的誓言。于是，他们选举\Person{埃瓦格里乌斯}为主教。\Person{埃瓦格里乌斯}在受任后不久便去世了，尽管\Person{弗拉维安}阻止了另一位主教的选举，那些与他断绝共融的人仍然继续分开举行集会。\par

大约这个时期，\Place{亚历山大}的主教，皇帝应其请求将狄奥尼索斯神庙赐予他，他将该建筑改建成教堂。雕像被移走，内殿被暴露；为了羞辱异教的神秘仪式，他举行了一场游行展示这些物品；\OriginalTerm{法罗斯}{phalli}以及内殿中隐藏的任何其他真正或看似荒谬的物品，他都公开展示。异教徒们对如此意想不到的暴露感到惊愕，无法忍受沉默，于是合谋攻击基督徒。他们杀死了许多基督徒，打伤了其他人，并占领了\OriginalTerm{塞拉皮斯神庙}{Serapion}，这座神庙以其美丽和巨大而引人注目，坐落在一座高地上。他们将此改造成临时堡垒；并将许多基督徒带到那里，施以酷刑，强迫他们献祭。拒绝服从的人被钉在十字架上，双腿被打断，或以某种残酷的方式处死。当叛乱持续了一段时间后，统治者们前来敦促人民遵守法律，放下武器，交出塞拉皮斯神庙。随后，\Place{埃及}的军事军团将军\Person{罗马努斯}和\Place{亚历山大}的总督\Person{埃瓦格里乌斯}到来。然而，他们试图使人民服从的努力完全徒劳，于是向皇帝报告了所发生的事。那些将自己关在塞拉皮斯神庙里的人准备更激烈的抵抗，因为他们害怕其大胆行为将面临的惩罚，并且一个名叫\Person{奥林皮乌斯}的人，穿着哲学家的服装，发表煽动性言论，进一步煽动他们叛乱，他告诉他们，他们应该为祖先的诸神而死，不可忽视它们。看到他们因偶像雕像被毁而士气低落，他向他们保证，这种情况并不足以让他们放弃宗教；因为雕像是由可朽坏的材料制成的，仅仅是图像，因此会消失；而曾经居住在其中的力量已飞升天堂。用这样的说法，他使群众留在塞拉皮斯神庙中。\par

当皇帝获悉这些事件后，他宣布被杀害的基督徒是有福的，因为他们被接纳享有殉道者的荣誉，并为信仰而受苦。他宽恕了杀害他们的人，希望借此仁慈之举使他们更愿意接受基督教；并下令拆除\Place{亚历山大}那些引发民众暴乱的神庙。据说，当这道皇帝谕旨被公开宣读时，基督徒发出大声欢呼，因为皇帝将所发生之事的责任归咎于异教徒。守卫塞拉皮斯神庙的人民听到这些欢呼声非常恐惧，以至于逃跑了，基督徒立即占领了该地，并一直保留至今。我听说，在此事发生的前一晚，\Person{奥林皮乌斯}在塞拉皮斯神庙中听到有人唱\OriginalTerm{阿肋路亚}{hallelujah}的声音。门都关着，一切都静悄悄的；他看不见任何人，只能听到歌唱者的声音，他立刻明白了这征兆的含义；他悄悄离开塞拉皮斯神庙，乘船前往\Place{意大利}。据说，当神庙被拆除时，发现了一些石头，上面有十字架形状的\OriginalTerm{象形文字}{hieroglyphic}，经学者鉴定，被解释为象征永生。这些字符导致几个异教徒皈依，同样，在同一地点发现的其他铭文也包含了对神庙毁灭的预言。塞拉皮斯神庙就是这样被占领的，不久之后被改建成教堂；并以皇帝\Person{阿尔卡狄乌斯}的名字命名。\par

许多城市中仍有异教徒，他们热心地为他们的神庙辩护；例如，\Place{阿拉伯}的\Place{佩特拉}和\Place{阿雷奥波利斯}的居民；\Place{巴勒斯坦}的\Place{拉菲}和\Place{迦萨}的居民；\Place{腓尼基}的\Place{赫里奥波利斯}的居民；以及\Place{叙利亚}的\Place{阿克西乌斯河}畔\Place{阿帕梅亚}的居民。我听说，最后提到的这座城市的居民经常武装\Place{加里肋亚}人和\Place{黎巴嫩}的农民来保卫他们的神庙；最后，他们甚至胆大妄为到杀害了一位名叫\Person{马尔塞鲁斯}的主教。这位主教曾命令拆除该市和所有村庄的神庙，认为否则他们很难从先前的宗教中皈依。他听说在\Place{阿帕梅亚}的一个地区\Place{奥隆}有一座非常宽敞的神庙，便率领一队士兵和\OriginalTerm{角斗士}{gladiators}前往那里。他站在远离冲突现场、箭矢射程之外的地方；因为他患有痛风，无法战斗、追击或逃跑。当士兵和角斗士正在攻击神庙时，一些异教徒发现他独自一人，便赶到他与战斗分离的地方；他们突然出现，抓住他，并活活烧死了他。犯下此罪的人当时未被发现，但后来被查出，\Person{马尔塞鲁斯}的儿子们决定为他复仇。然而，该省议会禁止他们执行这一计划，并宣布\Person{马尔塞鲁斯}的亲属或朋友不应寻求复仇；他们反而应当感谢天主，因为天主认为他配得上为如此事业而死。\par

% structure: p0050 source=h2 target=subsection reason=relative-heading-rank; base=h2

\subsection{第十六章：被任命主管施行补赎的司铎之职以何种方式、因何原因被废除。论施行补赎的方式。}

大约在这时期，\Person{纳克塔留斯}废除了主管施行补赎的司铎之职；这是教会内首次废止此职。各地区的诸位主教纷纷效法。关于此职的性质、起源及废除的原因，有各种不同的说法。我将陈述自己对此事的看法。无罪性是天主的属性，不属于人的本性；因此，天主决定即使在多次过犯之后，仍应宽免悔改者。在祈求宽恕时，必须告明己罪，看来司铎们从一开始就觉得在全体会众面前公开告明是令人厌烦的。于是他们任命一位至为圣洁、无疑最具明智的司铎来处理这些场合；忏悔者到他那里告明己过；他的职责是指明与每种罪过相称的补赎方式，然后在补赎完成后宣告赦免。由于补赎的习惯在诺瓦替安派中从未扎根，这类规定对他们自然是不必要的；但此习惯在所有其他异端者中盛行，甚至延续至今。西方教会，尤其是\Place{罗马}，极严格地遵守此习惯，那里有一处专为忏悔者设立的场所，他们在该处站立哀哭直到礼仪结束，因为他们不得参与奥迹；然后他们呻吟哀叹，俯伏在地。主教主持仪式，流泪并同样俯伏；全体会众也流泪并大声叹息。之后主教首先从地上升起，并扶起其他人；他为忏悔者献上祈祷，然后遣散他们。每位忏悔者私下自愿承受苦行，或通过禁食、或戒浴、或戒食多种肉类，或通过其他规定方式，直到主教指定的某个时期。届时，他便可从罪的后果中得释放，与民众一同在教会集会。罗马的司铎们从古至今一直谨慎遵守此习惯。在\Place{君士坦丁堡}的教会中，总是有一位司铎主管忏悔者，直到某位贵妇作证说，她作为忏悔者到司铎那里去禁食并向天主祈求，为此目的留在教堂内时，竟被一位执事强奸。此事传开后，民众极为不满，因为这将使教会蒙羞；司铎们尤其因此大感不安。\Person{纳克塔留斯}在犹豫许久不知该如何处理之后，罢免了那位执事；并听从某些人的建议，他们主张在参与神圣奥迹之前必须让每个人自我省察，于是他便废除了主管补赎的司铎之职。从那时起，补赎的实行便废弃了；在我看来，古老的严厉与严格就这样被极端的宽纵原则所取代。在古代制度下，我认为过犯更为罕见；因为人们因惧怕告明己过并将之暴露于严厉审判者的审查之下，而被阻止犯罪。我相信出于类似的考虑，皇帝\Person{德奥多西}——他始终热切促进教会的荣耀——颁布了一项法律，规定妇女除非已有子女且年逾六十，不得进入圣职，这是遵照宗徒\Person{保禄}的训导。该项法律还规定，剃了头的妇女应从教会中逐出；而接纳此类妇女的主教应从主教职上被罢免。\par

% structure: p0052 source=h2 target=subsection reason=relative-heading-rank; base=h2

\subsection{第十七章：\Person{欧诺米乌斯}被\Person{德奥多西}大帝流放。他的继任者\Person{德奥夫洛尼乌斯}；论\Person{欧提基乌斯}和\Person{多罗特乌斯}及其异端；论被称为撒提里派者；亚略派分裂为不同派别；\Place{君士坦丁堡}的派别更为有限。}

不过，上述这类问题最好交由个人判断去决定。\par

大约在这时期，皇帝判处欧诺米流放。这个异端者定居在君士坦丁堡郊区，经常在私人宅邸举行聚会，并宣读自己的著作。他引诱许多人接受他的观点，以致以他命名的宗派分子变得极为众多。他被流放后不久便去世了，安葬在他的出生地达科拉，这是卡帕多细亚的一个村庄，位于凯撒勒雅境内亚尔格山附近。德奥弗洛尼也是卡帕多细亚人，曾是他的门徒，继续传播他的学说。他通过亚里士多德的著作获得了一点皮毛，编写了一部研究其中三段论的导论，题为\Quote{\WorkTitle{Exercises for the Mind.}}。但我听说，他后来卷入许多无益的辩论，很快就不再局限于他老师的学说。由于他热衷于新奇事物，他试图从圣经所用的措辞中证明：虽然天主预知尚未存在之事，知道已存在之事，并记忆已发生之事，但天主对将来和现在的认知方式并不总是同一，而且祂对过去的认知也会改变。由于这个假设在欧诺米派看来极为荒谬，他们将他逐出教会；他自立为以他命名的\Quote{德奥弗洛尼派}的新宗派首领。不久之后，欧诺米派的欧提古在君士坦丁堡创立了另一个以他命名的宗派。因为有人提出问题：天主圣子是否知道最后的时辰？为解决这问题，人们引用了圣史的话，其中说那日子和时辰只有圣父知道。\BibleRef{玛 24:36}然而欧提古争辩说，这种知识也属于圣子，因为祂从圣父领受了一切。欧诺米派的会长们谴责了这一意见，他便脱离与他们的共融，前往流放地加入欧诺米。一名执事及另外几个人，被从君士坦丁堡派去控告欧提古，并在必要时与他辩论，他们先到达了目的地。欧诺米得知他们此行的目的后，表示赞同欧提古提出的观点；并于欧提古到达后，与他一同祈祷，尽管与任何未携带由神圣字符写成的、证明其共融身份的书信的人一同祈祷是不合法的。这次争论后不久欧诺米去世；君士坦丁堡的欧诺米派会长拒绝接受欧提古进入共融；因为他出于嫉妒而敌视欧提古，因为欧提古甚至没有神职身份，而且他无法反驳欧提古的论点，也找不到解决他问题的方法。因此，欧提古将赞同他观点的人分离出来，自成一派异端。许多人断言，他和德奥弗洛尼是最先提出欧诺米派关于神圣圣洗的独特观点的人。以上是我为简洁地提供导致欧诺米派内部分裂的原因所能给予的概述。如果我进一步详细说明，就会过于冗长；而且，这对我而言绝非易事，因为我缺乏这样的辩证法技巧。\par

与此同时，在君士坦丁堡的亚略派中，有人就以下问题争论不休：在圣子存在之前（他们认为圣子是出自虚无的），天主是否可被称为圣父？从安提约基雅被召来代替玛里诺治理他们的多罗特奥认为，在圣子存在之前，天主不能被称为圣父，因为圣父的名与圣子有着必然的联系。而玛里诺则坚持，即使在圣子不存在时，圣父仍是圣父；他提出这一意见，要么是出于信念，要么是出于争辩的欲望，以及对多罗特奥在教会中受偏爱的嫉妒。于是亚略派分裂为两派；多罗特奥及其追随者保留了祈祷所的产业，而玛里诺及其分离出来的人则建造了新建筑，用于举行他们自己的聚会。玛里诺的支持者被称为\Quote{普萨提里安派}和\Quote{哥特派}；普萨提里安，是因为一个\OriginalTerm{糕点贩子}{\GreekText{ψαθυροπώλης}}特奥克提斯图斯是他们观点的热心拥护者；哥特派，则是因为他们的观点得到了该民族的主教塞利努斯的赞同。几乎所有这些蛮族人都遵循塞利努斯的训导，并与玛里诺的追随者一起在教堂聚集。哥特人特别被塞利努斯所吸引，因为他曾是乌尔菲拉的秘书，并继任他为主教。他不仅能用本地语言教导，还能用希腊语在教堂中教导。\par

不久之后，玛里诺与阿加皮乌之间发生了有关优先权的争执，阿加皮乌是玛里诺亲自祝圣为厄弗所亚略派主教的；在随之而来的争吵中，哥特人支持阿加皮乌。据说，该城的许多亚略派神职人员因这两位主教表现的野心而极其愤怒，以致他们与天主教会共融。这就是亚略派分裂为两派的起源——这一分裂至今仍然存在；以致在每座城市，他们都有各自的教堂。然而，君士坦丁堡的亚略派在分裂三十五年后，由前执政官、骑兵与步兵将军普林塔斯（他在朝廷极具影响力）促成了彼此的和解。为防止他们之间旧纷争的复燃，禁止讨论曾导致分裂的问题。这些事情发生在稍后。\par

% structure: p0057 source=h2 target=subsection reason=relative-heading-rank; base=h2

\subsection{第十八章　另一种异端，即萨巴提安派的兴起，由诺瓦提安派所产生。他们在桑加鲁斯举行的会议。关于逾越节盛宴的更详细记载。}

在同一位皇帝统治期间，诺瓦替安派内部因庆祝复活节的问题发生了分裂；这场争论又引发了另一派别，称为萨巴提安派。萨巴提乌斯曾与狄奥克提斯托斯和玛卡里乌斯一起由马尔西安祝圣为司铎。他采纳了那些在瓦伦斯统治时期于帕祖科马召集的同僚司铎们的意见，坚持基督徒应当像犹太人一样庆祝逾越节（复活节）。他起初脱离教会是为了实践更严格的苦修，因为他声称要采取一种非常刻苦的生活方式。他还宣称，他脱离的一个动机是，他看来许多参与奥秘的人不配得此殊荣。然而，当他引入革新的意图被察觉后，马尔西安表示后悔祝圣了他，据说常听见他呼喊，他宁愿把手按在荆棘上也不愿按在萨巴提乌斯的头上。马尔西安意识到他教区的民众正在分裂为两派，便召集了他自己一派的所有主教到桑加鲁斯——这是比提尼亚海边的一个城镇，离埃莱诺波利斯城不远。当他们聚集后，他们召来萨巴提乌斯，请他陈述他抱怨的原因。他只是抱怨关于节日的分歧，他们怀疑他把这当作掩饰其贪爱首席的借口，便让他发誓永不接受主教职位。当他发了所要求的誓之后，所有人都意见一致，他们投票决定要共同维护教会，因为庆典逾越节的分歧绝不应成为分离共融的理由。他们决定每人都可按照自己的判断自由地庆祝该节日。他们为此制定了一条规则，称之为\Quote{\OriginalTerm{无区别}{\GreekText{ἁ} \GreekText{διάφορος}}规则}。这就是桑加鲁斯会议的经过。从那时起，萨巴提乌斯遵守犹太人的惯例；除非所有人恰好同时过节，他则按照惯例提前守斋，并独自按规定守逾越节。他从傍晚到规定时间守夜，念诵规定的祈祷文；次日与众人一同聚会，领受奥秘。这种过节方式起初并未引起人们注意；但随着时间推移，它开始引人注目并更为人知，他找到了许多模仿者，尤其在弗吕家和迦拉达，这些人视此节日庆祝为民族习俗。最终他公开脱离共融，成为那些接受他观点者的主教，我们将在适当的地方叙述。\par

我本人对撒巴提乌斯及其追随者企图引进这种新做法感到惊讶。据尤西比乌斯记载，依据斐洛、若瑟夫、阿里斯托布鲁等人的见证，古希伯来人在春分之后，太阳位于黄道第一宫（希腊人称为白羊宫），月亮处于天空对宫且为其年龄第十四日时，才献上祭品。甚至诺瓦替派中那些对此问题做过较精确研究的人也称，他们异端的创始人及其首批门徒并未遵循这一习俗；此习俗是由那些在帕祖科马集会的人首次引入的；而在古罗马，该教派成员仍与罗马人一样遵守相同做法，罗马人在这方面并未偏离其原始惯例，因为该习俗是由圣宗徒伯多禄和保禄传授给他们的。此外，严格遵行梅瑟法律的撒玛黎雅人，直到初熟之物成熟后才庆祝这一庆节；他们说，在法律中，这被称为\Quote{初熟之物的庆节}，在这些初熟之物出现之前，守节是不合法的；因此，春分必然先于它。由此我惊讶的是，那些自称在庆祝此庆节时采用犹太习俗的人，并未遵循犹太人的古老做法。除了上述民族和亚细亚的十四日派之外，我相信所有异端都以与罗马人和埃及人相同的方式庆祝逾越节。十四日派之所以得名，是因为他们像犹太人一样在月亮的第十四日庆祝此庆节，因此得名。诺瓦替派则庆祝复活之日。他们遵循犹太人和十四日派的习俗，除非月亮第十四日恰好落在一周的第一天，在这种情况下，他们会在犹太人之后的天数等于月亮第十四日与随后的主日之间的天数时庆祝庆节。蒙丹派（又称佩普齐派和弗吕家派）按照他们引入的一种奇怪方式庆祝逾越节。他们责备那些根据月亮运行确定庆节时间的人，并断言应完全遵循太阳的周期。他们规定每个月为三十天，并将春分后的第二天视为一年的第一天，按罗马计算法，这称为四月朔日前第九日。他们说，正是在这一天，创造了为指示节令和年岁而设立的两大光体。他们以每八年太阳和月亮在天上同一地点相遇这一事实来证明这一点。月亮的八年周期在九十九个月和二千九百二十二天内完成；在此期间，太阳运行八周，每周三百六十五天加四分之一天。因为他们认定，圣经中记载的太阳创造之日是月亮第十四日，发生在四月朔日前第九日之后，对应同月的月中日前第八日。他们总是在这一天庆祝逾越节，若这天恰好是复活之日；否则就在随后的主日庆祝；因为他们声称，根据经书，庆节可以在十四日至二十一日之间的任何一天举行。\par

% structure: p0060 source=h2 target=subsection reason=relative-heading-rank; base=h2

\subsection{章十九．一份值得研究的列表，由历史学家给出，关于不同民族和教会的习俗。}

我们现在已经描述了在庆祝逾越节时盛行的各种习惯。在我看来，罗马主教\Person{维克多}和斯米纳主教\Person{颇利加布}就他们之间产生的争议作出了非常明智的决定。因为西方的诸位主教认为没有必要违背由\Person{伯多禄}和\Person{保禄}传给他们的传统，而亚细亚的诸位主教则坚持遵循圣史\Person{若望}所定的规则，他们一致同意继续按照各自习惯遵守此庆节，而不彼此断绝共融。他们忠实而公正地认为，凡在敬拜的要素上一致者，不应因习惯而彼此分离。因为即使各教会持相同见解，在每个点上也都能找到完全相似的传统。例如，在\Place{斯基提亚}有许多城市，却总共只有一位主教；而在其他民族，一位主教甚至充当一个村庄的司祭，正如我本人在\Place{阿拉伯}、\Place{塞浦路斯}，以及在\Place{夫里基雅}的诺瓦替派和蒙丹派中所观察到的那样。再者，现今在\Place{罗马}也只有七位执事，恰好符合宗徒们所规定的数目，其中\Person{斯德望}是首位殉道者；而在其他教会，执事的人数则无关紧要。在\Place{罗马}，阿肋路亚每年只歌唱一次，即在逾越节节期的第一天；因此罗马人常以听到或歌唱这首赞美诗的事实来发誓。在那座城市，人民不受主教教导，也不受教会中任何人教导。在\Place{亚历山大}，只有该城的主教教导人民，据说这习惯自\Person{亚略}时代起就在那里盛行，\Person{亚略}虽仅是一位司铎，却提出了新的教义。在\Place{亚历山大}还有另一个奇怪的习惯，是我从未在其他地方见过或听过的，那就是，当诵读福音时，主教不从座位上起立。在这座城市，只有总执事诵读福音，而在有些地方由执事诵读，在许多教会中仅由司铎诵读；而在特殊日子则由主教诵读，例如在\Place{君士坦丁堡}，在复活节节期的第一天。在一些教会，这庆节前称为四旬期的期间，人民用于守斋，定为六周；\Place{伊利里亚}、西方地区、\Place{利比亚}、全\Place{埃及}和\Place{巴勒斯坦}的情况即是如此；而在\Place{君士坦丁堡}及邻近省份直到\Place{腓尼基}，则定为七周。在一些教会，人们在六或七周期间隔三周守斋，而在其他教会，他们在庆节前的三周连续守斋。有些人，如蒙丹派，只守斋两周。所有教会并非在相同的时间或方式举行集会。\Place{君士坦丁堡}的人民，以及几乎各处的人民，在安息日和一周的第一天都集会，这一习惯在\Place{罗马}或\Place{亚历山大}从未遵守。在\Place{埃及}有几个城市和村庄，与别处确立的习惯相反，人民在安息日晚上集会，并且尽管之前已用过餐，仍领受圣事。在所有教会中，同一场合并不诵念相同的祷文和圣咏，也不诵读相同的读经。因此，题为\Quote{\WorkTitle{伯多禄默示录}}的书，虽被古人完全视为伪作，但在\Place{巴勒斯坦}的一些教会仍在预备日诵读，那一天人民守斋以纪念救主的苦难。同样，题为\Quote{\WorkTitle{宗徒保禄默示录}}的作品，虽不被古人承认，但仍受到多数隐修士的珍视。有些人断言，本书是在当朝，通过天主的启示，在\Place{基里基雅}的\Place{塔尔索}的\Person{保禄}房屋地下埋藏的一个大理石箱子里被发现的。我被告知此说法是假的，消息来自\Person{基利克斯}，他是\Place{塔尔索}教会的司铎，年事甚高，白发可证；他说他们中无人知晓此事，并怀疑是否异端者编造了这故事。我在此主题上所说的话到此已足够。许多其他习惯仍可在城市和乡村见到；那些在这些习惯中长大的人，出于对创建并延续这些习惯的伟人的尊敬，会认为废除它们是不当的。对于在庆祝此庆节中持守不同习惯的人，也应归以相似的动机；正是这庆节引我们进入了这番长篇离题。\par

% structure: p0062 source=h2 target=subsection reason=relative-heading-rank; base=h2

\subsection{第二十章. 我们教义的扩展，以及偶像庙宇的彻底拆毁。尼罗河的泛滥。}

当异端者彼此分裂时，天主教会因着来自异己阵营的许多人加入而日益增长，这既由于他们之间的纷争，尤其由于众多外教人。皇帝观察到，偶像崇拜的习惯因进出庙宇的方便而大大助长，便下令关闭所有庙宇的入口；最终他命令拆毁许多这类建筑。当外教人发现自己被剥夺了他们自己的祈祷之所，便开始频繁出入我们的教会；因为他们不敢秘密地按外教形式献祭，因为这很危险，因为献祭要冒死刑和没收财产的风险。\par

据说埃及的河流今年在适当季节没有泛滥；埃及人愤怒地将这一情况归咎于依照祖传法律禁止向其献祭。该省总督担心普遍的不满会演变成叛乱，便就此向皇帝呈报。然而，皇帝远未将尼罗河带来的暂时肥沃看得比他应尽的天主忠信和宗教利益更为重要，答复如下：\Quote{让那条河停止流动，如果法术是保证其规律运行所必需的；或者如果它乐于接受祭献，或者如果必须用血来混合源自天主乐园的水。}不久之后，尼罗河泛滥得如此猛烈，以至最高的高地都被淹没。当它达到最远界限并几乎达到最大水时，水流并未因此减弱，反而继续上涨，使埃及人陷入了相反的恐惧。他们担心的是亚历山大城和利比亚的一部分会被淹没。亚历山大的异教徒对这一意外事件感到愤怒，在公共剧院里嘲笑说，河流像一个老人或傻瓜，无法控制自己的行为。许多埃及人因此被引导放弃祖先的迷信，皈依了基督信仰。这些事件是我所了解的。\par

% structure: p0065 source=h2 target=subsection reason=relative-heading-rank; base=h2

\subsection{第二十一章 发现我们主的前驱的光荣首级及其相关事件}

大约此时，洗者若翰的头颅——黑落狄亚曾向分封侯黑落德索要过若翰的头——被移往君士坦丁堡。据说它是由一些马其顿异端\OriginalTerm{马其顿异端}{Macedonian heresy}的隐修士发现的，他们最初住在君士坦丁堡，后来又定居在基里基雅。宫内第一总管马尔多尼乌斯在前朝统治期间将这一发现禀报了宫廷，瓦伦斯皇帝命令将该圣髑移往君士坦丁堡。负责运送的官员将它放在一辆公共马车上，一路行进到卡尔西顿境内的潘提基乌姆地区。在那里，马车的骡子突然停下；不论鞭打还是马夫的威吓，都无法驱使它们再前进一步。这一极不寻常的事件被所有人，甚至皇帝本人，都视为出于天主；因此圣头被存放在附近的科西劳斯村，该村属于马尔多尼乌斯。不久后，德奥多西皇帝因天主的感动或先知的催迫，前往该村。他决定迁移这位洗者的遗骸，据说除了圣女玛特洛娜——她曾负责守护这圣髑——之外，没有遇到其他反对。皇帝放下了所有权威和武力，经过多次恳求，终于从她那里勉强获得了移走头颅的同意；因为她铭记着瓦伦斯皇帝命令迁移时期所发生的事。皇帝将它连同存放的盒子放在自己的紫袍里，带入君士坦丁堡郊区一个称为赫布多莫斯的地方，在那里建造了一座宽敞而宏伟的圣殿。皇帝虽动用了恳求和许诺，却无法说服那位被指派守护圣髑的妇人放弃她的宗教立场；因为她似乎属于马其顿异端。一位持同样倾向的司铎，名叫文森特，他也负责守护这位先知棺椁并为其举行司祭职事，后来追随了皇帝的宗教意见，并与天主教会共融。马其顿异端者声称，他曾发誓绝不偏离他们的教义；但他后来公开宣称，如果洗者愿意跟随皇帝，他也要与他共融并分离。他是波斯人，与他的亲属阿达斯一起，于君士坦提乌斯统治期间离开祖国，为躲避当时波斯基督徒所受的迫害。到达罗马帝国后，他被列入圣职人员行列，并升任司铎之职。阿达斯结了婚，为教会做出了巨大贡献。他留下了一个儿子奥克森提乌斯，以其极为忠信的虔诚、对朋友的热忱、节制的生活、对文学的热爱以及在外教和教会文学上的深厚造诣而著称。他举止谦逊退隐，虽被允许与皇帝和朝臣亲近，并拥有非常显赫的职位。他的记忆至今仍受认识他的所有隐修士和热忱者的尊敬。曾受托保管圣髑的妇人，在科西劳斯度过了余生。她以虔诚和智慧极为出众，教导了许多圣洁的童女；我确实听说，至今仍有许多人体现出在玛特洛娜教导下所培养出的高尚品格。\par

% structure: p0067 source=h2 target=subsection reason=relative-heading-rank; base=h2

\subsection{第二十二章 在罗马的小瓦伦提尼安皇帝被绞死。僭主欧根尼乌斯。特巴依德隐修士若翰的预言。}

当\Person{狄奥多西}如此明智而和平地统治东方臣民并事奉天主时，传来消息说\Person{瓦伦提尼安}被勒死了。有些人说，他是被寝宫太监们处死的，这是在军事首领\Person{阿尔博加斯特}和一些朝臣的请求下进行的；这些朝臣不满意这位年轻君主开始效法其父治理国政，并违背他们所认可的意见。然而，另一些人断言，\Person{瓦伦提尼安}是亲手自尽的，因为他发现自己试图行不符合其年龄之事时受到阻碍；因此他认为不值得活下去；因为虽是皇帝，却不被允许做他愿意的事。据说这个孩子相貌高贵，在帝王风范上出类拔萃；若他活到成年，必显出治理帝国之才，并在宽宏和正义上超越其父。然而，虽具这些前途无量的品质，他却以上述方式死去了。\par

一个名叫\Person{欧根尼}的人，在基督教信仰表白上绝非真诚，却觊觎王权，并接受了帝王权力的标志。他希望通过这一尝试安全成功；因为他被一些人的预言引导，这些人自称通过观察动物内脏和肝脏以及星象来预见未来。罗马人中的最高阶层沉迷于这些迷信。当时的近卫总长\Person{弗拉维安}，一个有学问且看似有政治天赋的人，以通晓各种预言未来之术而闻名。他说服\Person{欧根尼}拿起武器，向他保证他注定要登上王位，他的军事行动必将胜利，基督教将被废除。被这些谄媚的表象所迷惑，\Person{欧根尼}招募了一支军队，占领了进入意大利的门户，罗马人称之为\OriginalTerm{尤利安阿尔卑斯山}{Julian Alps}，这是一片高耸险峻的山脉；他事先占领并加固了这些隘口，因为那里只有一条狭窄通道，两侧被悬崖和极高山峰封闭。\Person{狄奥多西}感到困惑，不知是该等待战争结果，还是最好先攻击\Person{欧根尼}；在这两难之际，他决定咨询\Person{约翰}，一位来自\Place{底比斯}的隐修士，如前所述，他以预知未来而闻名。于是他派遣宫中的太监、忠诚可靠的\Person{欧特罗皮乌斯}前往埃及，命令他尽可能将\Person{约翰}带到朝廷；但若他拒绝，则了解应当做什么。当他来到\Person{约翰}那里时，隐修士无法被说服去见皇帝，但通过\Person{欧特罗皮乌斯}传话说，战争将以\Person{狄奥多西}胜利、暴君被杀而告终；但胜利之后，\Person{狄奥多西}本人将死于意大利。这两个预言的真实性随后被事件所证实。\par

% structure: p0070 source=h2 target=subsection reason=relative-heading-rank; base=h2

\subsection{第23章。在安提约基雅征收贡金及摧毁皇帝雕像。由大司祭\Person{弗拉维安}率领的使团。}

此时，由于战争需要，负责官员认为有必要征收比往常更多的赋税；因此，叙利亚\Place{安提约基雅}的民众发生叛乱；皇帝和皇后的雕像被推倒，用绳索拖过全城，正如这类事件中常见的那样，狂怒的群众说出了激情所能激起的每一种侮辱性称呼。皇帝决定以处死许多安提约基雅市民来报复这一侮辱；人民在听到这消息时惊得哑口无言；市民的愤怒已经平息，代之以忏悔；仿佛已经受到威胁的惩罚，他们沉浸在哀叹和眼泪中，恳求天主转离皇帝的愤怒，并在他们的连祷中使用了一些哀歌。他们派遣主教\Person{弗拉维安}作为大使前往\Person{狄奥多西}处；但当他到达时，发现皇帝对所发生之事的怒气未消，于是他采取了以下计策。他让一些习惯在皇帝餐桌上唱歌的年轻人，用安提约基雅人的连祷文吟唱这些赞美诗。据说皇帝的仁慈被激发；他立刻被怜悯所征服；他的怒气被抑制，当他心中挂念那城市时，眼泪滴落在手中的杯子上。据说在叛乱发生的前一晚，曾有一个形如高大可怖妇人的幽灵，在城中街道上徘徊，用一根刺耳的鞭子抽打空气，类似于在公共剧院中用来驱赶出场的野兽的那种鞭子。可以推断，这场叛乱是由某个邪恶恶毒的魔鬼所煽动。毫无疑问，若不是皇帝因尊重这位司祭的祈求而止息怒气，必将发生大量流血事件。\par

% structure: p0072 source=h2 target=subsection reason=relative-heading-rank; base=h2

\subsection{第24章。\Person{狄奥多西}皇帝战胜\Person{欧根尼}。}

当他完成战争准备后，\Person{德奥多西乌斯}宣布他的幼子\Person{霍诺留斯}为皇帝，并将他留在\Place{君士坦丁堡}与先前已被立为皇帝的\Person{阿尔卡狄乌斯}共同统治，自己则率领军队从东方前往西方。他的军队不仅包括罗马士兵，还有来自\OriginalTerm{伊斯特尔河}{Ister}沿岸的蛮族队伍。据说他离开\Place{君士坦丁堡}后，来到\Place{第七里程碑}处，进入他为纪念\Person{圣若翰洗者}而建造的圣堂向天主祈祷；他奉\Person{圣若翰洗者}之名祈求罗马军队获胜，并恳求洗者本人相助。献上这些祈祷后，他继续向意大利进发，越过\Place{阿尔卑斯山}，占领了前哨阵地。从山巅下来时，他看见面前平原上布满了步兵和骑兵，同时注意到敌军的一些部队正埋伏在他身后的山间凹地中。他的先头部队攻击了驻扎在平原上的步兵，一场激烈而胜负难料的战斗随之爆发。此外，当军队包围了他时，他意识到自己已落入他人手中，即使有人愿意救他也无能为力，因为埋伏在他后方的人已占据了高地；他便俯伏在地，含泪祈祷，天主立即应允了他；因为埋伏在高地上的部队将领派人向他表示，愿意作为盟友为他效力，条件是他要在军队中赐予他们尊贵的职位。由于手头既无纸也无墨，他拿起一些蜡板，在上面写下他将赐予他们的崇高而适当的职位，只要他们履行对他的承诺。在这些条件下，他们来到皇帝面前。战局尚未偏向任何一方，平原上的战斗仍然势均力敌，这时一阵狂风迎面扑向敌人。这样的狂风我们从未记载过，它冲散了敌军的行列。射向罗马人的箭矢和标枪，仿佛被敌军自身反弹回来，转而刺向投掷者的身体；他们的盾牌也从手中被夺走，带着泥土和灰尘旋转着击向他们。他们这样暴露无遗、毫无防御地面对罗马人的武器，许多人丧命，而少数试图逃跑的人很快被俘。\Person{欧根尼乌斯}俯伏在皇帝脚下，恳求饶他一命；但正当他这样哀求时，一名士兵砍下了他的头颅。\Person{阿尔博加斯特}战后逃跑，自杀了。据说战斗进行时，一个附魔的人出现在\Place{赫布多摩斯}的天主圣殿里——皇帝出发时曾在那里祈祷——他侮辱\Person{圣若翰洗者}，讥讽他头被砍掉，并喊出以下话语：\Quote{你战胜我，并设计陷害我的军队。}当时在场且急切等待战争消息的人们大为惊奇，当日就将此事记录下来，事后证实这正是战斗发生的同一天。这些事迹的历史就是这样。\par

% structure: p0074 source=h2 target=subsection reason=relative-heading-rank; base=h2

\subsection{第二十五章 \Person{圣盎博罗削}在\Person{德奥多西乌斯}皇帝面前的无畏姿态。\Place{得撒洛尼}大屠杀。记述这位圣人的其他义行。}

\Person{欧根尼乌斯}死后，皇帝前往\Place{米兰}，走向教堂，要在其墙内祈祷。当他走近殿门时，城中的主教\Person{盎博罗修}迎面而来，抓住他的紫袍，当着人群对他说：\Quote{退后！一个被罪恶玷污、双手沾满无辜流血的人，若不悔改，不配进入这些神圣的场所，或分享神圣的奥迹。}皇帝被主教的胆量所震撼，开始反思自己的行为，满怀痛悔地转身退去。这罪的起因如下：当\Person{布特里克}是\Place{伊利里亚}的军队统帅时，一个战车御者看见他可耻地暴露在一家酒馆，试图侮辱他；御者被逮捕并拘押。一段时间后，赛马场要举行一些盛大的赛马，\Place{塞萨洛尼卡}的民众要求释放囚犯，认为他对比赛不可或缺。因他们的请求未受重视，他们起来暴动，最终杀死了\Person{布特里克}。皇帝听闻此事，立即被激起愤怒，命令处死一定数量的市民。城市充满许多无辜被杀的血；因为刚刚到达那里的前往其他地方的旅客也与其他受害者一同被牺牲。有许多值得同情的受苦事例，以下是其中之一：一个商人献出自己为他两个都被选为牺牲品的儿子代死，并向士兵承诺，只要他们同意交换，就把他所有的金子给他们。士兵们不能不同情他的不幸，同意他以一子代之，但声明不敢释放两个年轻人，因为那会使被杀的人数不足。父亲注视着他的儿子们，呻吟哭泣；他不能从死亡中救出任何一个，但由于对两者的爱相等，他犹豫不决，直到他们被杀。我还得知，一个忠信奴隶自愿代替主人去死，主人正被带往刑场。看来正是为这些和其他残忍行为，\Person{盎博罗修}斥责皇帝，禁止他进入教堂，并将他绝罚。\Person{狄奥多西}在教堂中公开忏悔自己的罪，在补赎期间，按哀悼者的习俗，不佩戴皇室的装饰。他还颁布了一项法律，命令负责执行皇帝命令的官员，在命令发布三十天之后，才可处死刑，以便皇帝的怒气有时间平息，为慈悲和悔改留出空间。\par

\Person{盎博罗修}无疑还做了许多其他配得上他司铎职份的事，这些事可能只有当地居民才知道。在他可称颂的事迹中，我知道以下这件。皇帝习惯坐在教堂集会的祭台围栏内，因此他远离其他民众而坐。\Person{盎博罗修}认为这习俗要么出于谄媚，要么出于纪律缺乏，便让皇帝坐在祭台栅栏之外，这样他就坐在民众前面，司铎后面。\Person{狄奥多西}皇帝赞同这一最佳传统，他的继任者也同样；我们被告知，从此以后这规矩被一丝不苟地遵守。\par

我认为有必要记录这位主教所做的另一件值得一提的事。一位显赫的异教徒侮辱了\Person{格拉提安}，宣称他不配作他父亲；因此他被判死刑。当他被带出执行时，\Person{盎博罗修}前往皇宫请求赦免。\Person{格拉提安}当时正观看私人狩猎表演，这是皇帝们常为私人娱乐而不是公众消遣而举行的。发现这种情况后，主教走到运进野兽的门；他藏起来，与管理动物的猎人们一同进入，尽管\Person{格拉提安}和随从们抵抗，但他没有停止，直到立即获得皇帝拯救的许可，释放了将被带去处死的人。\Person{盎博罗修}非常勤勉于遵守教会的法律，并维持其神职人员的纪律。我从他众多高尚事迹的记录中选取了以上两件事，以表明当涉及天主的服务时，他是以何等无畏的精神对掌权者说话。\par

% structure: p0078 source=h2 target=subsection reason=relative-heading-rank; base=h2

\subsection{第26章 圣\Person{多纳图}，\Place{尤罗亚}主教，以及\Place{斯基提亚}的大司祭\Person{德奥提莫}。}

当时在帝国各地有许多主教，以其圣德和卓越才干而备受赞誉，其中尤以\Place{伊庇鲁斯}\Place{欧罗亚}主教\Person{多纳图斯}为典型。该国民众传述了他所行的许多非凡奇迹，其中最有名的似乎是杀死一条巨蟒。它盘踞在大路上一个叫\Place{卡迈格菲拉}的地方，吞噬羊、山羊、牛、马和人。\Person{多纳图斯}遇见这只野兽，徒手攻击它，没有刀剑或矛；蟒蛇抬起头要冲向他时，\Person{多纳图斯}用手指在空中画十字，并朝蟒蛇吐唾沫。唾液进入它口中，它立刻毙命。它伸展在地上，大小不亚于印度著名的蛇。我得知当地人套了八对牛把尸体拖到附近田里烧掉，以免腐烂污染空气致病。这位主教的坟墓安放在一座以其命名的宏伟祈祷所中，位于一处多水泉源旁，这泉源是天主应他祈祷从干旱之地涌出的，那里原本无水。据说有一天他旅行经过此地，见同伴口渴，便用手翻土祈祷；祈祷未完，泉水涌出，至今未干。\Place{伊索里亚}（\Place{欧罗亚}领土内的村庄）居民见证此事为真。\par

\Place{托米}教堂乃至整个\Place{斯基提亚}的诸教会当时由\Person{德奥提摩}管辖。他受哲学培养；其德行令\Place{多瑙河}畔的蛮族匈人钦佩，称其为罗马人的天主，因他们目睹他行的神圣事迹。据说有一天他朝蛮族地区旅行，发现一些蛮族向\Place{托米}进军。随从痛哭，自认必死；他下马祈祷。结果蛮族走过却没有看见他、随从和他们的马。这些部落屡次掠夺\Place{斯基提亚}，他试图用食物和礼物驯服他们的凶残。一个蛮族由此认为他富有，决心俘虏他，倚着盾牌（这是他与敌人谈判的惯例）举起右手要向他抛套索，想把他拖走；但手悬在空中，那蛮族直到同伴恳求\Person{德奥提摩}为他向天主代祷才被释放。\par

据说\Person{德奥提摩}始终留着他初次献身哲学时所留的长发。他非常节制，没有固定用餐时间，只在饥饿口渴逼迫使然时才吃喝。我认为满足这些本能需求是出于必要而非感官享受，这才算哲学家的作为。\par

% structure: p0082 source=h2 target=subsection reason=relative-heading-rank; base=h2

\subsection{第二十七章。\Person{圣厄丕法尼}，\Place{塞浦路斯}主教，及其事迹详录。}

\Person{厄丕法尼}当时是\Place{塞浦路斯}主教省的首领。他不仅在生前因德行和奇迹著称，死后也因天主赐予的尊荣而闻名；据说在他墓前魔鬼被驱逐，疾病得痊愈。许多生前行的事迹归功于他，以下是我们所知最有名的一个。他对穷人非常慷慨，无论遭海难或其他灾祸的人；在把自己的全部家产用于周济后，又动用教会财富。这些财富因各省虔诚人士的捐赠而大增，他们因钦佩\Person{厄丕法尼}，生前委托他分发施舍，或死后遗赠给他。据说有一次，一位虔诚的司库发现教会收入几乎耗尽，库里所剩无几，便责备主教挥霍。但\Person{厄丕法尼}不顾劝阻，把剩下的小额也施舍了；这时一个陌生人到司库住的小屋，交给他一个装有大量金币的钱袋。由于馈赠者或送钱者都看不见，这显然是个奇迹，在一笔如此巨额的赠予中，一个人能隐藏自己身份，因而人人都认为这是天主的作为。\par

我也愿讲述另一桩归于\Person{埃丕法尼}的奇迹。我曾听说，类似的事迹也曾归于从前治理\Place{新凯撒勒雅}的\Person{额我略}；我看不出有任何理由怀疑记载的真实性；但这并不否定归于\Person{埃丕法尼}的奇迹的真实性。宗徒\Person{伯多禄}并非唯一曾使死者复活的人；福音作者\Person{若望}在\Place{厄弗所}行了类似的奇迹；\Person{斐理伯}的女儿们在\Place{耶辣颇里}也是如此。在不同的时代，天主的人曾行过类似的事迹。我想举出的奇迹如下：两个乞丐打听到\Person{埃丕法尼}何时会经过那条路，便约定用计谋从他那里多讨些施舍。一看到主教走近，其中一个乞丐便扑倒在地，装死；另一个站在旁边大声哀号，痛悼同伴的去世，并哀叹自己贫穷，无力为他埋葬。\Person{埃丕法尼}向天主祈祷，愿死者安息；他给了活着的那个人足够的钱用于埋葬，并对那哭泣者说：\Quote{吾儿，安排你同伴的安葬吧，不要再哭泣了；他现在不能从死中复活了；这灾祸是难免的，所以你应当顺受。}说完这些话，主教便离开了那里。一等到看不见一个人影，那曾对\Person{埃丕法尼}说话的乞丐用脚碰了碰仍伸躺在地上的同伴，对他说：\Quote{你演得很好；现在起来吧，因为靠你的辛劳，我们今天有了好口粮。}然而，那人仍那样躺着，既听不见任何呼喊，也感觉不到那正用尽全力推他的人；那另一个乞丐追上司铎，承认了他们的诡计，并扯着头发哀哭，恳求\Person{埃丕法尼}救活他的同伴。\Person{埃丕法尼}只是劝他耐心承受这灾祸，便打发他走了。天主并没有取消已发生的事，因我深信，祂的用意是要表明，那些欺骗祂仆人的人，其欺诈之罪如同是行在鉴察一切、倾听一切的主身上一样。\par

% structure: p0085 source=h2 target=subsection reason=relative-heading-rank; base=h2

\subsection{第二十八章 \Place{贝洛雅}的\Person{阿加丘}主教、\Person{则诺}、\Person{阿雅克斯}——德性杰出与著称的人物}

以下细节也是查询所得。\Person{阿加丘}在主教中出类拔萃；他此前已担任\Place{叙利亚}\Place{贝洛雅}的主教职务。当然，他的许多事迹都值得记录。他自幼便受培育于苦修的隐修生活，并严格遵行这种生活方式的一切规条。当他被提升为主教时，他显出了卓越德性的这一证据：他使主教府在一天中的任何时候都开放，以致市民和陌生人随时可以拜访他，即使在他用餐或休息时也不例外。依我之见，这种行为方式非常可钦敬；因为要么他生活的方式使他时刻能把握自己，要么他将此作为一种防备人性中邪恶的手段，这样，由于预料到会被人突然闯入，他就必须持续警惕，在职责上不犯错误，而更致力于履行盟约的行为。\par

\Person{则诺}和\Person{阿雅克斯}，两位著名的兄弟，约在同一时期活跃。他们献身于哲学（指修行）的生活，但并未定居旷野作隐修士，而是在\Place{加匝}（也称为\Place{玛犹玛}）这座沿海城市。他们都非常忠实地捍卫自己宗教的真理，并勇敢地宣认天主，以致经常受到外教人非常残忍和粗暴的对待。据说\Person{阿雅克斯}娶了一位非常可爱的女子，在与她同房总共三次后，生了三个儿子；此后他不再与她有来往，而是坚持隐修的操练。他将两个儿子培养去过奉献生活与独身，而允许第三个儿子结婚。他得体而卓越地治理了\Place{波托里乌姆}的教会。\par

\Person{则诺}自幼便弃绝世俗与婚姻，坚持不懈地忠诚事奉天主。据说——而我自己就是这断言真实性的见证——当他担任\Place{玛犹玛}教会的主教时，除非因病所困，他从未缺席清晨或傍晚的赞美诗或其他任何对天主的敬礼；而当时他已是一位老人，将近百岁。他继续在隐修哲学中度过一生，但通过从事织麻布的手艺，持续赚取所需，以自给并周济他人。他从未偏离此行径，直至生命终结，尽管他在该省所有司铎中年纪最长，并管理着最大教会的信众与财产。\par

我列举这些作为当时担任司铎者的榜样。若要一一列举所有善良者，那将是繁重的任务；他们中大多数是良善的，天主也因乐意垂听他们的祈祷并显行许多奇迹而为他们的生活作证。\par

% structure: p0090 source=h2 target=subsection reason=relative-heading-rank; base=h2

\subsection{第二十九章 发现\Person{哈巴谷}与\Person{米该亚}二位先知的遗骸。大帝\Person{德奥多西}皇帝驾崩。}

当整个教会处于这些卓越人物的治理之下时，神职人员和民众都被激励去效仿他们的德行和热忱。这个时代的教会不仅因这些杰出的虔诚典范而闻名，而且大约在此时，古先知哈巴谷以及稍后米该亚的遗骸也被发现。据我所知，天主通过梦境的神视向当时担任厄勒乌特罗波利斯教会主教的泽本努斯启示了这两具遗体的安放之处。哈巴谷的遗骸是在一个从前称为克依拉的城市\Place{克拉}发现的。米该亚的墓则是在距离克拉十斯塔狄昂的地方发现的，那个地方叫做\Place{贝拉特撒提亚}。当地人无知地将此墓称为\Quote{信徒的坟墓}，用他们的本地话来说就是\OriginalTerm{那弗撒米玛纳}{Nephsameemana}。这些发生在德奥多西皇帝统治期间的事件，足以使基督徒的宗教获得美誉。\par

在击败欧格尼乌斯之后，皇帝德奥多西在米兰逗留了一段时间，在那里他患上重病。他回想起隐修士\Person{若望}的预言，猜测自己的病会致死。他急忙派人从君士坦丁堡接来儿子奥诺留；见到儿子在身边后，他似乎好转，能够出席竞技场的比赛。然而，晚餐后他突然病情恶化，派人去请儿子主持观礼。他在次日夜间去世。此事发生在奥利布里和普罗比亚诺兄弟执政期间。\par

\SourceRecord{Translated by Chester D. Hartranft.；Nicene and Post-Nicene Fathers, Second Series；Vol. 2.；Edited by Philip Schaff and Henry Wace.；Buffalo, NY: Christian Literature Publishing Co.,；1890.；<http://www.newadvent.org/fathers/26027.htm>.}

% structure: p0001 source=h1 target=section reason=page-title

\section{教会史（第八卷）}

% structure: p0002 source=h2 target=subsection reason=relative-heading-rank; base=h2

\subsection{第一章  狄奥多西大帝的继位者。鲁菲努斯，近卫军长官，被杀。主要城市的首席司铎。异端者之间的分歧。关于诺瓦提安派主教西西尼乌斯的记述。}

狄奥多西之死正是如此，他曾如此有效地促进教会的壮大。他在六十岁、执政第十六年时去世。他留下两个儿子作为继位者。长子阿卡狄乌斯统治东方，霍诺留统治西方。二人均持守与他们父亲相同的宗教情感。\par

达玛苏斯已去世；此时，西里修是罗马教会的领袖；纳克塔里乌斯是君士坦丁堡教会的领袖；狄奥斐禄是亚历山大里亚教会的领袖；弗拉维安是安提阿教会的领袖；若望是耶路撒冷教会的领袖。此时，亚美尼亚和东部各省被野蛮的匈人蹂躏。东方长官鲁菲努斯被怀疑秘密邀请他们蹂躏罗马领土，以推进其野心计划；因传说他觊觎暴政。为此，他不久后被杀死；当军队从征服欧根尼乌斯返回时，阿卡狄乌斯皇帝按惯例从君士坦丁堡出城迎接他们；士兵们趁机杀害了鲁菲努斯。这些情况大大促进了宗教的扩展。皇帝们将其父的虔诚归因于暴君轻易被击败，以及鲁菲努斯夺取他们政府的阴谋被阻止；他们欣然确认了前任为教会利益颁布的所有法律，并额外地赐予自己的礼物。臣民以其为榜样，就连外教人也毫无困难地皈依了基督教，异端者也与天主教会合一。\par

由于阿里乌斯派和欧诺米派之间发生的纷争，我已在前面提及，这些异端者逐日减少。他们中的许多人，在思考自己一派中盛行的意见分歧时，断定天主的真理不可能与他们同在，于是归附于与皇帝持相同信仰的人。\par

君士坦丁堡的马其顿派的利益受到实质影响，因为他们在那时没有主教；自康斯坦提乌斯统治时期被欧多克西乌斯剥夺了教堂以来，他们一直只由长老管理，直到下一个统治时期。另一方面，诺瓦提安派虽然因逾越节的争论而动荡——这是萨巴提乌斯的一项创新——但他们中的大多数人仍平静地保有其教堂，且未曾受到针对其他异端者所颁布的任何惩罚或法律的骚扰，因为他们坚持天主圣三的三位是同一实体。他们领袖的品德也大大有助于在他们中间维持和睦。在阿格利乌斯的主席任职之后，他们由好人马尔西安管理；他逝世后，略早于现在所考虑的时候，主教职位传给了西西尼乌斯，一个非常雄辩的人，精通哲学和圣经的教义，并且如此善于辩论，以至于连欧诺米乌斯——他在这方面得到很好认可且工作有效——也经常拒绝与他辩论。他的生活方式谨慎且无可指摘；然而他纵情于奢华甚至过度，以致那些不认识他的人怀疑他能否在如此丰裕中保持节制。他在集会中举止优雅和蔼，因此受到天主教会主教、统治者及有学问之人的尊重。他的玩笑充满善意，能忍受嘲笑而不显出丝毫恼怒。他的应答迅捷而机智。有一次被问到，既然他是主教，为何每日沐浴两次，他回答说：\Quote{因为我不洗三次。}另一次，被一位天主教会成员嘲笑他穿白衣，他问哪里吩咐他该穿黑衣；当对方犹豫不答时，他继续说：\Quote{你无法为你的立场提供任何论证；但我向你们举荐\Person{所罗门}，最智慧的人，他说：\BibleQuote{愿你的衣服时常洁白。}况且，\WorkTitle{福音}中描述基督曾以白衣显现，\Person{梅瑟}和\Person{厄里亚}也以白衣显现给宗徒们。}依我看，以下回答也极为巧妙。加拉太地区安西拉的主教良提乌斯，在剥夺了他省内诺瓦提安派的教堂后，定居于君士坦丁堡。西西尼乌斯去见他，请求归还教堂；但他不但不依从，反而辱骂诺瓦提安派，说他们不配拥有教堂，因为他们废除了补赎的实践，阻断了天主的仁爱。对此，西西尼乌斯回答说：\Quote{没有人像我一样做补赎。}良提乌斯问他如何做补赎。\Quote{来见你}，西西尼乌斯反驳道。还有许多其他机智的言谈归于他，据说他甚至写过几部相当优雅的作品。但他的演说比他的著作获得更大的喝彩，因为他最擅长朗诵，并能通过声音、表情和愉悦的面容吸引听众。这段简要描述可作为这位伟大人物的气质和生活方式的证明。\par

% structure: p0007 source=h2 target=subsection reason=relative-heading-rank; base=h2

\subsection{第二章  伟大的金口圣若望的教育、训练、品行和智慧；他被擢升为教座；亚历山大里亚主教狄奥斐禄成为他确定的对手。}

\Person{Nectarius}逝世了大约在这个时候，关于继任者的祝圣展开了冗长的辩论。他们各自投票给不同的人，似乎不可能团结在一个人身上，时间沉重地流逝。然而，在\Place{Orontes}河上的\Place{Antioch}，有一位名叫\Person{John}的司铎，出身高贵，生活卓越，拥有如此令人惊叹的雄辩和说服力，以致叙利亚的诡辩家\Person{Libanius}宣称他超过了当代所有的演说家。当这位诡辩家临终时，他的朋友问他谁应接替他的位置。\Quote{本应是\Person{John}，}回答他，\Quote{要不是基督徒把他从我们这里夺走了。}许多在教堂里听了\Person{John}讲道的人因此被激发起对美德的热爱，并接受了他自己的宗教情感。因为通过过神性的生活，他将自己的美德所生的热忱传递给了听众。他产生了类似自己的信念，因为他不是通过修辞技巧和力量来强制推行它们，而是真实而诚恳地解释圣经。因为被行为装饰的话语通常会显得值得相信；但若没有这些，说话者就显得像个冒名顶替者和自己话语的叛徒，即使他热切地教导。\Person{John}在这两方面都值得称赞。他致力于谨慎的生活方式和严格的公共生涯，同时也使用清晰的语言，结合了言语的光彩。\par

他的天赋极好，并通过跟随最好的老师学习来提升它们。他向\Person{Libanius}学习修辞，向\Person{Andragathius}学习哲学。当人们期望他从事法律职业并加入辩护律师的行列时，他决定在圣经上操练自己，并按照教会的法律实践哲学。他在这哲学上的老师是\Person{Carterius}和\Person{Diodorus}，两位著名的苦修机构的主持。\Person{Diodorus}后来成为了\Place{Tarsus}教会的治理者，而且我被告知，他留下了许多自己写的书，在其中他解释了圣经文字的意义并避免了寓意。\Person{John}不仅自己接受了这些人的教导，还说服了\Person{Theodore}和\Person{Maximus}——他们曾是他跟随\Person{Libanius}学习时的同伴——与他同行。\Person{Maximus}后来成为\Place{Isauria}的\Place{Seleucia}的主教；\Person{Theodore}成为\Place{Cilicia}的\Place{Mompsuestia}的主教。\Person{Theodore}精通圣经以及修辞学家和哲学家的其他学科。在学习教规并常与圣人来往后，他对苦修生活方式充满了钦佩并谴责城市生活。但他没有坚持同样的目的，而是在改变后被拉回到以前的生活方式；并且为了证明自己的行为，他引用了许多他熟知的历史上的例子，然后回到了城市。听说他忙于事务并打算结婚后，\Person{John}写了一封书信，其语言和思想比人心所能产生的更为神圣，并寄给了他。读信后，他悔改了，立即放弃了财产，放弃了结婚的意图，被\Person{John}的建议所拯救，并回到了哲学生涯。在我看来，这是\Person{John}雄辩力量的一个显著例子；因为他很容易迫使一个习惯于说服和说服别人的人信服。凭借同样的雄辩，\Person{John}吸引了人们的钦佩；同时，他甚至在教堂里严厉地谴责罪人，并勇敢地对抗所有不公正的行为，就好像这些行为是对他自己犯下的一样。这种勇敢取悦了人民，但却使那些富人和有权势的人伤心，因为他们犯有他所谴责的许多恶习。\par

因此，\Person{John}被那些通过经验认识他的人，以及那些通过他人报告而认识他的人如此高度评价，以至于所有罗马帝国的臣民都认为他配得上言行一致地成为\Place{Constantinople}教会的主教。神职人员和民众一致选举他；他们的选择得到了皇帝批准，皇帝也派出了迎接他的使团；并且为了给他的祝圣增添更大的隆重性，召开了一次会议。皇帝的信函发出后不久，东方将军\Person{Asterius}派人请\Person{John}到他那里，好像他需要他一样。\Person{John}一到，他立刻让他上自己的战车，迅速将他送到一个名为\Place{Pagras}的军事驿站，在那里将他交给了皇帝派来寻找他的官员。\Person{Asterius}在\Place{Antioch}市民知道要发生什么事之前就派人请\Person{John}，这是非常谨慎的；因为他们很可能会煽动叛乱，并对他人造成伤害，或者自己遭受暴力行为，而不愿让\Person{John}从他们身边被夺走。\par

当\Person{若望}抵达\Place{君士坦丁堡}，众司铎聚集时，\Person{德奥斐禄}反对他的祝圣，并提议推举其教会中的一位\OriginalTerm{司铎}{presbyter}，名叫\Person{依西多尔}，他在\Place{亚历山大里亚}负责接待旅客和穷人。我从认识\Person{依西多尔}的人得知，他从年轻时起就在\Place{赛克斯}附近实践哲学美德。另有人说，他因参与并熟悉一项非常危险的事业而赢得了\Person{德奥斐禄}的友谊。相传在对抗\Person{马克西穆斯}的战争期间，\Person{德奥斐禄}托付\Person{依西多尔}分别致皇帝和僭主的礼物和信件，派他去\Place{罗马}，要他留在那里直到战争结束，然后把礼物和信件交给获胜者。\Person{依西多尔}遵照指示行事，但计谋被识破；他担心被捕，逃往\Place{亚历山大里亚}。从那时起，\Person{德奥斐禄}对他表现出极大的眷顾，并为了酬报他的服务，竭力将他擢升为\Place{君士坦丁堡}的主教。但无论这一传闻是否属实，还是\Person{德奥斐禄}因\Person{依西多尔}的卓越而想祝圣他，可以肯定的是，他最终屈服于那些支持\Person{若望}的人。他惧怕\Person{欧特罗比乌}，此人巧妙而急切地要求这次祝圣。\Person{欧特罗比乌}当时掌管皇室，据说他威胁\Person{德奥斐禄}，除非他与其余主教一致投票，否则他必须对那些想要控告他的人进行辩护；因为当时有许多书面的控告状在议会中针对他。\newline{}\par

% structure: p0012 source=h2 target=subsection reason=relative-heading-rank; base=h2

\subsection{第三章 若望迅速晋升主教，更加剧烈地处理事务。他在各地教会重建纪律。通过向罗马派遣使节，他消除了对弗拉维安的敌意。}

\Person{若望}一晋升主教职，首先致力于改革他神职人员的品行；他责备并纠正他们的生活方式、饮食以及各项事务的一切程序。他还将一些神职人员逐出教会。他天生倾向于责备他人的过失，并正义地对抗那些行为不义的人；在主教任上，他更表现出这些特点；因为他的本性，一旦获得权力，便使他的舌头趋向于责备，并更轻易地对敌人激起愤怒。他并没有将努力局限于自己教会的改革；作为一个善良而心胸宽广的人，他寻求纠正整个世界的弊端。他一进入主教职位，立即努力结束关于\Person{保利努斯}在西方和埃及主教与东方主教之间引起的纷争；因为为此，整个帝国的教会普遍分裂。他请求\Person{德奥斐禄}协助促成\Person{弗拉维安}与\Place{罗马}主教的和解。\Person{德奥斐禄}同意与他合作恢复和睦；\Place{贝雷亚}的主教\Person{阿卡基乌}和\Person{德奥斐禄}曾提议取代\Person{若望}祝圣的候选人\Person{依西多尔}被派往\Place{罗马}为使节。他们很快完成了使命，乘船返回\Place{埃及}。\Person{阿卡基乌}前往\Place{叙利亚}，带着来自\Place{埃及}和西方司铎致\Person{弗拉维安}追随者的和解信。各教会在长时间延迟后，终于再次放下纷争，彼此共融。在\Place{安提约基雅}，那些被称为\OriginalTerm{尤斯塔提派}{Eustathians}的人，尽管没有主教，仍继续举行分离的集会。\Person{保利努斯}的继任者\Person{埃瓦格里乌}，如我们所说，没有比他活得更久；我认为，由于没有人反对，主教们更容易实现和解。平信徒，如同民众的习惯，逐渐转向在\Person{弗拉维安}领导下集会的人；于是，随着时间的推移，他们越来越合一。\newline{}\par

% structure: p0014 source=h2 target=subsection reason=relative-heading-rank; base=h2

\subsection{第四章 哥特蛮族盖纳斯的企图。他所行的恶事。}

\Emph{一个蛮族}，名叫\Person{盖纳斯}，曾投奔\OriginalTerm{罗马人}{Romans}，并从军队最低阶层升至军事指挥权，图谋篡夺罗马帝国皇位。为此，他从自己的故乡召唤同族\OriginalTerm{哥特人}{Goths}来到罗马领土，并任命他的亲属为护民官和千夫长。他的亲属\Person{提尔宾吉鲁斯}在\Place{弗里吉亚}指挥一支大军，开始叛乱；明眼人都清楚他在铺路。\Person{盖纳斯}以愤慨许多\Place{弗里吉亚}城市被蹂躏为借口——这些城市曾被委托给他监管——转向援助它们；但当他到达时，大量蛮族已装备备战，他便暴露了先前隐藏的计划，洗劫了受命守卫的城市，并准备攻击其他城市。随后他前往\Place{比提尼亚}，在\Place{加采东}边界扎营，威胁战争。东方和亚洲的城市，以及居住在这些地区与\Place{攸克辛海}之间的所有人，都处于危险之中。皇帝及其谋士认为，在没有准备的情况下，冒险对抗这些已绝望的人是不安全的；因为皇帝宣称他愿意在各方面对他友善，并派人向\Person{盖纳斯}承诺他可能要求的任何东西。\newline{}\par

\Person{盖纳斯}要求将两位执政官——他怀疑对其怀有敌意的\Person{萨图尔尼努斯}和\Person{奥勒良}——交给他；当他掌控他们后，却赦免了他们。随后，他在\Place{加采东}附近、安放殉道者\Person{欧费米雅}墓穴的祈祷所与皇帝会晤；他与皇帝彼此以友谊之誓相互约束后，便放下武器，前往\Place{君士坦丁堡}。在那里，根据帝国敕令，他被任命为步兵与骑兵统帅。如此超出其应得的顺遂，他无法以节制之心承受；正如完全出乎意料地，他在先前的行动中取得了惊人的成功，于是他决意破坏天主教会的和平。他是一名基督徒，与其他蛮族一样，信奉\OriginalTerm{亚略异端}{Arian heresy}。或受该派首领怂恿，或出于自身野心的驱使，他向皇帝请求将城中一座教堂交与亚略派手中。他声称，作为罗马军队的统帅，若想在祈祷时却被迫退至城外，这是既不公正也不适当的。\Person{若望}得知这些行径后并未坐视不管。他召集当时城中所有主教，一同前往皇宫。他在皇帝与\Person{盖纳斯}面前长篇陈辞，斥责后者乃外邦人与逃亡者，并提醒他其性命乃由皇帝之父所救，他曾向皇帝之父宣誓效忠，亦向其子女、罗马人及他正试图使之失效的法律宣誓。讲完这番话后，他展示了\Person{德奥多西}所立、禁止异端在城内拥有教堂的法律。随后，\Person{若望}转向皇帝，劝勉他维护已制定的反对异端者的法律；并告诉他，丧失帝国也比因背叛天主的殿宇而陷于不敬更好。\Person{若望}就这样勇敢地像男子汉一样发言，在他照管的教会中未给革新留任何余地。然而，\Person{盖纳斯}不顾誓言，进攻城市。他的企图由一颗直接出现在城市上空的彗星预先宣告；这颗彗星异常巨大，据说比先前所见任何彗星都大，几乎触及地面。\Person{盖纳斯}意图首先夺取银行家的仓库，并希望聚集他们庞大的财富。但因其计划泄露，银行家们藏匿了现款，不再像往常一样在桌上公开陈列银钱。\Person{盖纳斯}于是夜间派遣一些蛮族去焚烧皇宫；但他们不熟练且充满恐惧，于是折返。因为他们接近宫殿时，仿佛看见一群全副武装的巨人，便回去报告\Person{盖纳斯}说，有增援部队刚刚抵达。\Person{盖纳斯}不相信他们的报告，因为他确信没有军队进入城中。然而，次日夜间他派去同样目的的其他人员带回同样的报告，于是他亲自出去亲眼目睹这一奇景。他想象面前的军队是从其他城市调来的士兵，这些部队夜间保卫城市和宫殿，白天隐藏起来；\Person{盖纳斯}假装被鬼魔附身，以祈祷为借口，前往皇帝之父为纪念\Person{若翰洗者}而在\Place{赫卜多莫斯}建造的教堂。一部分蛮族留在\Place{君士坦丁堡}，其他则随同\Person{盖纳斯}；他们秘密携带武器和装满标枪的罐子藏在妇女的马车里，但被发现后，他们杀死了试图阻止武器运出的城门守卫。于是城中充满混乱与喧嚣，仿佛突然被占领。这可怕时刻被一个良善的念头所支配：皇帝毫不延迟地宣布\Person{盖纳斯}为公敌，并下令杀死城中所有留下的蛮族。此令一出，士兵们立即扑向蛮族，杀死了其中绝大部分；然后他们放火烧了以哥特人命名的教堂；因为按照惯例，蛮族聚集在那祈祷所中，因为别无避难之处，城门已关闭。\Person{盖纳斯}听闻此灾，穿过\Place{色雷斯}，前往\Place{刻索尼苏斯}，意图渡过\Place{赫勒斯滂}；因为他认为若能征服对面的亚细亚海岸，便能轻易使帝国东部的所有省份归顺自己。然而这一切都出乎他的希望，因为罗马人得到了神力的眷顾。皇帝派遣的军队由\Person{弗拉维塔}率领，从陆路和海路同时到达；此人虽是蛮族出身，却是个好人，且是位能干的将军。蛮族没有船只，轻率地试图用木筏渡过\Place{赫勒斯滂}到对面大陆；突然狂风大作，猛烈地将他们吹散，并将他们冲向罗马船只。大部分蛮族及其马匹溺死；但许多被军队所杀。然而，\Person{盖纳斯}与少数随从逃脱；但不久后，当他们穿过\Place{色雷斯}逃跑时，又遭遇另一支罗马军队分队，\Person{盖纳斯}及其所有蛮族全军覆没。这就是\Person{盖纳斯}的狂妄图谋与生命的终结。\par

\Person{弗拉维塔}在这次战争中表现极为突出，因此被任命为执政官。在他与\Person{味增爵}执政期间，皇帝得一子。年轻的王子以其祖父之名命名，并在下一个执政官任期开始时被立为\OriginalTerm{奥古斯都}{Augustus}。\par

% structure: p0018 source=h2 target=subsection reason=relative-heading-rank; base=h2

\subsection{第五章 \Person{若望}以其教导感化民众。关于一位追随\Person{马其顿纽斯}的妇人，因她之故，饼变成石头。}

\Person{若望}以堪为典范的明智管理\Place{君士坦丁堡}教会，并引领许多外邦人和异端者与他合而为一。每日都有大批群众到他那里去，有些是为了听他的讲论而受到建树，另有些则是意图试探他。然而，他取悦并吸引了各阶层的人，引导他们接受与他相同的宗教情感。由于民众拥挤在他周围，仍不满足于他的言语，以致当他们被推来挤去、彼此挤压时，便陷入危险；而每个人都奋力向前，以便站得更近，更清楚地听\Person{若望}所说的话，他便自己站到读经者的讲台中央，坐下教训群众。在我看来，现在正是我的历史中插入记述\Person{若望}在世时所行一件奇迹的适宜之处。有一人属\Place{马其顿}异端，与一位信仰相同的妻子同住；他偶然听到\Person{若望}论及应持守何种关于天主本性的见解的话语，被他所听到的论证说服，便竭力劝妻子接受同样的见解。但她先前的思维习惯以及所认识其他妇人的谈话，使她不愿顺从丈夫的意愿；当他发现所有说服她的努力都徒劳无功时，便告诉她，\Quote{除非她在神圣之事上与他同心，否则就不能继续与她同住。}那妇人因此答应按他所要求的去做；但同时，她把这事告诉了一个她信赖其忠心的婢女，并利用她来欺骗丈夫。在举行奥秘庆典的季节（领受奥秘者会明白我的意思），这妇人保留着给她的东西，并低着头，仿佛在祈祷。站在她身后的婢女，将随身带来的一块面包放在她手中；但她刚把面包放入齿间，面包就变成了石头。既然有这样神圣的事发生在她身上，她非常害怕再有灾祸临头，便跑到主教那里，作了告解。她向主教出示那块带有齿痕的石头；这块石头由某种未知的物质构成，颜色极其奇特。她含泪恳求宽恕，此后便一直与丈夫持守相同的宗教信条。若有人以为这段记述难以置信，他可以去查看那块石头；因为它至今仍保存在\Place{君士坦丁堡}教会的宝库中。\par

% structure: p0020 source=h2 target=subsection reason=relative-heading-rank; base=h2

\subsection{第六章 \Person{若望}在\Place{亚细亚}和\Place{夫黎基雅}的行实。\Person{赫拉克利德}，\Place{厄弗所}主教，与\Person{革隆提乌}，\Place{尼苛美底亚}主教。}

\Person{若望}得知亚细亚及其邻近地区的教会由不称职的人管理，他们为了收入与馈赠而出卖司铎职，或因私情授予该尊位，便前往厄弗所，罢免了十三位主教，部分在吕西亚和夫黎基雅，部分在亚细亚本地，并另立他人接替。厄弗所的主教已去世，因此他祝圣了赫拉克利德为厄弗所教会的主教。赫拉克利德是塞浦路斯人，曾是若望属下的一名执事；他先前加入了斯塞提斯的隐修士团体，并曾师从隐修士厄瓦格留。若望也驱逐了尼科默底亚教会的主教格隆提。后者曾是米兰教会盎博罗削属下的执事；他宣称——我不知道为何，或是为了编造奇迹，或是他本人被魔鬼的伎俩与幻象所欺骗——他夜间抓住了一个形似驴子的东西(\OriginalTerm{驴形物}{\GreekText{ὀνοσκελίς}})，砍下了它的头，并扔进了一个磨坊。盎博罗削认为这种言论有失天主执事的身份，命令格隆提闭门思过，直到他以补赎洗清过错。然而，格隆提是一位医术精湛的医生；他善于言辞、有说服力，深谙交友之道；因此他嘲笑了盎博罗削的命令，前往君士坦丁堡。不久他就获得了宫廷中一些权贵的友谊；不久之后，他便被擢升为尼科默底亚主教。为他祝圣的是卡帕多细亚凯撒勒雅的主教赫拉迪乌，后者之所以更乐意为他行此礼，是因为格隆提通过自己在宫廷的人脉，为赫拉迪乌的儿子谋得了军队中的高位。盎博罗削闻知此事后，写信给君士坦丁堡教会的主席纳克塔留，要求他将格隆提从司铎职中驱逐出去，不容许司祭职与教会秩序受到如此践踏。纳克塔留虽愿遵从这一命令，却因尼科默底亚民众的坚决反抗而始终未能执行。若望罢免了格隆提，而祝圣了潘索斐，后者曾是皇后的教师，虽是一位虔诚、温和且性情温良的人，却不受尼科默底亚人喜爱。他们频频骚动，公开或私下列举格隆提的恩惠，以及他的医术所带来的慷慨益处，无论贫富均一视同仁；正如我们赞美所爱之人时常做的那样，他们还为他增添了许多其他美德。他们在本城和君士坦丁堡的大街小巷游行，仿佛发生了地震、瘟疫或其他天主义怒的降临，唱着圣咏，献上祈祷，恳求能继续拥有格隆提作他们的主教。他们最终被迫屈服于现实，含泪叹息地与格隆提分离，迎来了一位他们既恐惧又厌恶的主教。被罢免的主教们及其所有追随者纷纷谴责若望，称他是在教会中制造革新的始作俑者，违背祖传法律改变受圣职者的权利；出于愤懑，他们谴责了若望所做的一些事情，而这些事按照大多数人的看法是值得称赞的。在其他事项中，\newline{}他们指责他对欧特洛彼所采取的行动。\par

% structure: p0022 source=h2 target=subsection reason=relative-heading-rank; base=h2

\subsection{第七章　关于太监长\Person{欧特洛彼}和他所颁布的法律。他因被逐出教会而被处死。针对\Person{若望}的怨言。}

\Person{欧特洛彼}原是太监长，是我们所知所闻的第一位也是唯一一位获得执政官和贵族尊位的太监。当他被擢升到目前的权力时，他并未思及未来，也未想到人事的无常，竟下令将那些在教堂寻求避难的人驱逐出去。他如此对待了\Person{提玛修}之妻\Person{彭塔迪亚}。提玛修是军队中的一位将领，能干且令人畏惧；但欧特洛彼以他图谋篡位为借口，获得了一项谕旨，将他流放到埃及的帕西斯。据我所知，提玛修死于口渴，或是担心更糟的遭遇，他在那里的沙漠中被困，被发现时已死。欧特洛彼颁布了一项法律，规定无人可以在教堂寻求避难，已经逃到那里的也要被赶出去。然而，他却是第一个违反这条法律的人；因为法律颁布后不久，他得罪了皇后，立即离开宫廷，逃到教堂作恳求者。当他躺在祭台下时，\Person{若望}发表了一篇讲道，斥责权力的骄傲，并引导人们注意人间伟大的无常。于是，若望的敌人借此机会责备他，说他非但不同情一个遭遇厄运的人，反而斥责他。欧特洛彼不久后便因他邪恶的计划付出了代价，被斩首；他所颁布的法律也从公开碑文上被抹去。天主的义怒就这样迅速临到对教会所行的不义，教会恢复了繁荣，神圣敬拜也增加了。君士坦丁堡的民众比以前更加勤勉地参加早晚颂歌的咏唱。\par

% structure: p0024 source=h2 target=subsection reason=relative-heading-rank; base=h2

\subsection{第八章　\Person{若望}引入应答式圣歌以对抗亚略派。若望的教导大大增进了正统派信徒的利益，而富人则越来越愤怒。}

亚略派信徒在德奥多西皇帝在位期间被剥夺了他们在君士坦丁堡的教堂后，便在城墙外聚集他们的教会。他们先前在夜间聚集于公共柱廊，分成队列，以便对唱圣歌，因为他们编撰了某些反映他们自身教义的副歌，并在拂晓时列队游行，唱着这些圣歌，前往他们举行教会的地方。他们在所有隆重庆节，以及一周的第一天和最后一天，都这样做。这些颂歌中提出的情感很可能引发争端。例如，以下这句：\Quote{那些说三个位格构成一个权能的人在哪里？} 其他类似的尖刻评论也穿插在他们的作品中。若望担心自己教会中的信众看到这些表演而被引入歧途，因此命令他们也以同样的方式唱圣歌。正统信徒变得更加突出，短时间内，在人数和游行方面超过了反对的异端者；因为他们有银制十字架和点燃的蜡炬在前面抬着。皇后的宦官被指派来管理这些游行，支付所需的一切费用，并预备圣歌。因此，亚略派信徒或被嫉妒所驱，或被报复所驱，攻击了天主教会的成员。双方发生了大量流血事件。布里索（这是那位皇家宦官的名字）被投掷的石头击中了额头。皇帝的愤怒被点燃，他制止了亚略派信徒的集会。由于从上述方式和原因开始了唱圣歌的习俗，天主教会的成员并没有停止这种做法，而是保留至今。这些游行的设立以及他在教会中的服务使若望深受民众爱戴；但他却被圣职人员和权贵所憎恨，因为他的自由大胆，他一旦发现圣职人员有不义之举，必会斥责他们；当权贵滥用财富、行不敬之事或沉溺于享乐时，他也必会劝诫他们回归德行。\par

% structure: p0026 source=h2 target=subsection reason=relative-heading-rank; base=h2

\subsection{第九章。\Person{塞拉皮雍}，总执事，以及圣\Person{奥林比亚斯}。一些知名人士蛮横地攻击若望，诽谤他为固执和暴躁的人。}

圣职人员对若望的敌意因他的总执事\Person{塞拉皮雍}而大大增加。他是一名埃及人，天性易怒，总是侮辱对手。敌意因若望给予\Person{奥林比亚斯}的劝告而进一步加剧。\Person{奥林比亚斯}出身最为显赫，虽年轻守寡，却热忱地按照教会法规从事隐修生活，\Person{纳克塔里乌斯}已祝圣她为女执事。若望见她慷慨地将财物施予任何求乞者，且轻视一切，只事奉天主，便对她说：\Quote{我赞赏你的意图；但愿意你知道，那些渴望按照天主达到德行完美的人，应当有节制地施舍财富。然而，你却将财富赠予富人，这如同将财富投入大海一般无用。难道你不知道，你自愿为了天主将所有财产用于救济穷人。因此，你应当视你的财富属于你的主，并记住你须为它的分配负责。如果你听从我，今后你将根据请求救济者的需要来调节你的捐赠。这样，你将能够扩展你的慈善范围，而你的怜悯和极其热心的关怀必得自天主的赏报。}\par

若望与许多隐修士发生过几次争论，特别是与\Person{依撒格}。他高度称赞那些安静地留在隐修院内并在那里修行的人；他保护他们免受一切不义，并热切地供应他们可能需要的任何必需品。但那些走出隐修院、出现在城市里的隐修士，他斥责他们，并认为他们侮辱了修道生活。因为这些原因，他招致了圣职人员和许多隐修士的仇恨，他们称他为严厉、暴躁、孤僻和傲慢的人。因此，他们试图公开诽谤他的生活，断言（仿佛是真事一般）他不与任何人一同进餐，并拒绝每一个邀请他吃饭的邀请。我不知道有什么借口能引起这种断言，除了——如同一位真实可靠之人向我保证的——若望由于严格的苦修而导致头部和胃部疼痛，因此无法出席一些精选的宴会。然而，由此产生了针对他的最严重的指控。\par

% structure: p0029 source=h2 target=subsection reason=relative-heading-rank; base=h2

\subsection{第十章。\Person{塞维里安}，\Place{加巴勒}主教，以及\Person{安提奥古斯}，\Place{托勒迈斯}主教。\Person{塞拉皮雍}与\Person{塞维里安}之间的争论。皇后促成了他们之间的和解。}

若望也同样因叙利亚的加巴里主教塞味黎安的阴谋而招致皇后敌视。塞味黎安与腓尼基的托勒迈城主教安提约古都是有学问之人，且极适于在教会中教导。安提约古嗓音优美、口才出众，故有人称他为\OriginalTerm{金口}{Chrysostom}。而塞味黎安则语言带有叙利亚人的粗硬，然而在学识和圣经的明证方面，他被认为优于安提约古。似乎是安提约古首先造访君士坦丁堡；他的讲道获得极大赞誉，积聚了一些财产，然后返回本城。塞味黎安步其后尘，也去了君士坦丁堡。他与若望结为密友，常在教会中讲话，受人钦佩。他享有尊荣，并结识了许多权贵以及皇帝和皇后。若望前往亚细亚时，将教会托付给他照管；因为他被塞味黎安的奉承所蒙蔽，以致以为他是自己热心的朋友。然而塞味黎安只想着取悦听众，以讲道博取民心。若望得知此事后，满怀嫉妒；据说塞辣披盎的陈述进一步点燃了他的怨恨。若望从亚细亚返回后，塞辣披盎偶然看见塞味黎安经过；他非但没有起身致意，反坐在原位，以表示对那人的极度轻蔑。塞味黎安因这轻慢的表示而恼怒，喊道：\Quote{若塞辣披盎死时仍是神职人员，那么基督便未曾降生成人。}塞辣披盎报告了这些话；于是若望将塞味黎安逐出该城，因为他傲慢无礼，亵渎天主；因为有人作证证实他确实说了上述的话。塞辣披盎的一些朋友甚至进一步篡改塞味黎安的话，声称他曾宣布基督未曾降生成人。若望也责问塞味黎安：\Quote{若塞辣披盎死时不属于神职人员，难道就是基督未曾降生成人了吗？}皇帝之妻一经塞味黎安的朋友告知所发生的事，立即派人从加采东将他召来。若望不顾皇后的一切劝阻，坚决拒绝与他交往，直到皇后将她的儿子德奥多西放在名为宗徒的教堂前面跪下；她坚持不懈地恳求他，屡次恳请他，直到若望不情愿地同意与塞味黎安和好。以上是我所收到的关于这些事件的记述。\par

% structure: p0031 source=h2 target=subsection reason=relative-heading-rank; base=h2

\subsection{第11章。关于埃及所激起的质疑：天主是否具有形体的形状。提阿菲罗，亚历山大里亚主教，与奥力振的著作。}

在此时期，埃及激起了一个问题，该问题不久前已被提出：即相信天主具有人形是否合理。由于他们单纯地抓住圣经字句而不加任何质疑，该地区大部分隐修士都持此见解；他们认为天主有眼、有脸、有手，以及身体组织的其他肢体。但那些探究圣经词语隐义的人则持相反意见；他们认为否认天主的无形体性是亵渎。后一种见解被提阿菲罗采纳，并在教会中宣讲；在他按惯例所写的关于庆祝逾越节的书信中，他借机陈述天主应被视为无形体，与人的形状相异。当此事通知给埃及隐修士，说提阿菲罗提出了这些观点时，他们前往亚历山大里亚，聚集民众于一处，引起骚动，并决定杀死这位主教，认为他是恶人。然而提阿菲罗立即出现在暴乱者面前，对他们说：\Quote{当我注视你们时，仿佛看见了天主的容貌。}这番话足以平息众人；他们压下怒气，回答说：\Quote{那么，若你真正持守正统教义，为何不谴责奥力振的著作？因为读这些书的人被导向这样的见解。}\Quote{我早有此意，}他回答说，\Quote{我将照你们的建议行事；因为我不亚于你们，同样责备所有追随奥力振学说的人。}通过这些手段，他欺骗了弟兄们，瓦解了骚乱。\par

% structure: p0033 source=h2 target=subsection reason=relative-heading-rank; base=h2

\subsection{第12章。关于被称为\Quote{长人}的四个兄弟，他们是苦修者，提阿菲罗仇视他们；关于依西多禄，以及因这四人而发生的事件。}

这场争论很可能已经结束，若不是德奥斐洛本人因对阿蒙尼乌斯、狄奥斯库鲁斯、欧瑟比乌斯和欧提米乌斯（他们被称为\Quote{\OriginalTerm{高个子}{the long}}）的敌意而重新挑起。这些人是兄弟，正如我们先前所述，他们在\Place{赛格提斯}的哲学家当中声名显赫。他们一度被\Person{德奥斐洛}爱戴，胜于埃及所有其他隐修士；他寻求他们的陪伴，并常与他们同住。他甚至将\Place{赫尔莫波利斯}的主教职授予\Person{狄奥斯库鲁斯}。他对他们的仇恨因对\Person{伊西多尔}的敌意而加深，他曾在\Place{君士坦丁堡}\Person{内克塔里乌斯}之后试图祝圣\Person{伊西多尔}。有人说，一个属于摩尼教异端的妇人皈依了天主教会的信仰；\Person{德奥斐洛}斥责了总司铎（他对他还有其他怨恨的原因），因为他在她弃绝先前异端之前就允许她参与圣事。\Person{伯多禄}（这是总司铎的名字）坚称他是按照教会的法规并得到\Person{德奥斐洛}的同意才接纳这妇人进入共融的；并引\Person{伊西多尔}为证，证明他所述属实。\Person{伊西多尔}当时恰好在\Place{罗马}出使；但回来后，他证实\Person{伯多禄}的断言是真的。\Person{德奥斐洛}将此声明视为诽谤，并将他与\Person{伯多禄}一同逐出教会。这是某些人对这件事的描述。然而，我从一位确实可靠的人（他与上述隐修士非常亲密）那里听说，\Person{德奥斐洛}对\Person{伊西多尔}的敌意源于两个原因。原因之一与司铎\Person{伯多禄}所指出的相同，即\Person{伊西多尔}拒绝见证一份将遗产留给\Person{德奥斐洛}妹妹的遗嘱；另一个原因据此人称，是\Person{伊西多尔}拒绝交出托付给他用于救济穷人的某些款项，而这些钱\Person{德奥斐洛}想要用于建造教堂；\Person{伊西多尔}说，救济受苦之人的身体（这些身体更应被视为天主的殿宇）比为建筑墙壁而建造教堂更好。无论\Person{德奥斐洛}的敌意源于何因，\Person{伊西多尔}在绝罚后立即加入了他昔日的同伴，即\Place{赛格提斯}的隐修士们。\Person{阿蒙尼乌斯}和另外几人于是去见\Person{德奥斐洛}，恳求他恢复\Person{伊西多尔}与教会的共融。\Person{德奥斐洛}爽快地答应按他们所求而行；但随着时间的流逝，他们毫无进展，且明显看出\Person{德奥斐洛}在敷衍，他们就再次去找他，重新恳求，并敦促他信守承诺。\Person{德奥斐洛}非但没有遵从，反而将其中一位隐修士投入监狱，以恐吓其他人。但他错了。\Person{阿蒙尼乌斯}和他所有的同伴便前往监狱，狱卒以为他们是来给囚犯送食物的，便轻易让他们进入；但他们一进去，就拒绝离开监狱。\Person{德奥斐洛}听说他们自愿被监禁，便派人叫他们来见他。他们回答说，\Quote{他应当先亲自将他们带出监狱，因为他们已经遭受了公开的侮辱，不应该私下获得释放。}最后，他们终于让步，去见他。\Person{德奥斐洛}为所发生的事道歉，并打发他们离去，仿佛不再打算骚扰他们；但他独自一人时，咬牙切齿，非常恼怒，决心加害他们。然而，他不知道如何加害他们，因为他们没有财产，除了哲学之外轻视一切，直到他想起扰乱他们退隐的安宁。从以前与他们的交往中，他得知他们指责那些相信天主具有人形的人，并且他们坚持\Person{奥利振}的观点；于是他让他们与持有相反观点的众多隐修士发生冲突。隐修士中爆发了可怕的争执，因为他们不想用激烈争论有条理地说服对方，而是陷入了相互侮辱。他们将主张天主无形体的称为奥利振派，而持相反意见的则被称为神人同形论者。\par

% structure: p0035 source=h2 target=subsection reason=relative-heading-rank; base=h2

\subsection{第十三章。这四人因\Person{若望}的声望而投奔他；因此，\Person{德奥斐洛}被激怒，准备与\Person{若望}作战。}

\Person{狄奥斯柯若斯}、\Person{阿摩尼乌斯}与其他隐修士发现了\Person{提阿斐洛}的阴谋，便退至\Place{耶路撒冷}，又从那里前往\Place{斯基托波利斯}；他们认为那里有许多棕榈树，隐修士们常以棕榈叶进行劳作，因此居住在那里对他们有利。\Person{狄奥斯柯若斯}和\Person{阿摩尼乌斯}约有八十名其他隐修士陪同。与此同时，\Person{提阿斐洛}派遣信使前往\Place{君士坦丁堡}，控告他们，并阻止他们可能向皇帝呈上的任何请愿。\Person{阿摩尼乌斯}和隐修士们得知此事后，便乘船前往\Place{君士坦丁堡}，并带上了\Person{依西多禄}；他们请求在皇帝和主教面前审理此案，因为他们认为，\Person{若望}为人正直、敢作敢为，必能帮助他们伸张正义。\Person{若望}虽以善意接待他们，予以尊重，并未禁止他们在教堂祈祷，但拒绝让他们参与\OriginalTerm{奥迹}{\GreekText{μυστήρια}}，因为在调查之前这样做是不合法的。他写信给\Person{提阿斐洛}，希望他重新接纳他们共融，因为他们的神性见解是正统的；并请求他，若对他们的正统性有疑虑，可派人前来控告。\Person{提阿斐洛}对此信未予回复。此后不久，\Person{阿摩尼乌斯}及其同伴在皇后骑马出行时拦住她，控诉\Person{提阿斐洛}对他们的阴谋。皇后知道他们所遭遇的陷害；她站起来以示尊敬，从御辇上俯身点头，对他们说：\Quote{请为皇帝、为我、为我们的子女、并为帝国祈祷。至于我，我将很快召开一次公议会，届时\Person{提阿斐洛}也将被传唤。}随后，在\Place{亚历山大里亚}流传着一个虚假的报告，称\Person{若望}已接纳\Person{狄奥斯柯若斯}及其同伴共融，并尽全力给予他们一切帮助和鼓励。\Person{提阿斐洛}开始思考采取何种措施才能将\Person{若望}从主教职位上赶走。\par

% structure: p0037 source=h2 target=subsection reason=relative-heading-rank; base=h2

\subsection{第十四章 \Person{提阿斐洛}的乖戾。\Person{圣埃皮法尼乌斯}：他在\Place{君士坦丁堡}的居留及准备煽动民众反对\Person{若望}。}

德奥斐洛对若望的计谋保持尽可能秘密，并致信各城主教，谴责奥利振的著作。他又想到，若能争取到塞浦路斯撒拉米斯的主教厄丕法尼乌斯站在自己一边，则颇为有利；此人因其生活而受人尊敬，是同时代中最杰出的人物。于是他与他结为朋友，尽管此前他曾责备厄丕法尼乌斯主张天主具有人形。德奥斐洛如同后悔曾有任何不同见解，致信厄丕法尼乌斯，告知他自己如今持有与厄丕法尼乌斯相同的意见，并发动对奥利振著作的攻击，视之为如此邪恶教义的根源。厄丕法尼乌斯长久以来对奥利振的著作怀有特别的厌恶，因而轻易相信了德奥斐洛的书信。他不久后召集塞浦路斯的主教们，禁止审查奥利振的著作。他还致信其他主教，其中包括君士坦丁堡的主教，劝他们召集会议，作出同样的决定。德奥斐洛见追随厄丕法尼乌斯的榜样并无危险，而厄丕法尼乌斯为民众所称赞，因其生活的德行受人钦佩（无论其意见如何），便与所属主教一致通过了与厄丕法尼乌斯相似的决议。另一方面，若望对厄丕法尼乌斯和德奥斐洛的书信不大在意。那些与他为敌的权贵和神职人士觉察到德奥斐洛的计谋意在使若望被驱逐出主教之职，于是尽力促成在君士坦丁堡召开一次会议，以执行这一措施。德奥斐洛知道此事后，竭尽全力召集这次会议。他命令埃及的主教们乘船前往君士坦丁堡；他写信请求厄丕法尼乌斯和其他东方主教尽快前往该城，而他自己则从陆路出发。厄丕法尼乌斯首先从塞浦路斯出发，在君士坦丁堡的郊区赫布多摩斯上岸；他在那里的教堂祈祷后，便进入城中。为表敬意，若望带领全体神职人员出城迎接他。然而，厄丕法尼乌斯的行为清楚表明他相信了对若望的指控；因为，尽管受邀住在教会住所，他却不肯留在那里，并拒绝在那里与若望会面。他还私下召集所有住在君士坦丁堡的主教，向他们展示了他对奥利振言论所颁布的法令。他说服一些主教赞同这些法令，另一些则反对。斯基泰的主教德奥提摩斯强烈反对厄丕法尼乌斯的行动，告诉他说，侮辱一位早已去世之人的记忆是不对的；而且，攻击古人就此问题所达成的结论并废除他们的决定，也不免带有亵渎意味。他一边这样说，一边拿出一本随身携带的奥利振的著作，大声诵读了一段有益于教会教育的段落，并指出那些谴责这些言论的人行为荒谬，因为他们有侮辱这些言论所谈论对象本身的危险。若望仍对厄丕法尼乌斯保持尊敬，邀请他参加自己教会的聚会，并与他同住。但厄丕法尼乌斯宣称，除非若望谴责奥利振的著作并驱逐狄奥斯库若及其同伴，否则他既不与若望同住，也不与他公开祈祷。若望认为在案件作出判决之前不宜如此行事，便试图推迟此事。当会议即将在宗徒教堂举行时，那些对若望怀有恶意的人策划让厄丕法尼乌斯先去公开向民众谴责奥利振的著作，并斥责狄奥斯库若及其同伴为该作者的同伙；还要攻击该城主教为那些异端者的支持者。有些人参与了此事；因为以此手段，人们认为民众对主教的爱戴将会被疏离。次日，当厄丕法尼乌斯正要进入教堂以实施其计划时，他被塞拉皮雍拦住，这是奉若望之命，因为若望已得知这一阴谋。塞拉皮雍向厄丕法尼乌斯证明，他所筹划的计划本身是不义的，对他个人也无益；因为若它激起民众骚乱，他将被视为随后可能发生的暴行的责任人。厄丕法尼乌斯被这些论据说服，放弃了攻击。\par

% structure: p0039 source=h2 target=subsection reason=relative-heading-rank; base=h2

\subsection{第十五章。皇后的儿子与圣厄丕法尼乌斯。\Quote{长弟兄们}与厄丕法尼乌斯的会谈，及其重新登船前往塞浦路斯。厄丕法尼乌斯与若望。}

约在此时，皇后之子患重病，皇后忧惧后果，遣人恳请\Person{Epiphanius}为之祈祷。\Person{Epiphanius}回报说，若她能避免与异端者\Person{Dioscorus}及其同伙交往，病者即可存活。皇后对此回覆如下：\Quote{若天主愿意带走我的儿子，愿祂的旨意成就。赐我孩儿的主，也能将他收回。你并无能力复活死者，否则你的总执事便不会死了。}她暗指不久前去世的总执事\Person{Chrispion}。他是隐修士\Person{Fuscon}和\Person{Salamanus}的兄弟，我在详述\Person{Valens}统治期间事件时曾提及他们；他原为\Person{Epiphanius}的同伴，被任命为他的总执事。\Person{Ammonius}及其同伴在皇后允许下前往\Person{Epiphanius}处。\Person{Epiphanius}询问他们是谁，\Person{Ammonius}答道：\Quote{神父，我们是\OriginalTerm{长弟兄}{Long Brothers}；我们恭敬前来，想知道您是否读过我们的著作或我们弟子的作品？}\Person{Epiphanius}回答未曾见过，他继续说：\Quote{那么，您既无证据表明我们持何种见解，为何视我们为异端者？}\Person{Epiphanius}说他只是根据听闻判断；\Person{Ammonius}回答说：\Quote{我们的行为与您大不相同。我们与您的弟子交谈，常读您的著作，其中包括名为\WorkTitle{锚定}之作。我们遇到有人嘲笑您的见解，宣称您的著作充满异端时，我们为您辩护，将您视为神父。难道您应仅凭传闻就谴责不在场、且您毫无确知之人，或如此回报那些曾赞誉您的人？}\Person{Epiphanius}略有醒悟，便让他们离去。不久后他启程前往\Place{塞浦路斯}，或许因为他意识到赴\Place{君士坦丁堡}之行的徒劳，或者有理由相信，天主已启示他死期将近；因他在返回\Place{塞浦路斯}的航程中去世。据传他对陪同至登船处的主教们说：\Quote{我将城市、宫殿和舞台留给你们，因为我即将离去。}我从数人处得知，\Person{John}预言\Person{Epiphanius}将死于海上，而后者预言了\Person{John}的被废。因在他二人争执最烈时，\Person{Epiphanius}对\Person{John}说：\Quote{我愿你不会以主教身份而死，}\Person{John}回答：\Quote{我愿你永远回不了你的主教区。}\par

% structure: p0041 source=h2 target=subsection reason=relative-heading-rank; base=h2

\subsection{第十六章　皇后与\Person{John}之争。\Person{Theophilus}自埃及抵达。\Person{Cyrinus}，\Place{加尔西顿}主教。}

\Person{Epiphanius}离去后，\Person{John}照常在教堂讲道时，偶然痛斥女性易犯的恶习。众人以为他的批评影射皇帝之妻。主教的仇敌当然将此话按此意禀报皇后；她自觉受辱，向皇帝抱怨，力促\Person{Theophilus}尽快前来并召开会议。\Place{加巴拉}主教\Person{Severian}尚未放下对\Person{John}的旧怨，也配合推动此事。我并无足够资料断定当时流行之说是否真实，即\Person{John}上述讲道确指皇后，因他怀疑她煽动\Person{Epiphanius}反对自己。\Person{Theophilus}不久后抵达\Place{彼提尼亚}的\Place{加尔西顿}，许多主教随后赶到。有些主教应他邀请而来，另一些则遵皇帝之命。\Person{John}在亚洲所革职的主教们欣然前往\Place{加尔西顿}，凡对他怀有敌意者也同样前往。\Person{Theophilus}从埃及期待的船只已到\Place{加尔西顿}。当他们再次在同一地点聚集，商议如何明智推进对\Person{John}的图谋时，\Place{加尔西顿}教会领袖\Person{Cyrinus}——他是埃及人，\Person{Theophilus}的亲戚，且与\Person{John}有其他过节——对他大肆辱骂。然而，公义似乎很快追随了他：陪同主教们的\Place{美索不达米亚}本地人\Person{Maruthas}恰巧踩到他的脚；\Person{Cyrinus}因这一意外剧痛难忍，无法与其他主教一同前往\Place{君士坦丁堡}，尽管他的协助对实现针对\Person{John}的阴谋必不可少。伤口情况危急，外科医生不得不在腿上多次手术；最终坏疽发生，蔓延全身，甚至扩展到另一只脚。他不久后在剧痛中去世。\par

% structure: p0043 source=h2 target=subsection reason=relative-heading-rank; base=h2

\subsection{第十七章　\Person{Theophilus}与\Person{John}的控告者在\Place{Rufinianæ}举行会议。\Person{John}被传唤出席，未到会，因而被他们废黜。}

\Person{德奥斐洛}进入\Place{君士坦丁堡}时，没有一位圣职人员出去迎接他，因为他与主教为敌已是众所周知。然而，一些来自\Place{亚历山大里亚}的水手，恰巧在岸上，既有来自运粮船的，也有来自其他船只的，聚集在一起，以极大的欢呼迎接了他。他绕过圣堂，径直前往宫殿，那里已为他准备好住宿。他很快发现城中有许多人强烈反对\Person{若望}，并准备对他提出控告；于是，他采取相应措施，前往\Place{加尔西顿}郊区一个叫做\Quote{橡树园}的地方。这个地方现在以\Person{鲁菲努斯}的名字命名；因为他是一位执政官，在这里建造了一座宏伟的宫殿，并修建了一座大圣堂以尊崇宗徒\Person{伯多禄}和\Person{保禄}，因此将其命名为\OriginalTerm{宗徒堂}{Apostolium}；并指定了一个隐修士团体在该圣堂中履行圣职。当\Person{德奥斐洛}与其他主教在此集会商议时，他认为不宜再提及\Person{奥力振}的著作，并呼吁\Place{斯塞提斯}的隐修士们悔改，应许不再记念过犯，也不施以祸害。他的支持者热切地附和这一努力，并告诉他们，必须请求\Person{德奥斐洛}宽恕他们的行为；由于集会全体成员都同意这一请求，隐修士们感到不安，并认为必须按照这么多主教所期望的去做，于是使用了他们即使在受伤害时也惯常使用的话，说道：\Quote{宽恕我们吧}。\Person{德奥斐洛}欣然接纳他们，恢复他们与教会的共融；于是，关于对\Place{斯塞提斯}隐修士们所施伤害的问题就此了结。我深信，倘若\Person{狄奥斯库鲁斯}和\Person{阿摩尼乌斯}与其余隐修士一同在场，此事就不会如此迅速了结。然而\Person{狄奥斯库鲁斯}已于先前去世，被安葬在殉道者圣\Person{莫基乌斯}的圣堂内。\Person{阿摩尼乌斯}也在召集会议的准备工作进行时染病；尽管他坚持前往\Quote{橡树园}，但他的病情因此大大加重：他在旅行后不久去世，并在附近的隐修士中得到了隆重的安葬，至今安息在那里。据说，\Person{德奥斐洛}听到他的死讯时流下了眼泪，并宣称尽管他造成了许多困扰，但从未找到一位品性比\Person{阿摩尼乌斯}更高尚的隐修士。然而必须承认，这位隐修士的死极大地促进了\Person{德奥斐洛}计划的成功。\par

会议成员传唤\Place{君士坦丁堡}的所有圣职人员到他们面前，并威胁要罢免那些不服从传唤的人。他们传唤\Person{若望}前来应诉；同样也传唤了\Person{塞拉皮翁}、司铎\Person{提格里乌斯}和读经员\Person{保禄}。\Person{若望}通过\Place{皮西努斯}主教\Person{德默特琉斯}以及其他几位是他朋友的神职人员的媒介，告知他们：他不会逃避调查，但若将控告者的姓名和控告内容告知于他，他愿意在一个更大的会议面前为自己的行为辩护；因为他不愿被视为疯狂，也不愿承认那些显而易见的敌人为审判官。主教们对\Person{若望}的不服从表现出极大的愤慨，以致他派去会议的一些神职人员受到恐吓，没有回到他那里。\Person{德默特琉斯}以及那些将他的利益置于一切考虑之上的人离开了会议，回到了他身边。同日，从宫中派出一名信使和一名速记员，命令\Person{若望}到主教们那里去，并催促主教们不再拖延地裁决他的案件。在\Person{若望}被传唤了四次，并上诉至大公会议之后，除了他拒绝服从会议的传唤之外，没有其他指控能得到证实；基于这一理由，他们罢免了他。\par

% structure: p0046 source=h2 target=subsection reason=relative-heading-rank; base=h2

\subsection{第十八章 人民对德奥斐洛的叛乱；他们诽谤统治者。若望被召回，再次就任主教座。}

\Place{君士坦丁堡}的人民在傍晚时分得知了会议的谕令；他们立即起来叛乱。天一亮，他们便跑进圣堂，在其许多计划中呼喊，应当召开一个更大的会议来审理此事；他们阻止了皇帝派来将\Person{若望}押送流放的官员执行谕令。\Person{若望}担心会以他违抗皇帝命令或煽动民众叛乱为借口而对他提出另一项指控，于是在人群散去后，于他被罢免后第三天中午，悄悄从圣堂逃走了。当人民得知他已流亡时，叛乱变得严重，并发表了许多侮辱性的言论反对皇帝和会议；特别是针对\Person{德奥斐洛}和\Person{塞维里安}，因为他们被视为阴谋的策划者。\Person{塞维里安}恰好在这些事件发生时在圣堂中讲道；他借机赞扬对\Person{若望}的罢免，并声称，即使假设他在其他罪行上无罪，\Person{若望}也因其骄傲而应当被罢免；因为，尽管天主乐意赦免人的一切其他罪过，但他却反对骄傲的人。这一言论使人民在委屈中变得焦躁，重新燃起愤怒，陷入无法控制的叛乱。他们跑到圣堂、市场，甚至皇帝宫殿，以嚎哭和哀叹要求召回\Person{若望}。皇后最终被他们的迫切要求所打动；她劝说丈夫顺从人民的意愿。她迅速派出一位她信任的宦官，名叫\Person{布里索}，将\Person{若望}从\Place{比提尼亚}的\Place{普雷内图斯}城接回；并声明她并未参与针对他的阴谋，相反，她一直尊敬他，视他为司铎和其子女的入门引导者。\par

当\Person{若望}踏上归途，抵达属皇后城郊时，他在\Place{阿纳普卢斯}附近停下，拒绝重入城内，除非有更大规模的主教会议承认他受废黜的不公。但此一拒绝加剧了民众的骚动，并导致许多公开抨击皇帝和皇后的言论，于是他听从劝说进入城中。民众前去迎接他，唱着与当前境况相关的圣咏；许多人手持蜡烛。他们引导他前往教堂；虽然他拒绝，并屡次声称那些定他罪的人应先重新考虑他们的投票，但他们强迫他登上主教宝座，并按照司铎的惯例向民众致平安。随后他作了一次即席讲论，其中以巧妙的比喻声称，\Person{狄奥菲鲁斯}曾图谋伤害他的教会，正如\Place{埃及}王意图玷污圣祖\Person{亚巴郎}的妻子\Person{撒辣}，此事记载于希伯来人的书中：然后他称赞民众的热忱，并颂扬皇帝和皇后对他的善意；他激起民众对皇帝及其配偶的热烈掌声和称颂，以至于他不得不半途中断讲论。\par

% structure: p0049 source=h2 target=subsection reason=relative-heading-rank; base=h2

\subsection{第十九章　\Person{狄奥菲鲁斯}的顽固。\Place{埃及}人与\Place{君士坦丁堡}市民之间的敌意。\Person{狄奥菲鲁斯}的逃亡。隐修士\Person{尼拉孟}。关于\Person{若望}的主教会议。}

尽管\Person{狄奥菲鲁斯}想以\Person{若望}非法自行复任主教职为由提出控告，但因惧怕触怒皇帝而作罢——皇帝是被迫召回\Person{若望}以平息民众叛乱的。然而，\Person{狄奥菲鲁斯}在被告缺席期间接受了针对\Person{赫拉克利德}的控告，希望借此使已经发出的对\Person{若望}的定罪判决具有正当性。但\Person{赫拉克利德}的朋友们介入，宣称缺席定罪是不公正的，且违反教会法规。\Person{狄奥菲鲁斯}及其同党坚持相反观点：\Place{亚历山大里亚}和\Place{埃及}的人民支持他们，而\Place{君士坦丁堡}市民则反对他们。两派之间的冲突变得如此激烈，以致发生流血事件：许多人受伤，还有人在争斗中丧生。\Person{塞维里安}和当时在\Place{君士坦丁堡}所有不支持\Person{若望}事业的主教们因担心人身安全，匆忙离开了该城。\Person{狄奥菲鲁斯}也在初冬逃离该城，与隐修士\Person{依撒格}一同乘船前往\Place{亚历山大里亚}。一阵风将船吹至\Place{革辣}——距\Place{培琉喜阿姆}约五十斯塔狄亚的一座小城。该城的主教去世，据我所知，居民选举\Person{尼拉孟}主持他们的教会；他是个善人，已臻于隐修哲学之巅峰。他住在城外一间用石块砌上门的小室中。他拒绝接受司铎职的尊荣；因此\Person{狄奥菲鲁斯}亲自拜访他，劝他接受自己的祝圣。\Person{尼拉孟}屡次拒绝这一荣誉；但因\Person{狄奥菲鲁斯}不肯接受拒绝，他对他说：\Quote{父亲，明天您可按您的意愿行事；今天我需要安排我的事务。}第二天，\Person{狄奥菲鲁斯}来到隐修士的小室，命令打开门；但\Person{尼拉孟}喊道：\Quote{让我们先祈祷吧。}\Person{狄奥菲鲁斯}同意，开始祈祷。\Person{尼拉孟}也在小室内祈祷，并在祈祷中去世。\Person{狄奥菲鲁斯}和那些站在小室外的人当时对此一无所知；但当大半天过去，\Person{尼拉孟}的名字被反复大声呼喊却无人应答时，门上的石块被移开，发现隐修士已死。他们给他穿上必要的衣袍，为他举行了公开葬礼；居民在他的坟墓周围建了一座祈祷之所；直到今日，他们仍以非常隆重的方式纪念他的逝世之日。\Person{尼拉孟}就这样去世了——如果离开此生进入另一生命可以称为死亡的话——而不愿接受一个以非凡谦逊自认为不配的主教职。\par

返回\Place{君士坦丁堡}后，\Person{若望}似乎比以往更受人民爱戴。六十位主教在该城集会，废除了\Quote{橡树}会议的所有法令。他们确认\Person{若望}保有主教职，并规定他应担任司铎、施行祝圣、并履行通常由教会首领承担的一切教会职责。此时，\Person{色拉比翁}被任命为\Place{色雷斯赫拉克里亚}的主教。\par

% structure: p0052 source=h2 target=subsection reason=relative-heading-rank; base=h2

\subsection{第二十章　皇后的雕像；那里发生的事；\Person{若望}的教导；召集另一次反对\Person{若望}的主教会议；他的废黜。}

这些事件发生后不久，皇后的银质雕像——至今仍可见于教堂南边，正对\Place{大会堂}——被竖立在一根斑岩柱上，置于高台之上；此事在当时皇帝雕像竖立时惯常的舞蹈和哑剧等民众表演的喝彩声中庆祝。在一次对民众的公开讲论中，\Person{若望}指斥这些行为有损教会荣誉。此言使皇后忆起旧怨，她因受辱而异常愤怒，决心召集另一次主教会议。他并未让步，反而更公开地在教堂中抨击她，火上浇油；正是在这一时期，他发表了那篇著名的讲论，开头的话是：\Quote{\Person{黑落狄雅}再次发怒；再次跳舞；再次要求将\Person{若望}的头放在盆中。}\par

不久之后，有几位主教抵达\Place{君士坦丁堡}，其中有\Place{安西拉}主教\Person{良提乌}和\Place{贝雷亚}主教\Person{阿卡西}。那时主的圣诞临近，皇帝没有像往常一样前往教堂，反而派人通知\Person{若望}，除非他澄清所受的指控，否则不能与他共融。\Person{若望}勇敢地回答说，他愿意证明自己的清白；这让指控者十分畏惧，以致不敢继续追究。审判者们决定，既然他已被罢免一次，不应接受第二次审判。但他们要求\Person{若望}只在此点上为自己辩护：在他被罢免之后、在会议恢复他职位之前，他是否坐上了主教座位。\Person{若望}在辩护中援引了那些在\Quote{橡树}会议之后仍与他共融的主教们的决定。审判者们驳回了这一论点，理由是那些与\Person{若望}共融的人数少于罢免他的人，且有一项教会法已定他的罪。他们以此为借口罢免了他，尽管该法律是由异端者制定的；因为亚略派在利用各种诽谤将\Person{亚大纳削}逐出\Place{亚历山大}教会之后，出于对政局变化的忧虑，制定了这项法律，企图使针对他的判决不受调查。\par

% structure: p0055 source=h2 target=subsection reason=relative-heading-rank; base=h2

\subsection{第21章。\Person{若望}被逐后人民所受的灾祸。针对他的暗杀阴谋。}

\Person{若望}被免职后，不再在教堂举行聚会，而是安静地留在主教住所。四旬期结束时，正值每年纪念\Person{基督}复活的圣夜，\Person{若望}的追随者被士兵和他的敌人逐出教堂，士兵们袭击了正在庆祝神圣礼仪的民众。由于这一事件出乎意料，洗礼池旁发生了极大骚乱。妇女们哭泣哀叹，孩童尖叫；司铎和执事遭到殴打，被强行从教堂中赶出，他们当时还穿着主持礼仪的圣服。他们被指控犯下了那些已领受圣事者所能想象的种种无序行为，但我认为有必要对此缄默不言，以免我的著作落入未受启迪者手中。\par

当民众识破这阴谋后，次日他们未使用教堂，而是在以\Person{君士坦狄乌斯}皇帝命名的极其宽敞的公共浴室中庆祝逾越节庆典。主教和司铎，以及有权管理教会事务的其他人主持了礼仪。支持\Person{若望}的人与民众一同在场。然而，他们被逐出那里，随后聚集在城外一处空地，该地是\Person{君士坦丁}皇帝在城市建设前令人清理并用栅栏围起来的，用于举行赛马比赛。从那时起，民众分开聚会，有时只要可行就在那地点，有时在别处。他们被称为若望派。大约此时，有一个被鬼附身或假装被附身的人被抓获，身上带着匕首，意图刺杀\Person{若望}。他被民众当作被雇来参与这阴谋的人抓住，带到总督面前；但\Person{若望}派了他一方的几位主教，在审讯拷问之前将其释放。此后不久，司铎\Person{厄尔彼狄乌斯}的一名奴隶——他对执事公开怀有敌意——被看见拼命向主教住所跑去。一个路人试图拦住他，询问为何如此匆忙；但奴隶非但不回答，反而用匕首刺伤了他。另一个站在旁边的人，见人受伤而呼喊，也遭到奴隶类似刺伤；第三个旁观者同样受伤。附近所有民众看到所发生的事，大声呼喊要逮捕奴隶。他转身逃跑。追赶者向前方的人喊叫要求抓住逃犯，这时一个刚从浴室出来的人试图拦住他，被重伤倒地，当场死亡。最后，民众设法包围了奴隶。他们抓住他，将他押送到皇宫，声称他意图刺杀\Person{若望}，罪行应受惩罚。总督将罪犯关押起来，并保证会依法惩处，以此平息了民众的愤怒。\par

% structure: p0058 source=h2 target=subsection reason=relative-heading-rank; base=h2

\subsection{第22章。\Person{若望}被非法逐出主教职。随后的骚乱。教堂被天火焚毁。\Person{若望}被流放至\Place{库库苏斯}。}

从此时起，最热心的民众轮流守护\Person{若望}，昼夜在主教的住所周围站岗。宣判\Person{若望}的主教们指责这种行为违反教会法规，声称他们能为自己所定的判决之公正性负责，并断言除非将他逐出城市，否则民众无法恢复平静。一位信使将皇帝命他立即离开的谕旨传达给他，\Person{若望}服从了，在那些被指派看守他的人不注意时逃离了城市。他唯一责备的是：在没有合法审判或任何法律程序的情况下被流放，他受到的待遇比杀人犯、巫师和奸淫者更严厉。他被一艘小船送到\Place{比提尼亚}，然后立即继续行程。他的一些敌人担心民众得知他离开后会追赶并将他强行带回，因此命令关闭教堂的大门。当身处城市公共场所的民众得知所发生的事时，引发了极大的混乱；一些人跑向海边仿佛要追随他，另一些人四处奔逃，惊恐万分，因为皇帝的愤怒预计会因他们制造如此大的骚动和喧哗而降临到他们身上。那些在教堂内的人因互相拥挤着冲向门而进一步堵住了出口。他们费了很大劲才用暴力撞开门；一派人用石头砸门，另一派人将门向内拉，从而将人群向后挤入建筑物。与此同时，教堂突然从四面燃起大火。火焰向四处蔓延，教堂南侧相邻的元老院会议大厦也被烧毁。双方互相指责对方纵火。\Person{若望}的敌人断言他的党羽出于报复，因为会议通过了反对他的决议而犯下此罪行。另一方则声称他们被诬蔑，纵火是他们的敌人所为，意图将他们烧死在教堂里。大火从傍晚一直燃烧到早晨，并向尚未烧毁的材料蔓延时，负责看管\Person{若望}的官员将他送到\Place{亚美尼亚}的\Place{库库苏斯城}——皇帝通过谕旨指定这里为被定罪者的居住地。其他官员奉命逮捕所有支持\Person{若望}事业的主教和圣职人员，并将他们囚禁在\Place{加采东}。那些被怀疑效忠于\Person{若望}的市民被搜出并投入监狱，被迫宣布弃绝\Person{若望}。\par

% structure: p0060 source=h2 target=subsection reason=relative-heading-rank; base=h2

\subsection{第二十三章 \Person{阿尔萨基乌斯}被选继任\Person{若望}。对\Person{若望}追随者所行之恶。圣\Person{尼卡瑞特}。}

不久之后，\Person{阿尔萨基乌斯}——即此前管理主教区的\Person{内克塔里乌斯}的兄弟——被祝圣为\Place{君士坦丁堡}主教。他性情非常温和，虔诚至极；但他作为司铎所获得的声誉因一些被他委托权力的圣职人员的行为而受损，他们在他的名义下为所欲为；因为他们的恶行被归咎于他。然而，没有什么比针对\Person{若望}追随者的迫害更对他不利。他们拒绝与他共融，甚至拒绝与他一起祈祷，因为\Person{若望}的敌人与他联合；当他们如我们先前所述，坚持在城郊另立教堂时，他向皇帝控告了他们的行为。护民官奉命率兵攻击他们，很快用棍棒和石块驱散了他们。其中地位最显赫以及最热忱支持\Person{若望}的人被投入监狱。士兵们按照此类场合的惯例，越权行事，暴力剥去妇女的饰品，掠夺她们的锁链、金腰带、项链和戒指项圈；她们扯掉耳垂连同耳环。虽然全城充满了忧愁和哀哭，民众对\Person{若望}的热爱仍一如既往，他们避免在公共场合露面。许多人远离市场和公共浴场，另一些人觉得在自己家中也不安全，便逃离了城市。\par

在采纳后一措施的热忱男子和优秀女子中，有\Place{彼提尼雅}的\Person{尼卡莱特}女士。她出身显赫的贵族家庭，因永葆童贞和德行高尚而闻名。她是我们所认识的一切热忱女子中最谦逊者，在言行举止上皆井井有条，终其一生，她始终将事奉天主置于一切世俗考量之上。面对逆境突然逆转，她表现出勇敢而深思的承受能力；虽见自己丰厚的祖产大部分被不义地剥夺，却未显任何愤慨，而将余下的小部分管理得如此节俭，以致年事已高，仍能供给全家的需要，并慷慨济助他人。因她性喜人道，也为贫病者之需备办各种药物，常能治愈那些经普通医生诊治无果的病人。她以虔诚的力量获得最佳成效，紧闭双唇。总之，我们从未见过一位虔敬的妇女具有如此礼仪、庄重及一切美德。她虽如此卓越，却隐藏了自己大部分的本性和事迹；因她本性谦逊且富有哲学修养，常致力于隐晦。她不肯接受女执事之职，也不愿担任奉献于教会服务的贞女的训导者，因为她自觉不配，虽然\Person{若望}多次将此荣耀强加于她。\par

民众暴动平息后，该城长官公开露面，佯装调查火灾及议事厅焚毁的原因，并严厉惩罚了许多人；但因他是异教徒，竟嘲笑教会的灾难，并对它的不幸感到欣喜。\par

% structure: p0064 source=h2 target=subsection reason=relative-heading-rank; base=h2

\subsection{二十四、\Person{欧特洛彼}诵读人、\Person{真福奥琳皮安}及\Person{司铎提格里}因对\Person{若望}的依恋而受迫害。牧首。}

\Person{欧特洛彼}诵读人受命指认纵火焚烧教堂者；但尽管他遭受严刑鞭打，肋旁和面颊被铁钉撕裂，点燃的火把烫及他身上最敏感的部位，却未能从他口中逼出任何供认，尽管他年轻且体质纤弱。受此酷刑之后，他被投入地牢，不久便在那里去世。\par

一个关于欧特洛彼的梦似乎值得载入本史。\OriginalTerm{诺瓦替安派}{Novatians}的主教\Person{西西尼}在梦中看见一个容貌出众、身材高大的人站在诺瓦替安派为尊崇首位殉道者\Person{斯德望}而建造的教堂的祭台旁；那人抱怨好人稀少，说他找遍了全城，只找到一个好人，那就是\Person{欧特洛彼}。\Person{西西尼}对所见的景象感到惊奇，便将梦境告知他教会中最忠信的\OriginalTerm{司铎}{presbyters}，并命他无论欧特洛彼在哪里，都要去找他。那位司铎正确推测，此欧特洛彼必是无他，正是被长官野蛮折磨的那一位，于是一个监狱接一个监狱地寻找他。终于找到了他，并在交谈中告知了主教的梦，含着泪恳求他为自己祈祷。这就是我们所掌握的关于欧特洛彼的详情。\par

在这些灾难中，女执事\Person{奥琳皮亚}显示了极大的刚毅。为此她被拖到法庭，长官审问她纵火烧教堂的动机，她回答说：\Quote{我过去的生活理应消除一切嫌疑，因为我已将自己庞大的财产用于修缮天主的圣殿。}长官声称他非常了解她过去的生活。她继续说：\Quote{那么，你应当站在控告人的位置，让另一个人来审判我们。}由于对她的指控完全缺乏证据，长官发现没有正当理由指责她，便转而用较轻的罪名，仿佛想劝告她，责备她和其他妇女拒绝与他所支持的主教共融，尽管她们有可能忏悔并改变自己的处境。她们出于恐惧都听从了长官的劝告，但\Person{奥琳皮亚}对他说：\Quote{这不公平：我被人公开诽谤，法庭上没有任何针对我的指控被证实，却要我为自己开脱那些与本案完全无关的罪名。还是让我就原告对我提出的最初指控进行辩护吧。即使你动用非法强制手段，我也不愿与我应当分离的人共融，也不愿同意任何对虔诚者不合法的事。}长官见她无法被说服与\Person{阿尔撒奇}共融，便打发她去咨询律师。然而，另一天，他又派人将她召来，判处她巨额罚款，因为他以为这样会迫使她改变主意。但她对财产的损失完全置之不理，离开君士坦丁堡前往\Place{基齐科}。大约同一时期，司铎\Person{提格里}被剥去衣服，背受鞭打，手足被缚，并被吊上刑架。他本是蛮族，且是阉人，但并非天生。他原是一位权贵家中的奴隶，因忠心服役获得释放。后来被祝圣为司铎，以温良柔和，以及对外乡人和穷人的爱德而著称。这些就是发生在君士坦丁堡的事情。\par

其间，西里修在治理罗马主教职位十五年后逝世。阿纳斯塔修担任同一主教职位三年后去世，由依诺森继任。弗拉维盎因拒绝同意对若望的罢免也已去世；波菲里被任命接替他治理安提约基雅教会，并赞同那些谴责若望的人，叙利亚的许多信徒便脱离了安提约基雅教会；由于他们自行集会，遭受了许多残酷迫害。为强制与阿尔萨修、此波菲里以及亚历山大主教德奥斐洛共融，在宫中权贵的热心推动下，颁布了一条法律，规定正统信徒不得在教会之外集会，凡不与这些人共融者应被驱逐。\par

% structure: p0069 source=h2 target=subsection reason=relative-heading-rank; base=h2

\subsection{第二十五章. 既然教会中存在这些祸患，世俗事务也陷入混乱。论霍诺留的将军斯提利科之事。}

大约在这时期，震撼教会的纷争正如常有的情形，随之引起了国家中的骚乱和动荡。匈奴人越过伊斯特河，蹂躏了色雷斯。伊索里亚的强盗大量集结，劫掠了远至卡里亚和腓尼基的城邑村庄。霍诺留的将军斯提利科，一个若曾有任何人达到过巨大权势的人，统率着罗马和蛮族军队的精锐，对阿尔卡狄治下的官员心怀敌意，决意使两个帝国彼此为敌。他令霍诺留任命哥特人的领袖阿拉里克为罗马军队的将军，并派他进入伊利里亚；同时又派遣禁卫军长官约维乌斯前往那里，并许诺自己将率领罗马士兵在那里与他们会合，以便将那个行省并入霍诺留的领土。阿拉里克率领他的军队从达尔马提亚和潘诺尼亚交界处的蛮族地区出发，到达伊庇鲁斯；在那里等待了一段时间后，他返回了意大利。斯提利科未能履行与阿拉里克会合的约定，因为霍诺留给他送来的几封信阻止了他。这些事件就是这样发生的。\par

% structure: p0071 source=h2 target=subsection reason=relative-heading-rank; base=h2

\subsection{第二十六章. 罗马主教依诺森的两封书信，一封致金口若望，另一封致君士坦丁堡的圣职人士，论若望之事。}

罗马主教依诺森得知对若望采取的措施后，极其愤慨，并谴责了全部程序。随后他转而关注召集大公会议，并写信给若望，部分地也写给君士坦丁堡的圣职人士。以下附录这两封信，正如我找到它们时那样，从拉丁文译成了希腊文。\par

依诺森致可爱的弟兄若望。\par

\Quote{虽然一个自知自己无辜的人理应期望一切福佑，并向天主祈求怜悯，但我们认为，通过执事齐里亚古给您送去一封相称的信函，规劝您恒心忍耐，似乎是适宜的，以免加于您的侮辱在挫败您的勇气方面，比良心的证言在鼓励您抱有希望方面更有力量。无需教导您这位如此伟大子民的导师和牧者：天主总是考验最优秀的人，看他们是否愿意留在忍耐的高峰，并不向任何受苦的劳作让步；而且，良心是如何坚固地对抗一切不义地临到我们的事，除非人在这些不幸中被忍耐所感动，否则他就为恶意的猜疑提供了根据。因为那首先信赖天主、其次信赖自己良心的人，应当忍受一切。尤其当一位卓越而良善的人能够在忍耐中操练自己时，他不能被战胜；因为圣经护卫着他的思想，而我们对民众讲解的虔诚读经，则充满榜样。这些圣经向我们保证，几乎所有的圣人都以不同的方式持续受苦，并通过某种考验而被试炼，从而进入了忍耐的冠冕。愿您的良心鼓励您的爱心，极可敬的弟兄；因为那种官能在患难中拥有对德行的鼓励。既然基督，我们的主，正在观察，被净化的良心必将您安放在平安的港湾。}\par

依诺森主教致在若望主教治下的君士坦丁堡教会的司铎、执事、全体圣职人士和信众：问候你们，亲爱的弟兄们。\par

从你通过司铎日耳曼和执事卡西安转交给我的爱心书信中，我以忧心忡忡得知你展现在我们眼前的邪恶景象。我反复阅读时，屡次看到信德被何等灾难与劳苦所困。唯有忍耐之安慰才能医治此种情形。我们的天主要尽快终结这样的磨难，而这磨难终将有助于你们的益处。但我们欣然赞同你们在爱心书信开头所提出的主张，即：正是这安慰是必要的，并包含你们忍耐的许多证据；因为我们本应给予的安慰，你们已在信中预先表达了。我们的主惯常赐予受苦者这种忍耐，为使基督的仆人能于陷入磨难时互相鼓励；他们应当自忖：他们所遭受的，以往在圣人们身上也曾发生过。就连我们自己也从你们的书信中获得安慰，因为我们并非与你们的苦难无干，而是在你们内受到磨炼。事实上，谁又能忍受那些本应特别热切追求和平安宁与和谐一致的人所引入的错误呢？但他们不仅不维护和平，反而将无罪的司铎从他们自己教会的前席逐出。我们的弟兄和同工、你们的主教若望，是首位未经审讯便遭受此不义对待的人。没有控告，没有听证。什么主张被废弃了，以至于任何审判的表象都无法产生或被寻求？别人竟坐在仍在世的司铎的位置上，仿佛任何从如此纷争开始的人，在任何人看来都能正当地拥有或做任何事情。我们从未见过我们的先辈有过如此胆大妄为。他们反而禁止此类标新立异，拒绝给予任何人在另一人还在位时被祝圣取代的权力，因为凡是不义地取代者，不能成为主教。\par

关于遵守教规，我们声明：唯有在尼西亚所制定的教规，才应得到天主教会的服从与承认。若有人试图引进其他与尼西亚教规相悖、且由异端者汇编的教规，天主教会应当拒绝此类教规，因为将异端的发明加于天主教教规是不合法的。因为他们总是想通过对手和无法无天的人来贬低尼西亚教父的决定。因此，我们说，我们所谴责的教规不仅应被忽视，而且应与异端及裂教者的教义一同被谴责，正如它们先前在萨迪卡会议中由作为我们前任的主教们所谴责的那样。因为，至可敬的弟兄们，与其确认任何违反教规的行为，不如谴责这些行径更为妥当。\par

对于目前的情况下，我们应当采取何种措施来对付如此行径呢？必须进行一次教务会议式的调查，而我们早已说过应当召集一次教务会议。没有其他办法可以止息这场风暴的狂怒。为了达成这一目标，与此同时，那藉由伟大的天主及其基督、我们的主的旨意而来的治愈得以彰显，将是有益的。这样，我们将看到所有因魔鬼的嫉妒而激起的、并作为我们信德考验的灾祸都得以止息。如果我们恒心持守信德，就没有什么事不应从主那里期待。我们一直在寻找公会议的机会，藉着天主的旨意，这些烦扰的骚动可告终结。那么，让我们在期间忍耐，并以忍耐之墙为堡垒，信赖我们的天主在恢复一切事上的助佑。\par

\Quote{我们早已从我们的同僚主教德默特琉、基里亚古、欧吕修斯和帕拉迪乌斯得知了你们所讲述的一切关于你们的磨难，他们曾在不同时期到访罗马，现与我们同在；我们从他们那里经过全面查问得知了所有细节。}\par

% structure: p0080 source=h2 target=subsection reason=relative-heading-rank; base=h2

\subsection{第二十七章 从对待若望的行为所导致的可怕事件。皇后欧多西亚之死。阿尔撒修斯之死。以及关于宗主教阿提库斯，他的生平和性格。}

这份信件就是\Person{依诺森}所写，从中可以轻易推断出他对\Person{若望}所持的看法。大约同一时期，一些异常巨大的冰雹降落在\Place{君士坦丁堡}和城郊。四天后，皇帝的妻子去世了。许多人将这些事件视为天主义怒的征兆，因为\Person{若望}曾遭受迫害。因为\Place{加采东}主教\Person{西利诺}，\Person{若望}的主要诽谤者之一，早先就在剧烈的身体痛苦中结束了生命，起因是他的脚发生了意外，医生不得不截肢了他的腿。\Person{阿尔撒息}也在主持\Place{君士坦丁堡}教会仅很短一段时间后去世了。许多人被提名为他的继任者；他去世四个月后，\Person{阿提库斯}，\Place{君士坦丁堡}圣职人员中的一位司铎，也是\Person{若望}的敌人之一，被祝圣为主教。他是\Place{亚美尼亚}\Place{塞巴斯特}人。他自幼便由马其顿异端的隐修士们教导隐修哲学的原则。这些隐修士当时因哲学而在\Place{塞巴斯特}享有极高的声誉，他们属于\Person{厄弗斯大提}的会规，之前已经提到过他，他是那里的主教，也是最优秀的隐修士的领袖。当\Person{阿提库斯}成年后，他接受了天主教会的教义。他天性多于学识，参与事务，既擅长策划阴谋，也善于躲避他人的诡计。他性格非常迷人，被许多人喜爱。他在教堂中发表的讲道并不超出平庸；虽然并非完全缺乏学识，但听众认为它们不值得保存成文字。由于专注，只要有机会，他就在最受推崇的希腊作者中练习；但为了避免在谈论这些作家时显得无知，他经常隐藏自己所知。据说他为自己相同见解的人表现出极大的热忱，并使自己在异端者面前显得可畏。当他愿意时，他能轻易使他们惊慌；但他立刻转变自己，显得温顺。这就是了解此人的人所提供的信息。\par

\Person{若望}即使在被流放时也获得了极大的声誉。他拥有充足的金钱资源，而且女执事\Person{奥林比亚}和其他人也慷慨地供给他钱财，他以此从\Place{以扫利亚}强盗手中赎回了许多俘虏，并使他们回到家人身边。他还供给许多有急需之人的需要，并用亲切的话语安慰那些不需要金钱的人。因此，他不仅在他居住的\Place{亚美尼亚}，而且在邻近各国的人民，以及\Place{安提约基雅}和\Place{叙利亚}其他地区及\Place{基里基雅}的居民中深受爱戴，这些人经常寻求与他交往。\par

% structure: p0083 source=h2 target=subsection reason=relative-heading-rank; base=h2

\subsection{第二十八章. \Person{依诺森}，\Place{罗马}主教，通过大公会议召回\Person{若望}的努力。关于他所派去验证此事的人。\Person{金口若望}之死。}

\Place{罗马}主教\Person{依诺森}非常渴望促成\Person{若望}的召回，正如他之前的信件所表明的。他派遣了\Place{罗马}教会的五位主教和两位司铎，连同从东方被委派为使节的主教们，一同去见皇帝\Person{霍诺留}和\Person{阿尔卡狄}，请求召开大公会议，并恳请他们指定时间和地点。\Place{君士坦丁堡}的\Person{若望}的敌人捏造了一项指控，声称这些事情是为了侮辱东方皇帝，并导致使节们可耻地被驱逐，仿佛他们侵犯了外国政府。与此同时，\Person{若望}被一道皇帝敕令判处流放到更偏远的地方，并派士兵将他押送到\Place{毕丢斯}；士兵们很快就到达，并执行了转移。据说在这段旅程中，殉道者\Person{巴西利斯库}在\Place{亚美尼亚}的\Place{科马尼}向他显现，并告知他死亡的日子。他因头痛发作，且无法忍受太阳的热度，无法继续行程，就在那个城镇结束了生命。\par

\SourceRecord{Translated by Chester D. Hartranft.；Nicene and Post-Nicene Fathers, Second Series；Vol. 2.；Edited by Philip Schaff and Henry Wace.；Buffalo, NY: Christian Literature Publishing Co.,；1890.；<http://www.newadvent.org/fathers/26028.htm>.}

% structure: p0001 source=h1 target=section reason=page-title

\section{\WorkTitle{Ecclesiastical History (Book IX)}}

% structure: p0002 source=h2 target=subsection reason=relative-heading-rank; base=h2

\subsection{第一章\newline{} \Person{阿卡狄乌斯}去世，小\Person{狄奥多西}皇帝执政。他的姐妹。\Person{普尔喀丽娅}公主的虔诚、美德与童贞；她蒙天主喜爱的工作；她合宜地教育了皇帝。}

以上是关于\Person{若望}所传下来的细节。在他去世后不久，即\Person{阿提库斯}升任\Place{君士坦丁堡}主教后三年，在\Person{巴苏斯}与\Person{菲利普}担任执政官期间，\Person{阿卡狄乌斯}驾崩。他留下儿子\Person{狄奥多西}——当时刚断奶——作为帝国继承人。他还留下三个幼龄女儿，名为\Person{普尔喀丽娅}、\Person{阿尔卡狄娅}与\Person{玛丽娜}。\par

在我看来，天主有意藉这一时期的事件表明：单凭虔诚便足以保障君主的救恩；没有虔诚，军队、强大的帝国以及一切其他资源都无济于事。守护宇宙的天主上智预见到这位皇帝将因虔诚而卓越，因此决定让他的姐妹\Person{普尔喀丽娅}成为他及其政府的保护者。这位公主尚未满十五岁，却已拥有超越年龄的极有智慧与神圣的心智。她首先将童贞奉献于天主，并教导她的姐妹们走同样的生活道路。为避免一切嫉妒与阴谋的缘由，她不允许任何男子进入她的宫殿。为坚定她的决心，她以天主、司铎及罗马帝国全体臣民为见证，证明她的自我奉献。为表明她的童贞以及她兄弟的元首地位，她在\Place{君士坦丁堡}的圣堂中祝圣了一张圣桌——一件非凡的工艺，极为美丽可观；它以黄金和宝石制成，并在桌面上刻下这些事，使众人皆知。在安静地执掌国事后，她卓越而极有条理地治理罗马帝国；她筹划如此周密，以致应办之事迅速被议决并完成。她能以完美的准确性用希腊语和拉丁语写作和交谈。她使一切事务都以她兄弟的名义办理，并非常用心地以最佳方式将他培养成一位君主，传授适合其年龄的知识。她让最干练的人教导他骑术、兵器操练和文字。但系统地教导他举止有序而有君主风度的，是他的姐姐；她教他如何整理袍服、如何就座、如何行走；她训练他克制笑容，视场合需要而显出温和或威严的面容，并亲切地询问前来呈递请愿者的案件。但她主要努力引导他度虔诚的生活、不断祈祷；她教导他定期前往圣堂，并以礼品和珍宝尊重祈祷之所；她激发他对司铎及其他善人，以及那些依照基督宗教法律献身于哲学的敬畏之心。她热忱而明智地防范宗教因虚假教义的革新而受到危害。在我们时代没有新异端得势，我们将发现这尤其要归功于她，正如我们随后将看到的。她以何等畏惧敬拜天主，任何人若要详述都会花费很长时间；她建造了多少华丽的祈祷之所，建立了多少客栈与隐修团体，以及为它们永久维持的开支安排和院内供应——这一切，若有人愿意从事情本身考察真相，而不被我的言语说服，只要他查阅她家管家所写的证明文件，并从真实记录中查核事实是否与我的历史相符，便会知道我的描述并非为了自我偏爱而虚假。若仅凭这些证据仍不能使他满意而相信，那么让天主亲自说服他——天主因她的行为而完全并处处垂顾她，以致迅速俯听她的祈祷，并时常预先指示应行之事。这种圣爱的标记，若非人以其行为配得，是不会赐予人的。但我情愿暂时略过赐予皇帝姐姐的许多个别圣宠显现——它们证明她是天主所爱的，以免有人责备我本从事别的事，却转而使用赞词。然而，有一件与她相关的事似乎本身以及对我的教会历史而言都如此恰当，又如此明显地为她爱天主的明证，我决定在此叙述，尽管它发生在后来一段时期。此事如下：——\par

% structure: p0005 source=h2 target=subsection reason=relative-heading-rank; base=h2

\subsection{第二章\newline{} 发现四十位神圣殉道者遗骸。}

\Emph{一位名叫尤西比亚的女子}，是马其顿派的女执事，在\Place{君士坦丁堡}城外有一处房屋和花园，其中保存着四十位士兵的圣髑，这些士兵曾在\Person{李锡尼}统治下在\Place{亚美尼亚}的\Place{塞巴斯特}殉道。当她感到死亡临近时，她把上述地点遗赠给几位正统派隐修士，并要他们起誓将她埋葬在那里，并在棺材顶端她的头部上方凿出一个单独的地方，将殉道者的圣髑与她一同存放，且不得告知任何人。隐修士照做了；但为了按照与\Person{尤西比亚}的约定秘密向殉道者表达合宜的敬意，他们在其墓旁建造了一座地下祈祷所。但在显眼之处，地基上建起了一座用烧砖围砌的建筑物，并设有从该处通往殉道者的秘密通道。不久后，\Person{凯撒}——一位权贵，曾晋升至执政官和总督之尊——丧妻，他将妻子安葬在\Person{尤西比亚}墓旁；因为两位女士因最亲密的友谊而紧密相连，在一切教义和宗教问题上意见一致。\Person{凯撒}因此被诱导买下此地，以便能安葬在妻子旁边。上述隐修士定居别处，没有透露任何关于殉道者的事。此后，当建筑被拆毁，泥土和瓦砾被四处丢弃，整个地方被整平。因为\Person{凯撒}本人亲自在那里建造了一座宏伟的圣殿献给天主，以光荣殉道者\Person{蒂尔索}。看来很可能是天主有意让上述地方消失，并让如此长的时间过去，以使殉道者的发现被视为更神奇、更显著的事件，并作为天主眷顾发现者的证明。实际上，发现者正是皇帝之姊——皇后\Person{普尔喀利亚}本人。可敬的\Person{蒂尔索}向她显现了三次，向她揭示了那些隐藏在地下的圣髑；并命令将它们存放在他的墓旁，以便共享同等的位置和尊荣。四十位殉道者自己也向她显现，身披闪耀长袍。但这事件似乎太过神奇而难以置信，且完全不可能；因为该地区的年长神职人员屡次探究之后，未能指明殉道者的位置，其他人也未能做到。最后，当一切无望时，一位名叫\Person{波利克罗尼乌斯}的长老（他曾是\Person{凯撒}府中的仆人）受天主提醒，得知所述地点曾曾有隐修士居住。于是他前往马其顿派的神职人员处打听他们的情况。所有隐修士都已去世，除了一位，他似乎被保全生命，专为指明神圣殉道者圣髑隐藏的地点。\Person{波利克罗尼乌斯}就此仔细询问他，发现因与\Person{尤西比亚}的约定，他的回答有些犹豫不决，便将天主启示和皇后的焦虑，以及她一切尝试的失败告诉了他。隐修士于是承认，天主已向皇后宣告了真理；因为当他还是一个长大的男孩，由年长的领导者教导隐修生活时，他确切记得殉道者的圣髑曾被存放在\Person{尤西比亚}墓旁；但后来的时间流逝和该地方发生的变化，使他无法回想起圣髑是存放在圣堂之下还是任何其他地方。\Person{波利克罗尼乌斯}接着说道：\Quote{我并未遭受类似的记忆失灵，因为我记得我曾出席\Person{凯撒}妻子的葬礼，并且根据大路的相对位置，我能判断她必定被埋葬在读经台之下}；读经台是供读经者使用的讲台。\Quote{因此，}隐修士补充说，\Quote{必须在\Person{凯撒}妻子遗骸附近寻找\Person{尤西比亚}的坟墓；因为两位女士曾以最亲密的友谊和交情生活，并相互约定埋葬在彼此旁边。}当需要根据上述提示挖掘并追踪神圣圣髑，而皇后得知实情后，她下令开始动工。在读经台旁掘土时，按照\Person{波利克罗尼乌斯}的推测发现了\Person{凯撒}妻子的棺材。在侧旁不远处，他们发现烧砖铺地，以及一块同等尺寸的大理石板，每块都与砖同大，下面露出了\Person{尤西比亚}的棺材；近旁是一座小圣堂，优雅地用白色和紫色大理石围砌。墓盖呈圣桌形状，在存放圣髑的顶端可见一个小孔。一位宫廷侍从碰巧站在旁边，将手中的手杖插入小孔；拔出手杖后，他将手杖凑到鼻子前，吸入一股没药的甜香，这使工匠和旁观者重振信心。他们急切地打开棺材，发现了\Person{尤西比亚}的遗骸，在她头部附近是墓穴的突出部分，精确地制成箱形，由其自身的盖子掩盖于内；边缘两侧包裹的铁器用铅牢牢固定。中央再次出现同样的小孔，更清楚地揭示了圣髑隐藏在内的事实。发现一经宣布，他们便跑到殉道者的圣堂，派人叫来铁匠打开铁栓，轻易地取下了盖子。下面发现了大量香料，在香料中找到两只银匣，里面安放着神圣圣髑。随后公主感谢天主，因天主认为她堪当如此伟大的显现，并得以寻获神圣圣髑。此后，她以最珍贵的宝匣尊荣了殉道者；在一场公开庆典结束时（庆典以合宜的荣耀和伴随圣咏的游行举行，我曾在场），圣髑被安放在神一般的\Person{蒂尔索}旁边。其他在场的人也能作证，这些事是按所述方式完成的，因为他们几乎都还健在。这一事件发生得晚得多，当时\Person{普罗克鲁斯}治理\Place{君士坦丁堡}教会。\par

% structure: p0007 source=h2 target=subsection reason=relative-heading-rank; base=h2

\subsection{第三章 普尔喀丽亚的德性；她的姊妹们。}

据说天主在许多其他情况下常常向公主启示将要发生的事，而且大多数都发生在她和她的姊妹们身上，作为天主之爱的见证。她们都追求同样的生活方式；她们殷勤对待司铎和祈祷之所，对贫困的客旅和穷人大方施舍。这些姊妹通常一起用餐和散步，白日黑夜相伴，歌颂天主的赞美。如同模范女子的习惯，她们致力于纺织和类似的工作。虽然身为公主，生在宫中长在宫中，她们却避免轻浮和懒惰，认为这与任何宣认童贞者不相称，因此将这种懈怠远远地弃于自己的生活之外。因此，天主的仁慈彰显出来，并为她们的家族而得胜；因为天主使皇帝年岁增长、治理稳固；一切针对他的阴谋和战争都自行瓦解。\par

% structure: p0009 source=h2 target=subsection reason=relative-heading-rank; base=h2

\subsection{第四章 与波斯休战。霍诺留和斯提利科。在罗马和达尔马提亚的事务。}

虽然波斯人已准备拿起武器，但他们被说服与罗马人缔结了一个为期百年的休战协定。\par

斯提利科，霍诺留的军队统帅，被怀疑曾阴谋宣布他的儿子欧赫里乌斯为东方皇帝，因而在拉文纳被军队杀死。早先，当阿尔卡狄乌斯还在世时，他就已对部下怀有强烈的敌意，因此促使两个帝国发生冲突。他让哥特人的领袖阿拉里克获得罗马将军的职位，并劝他夺取伊利里亚；又派遣被任命的总督约维安先行，并约定很快会合罗马军队，将其臣民置于霍诺留的统治之下。阿拉里克离开了靠近达尔马提亚和潘诺尼亚的蛮族地区——他之前就住在那——率领军队进军埃皮鲁斯；在那里停留一段时间后，他未达成任何目的便撤回意大利。因为他本来要根据协议迁移，却被霍诺留的书信阻止了。阿尔卡狄乌斯去世后，霍诺留计划前往君士坦丁堡代其侄子任命忠于他安全与帝国的官员；因为他视侄儿如子，担心那孩子因年幼而受苦，易遭阴谋；但当霍诺留即将启程时，斯提利科劝阻了他，向他证明他必须留在意大利，以压制君士坦丁的图谋——君士坦丁正试图夺取阿尔勒的统治权。斯提利科随后拿起罗马人称为拉布兰旗的权杖之一，从皇帝那里取得几封书信，率领四个军团出发向东作战；但有人散布谣言说，他阴谋反对皇帝，并与当权者合谋扶植其子登基，于是军队哗变，杀死了意大利和高卢的禁卫军长官、军事统帅和宫廷要员。斯提利科本人在拉文纳被士兵杀死。他几乎获得了绝对的权力，可以说所有的人，无论是罗马人还是蛮族，都受他控制。斯提利科就这样因涉嫌阴谋反对皇帝而丧命。他的儿子欧赫里乌斯也被杀。\par

% structure: p0012 source=h2 target=subsection reason=relative-heading-rank; base=h2

\subsection{第五章 不同的民族拿起武器反抗罗马人，其中一些因天主眷顾而被击败，另一些则被带至和好之约。}

大约同时，驻扎在色雷斯的匈人可耻地撤退，损失了很多人，尽管他们既未受到攻击也没有被追击。住在伊斯特河附近的蛮族部落首领乌尔迪斯率领大军渡河，在色雷斯边境扎营。他背信弃义地夺取了默西亚的一座城市——战神营，并从那里侵入色雷斯其余地区，傲慢地拒绝与罗马人结盟。色雷斯士兵的长官向他提出和平建议，但他指着太阳回答说，只要他愿意，征服这个发光体照耀下的每一个地区对他来说都很容易。然而，当乌尔迪斯发出此类威胁，并随心所欲地要求大量贡品，以此作为与罗马人媾和或继续开战的条件——在境况如此无望之时，天主对当前的统治显出了特殊眷顾的明显证据；因为不久之后，乌尔迪斯的贴身随从和部落首领们开始讨论罗马的政体、皇帝的人道、以及他奖赏最优和善之人的迅速与慷慨。这不是没有天主的安排：他们转而喜爱这些讨论的要点，并投奔了罗马人，连同自己统领的部队一起加入了他们的营地。乌尔迪斯发现自己被抛弃，勉强逃到河对岸。他的许多士兵被杀，其中整个野蛮的斯基里部落也覆灭了。这个部落在此之前人数众多，十分强大。一些人被杀，另一些人被俘，戴上锁链解送到君士坦丁堡。执政者们认为，如果让他们聚在一起，可能会发生叛乱。因此，一些人被低价出卖；另一些人被当作奴隶赠送，条件是他们永远不得返回君士坦丁堡或欧洲，而是被大海与熟悉的地方隔开。其中还有一些未售出，奉命定居在不同的地方。我在比提尼亚的奥林匹斯山附近见过很多人，彼此分开居住，耕种那地区的山丘和山谷。\par

% structure: p0014 source=h2 target=subsection reason=relative-heading-rank; base=h2

\subsection{第六章 哥特人阿拉里克。他袭击了罗马，并以战争围困之。}

东部帝国就这样出乎所有人的意料，因其统治者尚且年幼，而免于战祸且治理有序。与此同时，西部帝国却因许多僭主崛起而陷入混乱。\Person{斯提利科}死后，哥特人的首领\Person{阿拉里克}派遣使节到\Person{霍诺留}处商议和平，但无果。他进军\Place{罗马}，围攻之；并在\Place{台伯河}沿岸驻扎大量蛮族军队，有效地阻止了所有来自\Place{波图斯}的粮食供应入城。围攻持续了一段时间后，城内发生可怕的饥荒和瘟疫，许多奴隶和大多数城内的蛮族出身的居民都逃到阿拉里克那边。那些仍固守异教迷信的元老提议在\Place{卡皮托利山}和其他神庙献祭；而由城守召来的某些\Place{托斯卡纳}人，许诺以雷电驱走蛮族；他们自夸曾在\Place{托斯卡纳}的\Place{拉尼亚}城（阿拉里克曾经过该城前往罗马而未攻占）行过类似的事。然而事实证明，这些人不能为城市带来任何好处。所有有识之士都明白，这次围攻给罗马人带来的灾难，是天主愤怒的征兆，惩罚他们的奢侈、荒淫以及彼此之间和对陌生人的诸多不义行为。据说，当阿拉里克向罗马进军时，\Place{意大利}的一位善心的隐修士恳求他饶恕该城，不要成为如此多灾难的肇因。阿拉里克回答说，他并不打算开始围攻，但有一种不可抗拒的力量推动并命令他去攻打罗马；最终他这样做了。当他围攻该城时，居民向他献了许多礼物，一度他解除了围攻，因为罗马人同意说服皇帝与他订立和平条约。\par

% structure: p0016 source=h2 target=subsection reason=relative-heading-rank; base=h2

\subsection{第七章 \Person{依诺增爵}，\Place{罗马}长老团的主教。他派遣使团到\Person{阿拉里克}处。\Person{约维乌斯}，意大利总督。派遣使团到皇帝处。关于阿拉里克的事件。}

虽然派遣了使节去商议和平，但皇帝朝廷中\Person{阿拉里克}的敌人却竭力阻止与他缔结任何和约。但此后，当\Place{罗马}主教\Person{依诺增爵}派遣使团到他那里，并且皇帝用书信召见阿拉里克时，他前往\Place{阿里米努姆}城，该城距\Place{拉文纳}二百一十斯塔迪亚。\par

他在城外扎营；\Place{意大利}总督\Person{约维乌斯}与他会谈，并将他的要求转达给皇帝，其中一项是：他应被谕令任命为骑兵和步兵将军。皇帝授予约维乌斯全权，可以赐给\Person{阿拉里克}他想要的任何金钱和物资，但拒绝赐予他这一尊位。约维乌斯轻率地在阿拉里克的营地等候宫廷的使者；并下令当着所有蛮族的面宣读皇帝的决定。当发现尊位被拒绝时，阿拉里克被结果激怒，命令吹响号角，向\Place{罗马}进军。约维乌斯担心被皇帝怀疑与阿拉里克勾结，做了一个更不谨慎的举动：他以皇帝的安全起誓，并迫使主要官员发誓绝不与阿拉里克达成任何和平条款。然而，这位蛮族首领不久后又改变了主意，传话说他不想要任何职位尊严，而是愿意作为罗马人的盟友，只要他们赐给他一定数量的谷物和一些对他们而言次要的领土，以便他可以安居。\par

% structure: p0019 source=h2 target=subsection reason=relative-heading-rank; base=h2

\subsection{第八章 \Person{阿塔鲁斯}及其将军\Person{赫拉克良}的叛乱；以及他最终如何跪在\Person{霍诺留}脚下乞求宽恕。}

在先后两次派遣几位主教作为使者就此事进行商议，均无果而返之后，阿拉里克返回罗马，包围了该城；他占领了波尔图的一部分，并迫使罗马人承认当时担任城长的阿塔卢斯为他们的国王。当罗马人被任命担任其他官职时，阿拉里克被任命为骑兵与步兵统帅，其妻弟阿陶尔夫被提升为所谓禁卫骑兵的指挥官。阿塔卢斯召集了元老院，向他们发表了一篇冗长而精心准备的演说，其中他许诺恢复元老院的古老习俗，并将埃及及其他东方省份置于意大利的统治之下。这是一个注定连一年王位都无法保持之人的狂妄之语。他被一些占卜者的说法所欺骗，这些人向他保证他可以兵不血刃地征服阿非利加；他不听从阿拉里克的劝告，后者敦促他派遣一支适中的部队前往迦太基，若奥诺留的官员试图抵抗，便将他们处死。他也拒绝听从约翰的建议——约翰被他授予王室近卫军的指挥权——约翰建议他在康斯坦斯即将启程前往利比亚时，交给他一份他们称为《法令》的文件，仿佛是由奥诺留发出，凭此可使赫拉克良被罢免职务；赫拉克良曾被委托掌管阿非利加的军队。若采用此计，很可能成功，因为阿塔卢斯的图谋在利比亚尚无人知晓。但康斯坦斯一按照占卜者的建议启航前往迦太基，阿塔卢斯就心智软弱到认为此事毫无悬念，相信阿非利加人会如占卜者预言的那样成为他的臣民，并亲率大军向拉文纳进发。当消息传来，阿塔卢斯已抵达阿里米努姆，率领着一支由罗马人和蛮族混合组成的军队，奥诺留写信给他，承认他为皇帝，并派遣宫廷最高官员前去拜见他，提出与他分享帝国。然而，阿塔卢斯拒绝与他人共享权力，并传话给奥诺留，说他可以选一个岛屿或任何他喜欢的地方作为私人居所，并会获准享有所有皇帝的待遇。奥诺留的处境已危急到如此地步，以至于船只已备好，准备在必要时送他到他外甥那里；这时，一支四千人的军队从西方出发，于夜间出乎意料地抵达拉文纳；奥诺留用这支增援部队守卫城墙，因他不信任本土部队，认为他们有叛变的倾向。\par

与此同时，赫拉克良处死了康斯坦斯，并在阿非利加的港口和海岸沿线部署军队，以阻止商船前往罗马。结果，罗马发生了饥荒，他们派遣代表向阿塔卢斯报告此事。阿塔卢斯不知如何是好，返回罗马征询元老院的意见。饥荒如此严重，以至于人们用栗子代替粮食，有些人甚至被怀疑食用了人肉。阿拉里克建议派遣五百名蛮族前往阿非利加对抗赫拉克良，但元老院和阿塔卢斯反对，认为不应将阿非利加托付给蛮族。于是，阿拉里克清楚地看出天主不赞同阿塔卢斯的统治；并意识到为一桩力所不及之事徒劳无益，在获得某些保证之后，他与奥诺留达成协议，剥夺阿塔卢斯的王权。所有相关方在城外集结，阿塔卢斯丢弃了皇权的标志。他的官员们也解下了腰带，奥诺留对所有这些事件表示宽恕，每个人都保留他们最初拥有的荣誉和职位。阿塔卢斯与他的儿子退隐到阿拉里克那里，因为他认为，若继续留在罗马人中，自己的生命尚不安全。\par

% structure: p0022 source=h2 target=subsection reason=relative-heading-rank; base=h2

\subsection{第九章　关于阿塔卢斯在希腊人与基督徒中引起的骚动；勇猛的萨罗斯；阿拉里克用计夺取罗马，并保护了宗徒伯多禄的圣所。}

阿塔卢斯计划受挫，令外教人和阿里乌斯异端的基督徒深为不满。外教人根据阿塔卢斯众所周知的倾向和早年教育推断，他会公开维护他们的迷信，恢复他们古老的庙宇、节日和祭台。阿里乌斯派想象，一旦阿塔卢斯发现自己的统治稳固，就会重新使他们获得在君士坦提乌斯和瓦伦斯统治时期所享有的教会首席地位；因为他已由哥特人主教西格萨里乌斯施洗，这使阿拉里克和阿里乌斯派极为满意。\par

不久之后，阿拉里克驻扎在阿尔卑斯山，距拉文纳约六十斯塔迪昂，并与皇帝就缔结和平一事举行会谈。萨罗斯，一个蛮族出身、极富战争经验的人，手下仅有约三百人，但都是忠诚且极精干的士兵。他因过去的仇怨而猜疑阿拉里克，并推断罗马人与哥特人之间的条约对他毫无益处。他突然率部挺进，杀死了一些蛮族人。阿拉里克因这一事件既愤怒又恐惧，便退兵返回罗马，用计夺取了该城。他允许每个部下尽可能多地夺取罗马人的财富，并洗劫所有房屋；但出于对宗徒伯多禄的尊敬，他命令将建造在他坟墓周围的那座大而宽敞的教堂作为庇护所。这是唯一阻止罗马全城毁灭的原因；那些在那里获救的人——人数众多——重建了城市。\par

% structure: p0025 source=h2 target=subsection reason=relative-heading-rank; base=h2

\subsection{第十章　一位行谦逊之事的罗马贵妇。}

显然，像\Place{罗马}这样的大城市被攻陷，必然伴随着许多非凡的事件。因此，我现在将叙述那些在教会史中值得记载的事件。我将讲述一个野蛮人所行的虔诚之举，并记录一位罗马妇女为保全其贞洁而表现出的勇敢。这个野蛮人和这位妇女都是基督徒，但并非属于同一异端：前者是\OriginalTerm{亚略派信徒}{Arian}，后者是\OriginalTerm{尼西亚信理}{Nicene doctrines}的狂热追随者。\Person{阿拉里克}的一名年轻士兵看见这位非常美丽的妇女，被她的美貌征服，试图强行与她发生关系；但她退缩，竭力避免被玷污。他拔出剑，威胁要杀死她；但他因对她的情欲而克制，只是在她脖子上划了一道轻伤。鲜血大量涌出，她把脖子迎向剑刃；她宁愿死在贞洁中，也不愿在合法地与丈夫结合后，再被另一个男人侵犯而存活下来。当这个野蛮人重复他的意图，并以更可怕的威胁紧随其后时，他未能得逞；他对她的贞洁感到惊奇，便将她带到宗徒\Person{伯多禄}的教堂，给了守卫教堂的官员六个金币作为她的生活费用，并命令他们为她丈夫看护她。\par

% structure: p0027 source=h2 target=subsection reason=relative-heading-rank; base=h2

\subsection{第11章。当时在西方反抗\Person{霍诺里乌斯}的僭主。他们因皇帝对天主的爱而全被消灭。}

在此期间，许多僭主在西方政府中反抗\Person{霍诺里乌斯}。有些彼此相攻，另一些则以奇妙的方式被擒，从而显明天主对\Person{霍诺里乌斯}的眷顾非同寻常。\Place{不列颠}的士兵首先起来叛乱，他们拥立\Person{马克}为僭主。但后来他们又杀了\Person{马克}，拥立了\Person{格拉提安}。四个月之内，他们又杀了\Person{格拉提安}，改立\Person{君士坦丁}，以为凭着这个名字，他就能将帝国牢牢地控制在自己手中；而正是出于这个原因，其他几个同名的人也当选了僭主。\Person{君士坦丁}从不列颠渡海至高卢滨海城市\Place{博诺尼亚}，在诱使\Place{高卢}和\Place{阿基坦}的所有军队支持他的事业后，他使从分隔\Place{意大利}与\Place{高卢}的山脉（罗马人称之为\Place{科蒂安阿尔卑斯}）直到那片地区的居民都服从他。然后他派其长子\Person{君士坦斯}（他先前已任命他为\OriginalTerm{凯撒}{Cæsar}，后来又立他为皇帝）前往\Place{西班牙}。\Person{君士坦斯}在控制了这个行省并任命了自己的总督后，命令将\Person{霍诺里乌斯}的亲属\Person{狄迪莫斯}和\Person{维里尼安}戴上镣铐，带到他面前。\Person{狄迪莫斯}和\Person{维里尼安}起初彼此不和，但当他们发现自己面临同样的危险时，便和解了。他们联合了主要由武装农民和奴隶组成的军队，共同进攻\Place{卢西塔尼亚}，杀死了许多僭主派来捉拿他们的士兵。\par

% structure: p0029 source=h2 target=subsection reason=relative-heading-rank; base=h2

\subsection{第12章。——\Person{狄奥多西奥卢斯}和\Person{拉戈迪乌斯}。汪达尔人和苏维汇人的族群。\Person{阿拉里克}之死。僭主\Person{君士坦丁}和\Person{君士坦斯}的逃亡。}

不久之后，\Person{君士坦斯}的军队得到了增援，\Person{狄迪莫斯}和\Person{维里尼安}及其妻子被俘，最终被处死。他们的兄弟\Person{狄奥多西奥卢斯}和\Person{拉戈迪乌斯}住在其他行省，逃离了本国；前者逃往\Place{意大利}，到皇帝\Person{霍诺里乌斯}那里；后者逃往东方，到\Person{狄奥多西}那里。在这些事件之后，\Person{君士坦斯}回到他父亲那里，他已在通往\Place{西班牙}的路上部署了自己的士兵作为守卫；因为他没有允许西班牙人按照古老的习俗担任守卫，那是他们所请求的特权。这一预防措施后来很可能成了国家覆亡的原因；因为当\Person{君士坦丁}被剥夺权力时，野蛮的汪达尔人、苏维汇人和阿兰人族群壮了胆，征服了那条路，占领了\Place{西班牙}和\Place{高卢}的许多堡垒和城市，并逮捕了僭主的官员。\par

与此同时，\Person{君士坦丁}仍然认为事情会如他所愿，便立他的儿子为皇帝以取代\OriginalTerm{凯撒}{Cæsar}之位，并决心占领\Place{意大利}。为此，他越过了\Place{科蒂安阿尔卑斯}，进入\Place{利古里亚}的\Place{利韦罗纳}城。他正准备渡过\Place{波河}时，得知\Person{阿拉维库斯}死亡的消息，被迫折返。这位\Person{阿拉维库斯}曾是\Person{霍诺里乌斯}军队的统帅，因被怀疑密谋将整个西方政府置于\Person{君士坦丁}的统治之下，在照例应在皇帝前面行进的凯旋游行返回途中被杀害。事件发生后，皇帝立即下马，公开感谢天主将他从公然密谋反对他的人手中解救出来。\Person{君士坦丁}逃走并占领了\Place{阿尔勒}，他的儿子\Person{君士坦斯}从\Place{西班牙}赶来，在同城中寻求庇护。\par

在\Person{君士坦丁}的势力衰落时，汪达尔人、苏维汇人和阿兰人听说\Place{比利牛斯山}是一个繁荣且物产极其丰富的地区，便急切地攻取了它。由于\Person{君士坦斯}曾委托看守山口的守卫玩忽职守，入侵者便进入了\Place{西班牙}。\par

% structure: p0033 source=h2 target=subsection reason=relative-heading-rank; base=h2

\subsection{第13章。论及\Person{杰伦提乌斯}、\Person{马克西穆斯}和\Person{霍诺里乌斯}的军队。\Person{杰伦提乌斯}及其妻子的被俘；他们的死亡。}

\Emph{与此同时}，\Person{杰伦提乌斯}曾是\Person{君士坦丁}将领中最得力的一位，如今却成了他的敌人。他相信他的密友\Person{马克西穆斯}很有资格当僭主，便为他披上皇袍，允许他住在\Place{塔拉戈纳}。然后\Person{杰伦提乌斯}进军攻打\Person{君士坦丁}，并设法在\Place{维也纳}杀死了\Person{君士坦丁}的儿子\Person{君士坦斯}。\par

君士坦丁一听到马克西穆斯篡位，便派遣他的将军之一埃多维库斯渡过莱茵河，征集法兰克人和阿勒曼尼人的军队；又派其子君士坦斯守卫维也纳及邻近城镇。随后格隆提乌斯进军阿尔勒并围攻该城；但不久，当奥诺留的军队在君士坦提乌斯（其后成为罗马皇帝的瓦伦提尼安之父）的指挥下前来讨伐僭主时，格隆提乌斯带着少数士兵仓促撤退；因为他的大部分部队已投奔君士坦提乌斯的军队。西班牙士兵因格隆提乌斯撤退而对他极度蔑视，便商议如何杀他。他们密集列队，夜间袭击他的住所；但他与一位朋友阿拉努斯及少数仆从登上屋顶，用箭矢射击，致使不下三百名士兵丧命。箭矢用尽后，仆从们悄悄从建筑物上缒下逃脱；格隆提乌斯虽也可同样获救，却因对妻子农尼基亚的爱恋而不愿如此。次日黎明，士兵向房屋纵火；他见逃生无望，便应朋友阿拉努斯之愿砍下其头。此后，他的妻子哀泣着，泪流满面地以剑自逼，恳求在落入他人之手前死于丈夫之手，并乞求他赐予这最后的恩典。这女子以其勇气证明了自己配得上她的信仰，她是一名基督徒，就这样仁慈地死去；她为自己留下了永不磨灭的纪念。格隆提乌斯随后以剑自刺三下；但见未受致命伤，便拔出腰间匕首，刺入心脏。\par

% structure: p0036 source=h2 target=subsection reason=relative-heading-rank; base=h2

\subsection{第十四章　君士坦丁。奥诺留的军队及其将军埃多维库斯。埃多维库斯被君士坦丁的将军乌尔菲拉斯击败。埃多维库斯之死。}

尽管阿尔勒城被奥诺留的军队严密包围，君士坦丁仍坚持抵抗围攻，因为消息称埃多维库斯即将率众多盟友抵达。这使奥诺留的将军们极度恐慌。于是他们决定返回意大利，在那里继续战争。当他们商定此计划时，埃多维库斯已被报出现在附近，于是他们渡过罗讷河。步兵指挥官君士坦提乌斯静候敌军逼近，而君士坦提乌斯的同僚将军乌尔菲拉斯率骑兵埋伏在不远处。敌军经过乌尔菲拉斯的部队，正要与君士坦提乌斯的军队接战时，信号发出，乌尔菲拉斯突然出现，从后方袭击敌军。敌军立即溃逃。有的逃脱，有的被杀，另一些人则放下武器求饶，得以幸免。埃多维库斯上马逃往一位名叫埃克迪西乌斯的地主处，他曾为后者立下大功，因此以为他是朋友。然而，埃克迪西乌斯砍下他的头，献给奥诺留的将军们，希望获得重赏和荣誉。君士坦提乌斯收到首级后，宣布公开感谢埃克迪西乌斯，但乌尔菲拉斯的功劳却被归于他；当埃克迪西乌斯急切要陪伴他时，他命令他离开，因为他认为与一个恶毒的主人同行对自己和军队都不好。这个胆敢对身处困境的朋友和客人犯下最不圣洁谋杀的人——此人正如谚语所说，空张着嘴离去了。\par

% structure: p0038 source=h2 target=subsection reason=relative-heading-rank; base=h2

\subsection{第十五章　君士坦丁抛弃帝王权柄的象征，被祝圣为司铎；其后之死。其他阴谋反对奥诺留的僭主之死。}

此胜之后，奥诺留的军队再次围攻该城。君士坦丁听到埃多维库斯的死讯，便脱下紫袍和帝王饰物，前往教堂，让人祝圣自己为司铎。城内之人先得到誓言保证，然后打开城门，保全了性命。从此，整个省份重新效忠奥诺留，并一直服从他任命的地方长官。君士坦丁及其子尤利安被押送意大利，但途中遭伏击被杀。不久之后，上述僭主约维亚努斯、马克西穆斯、萨罗斯以及许多其他阴谋反对奥诺留的人意外被杀。\par

% structure: p0040 source=h2 target=subsection reason=relative-heading-rank; base=h2

\subsection{第十六章　统治者奥诺留，爱天主者。奥诺留之死。其继承人瓦伦提尼安及其女奥诺利亚；当时普世和平。}

此处不宜详细叙述暴君们的死因；但我认为有必要提及此事，以表明要确保帝国权力的稳固，皇帝只需虔诚地事奉\OriginalTerm{天主}{God}即可，洪诺留正是这样做的。他的同胞妹妹\Person{加拉·普拉西狄亚}与他同住，也同样以维护宗教和教会的真诚热心而著称。在勇将\Person{君士坦修斯}消灭了暴君\Person{君士坦丁}之后，皇帝将妹妹许配给他作为奖赏；还赐予他白鼬皮和紫袍，并让他分享统治权。君士坦修斯不久后去世，留下两个孩子：\Person{瓦伦提尼安}（后继承洪诺留）和\Person{霍诺利亚}。此时东方帝国太平无事，出乎所有人的意料，其事务被治理得井井有条，因为统治者尚是一个青年。看来\OriginalTerm{天主}{God}公开显示了祂对当今皇帝的眷顾，不仅以出人意料的方式处置了战事，而且还启示了许多古时最为敬虔之人的神圣遗体；其中发现了非常古老的先知\Person{匝加利亚}和由宗徒们祝圣为执事的\Person{斯德望}的遗骸。我似乎有责任描述这一方式，因为每一次发现都是奇妙而神圣的。\par

% structure: p0042 source=h2 target=subsection reason=relative-heading-rank; base=h2

\subsection{第十七章 发现先知\Person{匝加利亚}及\OriginalTerm{首位殉道者}{Proto-Martyr}\Person{斯德望}的遗骸。}

我先谈先知的遗骸。\Place{加法尔-匝加利亚}是巴勒斯坦城市\Place{厄勒乌特罗波利斯}地区的一个村庄。这块土地由一名叫\Person{卡莱梅鲁斯}的农奴耕种；他对主人很好，但对邻近的农民却刻薄不满且不公正。尽管他有这些性格缺陷，先知在梦中向他显现，指着一个特定的花园对他说：\Quote{去，在那花园中挖，距离通往彼瑟里比斯城的路旁篱笆两肘远。你会在那里找到两口棺材，里面的是木制的，外面的是铅制的。在棺材旁边，你会看到一个装满水的玻璃瓶，还有两条中等大小的蛇，但很驯服且完全无害，似乎习惯于被人触碰。}卡莱梅鲁斯按照先知的指示在指定地点热心地着手工作。当神圣的安放处以上述迹象被揭示时，神圣的先知穿着白色长巾向他显现，这使我认为他是\OriginalTerm{司铎}{priest}。在他的脚边棺材外面，躺着一个享有皇家葬礼的孩子；因为他的头上戴着金冠，脚上穿着金凉鞋，身上穿着昂贵的袍子。当时的智者和司铎对这个孩子是谁、来自何处以及为何如此穿戴感到十分困惑。据说在\Place{革拉里}的一个隐修院院长\Person{匝加利亚}找到了一份用希伯来语写的古老文件，它未被收入\OriginalTerm{正典书卷}{canonical books}中。该文件记载：当先知\Person{匝加利亚}被犹大王\Person{约阿士}处死后，王室很快遭到一场可怕的灾祸；因为在先知死后第七天，约阿士的一个他十分疼爱的儿子突然夭折。国王判断这场灾祸是\OriginalTerm{天主义怒}{Divine wrath}的特殊显示，便下令将儿子葬在先知脚边，作为对先知所犯罪行的某种\OriginalTerm{赎罪}{atonement}。这就是我所查明的关于此事的细节。\par

尽管先知在地下躺了那么多代，他看起来完好无损；他的头发剪得很短，鼻子挺直，胡须适中，头相当短，眼睛略显凹陷，被眉毛遮掩着。\par

\SourceRecord{Translated by Chester D. Hartranft.；Nicene and Post-Nicene Fathers, Second Series；Vol. 2.；Edited by Philip Schaff and Henry Wace.；Buffalo, NY: Christian Literature Publishing Co.,；1890.；<http://www.newadvent.org/fathers/26029.htm>.}
